%% Adapted copy for submission to:
%%   The Annals of Statistics [AOS]
%% Generated from the repository manuscript without editing the original .tex file.
%% Please do not edit imsart.cls or imsart.sty.

\documentclass[aos]{imsart}

%% Packages
\RequirePackage[english]{babel}
\RequirePackage{amsthm,amsmath,amsfonts,amssymb}
\RequirePackage[authoryear]{natbib}
\RequirePackage{graphicx}
\RequirePackage{subcaption}
\RequirePackage{xcolor}
\definecolor{forestgreen}{rgb}{0.0, 0.43, 0.13}
\RequirePackage{comment}
\RequirePackage{enumitem}
\RequirePackage{stackrel}
\RequirePackage{textcomp}
\RequirePackage{bbm}
\RequirePackage{booktabs}
\RequirePackage{array}
\RequirePackage{multirow}
\RequirePackage[colorlinks=true,citecolor=blue,urlcolor=blue,linkcolor=blue]{hyperref}

\newif\ifincludeimages
% \includeimagestrue  % Change to \includeimagesfalse to deactivate images, \includeimagestrue to activate them
 \includeimagesfalse  % Change to \includeimagesfalse to deactivate images, \includeimagestrue to activate them

\allowdisplaybreaks

\startlocaldefs
% Macros
\newcommand{\lp}{\left(}
\newcommand{\rp}{\right)}
\newcommand{\lc}{\left[}
\newcommand{\rc}{\right]}
\newcommand{\lb}{\left\{}
\newcommand{\rb}{\right\}}
\newcommand{\lf}{\left.}
\newcommand{\ri}{\right.}
%
\newcommand{\G}{\mathbb{G}}
\newcommand{\R}{\mathbb{R}}
\newcommand{\T}{\mathbb{T}}
\newcommand{\Sp}{\mathbb{S}}
\newcommand{\Z}{\mathbb{Z}}
\newcommand{\N}{\mathbb{N}}
%
\newcommand{\rd}{\mathrm{d}}
\newcommand{\cL}{\mathcal{L}}
\newcommand{\cA}{\mathcal{A}}
\newcommand{\cY}{\mathcal{Y}}
\newcommand{\cF}{\mathcal{F}}
\newcommand{\bx}{\boldsymbol{x}}
\newcommand{\be}{\boldsymbol{e}}
\newcommand{\bm}{\boldsymbol{m}}
\newcommand{\bn}{\boldsymbol{n}}
\newcommand{\by}{\boldsymbol{y}}
\newcommand{\bX}{\boldsymbol{X}}
\newcommand{\bY}{\boldsymbol{Y}}
\newcommand{\bZ}{\boldsymbol{Z}}
\newcommand{\bJ}{\boldsymbol{J}}
\newcommand{\bU}{\boldsymbol{U}}
\newcommand{\bR}{\boldsymbol{R}}
\newcommand{\bF}{\boldsymbol{F}}
\newcommand{\bmu}{\boldsymbol\mu}
\newcommand{\bu}{\boldsymbol{u}}
\newcommand{\bv}{\boldsymbol{v}}
\newcommand{\bb}{\boldsymbol{b}}
\newcommand{\bba}{\boldsymbol{a}}
\newcommand{\br}{\boldsymbol{r}}
\newcommand{\bs}{\boldsymbol{s}}
\newcommand{\bt}{\boldsymbol{t}}
\newcommand{\bz}{\boldsymbol{z}}
\newcommand{\bh}{{\boldsymbol{h}}}
\newcommand{\bc}{\boldsymbol{c}}
\newcommand{\bd}{\boldsymbol{d}}
\newcommand{\bk}{\boldsymbol{k}}
\newcommand{\bp}{\boldsymbol{p}}
\newcommand{\bpsi}{\boldsymbol{\psi}}
\newcommand{\one}{\mathbf{1}}
\newcommand{\zero}{\mathbf{0}}
\newcommand{\two}{\mathbf{2}}
\newcommand{\bxi}{\boldsymbol\xi}
\newcommand{\bphi}{\boldsymbol\varphi}
\newcommand{\ba}{\boldsymbol\alpha}
\newcommand{\bbeta}{\boldsymbol\eta}
\newcommand{\bbe}{\boldsymbol\beta}
\newcommand{\bzeta}{\boldsymbol\zeta}
\newcommand{\bga}{\boldsymbol\gamma}
\newcommand{\bthe}{\boldsymbol\theta}
\newcommand{\btheta}{\boldsymbol\theta}
\newcommand{\bkappa}{\boldsymbol\kappa}
\newcommand{\bTheta}{\boldsymbol\Theta}
\newcommand{\Ical}{\mathcal{I}}
\newcommand{\bsigma}{\boldsymbol\sigma}
\newcommand{\bSigma}{\boldsymbol\Sigma}
\newcommand{\bLambda}{\boldsymbol\Lambda}
\newcommand{\bGamma}{\boldsymbol\Gamma}
\newcommand{\bDelta}{\boldsymbol\Delta}
\newcommand{\bomega}{\boldsymbol\omega}
\newcommand{\bO}{\mathbf{0}_q}
\newcommand{\bOd}{\mathbf{0}_{q-1}}
\newcommand{\Hcal}{\mathcal{H}}
\newcommand{\bHcal}{\boldsymbol{\mathcal{H}}}
\newcommand{\X}{\boldsymbol{\mathcal{X}}_{\bx,p}}
\newcommand{\W}{\boldsymbol{\mathcal{W}}_{\bx}}
\newcommand{\bnabla}{\boldsymbol\nabla}
\newcommand{\blambda}{\boldsymbol\lambda}
\newcommand{\bB}{\boldsymbol{B}}
\newcommand{\bK}{\boldsymbol{K}}
\newcommand{\bH}{\boldsymbol{H}}
\newcommand{\bA}{\boldsymbol{A}}
\newcommand{\bV}{\boldsymbol{V}}
\newcommand{\bI}{\boldsymbol{I}}
\newcommand{\bL}{\boldsymbol{L}}
\newcommand{\bS}{\boldsymbol{S}}
\newcommand{\bP}{\boldsymbol{P}}
\newcommand{\bQ}{\boldsymbol{Q}}
\newcommand{\bT}{\boldsymbol{T}}
\newcommand{\bM}{\boldsymbol{M}}
\newcommand{\bW}{\boldsymbol{W}}
\newcommand{\pr}{^{\prime}}
\newcommand{\cqfd}{\hfill $\square$}
%
\newcommand{\equald}{\stackrel{d}{=}}
\newcommand{\inlaw}{\stackrel{d}{\rightsquigarrow}}

% More macros
\newcommand{\Xb}{\bX}
\newcommand{\Sb}{\mathbf{S}}
\newcommand{\thetab}{\boldsymbol\theta}
\newcommand{\etab}{\boldsymbol\eta}
\newcommand{\bnab}{\boldsymbol\nabla}
\newcommand{\xib}{\boldsymbol\xi}
\newcommand{\Vb}{\mathbf{V}}
\newcommand{\Qb}{\mathbf{Q}}
\newcommand{\Rb}{\mathbf{R}}
\newcommand{\Zb}{\mathbf{Z}}
\newcommand{\ab}{\mathbf{a}}
\newcommand{\ub}{\mathbf{u}}
\newcommand{\Tb}{\mathbf{T}}
\newcommand{\vb}{\mathbf{v}}
\newcommand{\xb}{\bx}
\newcommand{\yb}{\by}
\newcommand{\wb}{\mathbf{w}}
\newcommand{\Ab}{\mathbf{A}}
\newcommand{\Bb}{\bB}
\newcommand{\Hb}{\mathbf{H}}
\newcommand{\Ob}{\mathbf{O}}
\newcommand{\Cb}{\mathbf{C}}
\newcommand{\Db}{\mathbf{D}}
\newcommand{\Ub}{\mathbf{U}}
\newcommand{\Mb}{\mathbf{M}}
\newcommand{\Wb}{\mathbf{W}}
\newcommand{\Yb}{\mathbf{Y}}
\newcommand{\Lb}{\mathbf{L}}
\newcommand{\Gb}{\mathbf{G}}
\newcommand{\sigmad}{\sigma_{d}}
\newcommand{\sigmar}{\sigma_{\bd}}
\newcommand{\sigmai}{\sigma_{d_i}}
\newcommand{\sigmaim}{\sigma_{d_i-1}}
\newcommand{\sigmaj}{\sigma_{d_j}}
\newcommand{\sigmamj}{\sigma_{-d_j}}
\newcommand{\sigmamjk}{\sigma_{-(d_j,d_k)}}
\newcommand{\sigmak}{\sigma_{d_k}}
\newcommand{\sigmaell}{\sigma_{d_\ell}}
\newcommand{\sigmajm}{\sigma_{d_j-1}}
\newcommand{\sigmarm}{\sigma_{\bd-1}}
\newcommand{\sigmarmij}{\sigma_{d_1-1,\stackrel{i\neq j}{\ldots},d_r-1}}
\newcommand{\Sdr}{(\mathbb{S}^{d})^r}
\newcommand{\Sr}{\mathbb{S}^{\bd}}
\newcommand{\Srm}{\mathbb{S}^{\bd-1}}
\newcommand{\Sd}{\mathbb{S}^{d}}
\newcommand{\Sdm}{\mathbb{S}^{d-1}}
\newcommand{\Sj}{\mathbb{S}^{d_j}}
\newcommand{\Sk}{\mathbb{S}^{d_k}}
\newcommand{\Sell}{\mathbb{S}^{d_\ell}}
\newcommand{\Sjk}{\mathbb{S}^{d_j,d_k}}
\newcommand{\Smj}{\mathbb{S}^{\bd}_{-j}}
\newcommand{\Smjk}{\mathbb{S}^{\bd}_{-(j,k)}}
\newcommand{\Sjm}{\mathbb{S}^{d_j-1}}
\newcommand{\Sim}{\mathbb{S}^{d_i-1}}
\newcommand{\Hd}{\mathbb{H}^{d}}

% Commands with arguments
\newcommand{\lrp}[1]{\left(#1\right)}
\newcommand{\lrc}[1]{\left[#1\right]}
\newcommand{\lrb}[1]{\left\{#1\right\}}
\newcommand{\lrpbig}[1]{\big(#1\big)}
\newcommand{\lrcbig}[1]{\big[#1\big]}
\newcommand{\lrbbig}[1]{\big\{#1\big\}}
\newcommand{\blrp}[1]{\big(#1\big)}
%
%\newcommand{\Prob}[1]{\mathbb{P}\lb #1\rb}
%\newcommand{\E}[1]{\E\lp #1\rp}
\newcommand{\V}[1]{\mathbb{V}\mathrm{ar}\lc #1\rc}
\newcommand{\Es}[2]{\E_{#2}\lc #1\rc}
\newcommand{\Vs}[2]{\mathbb{V}\mathrm{ar}_{#2}\lc #1\rc}
\newcommand{\Cov}[2]{\mathbb{C}\mathrm{ov}\lc #1,#2\rc}
\newcommand{\supp}[1]{\mathrm{supp}\lp #1\rp}
\newcommand{\diag}[1]{\mathrm{diag}\lp #1\rp}
\newcommand{\cmod}[1]{\mathrm{cmod}\left(#1\right)}
%
\newcommand{\mse}[1]{\mathrm{MSE}\lrc{#1}}
\newcommand{\mise}[1]{\mathrm{MISE}\lrc{#1}}
\newcommand{\amise}[1]{\mathrm{AMISE}\lrc{#1}}
%
\newcommand{\df}[2]{\frac{\rd #1}{\rd #2}}
\newcommand{\dftwo}[2]{\frac{\rd^2 #1}{\rd #2^2}}
\newcommand{\pf}[2]{\frac{\partial #1}{\partial #2}}
\newcommand{\pftwo}[2]{\frac{\partial^2 #1}{\partial #2^2}}
\newcommand{\pfmix}[3]{\frac{\partial^2 #1}{\partial #2\partial #3}}
\newcommand{\norm}[1]{\left|\left| #1\right|\right|}
\newcommand{\abs}[1]{\left| #1\right|}
\newcommand{\tr}[1]{\mathrm{tr}\left[#1\right]}
\newcommand{\vect}[1]{\mathrm{vec}\left(#1\right)}
\newcommand{\vech}[1]{\mathrm{vech}\left(#1\right)}
\newcommand{\scalprod}[2]{\langle#1,#2\rangle}
%
\newcommand{\vlinel}[1]{\multicolumn{1}{|c}{#1}}
\newcommand{\vliner}[1]{\multicolumn{1}{c|}{#1}}
%
\newcommand{\alert}[1]{\textcolor{red}{#1}}

% Orders
\newcommand{\Om}[1]{\Omega_{#1}}
\newcommand{\om}[1]{\omega_{#1}}
\newcommand{\Iq}[2]{\int_{\Omega_{q}} #1\,\omega_{q}(d #2)}
\newcommand{\Iqr}[3]{\int_{\Omega_{q}\times\R} #1\,d #3\,\omega_{q}(d #2)}
\newcommand{\Iqp}[3]{\int_{\Omega_{q}\times\R^p} #1\,d #3\,\omega_{q}(d #2)}
\newcommand{\Iqm}[2]{\int_{\Omega_{q-1}} #1\,\omega_{q-1}(d #2)}
\newcommand{\Ir}[2]{\int_{\R} #1\,d #2}
\newcommand{\Iqq}[3]{\int_{\Omega_{q_1}\times\Omega_{q_2}} #1\,\omega_{q_2}(d #3)\,\omega_{q_1}(d #2)}
\newcommand{\Irr}[3]{\int_{\R\times\R} #1\,d #3\,d #2}
%
\DeclareFontFamily{OT1}{pzc}{}
\DeclareFontShape{OT1}{pzc}{m}{it}{<-> s * [1.10] pzcmi7t}{}
\DeclareMathAlphabet{\mathpzc}{OT1}{pzc}{m}{it}
\newcommand{\order}[1]{\mathpzc{o}\lp#1\rp}
\newcommand{\Order}[1]{\mathcal{O}\lp#1\rp}
\newcommand{\orderp}[1]{\mathpzc{o}_\mathbb{P}\lp#1\rp}
\newcommand{\Orderp}[1]{\mathcal{O}_\mathbb{P}\lp#1\rp}

% New counter
%\newcommand{\co}{\addtocounter{align}{1}\arabic{align}}

% New theorems
\theoremstyle{plain}
\newtheorem{theorem}{Theorem}[section]
\newtheorem{corollary}{Corollary}[section]
\newtheorem{proposition}{Proposition}[section]
\newtheorem{lemma}{Lemma}[section]

\theoremstyle{definition}
\newtheorem{definition}{Definition}[section]
\newtheorem{remark}{Remark}[section]
\newtheorem{algo}{Algorithm}[section]
\newtheorem{example}{Example}[section]

% Roman enumerate
\renewcommand{\theenumi}{\textit{\roman{enumi}}}

% Percentage sign without typing \%
\newcommand{\pct}{\%}

% Definition
\newcommand{\defin}{:=}

% Random forests
\DeclareMathOperator*{\argmin}{arg\,min}
\DeclareMathOperator*{\argmax}{arg\,max}
\DeclareMathOperator*{\arccosh}{arccosh}
\DeclareMathOperator{\arctantwo}{arctan2}
\newcommand{\convd}{\stackrel{d}{\rightarrow}}
\newcommand{\convl}{\stackrel{\mathcal{L}}{\rightarrow}}
\newcommand{\convp}{\stackrel{\mathsf{P}}{\rightarrow}}
\newcommand{\convas}{\stackrel{a.s.}{\longrightarrow}}
\newcommand{\convn}{\stackrel[n\to\infty]{}{\longrightarrow}}
\newcommand{\Prob}[1]{\mathsf{P}\left(#1\right)}
\newcommand{\Probbig}[1]{\mathsf{P}\big(#1\big)}
\newcommand{\Probstar}[1]{\mathsf{P}_*\left(#1\right)}
\newcommand{\Probbigstar}[1]{\mathsf{P}_*\big\{ #1 \big\}}
\newcommand{\Condprob}[2]{\mathsf{P}\left(#1 \mid #2\right)}
\newcommand{\Condprobbig}[2]{\mathsf{P}\big(#1 \mid #2\big)}
\newcommand{\E}[1]{\mathsf{E}( #1 )}
\newcommand{\Elr}[1]{\mathsf{E}\left( #1 \right)}
\newcommand{\Ebig}[1]{\mathsf{E}\big( #1 \big)}
\newcommand{\EBig}[1]{\mathsf{E}\Big( #1 \Big)}
\newcommand{\CondE}[2]{\mathsf{E}\left(#1 \mid #2\right)}
\newcommand{\CondEbig}[2]{\mathsf{E}\big(#1 \mid #2\big)}
\newcommand{\dY}{d_\mathcal{Y}}
\newcommand{\dH}{d_\mathcal{H}}
\newcommand{\Roob}{\widehat{R}^{\mathrm{oob}}}
\newcommand{\Var}[1]{\mathsf{Var}( #1 )}

\endlocaldefs

\begin{document}

\begin{frontmatter}
\title{Goodness-of-fit for distributions on metric spaces}
\runtitle{Goodness-of-fit for distributions on metric spaces}
\runauthor{Serrano, García-Portugués and Van Keilegom}

\begin{aug}
\author[A]{\fnms{Diego}~\snm{Serrano}}
\author[A]{\fnms{Eduardo}~\snm{García-Portugués}}
\author[B]{\fnms{Ingrid}~\snm{Van Keilegom}}

\address[A]{Department of Statistics, Carlos III University of Madrid}
\address[B]{ORSTAT, KU Leuven}
\end{aug}

\begin{abstract}
We propose a general goodness-of-fit framework for distributions on separable metric spaces. Under suitable identifiability conditions, probability distributions are characterized by distance profiles, which motivates their use in goodness-of-fit testing, for simple and composite null hypotheses. For composite null hypotheses, parameter estimation is incorporated via a Bahadur-type expansion, and the asymptotic distribution of the empirical process for distance profiles is obtained under the null. We define test statistics based on this empirical process and derive their asymptotic null distributions. We further study the behavior of the proposed tests under fixed and local alternatives, establishing consistency results. The methodology is illustrated with simulation studies and real-data experiments.
\end{abstract}

\begin{keyword}[class=MSC]
\kwd[Primary ]{62G10}
\kwd[; secondary ]{62R20}
\end{keyword}

\begin{keyword}
\kwd{Goodness-of-fit}
\kwd{Metric spaces}
\kwd{Distance profiles}
\end{keyword}

\end{frontmatter}
%---------------------------%
\section{Basic concepts and notation}
%---------------------------%
Consider a separable metric space $(\Omega, d)$. Let $(S, \Sigma, \mathsf{P})$ be a probability space, where $S$ is a sample space, $\Sigma$ is a sigma algebra of subsets of $S$, and $\mathsf{P}$ is a probability measure. An $\Omega$-valued random variable is called a random object, which is a measurable map $X: S \rightarrow \Omega$. Suppose that $X$ is distributed according to a Borel probability measure $\mu$, that is, $\mu(A)=\Prob{\{s \in S: X(s) \in A\}}=: \Prob{X \in A} = \mathsf{P}\left(X^{-1}(A)\right)$, for any Borel measurable set $A \subseteq \Omega$. 
For any $t \geq 0$, we define the \emph{distance profile} at $\omega \in \Omega$ of $X$ as
$$
F_\omega^\mu(t)=\Prob{d(\omega, X) \leq t}.
$$
To estimate the distance profiles $F_\omega^\mu$ from a sample $X_1, \ldots, X_n$ of independent realizations of $X$, $n\ge1$, we use empirical estimates
$$
F_{n,\omega}^\mu(t)\defin\frac{1}{n} \sum_{i=1}^n 1_{\left\{d\left(\omega, X_i\right) \leq t\right\}}, \quad t \geq 0.
$$

Distance profiles characterize the underlying probability measures under suitable conditions. 
Consider two Borel probability measures $\mu_1$ and $\mu_2$ on $(\Omega,d)$. \cite[][Proposition 1]{DubeyChenMuller2024} provide a sufficient condition for $\mu_1=\mu_2$ in terms of distance profiles: if there exists $\theta>0$ such that $(\Omega,d^\theta)$ is of strong negative type \citep{Lyons2013}, where $d^\theta(\omega,\omega^\prime):=d(\omega,\omega^\prime)^\theta$, then 
\begin{align}\label{eq:characterization_Omega}
\mu_1=\mu_2
\quad\Longleftrightarrow\quad
F_{\omega}^{\mu_1}(t)=F_{\omega}^{\mu_2}(t)
\ \ \text{for all } \omega\in\Omega,\ t\ge 0.
\end{align}


\cite{ChenDubey2025} obtain an alternative characterization under a condition imposed on the probability measures. 
Let $\mu$ be a Borel probability measure on $(\Omega,d)$. If there exists $L_\mu<\infty$ such that, for every $\omega\in \supp{\mu}$,
\begin{align}\label{eq:doubling}
\limsup_{r\downarrow 0}\frac{\mu\big(B(\omega,2r)\big)}{\mu\big(B(\omega,r)\big)}\le L_\mu,
\end{align}
then $\mu$ is said to satisfy the doubling condition.
If $\mu_1,\mu_2$ satisfy this condition, then Theorem 1 in \cite{ChenDubey2025} shows that
\begin{align}\label{eq:characterization_support}
\mu_1=\mu_2
\quad\Longleftrightarrow\quad
\mu_1\big(\overline B(\omega,t)\big)=\mu_2\big(\overline B(\omega,t)\big)
\ \ \text{for all } \omega\in \supp{\mu_1},\ t>0,
\end{align}
i.e., $\mu_1=\mu_2$ whenever $F_\omega^{\mu_1}(t)=F_\omega^{\mu_2}(t)$ for all $\omega\in\supp{\mu_1}$ and $t\ge 0$. The same result holds with $\supp{\mu_1}$ replaced by $\supp{\mu_2}$. 
Here, $\supp{\mu}$ denotes the support of $\mu$, which is the closure of the set of all points to which $\mu$ assigns positive mass in every open neighborhood, i.e.,
$$
\supp{\mu}=\operatorname{cl}(\{x \in \Omega: \mu(U)>0 \text { for all open sets } U \ni x\}).
$$


\citet[][p. 6]{ChenDubey2025} further discuss the set of values of $\omega$ for which \eqref{eq:characterization_Omega} and \eqref{eq:characterization_support} hold. 
In ibid., it is emphasized that, under the doubling condition \eqref{eq:doubling},
$\mu_1$ and $\mu_2$ coincide as long as they agree on all closed balls centered in 
\emph{either} of their supports \eqref{eq:characterization_support}. 
Under the strong negative type condition, they claim that the characterization in \eqref{eq:characterization_Omega} holds if $\mu_1$ and $\mu_2$ agree on all closed balls centered in \emph{both} of their supports,
i.e., $\omega\in\supp{\mu_1}\cup \supp{\mu_2}$.
However, in Proposition~1 of \cite{DubeyChenMuller2024}, condition \eqref{eq:characterization_Omega} is written in terms of $\omega \in \Omega$ (the whole metric space),
so the claim is false unless $\Omega\subset \supp{\mu_1}\cup \supp{\mu_2}$. 


In this manuscript, we assume that at least one of the following two conditions holds:
\begin{enumerate}[label=(\textit{a.\arabic*}), ref=(\textit{a.\arabic*})]
    \item $(\Omega, d)$ is of strong negative type. \label{a1} %and $\Omega\subset \supp{\mu_1}\cup \supp{\mu_2}$.
    
    %Examples: Spheres with the geodesic metric and the space of $n\times n$ network adjacency matrices with the Frobenius metric with bounded entries. Other examples of separable metric spaces of strong negative type include the space of one-dimensional distributions with finite second moments endowed with the $2$-Wasserstein metric, and the space of multivariate distributions with $L^2$ metric between cumulative distribution functions, but they are not totally bounded in general unless, for example, the support of the distributions is compact.  See Section 3 of \citet{Lyons2013} for a detailed discussion of examples and counterexamples of metrics that are of (strong) negative type. Examples of metric spaces that are not of strong negative type include Grassmannian manifolds and cylinders with their geodesic metrics.
    
    \item The probability measures satisfy the doubling condition. \label{a2}
    
    %Examples: Let $f: \R^n \rightarrow \R$ be Lebesgue measurable with $0<c \leq f \leq C$ a.e. and define $\mu(A)=\int_A f \, \rd \mathcal{L}^n$. Then $\mu$ is doubling (and $\R^n$ with the Euclidean metric is separable). The volume measure on a compact Riemannian manifold (which is a separable space with the geodesic distance) and the Lebesgue measure on the Heisenberg group (also separable with the Carnot–Carathéodory metric) are doubling. The Dirac measure and the volume in a hyperbolic space are not globally doubling. 
\end{enumerate}

Assumption \ref{a1} is a condition on the metric space, whereas \ref{a2} is an assumption on the probability measures.

Under certain conditions, \cite[][Theorem 5.1]{DubeyChenMuller2024} and \citet[][Theorem 3]{ChenDubey2025} establish that the distance profiles satisfy the Donsker condition. 
Let
\begin{align}\label{eq:def_Y}
\cY \defin \{y_{\omega, t}:(\omega, t)\in\Omega\times\mathbb R_{\ge0}\}, \ \text{ where } y_{\omega, t}(x):=1_{\{d(\omega,x)\le t\}}.
\end{align}
Then, $\cY$ is $\mu$-Donsker.
The conditions under which this property is established for distance profiles differ between \cite{DubeyChenMuller2024} and \cite{ChenDubey2025}. 


\cite{DubeyChenMuller2024} require the following two conditions:
\begin{enumerate}[label=(\textit{b.\arabic*}), ref=(\textit{b.\arabic*})]
\item\label{b1}
Let $N(\varepsilon,\Omega,d)$ be the covering number of $(\Omega,d)$, for $\varepsilon>0$. Then
$$
\varepsilon \log N(\varepsilon,\Omega,d)\ \to\ 0
\text{ as } \varepsilon\to 0^+.
$$
\item\label{b2}
For every $\omega\in\Omega$, $F_\omega^\mu$ is absolutely continuous with continuous density $f_\omega$. Defining
$$
\underline\Delta_\omega \defin \inf_{t\in \supp{F_\omega}} f_\omega(t),
\qquad
\bar\Delta_\omega \defin \sup_{t\in\mathbb R} f_\omega(t),
$$
it holds that $\underline\Delta_\omega>0$ for each $\omega\in\Omega$ and there exists $\bar\Delta<\infty$ such that $\sup_{\omega\in\Omega}\bar\Delta_\omega \le \bar\Delta$.
\end{enumerate}
%---------------------------%
\begin{remark}
%---------------------------%
    Condition~\ref{b1} implies that $(\Omega,d)$ is totally bounded. The assumption that $d(\omega,X)$ has a continuous distribution for every $\omega\in\Omega$, which is implied by \ref{b2}, dates back to \cite{Szabados1987} (page 383), who replaced the supremum over all $\omega \in \Omega$ in the Kolmogorov--Smirnov test statistic with maximization over a finite sample, as in our Proposition \ref{prop:sup_max_ks_simple}. 
    To achieve this, \cite{Szabados1987} restricted the analysis to \emph{continuous measures}, i.e., measures $P$ such that $P(B(\omega,t))$ depends continuously on $t>0$ for any $\omega \in \Omega$.
\end{remark}
%---------------------------%


\cite[][Assumption 1]{ChenDubey2025} assume the bracketing condition
\begin{align}\label{eq:bracketing_assumption_1}
\varepsilon \log N_{[]}\big(\varepsilon,\cY, L^1(P)\big)\ \to\ 0
\quad \text{as } \varepsilon\to 0^+,
\end{align}
where $N_{[]}(\varepsilon,\cY,L^1(P))$ denotes the $L^1(P)$ bracketing number of $\cY$. 
The condition given in \eqref{eq:bracketing_assumption_1} holds if both \ref{b1} and a modification of \ref{b2} are satisfied \citep[][Proposition~4]{ChenDubey2025}. 
The modification is:
\begin{enumerate}[label=(\textit{b.\arabic*'}), ref=(\textit{b.\arabic*'})]
\setcounter{enumi}{1}
\item\label{bp2}
There exists a constant $L_G>0$ such that
$$
\sup_{\omega\in \Omega}\big|F_\omega^\mu(r)-F_\omega^\mu(s)\big|
\ \le\ L_G|r-s|,
$$
\end{enumerate}
for every $r,s\in \R_{\ge0}$.


Assumption 1 in \cite{ChenDubey2025} does not require the ambient spaces to be totally bounded. 
A sufficient condition for this assumption to hold is the following. Assume that for any $\varepsilon>0$ there exists a closed ball $\overline{B}(r_\varepsilon)\subset \Omega$ of radius $r_\varepsilon>0$, such that
$\Prob{X\notin \overline{B}(r_\varepsilon)}\le \varepsilon.$ 
If \ref{bp2} holds and
\begin{align}\label{eq:bracketing_assumption_1_modified}
\varepsilon \log N\big(\varepsilon, \overline{B}(r_\varepsilon), d\big) \to 0
\quad \text{as }\varepsilon\to 0^+,
\end{align}
then Assumption 1 in \cite{ChenDubey2025} is satisfied. 
Examples of distributions on Euclidean spaces satisfying \eqref{eq:bracketing_assumption_1_modified} and \ref{bp2} include sub-Gaussian and sub-exponential distributions. 

In this manuscript, we will assume \ref{b1} and \ref{b2} hold.
\begin{comment}
We summarize where Assumptions \ref{b1} and \ref{b2} are used.

\begin{itemize}
    \item Assumption \ref{b1} implies total boundedness. This is used whenever we assume $\mathrm{diam}(\Omega)<\infty$, as in Lemma~\ref{ap:lem_bound_int_g_epsilon}, which is invoked in the proofs of Theorems~\ref{thm:test_statistics_simp}, \ref{thm:test_statistics_comp}, \ref{thm:local_alt_CvM_simple}, and \ref{thm:local_alt_Gn_composite}. 
\item Assumption \ref{b2} is omnipresent in this manuscript, although it could be explored whether it can be replaced with \ref{bp2}. It is used in the proofs of:
\begin{itemize}
    \item Theorems~\ref{thm:test_statistics_simp}, \ref{thm:test_statistics_comp}, \ref{thm:fixed_alternative_divergence_cvm}, and \ref{thm:fixed_alt_divergence_cvm_composite}.
    \item Lemmas~\ref{lem:delta_regularity}, \ref{lem:robust_separation}, \ref{lem:positive_mass_region}, and \ref{ap:lem_DCT_P_hat}.
\end{itemize} 
\end{itemize}
\end{comment}

%---------------------------%
\section{Distance profiles on particular metric spaces}\label{sec:dp_examples}
%---------------------------%

This section aims to provide closed formulas for the distance profile of a random variable $X$ following a specific distribution on a given metric space, with probability measure denoted by $\mu$.


%-------------------------------%
\begin{example}[multivariate normal distribution on Euclidean space]\label{ex:dp_mvn}
%-------------------------------%
Consider a random variable $\bX\sim \mathcal{N}_q(\bmu, \bSigma)$ taking values in $(\R^q, d)$, with $\bSigma$ positive definite and $d$ the Euclidean distance in $\R^q$. The distribution of $d(\bomega,\bX)^2=(\bomega-\bX)^\top(\bomega-\bX)$ characterizes %the distance profile at $\bomega\in\R^q$.
$F_{\bomega}$ at $\bomega\in\R^q$. It can be shown that
$$
(\bomega-\bX)^\top(\bomega-\bX) = \blambda^\top\bP,
$$
where $\bSigma=\bU\bLambda\bU^\top$ is the spectral decomposition of $\bSigma$, with $\bLambda=\mathrm{diag}(\lambda_1,\ldots,\lambda_q)$, $\boldsymbol{\eta}\defin \bU^\top(\bX-\bomega)$, $\boldsymbol{\nu}\defin \bU^\top(\bmu-\bomega)$, and $P_i\defin \eta_i^2/\lambda_i$. It holds that $P_i$ are independent, and $P_i\sim \chi_1^2\left(\nu_i^2/\lambda_i\right)$, $i=1,\ldots,q$. In the special case when $\bomega = \bmu$, then $P_i \sim \chi_1^2$.
\end{example}

%-------------------------------%
\begin{example}[\texorpdfstring{the Brownian motion in $L^2[0,1]$}{the Brownian motion in L2[0,1]}]\label{ex:dp_brownian}
%-------------------------------%
We study the distance profile of the Brownian motion $B(t)$, $t \in [0,1]$, as an element in the metric space $(L^2[0,1], d_{L^2})$, where $d_{L^2}$ denotes the distance induced by the $L_2$ norm $\|\cdot\|_{L^2}$.
The distribution of $\|B-\omega\|_{L^2}^2$ characterizes $F_{\omega}$ at $\omega\in L^2[0,1]$. It can be shown that
$$
\|B-\omega\|_{L^2}^2 = \sum_{i=1}^\infty X_i^2,
$$
where $X_i$ are independent, with $X_i\sim\mathcal{N}(c_i,\lambda_i)$, $c_i\defin\langle \omega,e_i\rangle$, $\lambda_i=(i-\frac{1}{2})^{-2}\pi^{-2}$. Here, $\{e_i\}_{i=1}^\infty$ are the eigenfunctions of the covariance kernel of $B$. In particular, when $\omega(t)=0$ for all $t\in[0,1]$, then $X_i^2 = \lambda_i Q_i,$ where $(Q_i)_{i\ge1}$ are iid and $Q_i\sim\chi_1^2$.
\end{example}

%-------------------------------%
\begin{example}[von Mises--Fisher (vMF) distribution on the sphere]\label{ex:dp_vmf}\label{subsec:dp_vmf}
%-------------------------------%
Consider the unit sphere $\mathbb{S}^{q}\defin\{\bx \in \R^{q+1}: \, \| \bx \|_2 = 1\}$.
Assume $\bX \sim \mathrm{vMF}(\bmu, \kappa)$, for some given $\bmu \in \mathbb{S}^{q}$, $\kappa>0$.
For $\bx \in \mathbb{S}^{q}$, the density of $\bX$ is given by $f_\mathrm{vMF}(\bx;\bmu,\kappa) \propto e^{\kappa\bmu^\top\bx}.$ We denote the normalization constant by $c_{q}^\mathrm{vMF}(\kappa)$.
% Here, $c_{q}^\mathrm{vMF}(\kappa) = \frac{\kappa^{\frac{q-1}{2}}}{ (2\pi)^\frac{q+1}{2} \mathcal{I}_{\frac{q-1}{2}}(\kappa)}$, where $\mathcal{I}_{q}$ denotes the modified Bessel function of the first kind of order $q$.

Proposition \ref{prop:dp_examples} characterizes the distance profile of $\bX$ at $\bomega\in \mathbb{S}^{q}$ under two common distances on $\mathbb{S}^{q}$: the chordal distance, given by $ d(\bX, \bomega)^2 = 2(1-\bomega^\top\bX)$; and the geodesic distance, given by $ d(\bX, \bomega) = \cos^{-1}(\bomega^\top\bX)$.
To characterize $F_{\bomega}$, the distribution of $T=\bomega^\top\bX$ plays a central role. Its density is given by
\begin{align}\label{eq:density_t_vmf}
    f_T(t) = \frac{c_{q}^\mathrm{vMF}(\kappa) e^{\kappa t \bmu^\top\bomega } (1-t^2)^{\frac{q-2}{2}}}{ c_{q-1}^\mathrm{vMF}(\kappa\sqrt{(1-t^2)(1-(\bmu^\top\bomega)^2)})}, \quad  -1 < t < 1.
\end{align}
\end{example}

%-------------------------------%
\begin{example}[hyperboloid von Mises--Fisher distribution (HvMF) on the hyperboloid]\label{ex:dp_hvmf}
%-------------------------------%
Consider the Minkowski pseudo-inner product $(\bx, \by)\defin -x_{1} y_{1}+\sum_{i=2}^{q+1} x_i y_i$ on the unit hyperboloid $\mathbb{H}^q\defin\{\bx \in \R^{q+1}: \allowbreak (\bx, \bx)=-1,\, x_{1}>0\}$. In this surface, distances are calculated as $d_{\mathbb{H}^q}(\bx, \by)=\arccosh \left( -(\bx, \by)\right)$.

The analogous distribution to the vMF on $\mathbb{H}^q$ has density $f_\mathrm{HvMF}(\by ; \bmu, \kappa) \propto \exp \left\{\kappa(\by, \bmu)\right\}$, for $\by, \bmu \in \mathbb{H}^q$, and $\kappa>0$, with respect to the Lebesgue measure on $\mathbb{H}^q$ \citep{Barndorff-Nielsen1978,Jensen1981}. We denote the normalization constant by $c_q^{\mathrm{HvMF}}$. For $\bX \sim \mathrm{HvMF}(\bmu, \kappa)$ and a fixed $\bomega \in \mathbb{H}^q$, set $\rho\defin d_{\mathbb{H}^q}(\bomega,\bmu)$. An important quantity in the characterization of $F_{\bomega}$ is $Z_r \defin d_{\mathbb{H}^q}(\bX,\bomega)$, with density given by
\begin{align}\label{eq:density_zr_hvmf}
    f_{Z_r}(z_r) = \frac{c_{q}^\mathrm{HvMF}(\kappa)}{c_{q-1}^\mathrm{vMF}(\kappa\sinh(z_r)\sinh(\rho))} (\sinh(z_r))^{q-1} e^{\kappa(\bmu, \bomega) \cosh(z_r)}, \quad z_r \ge 0.
\end{align}
\end{example}

%The tangent-normal decomposition on the unit hyperboloid allows expressing any $\bx \in \mathbb{H}^q$ as
%$$
%\bx = \cosh(r) \bomega + \sinh(r) \bu, \quad \text{where } (\bu, \bu) = 1 \text{ and } (\bomega, \bu) = 0,
%$$
%where $r \defin \arccosh(-(\bx, \bomega))$, that is, $(\bx, \bomega) = -\cosh(r)$.
%
%\textcolor{red}{Key aspect to discuss: We have that $(\bu, \bu)=1$. This is not a unit sphere, so it is not straightforward that we should integrate $\bu$ in $\mathbb{S}^{q-1}$. A quick research says that the induced Riemannian metric on the tangent space (the restriction of the Minkowski inner product) is positive definite, hence it gives a standard Riemannian volume on the points in the tangent space such that $(u,u)=1$. An isometry can be defined so that integration over $\mathbb{S}^{q-1}$ is actually justified. I would need to learn more about this; perhaps it is a well-known fact? I will now focus on creating the R code for validation, and see if the derived density is correct.}
%
%%$-(\bmu, \bomega) \ge 1$ (Cauchy-Schwarz)
%Now, decompose $\bmu$ as
%$$
%\bmu = -(\bmu, \bomega) \bomega + \sqrt{(\bmu, \bomega)^2-1} \, \bv,
%$$
%where $(\bv, \bv)=1$ and $(\bomega, \bv) = 0$. Then
%$$
%(\bmu, \bx) = (\bmu, \bomega) \cosh(r) + \sinh(r) \sqrt{(\bmu, \bomega)^2 - 1}  \, (\bv, \bu).
%$$
%
%The induced Riemannian measure on $\mathbb{H}^q$ in these coordinates \citep[see][]{Jensen1981} is
%$$
%d\sigma_{\mathbb{H}^q}(\bx) = (\sinh(r))^{q-1} dr \, d\sigma_{\mathbb{S}^{q-1}}(\bu).
%$$

\begin{comment}
\textcolor{blue}{I think we should remove Examples \ref{ex:dp_wishart} and \ref{ex:dp_dirichlet}}.
%-------------------------------%
\begin{example}[the \textcolor{blue}{Wishart distribution on the space of symmetric positive definite matrices}]\label{ex:dp_wishart}
%-------------------------------%
\textcolor{blue}{We focus on $\mathcal{S}_q^+$, the space of Symmetric Positive Definite (SPD) matrices with real entries of size $q\times q$, $q\in \mathbb{Z}_{>0}$. The space $\mathcal{S}_q^+$ can be endowed with the affine-invariant distance \citep{Pennec2006}}
$$
\textcolor{blue}{d_{\mathrm{AI}}(\bS_1, \bS_2)^2 \defin \| \log\big( \bS_1^{-1/2} \bS_2 \bS_1^{-1/2} \big) \|^2_\mathrm{F} = \sum_{i=1}^{q} \log(\lambda_i)^2,}
$$
\textcolor{blue}{where $\|\cdot\|_\mathrm{F}$ denotes the Frobenius norm, $\log(\cdot)$ stands for the matrix logarithm, and the $\lambda_i$'s are the eigenvalues of $\bS_1^{-1} \bS_2$.}

\textcolor{blue}{Consider a Wishart-distributed random matrix in $\mathcal{S}_q^+$, $\bS \sim \mathrm{Wishart}_q(d, \bSigma)$, for $\bSigma \in \mathcal{S}_q^+$ and $d \in \mathbb{Z}_{>0}$.}
\textcolor{blue}{For a fixed $\bW \in \mathcal{S}_q^+$, let $\{\lambda_i\}_{i=1}^q$ denote the eigenvalues of $\bW^{-1}\bS$. Proposition~\ref{prop:dp_examples} expresses the distance profile in terms of these eigenvalues. Moreover, \cite{Zanella2009} gives the joint density of the ordered eigenvalues of a Wishart-distributed random matrix, which characterizes the density of $\sum_{i=1}^q \log(\lambda_i)^2$. However, obtaining an explicit form for $F_{\bW}(t)$ is challenging.}
\end{example}

%-------------------------------%
\begin{example}[the Dirichlet distribution on the simplex]\label{ex:dp_dirichlet}
%-------------------------------%
Let $\bX \sim \mathrm{Dir}(\boldsymbol{\alpha})$ be a random vector on the simplex $\Delta^{q} = \left\{ \bx \in \R^{q+1} : x_i \ge 0, \ \sum_{i=1}^{q+1} x_i = 1 \right\}$,
with parameters $\boldsymbol{\alpha} = (\alpha_1, \dots, \alpha_{q+1})$, $\alpha_i > 0$. Let $\alpha_0 = \sum_{i=1}^{q+1} \alpha_i$.
The density of $\bX$ is given by
$$
f_{\bX}(\bx) = \frac{1}{B(\boldsymbol{\alpha})} \prod_{i=1}^{q+1} x_i^{\alpha_i - 1},
\qquad B(\boldsymbol{\alpha}) = \frac{\prod_{i=1}^{q+1} \Gamma(\alpha_i)}{\Gamma(\alpha_0)}.
$$
The Dirichlet distribution can be expressed in terms of gamma random variables. Specifically, consider a sequence of independent random variables $G_i$, where $G_i \sim \Gamma(\alpha_i, 1)$, and let $S = \sum_{i=1}^{q+1} G_i$. If we define $X_i = \frac{G_i}{S}$, then $\bX$ follows a Dirichlet distribution $\mathrm{Dir}(\boldsymbol{\alpha})$, and $S \sim \Gamma(\alpha_0, 1)$ is independent of $\bX$.

For a fixed $\bomega \in \Delta^{q}$, we study the distance profile
$$
F_{\bomega}^{\mu}(t) = \Probbig{d(\bX, \bomega) \le t}, \qquad t \ge 0,
$$
under two different distances: the Euclidean distance, and the Aitchison distance (induced by the central log-ratio transform).

The Euclidean distance on the simplex is defined by
$$
d_2(\bx, \bomega) = \|\bx - \bomega\|_2.
$$
The Aitchison distance between $\bx$, $\bomega \in \Delta^{q}$ is
$$
d_A(\bx, \bomega)
= \|\mathrm{clr}(\bx) - \mathrm{clr}(\bomega)\|_2,
$$
where the centered log-ratio transform of $\bx$, denoted by $\mathrm{clr}(\bx)$, is defined by
$$
\mathrm{clr}(\bx)
= \big( \log x_1 - \bar{\ell}, \dots, \log x_{q+1} - \bar{\ell} \big),
\quad
\text{ where } \quad  \bar{\ell} = \frac{1}{q+1} \sum_{k=1}^{q+1} \log x_k.
$$
The clr map is a bijection between $\Delta^{q}$ and the hyperplane
$H\defin\{\bz \in \R^{q+1} : \sum_{i=1}^{q+1} z_i = 0\}$. %The Aitchison distance in the simplex is isometric to the Euclidean norm restricted to $\bH$.

In the special case $q=1$, writing $\bX=(X_1,1-X_1)$, one has $X_1\sim \mathrm{Beta}(\alpha_1,\alpha_2)$. Then
$$
d_2(\bX,\bomega)=\sqrt{2}\,|X_1-\omega_1|.
$$
For the Aitchison distance in the same case $q=1$,
$$
\mathrm{clr}(\bX)
= \left( \frac{1}{2} \log \left(\frac{X_1}{1-X_1}\right), -\frac{1}{2} \log\left(\frac{X_1}{1-X_1}\right) \right),}
$$
$$
d_A(\bX, \bomega)
= \frac{\sqrt{2}}{2} \left| \log \left( \frac{X_1/(1-X_1)}{\omega_1/(1-\omega_1)} \right) \right|,}
$$
where $\mathrm{logit}(x) = \log\big(\frac{x}{1-x}\big)$, and $\mathrm{logit}^{-1}(x) = \frac{e^x}{1 + e^x}$.
\end{example}
\end{comment}

%-------------------------------%
\begin{example}[logistic Gaussian distribution on the simplex]\label{ex:dp_logistic_gaussian}
%-------------------------------%
Consider the simplex $\Delta^{q} = \{ \bx \in \R^{q+1} : x_i > 0, \ \sum_{i=1}^{q+1} x_i = 1 \}$, endowed with the Aitchison distance $d_A(\bx, \bomega) \defin \|\mathrm{clr}(\bx) - \mathrm{clr}(\bomega)\|_2.$ 
The centered log-ratio transform is defined by
$
\mathrm{clr}(\bx)
= \big( \log x_1 - \bar{\ell}, \dots, \log x_{q+1} - \bar{\ell} \big), 
\bar{\ell} \defin \frac{1}{q+1} \sum_{k=1}^{q+1} \log x_k.
$
Let $\mathbf{Y} \sim \mathcal{N}_{q+1}(\boldsymbol{\mu}, \boldsymbol{\Sigma})$, with $\boldsymbol{\Sigma}$ positive definite, and consider its softmax transformation
$$
\mathbf{X} \defin \bsigma(\mathbf{Y}) = \left( \frac{e^{Y_1}}{\sum_{j=1}^{q+1} e^{Y_j}}, \dots, \frac{e^{Y_{q+1}}}{\sum_{j=1}^{q+1} e^{Y_j}} \right).
$$
Then, $\mathbf{X}$ follows a logistic Gaussian distribution on $\Delta^{q}$,
with parameters $\boldsymbol{\mu}$ and $\boldsymbol{\Sigma}$.

Denote by $\bH = \bI_{q+1} - \frac{1}{q+1}\boldsymbol{1}_{q+1}\boldsymbol{1}_{q+1}^\top$ the centering matrix projecting onto $\{\bz \in \R^{q+1} : \sum_{i=1}^{q+1} z_i = 0\}$. Let $\bm \defin \bH\bmu - \mathrm{clr}(\bomega)$ and $\boldsymbol{C} \defin \bH\bSigma \bH^\top$, and let $\boldsymbol{C} = \bU \bA \bU^\top$ be a spectral decomposition with $\bA = \mathrm{diag}(\alpha_1,\ldots,\alpha_q,0)$, $\alpha_1,\ldots,\alpha_q>0$. Define $\delta_i \defin (\bu_i^\top\bm)^2\alpha_i^{-1}$, $i=1,\ldots,q$, where $\bu_i$ denotes the $i$-th column of $\bU$. It can be shown that $d_A(\bX, \bomega)^2 = \boldsymbol{\alpha}^\top \bQ,$ where $\boldsymbol{\alpha}=(\alpha_1,\ldots,\alpha_q)^\top$, and $(Q_i)_{i=1}^q$ are independent with $Q_i\sim \chi_1^2(\delta_i)$.
%---------------------------%
\begin{remark}[connection between isometric and central log-ratio transformations]\label{rmk:ilr_clr_equivalence}
%---------------------------%
The Aitchison distance can be defined using the isometric log-ratio (ilr) transform \citep[see, e.g.,][]{Egozcue2003} instead of clr, that is, $d_A(\bx, \bomega) = \| \mathrm{ilr}(\bx) - \mathrm{ilr}(\bomega) \|_2$. See the proof of Proposition~\ref{prop:dp_examples} for this equivalence.
\end{remark}
\end{example}

%---------------------------%
\begin{proposition}\label{prop:dp_examples}
%---------------------------%
    We provide distributional characterizations of the distance profiles, under the notation of the examples in Section \ref{sec:dp_examples}:
\begin{itemize}
\item If $\bX\sim \mathcal{N}_q(\bmu,\bSigma)$, with $\bSigma$ positive definite, then %with the notation of Example~\ref{ex:dp_mvn}, 
$F_{\bomega}(t)=\Prob{\blambda^\top\bP \le t^2}$.
%where $P_i$ are independent, and $P_i\sim \chi_1^2\left(\mu_i^2/\lambda_i\right)$.
% If additionally, $\bSigma=\bI_q$, where $\bI_q$ is the $q\times q$ identity matrix, then $d(\bmu,\bX)$ follows a $\chi_q$ distribution.

\item 
If $B$ is a Brownian motion in $L^2[0,1]$, then %with the notation of Example~\ref{ex:dp_brownian},
$F_{\omega}(t)=\Prob{\sum_{i=1}^\infty X_i^2 \le t^2}$.

\item If $\bX\sim \mathrm{vMF}(\bmu,\kappa)$, for the chordal distance,
$$
F_{\bomega}(t)=
\begin{cases}
1 - F_T\big(1-\frac{t^2}{2}\big), & 0\le t\le 2,\\
1, & t>2.
\end{cases}
$$
Under the geodesic distance,
$$
F_{\bomega}(t)=
\begin{cases}
1 - F_T\left(\cos(t)\right), & 0\le t\le \pi,\\
1, & t>\pi.
\end{cases}
$$
The density of $T$ is given in \eqref{eq:density_t_vmf}.

\item If $\bX \sim \mathrm{HvMF}(\bmu,\kappa)$, then $F_{\bomega}(t)=\Prob{Z_r\le t}.$ The density of $Z_r$ is given in \eqref{eq:density_zr_hvmf}. 

\begin{comment}
\item \textcolor{blue}{If $\bS\sim \mathrm{Wishart}_q(d,\bSigma)$, then with the notation of Example~\ref{ex:dp_wishart},}
\textcolor{blue}{$F_{\bW}(t)= \mathsf{P}\Big\{\sum_{i=1}^q \log(\lambda_i)^2 \le t ^2\Big\}$.}
\item Example~\ref{ex:dp_dirichlet}.
If $\bX\sim \mathrm{Dir}(\boldsymbol{\alpha})$, under the Euclidean distance, the corresponding CDF can be expressed as the push-forward of the Dirichlet density:
$$
F_{\bomega}^{\mu}(t)
= \int_{\{\bx \in \Delta^{q} : \|\bx - \bomega\|_2 \le t\}}
    \frac{1}{B(\boldsymbol{\alpha})} \prod_{i=1}^{q+1} x_i^{\alpha_i - 1} \, \rd x.
$$
For the Euclidean distance, if $q=1$,
$$
F_{\bomega}^{\mu}(t)
= \Probbig{ |X_1 - \omega_1| \le t / \sqrt{2} }
= F_{X_1}\left(\omega_1 + \frac{t}{\sqrt{2}}\right)
- F_{X_1}\left(\omega_1 - \frac{t}{\sqrt{2}}\right),
$$
where $F_{X_1}$ denotes the CDF of $X_1$. For the Aitchison distance, if $q=1$,
$$
F_{\bomega}^{\mu}(t)
= \Prob{
\left| \log\left(\frac{X_1/(1-X_1)}{\omega_1/(1-\omega_1)}\right) \right| \le \sqrt{2} t
}
= F_{X_1}\big(\mathrm{logit}^{-1}\big(\mathrm{logit}(\omega_1) + \sqrt{2}t\big)\big) - F_{X_1}\big(\mathrm{logit}^{-1}\big(\mathrm{logit}(\omega_1) - \sqrt{2}t\big)\big),
$$
\end{comment}

\item If $\bX$ follows a logistic Gaussian distribution on $\Delta^q$ with positive definite covariance matrix $\boldsymbol{\Sigma}$, then 
$F_{\bomega}(t)=\Prob{\boldsymbol{\alpha}^\top \bQ \le t^2}.$
\end{itemize}
\end{proposition}
The proof of this result is given in Appendix~\ref{ap:proof_prop_dp_examples}.
See Section~\ref{sec:numerical_experiments} for numerical illustrations of the distance profiles covered by Proposition~\ref{prop:dp_examples}.

%-----------------------%
\section{Null asymptotics}\label{sec:null_asymptotics}
%-----------------------%
Consider a random object $X$ taking values in a separable metric space $(\Omega, d)$ and denote by $\mu$ the underlying probability measure of $X$. 
In this section, we study the asymptotic distribution of $\sqrt{n}\left(F_{n,\omega}^{\mu}(t)-F_{\omega}^{\mu_0}(t)\right)$ under the null hypothesis $H_0$, for simple and composite null hypotheses, and use it to prove the consistency of the Kolmogorov--Smirnov and Cramér--von Mises type tests based on $\sqrt{n}\left(F_{n,\omega}^{\mu}(t)-F_{\omega}^{\mu_0}(t)\right)$.



%------------------------%
\subsection{Simple null hypothesis}\label{sec:null_asymptotics_simple}
%------------------------%
The goal is to test $\mu = \mu_{0}$, where $\mu_{0}$ is a given probability measure on $(\Omega, d)$.
Without loss of generality, we assume condition \ref{a2}, so that distance profiles characterize distributions on $(\Omega,d)$. 
Hence, testing $\mu = \mu_{0}$ is equivalent to:
$$
H_0: F_\omega^\mu=F_{\omega}^{\mu_0}, \text{ for every } \omega \in \supp{\mu_{0}}\subset\Omega.
$$
The alternative hypothesis is thus
$$
H_1: \text{ there exists } \omega \in \supp{\mu_{0}} \text{ such that } F_\omega^\mu\neq F_{\omega}^{\mu_0}.
$$
Alternatively, one could assume condition \ref{a1} and replace $\supp{\mu_0}$ by $\Omega$ in the above hypotheses.
Unless otherwise specified, all the results hold under either condition.

Consider a sample $X_1, \ldots, X_n$ of iid random objects taking values in $(\Omega, d)$, and following a common distribution $\mu$. 
Denote by $F_{n,\omega}^{\mu}$ the empirical distance profile associated to this sample. 
Define the random field 
\begin{align*}
    \G_{n,\mu}^{\mu_0}(\omega, t) \defin  \sqrt{n}\left(F_{n,\omega}^{\mu}(t)-F_{\omega}^{\mu_0}(t)\right), \quad (\omega, t) \in \left(\Omega \times \R_{\geq 0}\right).
\end{align*}
Observe that $\G_{n,\mu}^{\mu_0}(\omega, t)$ is the empirical process indexed by the class of functions $\cY$ given in \eqref{eq:def_Y}.  In other words, we can write
$\G_{n,\mu}^{\mu_0}(\omega, t) = \sqrt n(P_n y_{\omega, t} - \mu_0 y_{\omega, t})$
where $P_n$ is the empirical measure associated to the sample $X_1, \ldots, X_n$, i.e., $P_n = \frac{1}{n} \sum_{i=1}^n \delta_{X_i}$.
Thus, for notational convenience, we will also write $\G_{n,\mu}^{\mu_0}(y_{\omega,t})$ to denote $\G_{n,\mu}^{\mu_0}(\omega,t)$ in the proofs.
% This is more convenient when we want to explicitly mention the underlying class of functions when more precision is needed.
Theorem 5.1 in \cite{DubeyChenMuller2024} and Theorem 3 in \cite{ChenDubey2025} prove that, under $H_0$, the class $\cY$ is $\mu_0$-Donsker, i.e.,
\begin{align}\label{eq:thm_5.1_chendubeymuller}
\G_{n,\mu}^{\mu_0} \convl \G_{0} \text { in } \ell^{\infty}(\cY), 
\end{align}
as $n\to \infty$, where $\G_{0}$ is a zero mean Gaussian process with covariance given by
\begin{align}\label{eq:cov_simp}
\mathcal{C}_{(\omega_1,t_1), (\omega_2, t_2)} \defin \Cov{y_{\omega_1, t_1}}{y_{\omega_2, t_2}} = \Prob{d(\omega_1, X) \le t_1, d(\omega_2, X) \le t_2} - F_{\omega_1}^{\mu_0}(t_1)F_{\omega_2}^{\mu_0}(t_2),
\end{align}
for every $(\omega_1, t_1), (\omega_2, t_2) \in \Omega \times \R_{\geq 0}$.




%------------------------%
\subsection{Composite null hypothesis}
%------------------------%
Assume that $\mu$ belongs to a parametric family of probability measures on $(\Omega, d)$, indexed by a parameter $\btheta$ in an open subset $\Theta$ of $\R^p$. 
One may be interested in testing $\mu \in \{\mu_{\btheta} : \btheta \in \Theta\}$. Denote by $\btheta_0\in \Theta$ the true parameter value, i.e., $\mu = \mu_{\btheta_0}$.
In this case, the (composite) null hypothesis can be expressed as
$$
H_0: F_\omega^\mu=F_{\omega}^{\mu_{\btheta_0}}, \text{ for every } \omega \in \mathrm{supp}(\mu_{\btheta_0})\subset\Omega, \ \btheta_0\in\Theta\subset\R^p \text{ unknown}.
$$

The construction of test statistics for the composite null hypothesis requires estimating the unknown parameter $\btheta_0$. 
\citet[Chapter 5, Section 5]{Shorack2009} provides an overview of goodness-of-fit tests with estimated parameters. 
Let $\hat{\btheta}$ be a measurable estimator of $\btheta_0$ based on the sample $X_1, \ldots, X_n$. 
Let $\mu_{\hat{\btheta}}$ be the probability measure corresponding to $\hat{\btheta}$, and define $F_{\omega}^{\mu_{\hat{\btheta}}}$ as the distance profile of $\mu_{\hat{\btheta}}$ at $\omega$.
We can then define the random field
\begin{align*}
    \G_{n,\mu}^{\mu_{\hat{\btheta}}}(\omega, t) \defin \sqrt{n}\left(F_{n,\omega}^{\mu}(t)-F_{\omega}^{\mu_{\hat{\btheta}}}(t)\right), \quad (\omega, t) \in \left(\Omega \times \R_{\geq 0}\right).
\end{align*}

We study the asymptotic behavior of $\G_{n,\mu}^{\mu_{\hat{\btheta}}}$. 
The quantity $(\hat{\btheta} - \btheta_0)$ will play a central role in deriving this asymptotic distribution. 
For estimators based on $n$ replications of an experiment, the order at which $(\hat{\btheta} - \btheta_0)$ converges to zero is often $\sqrt{n}$ \citep[][p. 51]{vandervaart1998}, and its asymptotic distribution is usually normal. 
Under certain regularity conditions, every consistent estimator $\hat{\btheta}$ of $\btheta_0$ satisfies the so-called Bahadur's representation:
\begin{align}\label{eq:classical_bahadur}
\sqrt{n}\big(\hat{\btheta}-\btheta_0\big)= -\Big(\Ebig{\dot{\bpsi}_{\btheta_0}(X)}\Big)^{-1} \frac{1}{\sqrt{n}}\sum_{i=1}^n \bpsi_{\btheta_0}(X_i)+o_{P}(1).
\end{align}
In particular, the sequence $\sqrt{n}(\hat{\btheta}-\btheta_0)$ is asymptotically normal with mean zero and covariance matrix $\Big(\Ebig{\dot{\bpsi}_{\btheta_0}(X)}\Big)^{-1} \Ebig{ \bpsi_{\btheta_0}(X) \bpsi_{\btheta_0}^T(X)} \bigg(\Big(\Ebig{\dot{\bpsi}_{\btheta_0}(X)}\Big)^{-1}\bigg)^\top$. 
See Theorem 5.41 in \cite{vandervaart1998}. See also, e.g., Corollary 10.16 in \citet{Henze2024} for this representation in the context of maximum likelihood.

We now list the regularity conditions needed for \eqref{eq:classical_bahadur} to hold. The so-called ``classical conditions'' are the following:
\begin{enumerate}[label=(\textit{c.\arabic*}), ref=(\textit{c.\arabic*})]
    \item For each $\btheta$ in an open subset of Euclidean space, the mapping $\btheta \mapsto \bpsi_{\btheta}(x)$ is twice continuously differentiable for every $x$.\label{c1}
    \item $\Ebig{ \bpsi_{\btheta_0}(X) }=0$ \label{c2}
    \item $\Ebig{\big\|\bpsi_{\btheta_0}(X)\big\|^2}<\infty$. \label{c3}
    \item The matrix $\Ebig{\dot{\bpsi}_{\btheta_0}(X)}$ exists and is non-singular. \label{c4}
    \item The second-order partial derivatives with respect to $\btheta$ exist for every $x$ and are dominated by a fixed integrable function $M(x)$ for every $\btheta$ in a neighborhood of $\btheta_0$; that is,
    $$
    \left\| \frac{\partial^2 \bpsi_{\btheta}(x)}{\btheta^2}\right\| \le M(x)
    $$
    for some integrable function $M(x)$.\label{c5}
    \item $\hat{\btheta} \convp \btheta_0$. \label{c6}
    \item For every $n$, $\frac{1}{n} \sum_{i=1}^n \bpsi_{\hat{\btheta}}(X_i) = 0$. \label{c7}
\end{enumerate}

The classical conditions are too stringent and require the existence of third derivatives \citep{vandervaart1998}. This is why it is sometimes convenient to replace these assumptions using a Lipschitz condition. The alternative regularity conditions are:
\begin{enumerate}[label=($\tilde{c}$\textit{.\arabic*}), ref=($\tilde{c}$\textit{.\arabic*})]
    \item For each $\btheta$ in an open subset of Euclidean space, the mapping $x \mapsto \bpsi_{\btheta}(x)$ is a vector-valued function such that, for every $\btheta_1$ and $\btheta_2$ in a neighborhood of $\btheta_0$ and a measurable function with $\E{L(X)^2} < \infty$,
    $$
    \| \bpsi_{\btheta_1}(x) - \bpsi_{\btheta_2}(x)\| \le L(x) \|\btheta_1 - \btheta_2 \|.
    $$ \label{cp1}
    \item $\Ebig{\left\|\bpsi_{\btheta_0}(X)\right\|^2}<\infty$. \label{cp2}
    \item The map $\btheta \mapsto \Ebig{\bpsi_{\btheta}(X)}$ is differentiable at a zero $\btheta_0$, with non-singular derivative matrix $\bV_{\btheta_0}$.\label{cp3}
    \item $\hat{\btheta} \convp \btheta_0$. \label{cp4}
    \item $\frac{1}{n} \sum_{i=1}^n \bpsi_{\hat{\btheta}}(X_i) = o_P(n^{-1/2})$. \label{cp5}
\end{enumerate}
Under conditions \ref{cp1}--\ref{cp5}, Bahadur's representation becomes
\begin{align}\label{eq:bahadur}
\sqrt{n}\left(\hat{\btheta}-\btheta_0\right)= - \bV_{\btheta_0}^{-1} \frac{1}{\sqrt{n}}\sum_{i=1}^n \bpsi_{\btheta_0}(X_i)+o_{P}(1).
\end{align}
In particular, the sequence $\sqrt{n}(\hat{\btheta}-\btheta_0)$ is asymptotically normal with mean zero and covariance matrix $\bV_{\btheta_0}^{-1} \Ebig{\bpsi_{\btheta_0}(X) \bpsi_{\btheta_0}(X)^T}\left(\bV_{\btheta_0}^{-1}\right)^T$. See Theorem 5.21 in \cite{vandervaart1998} for this result.

\textcolor{blue}{Diego: As of now, I include both sets of conditions for completeness. But I think we should keep only \ref{cp1}--\ref{cp5} in the future.}

Proposition~\ref{prop:bahadur_vmf} provides Bahadur's representation for the von Mises--Fisher distribution on the sphere.

%---------------------------%
\begin{proposition}[a Bahadur representation, $\mathrm{vMF}(\bmu_0, \kappa_0)$]\label{prop:bahadur_vmf}
%---------------------------%
Let $\bX_1,\ldots,\bX_n$ be iid with common distribution $\mathrm{vMF}(\bmu_0,\kappa_0)$ on $\mathbb{S}^q$, where $\kappa_0>0$, and let $\bxi_0=\kappa_0\bmu_0$ be the canonical parameter.
Let $\hat{\bxi}$ be the Maximum Likelihood Estimator (MLE) of $\bxi_0$. Then
\begin{align}\label{eq:bahadur_vMF}
\sqrt{n}\big(\hat{\bxi}-\bxi_0\big)
=
-\big(\dot{\psi}_{\bxi_0}\big)^{-1}
\frac{1}{\sqrt{n}}\sum_{i=1}^n \psi_{\bxi_0}\left(\bX_i\right)
+o_{\mathsf{P}}(1),
\end{align}
where
\begin{equation}\label{eq:vMF_score}
\psi_{\bxi}(\bx)
\defin
\bx-A_q(\|\bxi\|)\frac{\bxi}{\|\bxi\|}
\end{equation}
with $A_q(r)\defin\mathcal{I}_{\frac{q+1}{2}}(r)/\mathcal{I}_{\frac{q-1}{2}}(r)$, $r>0$,
and, writing $\bxi=\kappa\bmu$,
\begin{align}\label{eq:vMF_der_psi}
\dot{\psi}_{\bxi}(\bx)
=
-\left(
A_q^\prime(\kappa)\bmu\bmu^\top
+
\frac{A_q(\kappa)}{\kappa}
\left(\bI_{q+1}-\bmu\bmu^\top\right)
\right).
\end{align}
\end{proposition}
The proof of Proposition~\ref{prop:bahadur_vmf} is given in Section~\ref{ap:sec_bahadur_vmf}.
See Section~\ref{sec:numerical_experiments} for a numerical validation of the Bahadur representation \eqref{eq:bahadur_vMF}.


We will use the notation $\dot{F}_{\omega}^{\btheta_0}(t) \defin \left.\frac{\partial}{\partial \btheta} F_{\omega}^{\mu_{\btheta}}(t)\right|_{\btheta=\btheta_0}$ for the partial derivative of $F_{\omega}^{\mu_{\btheta}}(t)$ with respect to~$\btheta$ at $\btheta_0$. 
The following differentiability condition is required in the derivation of the asymptotic distribution of $\G_{n,\mu}^{\mu_{\hat{\btheta}}}$:
\begin{enumerate}[label=(\textit{d}), ref=(\textit{d})]
    \item The map $\boldsymbol{\theta}\mapsto F^{\mu_{\boldsymbol{\theta}}}$ is differentiable at $\boldsymbol{\theta}_0$ as a map from $\Theta$ into $\ell^\infty(\Omega\times\mathbb{R}_{\ge0})$; that is,
    $$
    \sup_{\omega\in\Omega, t\ge0}
    \left|
    F_{\omega}^{\mu_{\boldsymbol{\theta}}}(t)
    -
    F_{\omega}^{\mu_{\boldsymbol{\theta}_0}}(t)
    -
    \dot F_{\omega}^{\boldsymbol{\theta}_0}(t)^\top
    (\boldsymbol{\theta}-\boldsymbol{\theta}_0)
    \right|
    =
    o(\|\boldsymbol{\theta}-\boldsymbol{\theta}_0\|)
    \quad
    \text{as }
    \boldsymbol{\theta}\to\boldsymbol{\theta}_0.
    $$
    In particular, for every $(\omega,t)\in\Omega\times\mathbb R_{\ge0}$, the map $\boldsymbol{\theta}\mapsto F_{\omega}^{\mu_{\boldsymbol{\theta}}}(t)$ is differentiable at $\boldsymbol{\theta}_0$, and $\sup_{\omega\in\Omega, t\ge0}\|\dot F_{\omega}^{\boldsymbol{\theta}_0}(t)\|<\infty$. \label{d}
\end{enumerate}

%---------------------------%
\begin{proposition}\label{prop:verification_assumptions_specific_models}
%---------------------------%
Assumption~\ref{d} holds for the following parametric models: the normal distribution $\mathcal{N}(\mu,\sigma^2)$ on $\mathbb{R}$, when either $\mu$, $\sigma$, or both parameters are unknown; and the $\mathrm{vMF}$ distribution on $\mathbb{S}^q$ with canonical parameter $\bxi=\kappa\bmu\in\R^{q+1}\setminus\{\boldsymbol{0}\}$.
\end{proposition}
The proof of this result is given in Section~\ref{ap:sec_verification_assumptions_specific_models}.


The following result characterizes the asymptotic behavior of $\G_{n,\mu}^{\mu_{\hat{\btheta}}}$.
The result follows by proving that the class of functions 
\begin{align}\label{eq:def_F}
\cF \defin \{f_{\omega,t} : \omega \in \Omega, t \in \R_{\geq 0}\}, \text{ where } f_{\omega,t}(x) \defin y_{\omega, t}(x) - F_{\omega}^{\mu_{\btheta_0}}(t) + \dot{F}_{\omega}^{\btheta_0}(t)^\top \bV_{\btheta_0}^{-1} \bpsi_{\btheta_0}(x),
\end{align}
is Donsker.

\textcolor{blue}{Diego: Something key here, and that has motivated how I have written the Donsker results, is that there is an extra term $o_P(1)$ that makes $\G_{n,\mu}^{\mu_{\hat{\btheta}}}$ not be exactly the empirical process indexed by the class $\cF$, but only asymptotically. This is why the Donsker result implies the weak convergence of $\G_{n,\mu}^{\mu_{\hat{\btheta}}}$, but it is not an equivalence, as usually.}
%---------------------------%
\begin{theorem}[asymptotic distribution of $\G_{n,\mu}^{\mu_{\hat{\btheta}}}$ under the null hypothesis]\label{thm:convergence_G_n_composite}
%---------------------------%
%Consider a sample $X_1, \ldots, X_n$ of iid random objects taking values in a complete separable metric space $(\Omega, d)$, and following some unknown distribution $\mu$. 
%Suppose that $\mu$ belongs to some parametric family of distributions indexed by $\btheta$, with $\btheta \in \bTheta \subset \R^p$, $p<\infty$, and let $\btheta_0$ denote the true parameter value. 
Consider an estimator $\hat{\btheta}$ satisfying the conditions \ref{cp1}--\ref{cp5}, and assume that \ref{d} holds. 
Then, under $H_0$, the class of functions $\cF$ defined in \eqref{eq:def_F} is Donsker. Therefore,
$$
\G_{n,\mu}^{\mu_{\hat{\btheta}}} \convl \tilde{\G}_0 \text { in } \ell^{\infty}(\mathcal{F}),
$$
where $\tilde{\G}_0$ is a zero mean Gaussian process with covariance given by
\begin{equation}\label{eq:cov_comp}
\begin{aligned}
\tilde{\mathcal{C}}_{(\omega_1,t_1), (\omega_2, t_2)}
 & = \mathcal{C}_{(\omega_1,t_1), (\omega_2, t_2)} + \dot{F}_{\omega_2}^{\btheta_0}(t_2)^\top \big(\bV_{\btheta_0}\big)^{-1} \mathsf{E}\left(y_{\omega_1, t_1}(X) \bpsi_{\btheta_0}(X)\right)  
 \\
& \quad + \dot{F}_{\omega_1}^{\btheta_0}(t_1)^\top \bV_{\btheta_0}^{-1} \mathsf{E}\left(y_{\omega_2, t_2}(X) \bpsi_{\btheta_0}(X)\right)
\\
 & \quad + \dot{F}_{\omega_1}^{\btheta_0}(t_1)^\top  \bV_{\btheta_0}^{-1}\mathsf{E}\left(\bpsi_{\btheta_0}(X) \bpsi_{\btheta_0}(X)^\top\right)  \bV_{\btheta_0}^{-1}\dot{F}_{\omega_2}^{\btheta_0}(t_2) ,
\end{aligned}
\end{equation}
with $\mathcal{C}_{(\omega_1,t_1), (\omega_2, t_2)}$ the covariance process under the simple null, given in \eqref{eq:cov_simp}.
\end{theorem}
See Section \ref{ap:sec_proofs_null_asymptotics_thm_convergence_G_n_composite} in the Appendix for the proof of this result.

%---------------------------%
\begin{proposition}[covariance process for specific parametric models]\label{prop:covariance_specific_models}
%---------------------------%
The following expressions hold when the unknown parameters are estimated by maximum likelihood:
\begin{enumerate}[label=(\textit{\roman*}), ref=(\textit{\roman*})]
    \item If $X \sim \mathcal{N}(\mu_0,\sigma^2)$, with $\sigma>0$ fixed and $\mu_0$ unknown, then \label{prop:covariance_specific_models_i}
    $$
    \tilde{\mathcal{C}}_{(\omega_1,t_1),(\omega_2,t_2)}
    =
    \mathcal{C}_{(\omega_1,t_1),(\omega_2,t_2)}
    -
    \sigma^2 \dot{F}_{\omega_1}^{\mu_0}(t_1)\dot{F}_{\omega_2}^{\mu_0}(t_2),
    $$
    where
	    \begin{equation}\label{eq:prop_cov_normal_mu_dotF}
	    \dot{F}_{\omega}^{\mu_0}(t)
	    =
	    - \frac{1}{\sigma}\phi\left(\frac{\omega - \mu_0 + t}{\sigma}\right)
	    +
	    \frac{1}{\sigma}\phi\left(\frac{\omega - \mu_0 - t}{\sigma}\right).
	    \end{equation}
    \item Take $X \sim \mathcal{N}(\mu,\sigma_0^2)$, with $\mu\in\R$ fixed and $\sigma_0$ unknown, then \label{prop:covariance_specific_models_ii}
    $$
    \tilde{\mathcal{C}}_{(\omega_1,t_1),(\omega_2,t_2)}
    =
    \mathcal{C}_{(\omega_1,t_1),(\omega_2,t_2)}
    -
    \frac{\sigma_0^2}{2}\dot{F}_{\omega_1}^{\sigma_0}(t_1)\dot{F}_{\omega_2}^{\sigma_0}(t_2),
    $$
    where
	    \begin{equation}\label{eq:prop_cov_normal_sigma_dotF}
	    \dot{F}_{\omega}^{\sigma_0}(t)
	    =
	    - \frac{\omega - \mu + t}{\sigma_0^2}\phi\left(\frac{\omega - \mu + t}{\sigma_0}\right)
	    +
	    \frac{\omega - \mu - t}{\sigma_0^2} \phi\left(\frac{\omega - \mu - t}{\sigma_0}\right).
	    \end{equation}
    \item If $X \sim \mathcal{N}(\mu_0,\sigma_0^2)$, with both $\mu_0$ and $\sigma_0$ unknown, then \label{prop:covariance_specific_models_iii}
    $$
	    \tilde{\mathcal{C}}_{(\omega_1,t_1),(\omega_2,t_2)}
	    =
	    \mathcal{C}_{(\omega_1,t_1),(\omega_2,t_2)}
	    -
	    \big(\dot{F}_{\omega_1}^{\mu_0,\sigma_0}(t_1)\big)^\top
	    \mathcal{I}_{\mu_0,\sigma_0}^{-1}
	    \dot{F}_{\omega_2}^{\mu_0,\sigma_0}(t_2),
	    $$
	    where $\mathcal{I}_{\mu_0,\sigma_0}=\mathrm{diag}\left(\frac{1}{\sigma_0^2},\frac{2}{\sigma_0^2}\right)$
	    and $\dot F_\omega^{\mu_0,\sigma_0}(t)\defin\big(\dot F_\omega^{\mu_0}(t),\dot F_\omega^{\sigma_0}(t)\big)^\top$, with $\dot F_\omega^{\mu_0}(t)$ and $\dot F_\omega^{\sigma_0}(t)$ given in \eqref{eq:prop_cov_normal_mu_dotF} and \eqref{eq:prop_cov_normal_sigma_dotF}, respectively, with $\sigma$ replaced by $\sigma_0$ and $\mu$ replaced by $\mu_0$.
	    \item Take $\bX \sim \mathrm{vMF}(\bmu_0,\kappa_0)$ on $\mathbb{S}^q$, and the canonical parameter $\bxi_0=\kappa_0\bmu_0$ is considered, then \label{prop:covariance_specific_models_iv}
	    $$
	    \tilde{\mathcal{C}}_{(\bomega_1,t_1),(\bomega_2,t_2)}
    =
	    \Prob{d(\bomega_1, \bX) \le t_1, d(\bomega_2, \bX) \le t_2}
	    -
	    F_{\bomega_1}^{\bxi_0}(t_1)F_{\bomega_2}^{\bxi_0}(t_2)
	    \big(1+\bm_1^\top\mathsf{Var}(\bX)^{-1}\bm_2\big),
	    $$
    where $\bm_i=\mathsf{E}(\bX\mid d(\bX,\bomega_i)\le t_i)-A_q(\kappa_0)\bmu_0$, $i=1,2$.
\end{enumerate}
\end{proposition}
The proof of Proposition \ref{prop:covariance_specific_models} is given in Section~\ref{ap:sec_covariance_specific_models} of the Appendix.



%-------------------%
\subsection{Test statistics}
%-------------------%
%From here, one could consider any continuous functional as a test statistic. 
First, consider the Kolmogorov--Smirnov test statistic, which we denote by either $T_n^{\mathrm{KS}}$ (simple null hypothesis), or $\tilde{T}_n^{\mathrm{KS}}$ (composite). These are defined as

\begin{align}\label{eq:ks_test_statistic_simple}
    T_n^{\mathrm{KS}} \defin \sup_{\omega \in \Omega, t \ge0}\Big|\G_{n,\mu}^{\mu_0}(\omega, t)\Big|, \quad \quad \tilde{T}_n^{\mathrm{KS}} \defin \sup_{\omega \in \Omega, t\ge0}\Big|\G_{n,\mu}^{\mu_{\hat{\btheta}}}(\omega, t)\Big|.
\end{align}

\begin{comment}
The empirical version of this test statistic is given by:
$$
\frac{1}{n^2} \sum_{i=1}^n \sum_{j=1}^n \left( \hat{F}_{X_i}(d(X_i, X_j)) - F_{X_i}^{\theta_0}(d(X_i,X_j)) \right)
$$
\end{comment}

Given a sample $X_1, \dots, X_n$, define the empirical probability distribution $P_n$ as the discrete measure that assigns mass $\frac{1}{n}$ to each of the points $X_1,\dots,X_n$. 
Consider the Cramér--von Mises test statistic, which we denote by either $T_n^{\mathrm{CM}}$ (simple null hypothesis), and $\tilde{T}_n^{\mathrm{CM}}$ (composite). These are defined, respectively, as
\begin{align}\label{eq:test_statistics_cvm}
T_n^{\mathrm{CM}} \defin \int_{\Omega\times\R_{\ge0}} (\G_{n,\mu}^{\mu_0}(\omega, t))^2 \, \rd F_{n,\omega}^{\mu}(t) \, \rd P_n(\omega), \quad \tilde{T}_n^{\mathrm{CM}} \defin \int_{\Omega\times\R_{\ge0}} (\G_{n,\mu}^{\mu_{\hat{\btheta}}}(\omega, t))^2 \, \rd F_{n,\omega}^{\mu}(t) \, \rd P_n(\omega).
\end{align}
Observe that this is an empirical version of the Cramér--von Mises test statistic, since the integration is with respect to the empirical measure $\rd F_{n,\omega}^{\mu}(t) \, \rd P_n(\omega)$ instead of the population version $\rd F_{\omega}^{\mu_0}(t) \, \rd \mu_0(\omega)$. 
The use of the empirical measure is motivated by the fact that $\rd F_{\omega}^{\mu_0}(t) \, \rd \mu_0(\omega)$ depends on the distribution under the null $\mu_0$ (and, under a composite null, on unknown parameters). 
Replacing it by its empirical counterpart yields a statistic whose integration weight is model-free and fully data-driven.

Theorem \ref{thm:test_statistics_simp} proves the asymptotic distributions under the simple null hypothesis of the Kolmogorov--Smirnov and Cramér--von Mises test statistics given in \eqref{eq:ks_test_statistic_simple} and \eqref{eq:test_statistics_cvm}, respectively.
%---------------------------%
\begin{theorem}[asymptotic distributions of $T_n^{\mathrm{KS}}$ and $T_n^{\mathrm{CM}}$ under $H_0$]\label{thm:test_statistics_simp}
%---------------------------%
The following statements are true under the simple null hypothesis:
\begin{enumerate}[label=(\textit{\roman*}), ref=(\textit{\roman*})]
\item For the Kolmogorov--Smirnov test statistic, as $n$ diverges to infinity, it holds that
\label{thm:test_statistics_simp_i}
\begin{align*}
T_n^{\mathrm{KS}} \convl \sup_{\omega \in \Omega, t\ge0}\big|\G_{0}(\omega, t)\big|,
\end{align*}

\item For the Cramér--von Mises test statistic, one has
\label{thm:test_statistics_simp_ii}
$$
T_n^{\mathrm{CM}}\convl \int_{\Omega\times\R_{\ge0}} (\G_{0}(\omega, t))^2 \rd F_{\omega}^{\mu_0}(t) \, \rd \mu_0(\omega), 
$$
as $n\to \infty$.
\end{enumerate}
\end{theorem}
See Section \ref{ap:sec_proofs_null_asymptotics_thm_test_statistics_simp} in the Appendix for the proof of Theorem \ref{thm:test_statistics_simp}.
See Section~\ref{sec:numerical_experiments} for numerical illustrations of the convergence of $T_n^{\mathrm{KS}}$.
%---------------------------%
\begin{remark}\label{rem:separability_G_0}
%---------------------------%
    We know that the class $\cY$ is Donsker, thus $\G_{0}(\omega, t)$ is tight in $\ell^\infty(\cY)$. Tightness is equivalent to there being a $\sigma$-compact set that has probability one (see, e.g., page 15 in \citealp{VanderVaartWellner2023}, and page 105 in \citealp{Kosorok2008}). Every $\sigma$-compact set is Lindelöf, and a metric space is Lindelöf if and only if it is separable \citep[see, e.g., Theorem 16.11 of][]{Willard2004}. Since there is a measurable and separable set with probability one, then $\G_{0}(\omega, t)$ is separable.  
\end{remark}
%---------------------------%
\begin{remark}
%---------------------------%
    The random variable $\sup_{\omega, t} \G_0(\omega, t)$ may fail to be measurable, since $\Omega \times \R_{\ge 0}$ is uncountable \citep[see, e.g., page 15 in][]{Adler2007}. Separability solves the measurability issues, since it guarantees the existence of a countable subset $D \subset \Omega \times \R_{\ge 0}$ such that $\sup_{(\omega,t) \in \Omega \times \R_{\ge 0}} \G_0(\omega, t) = \sup_{(\omega,t) \in D} \G_0(\omega, t)$ almost surely, and the supremum of a countable set of measurable random variables is measurable.
    Thus, it is sensible to speak of weak convergence of $\sup_{\omega, t} \G_{n,\mu}^{\mu_0}(\omega, t)$ to $\sup_{\omega, t} \G_0(\omega, t)$, since the limiting process is required to be separable. Observe that this also implies that $\sup_{\omega, t} \G_{n,\mu}^{\mu_0}(\omega, t)$ is asymptotically measurable by Lemma 1.3.8 in \cite{VanderVaartWellner2023}.
% Something interesting here is Exercise 6 in page 64 of \cite{VanderVaartWellner2023}. But I am not sure about the process being "indexed by a directed set that is trivial defined on complete probability spaces".
\end{remark}
\begin{comment}
%---------------------------%
\begin{remark}
%---------------------------%
    Theorem 2.6 in \cite{Potthoff2009} ensures that any real valued random field indexed by a separable metric space has a separable modification.  By assumption, $(\Omega, d)$ is separable and $(\R_{\ge0}, |\cdot|)$ is also separable, then the product metric space $\Omega \times \R_{\ge0}$ is separable when endowed with a metric that induces the product topology, such as $d_{\Omega \times \R_{\ge0}}((\omega_1, t_1), (\omega_2, t_2)) = d(\omega_1, \omega_2) + |t_1 - t_2|$. Hence, both $\G_n$ and $\G_{0}$ have separable modifications.
\end{remark}
\end{comment}

Theorem \ref{thm:test_statistics_comp} is the analog of Theorem \ref{thm:test_statistics_simp} for the composite null hypothesis.
%---------------------------%
\begin{theorem}[asymptotic distributions of $\tilde{T}_n^{\mathrm{KS}}$ and $\tilde{T}_n^{\mathrm{CM}}$ under $H_0$]\label{thm:test_statistics_comp}
%---------------------------%
Assume the conditions~\ref{cp1}--\ref{cp5} and~\ref{d}. The following statements are true under the composite null hypothesis:
\begin{enumerate}[label=(\textit{\roman*}), ref=(\textit{\roman*})]
\item For the Kolmogorov--Smirnov test statistic, as $n$ diverges to infinity, it holds that \label{thm:test_statistics_comp_i}
\begin{align*}
\tilde{T}_n^{\mathrm{KS}} \convl \sup_{\omega \in \Omega, t\ge0}\big|\tilde{\G}_{0}(\omega, t)\big|.
\end{align*}

\item If, in addition, the map $(\omega,t)\mapsto \dot F_{\omega}^{\btheta_0}(t)$ is uniformly continuous with respect to
$d_{\Omega\times\R_{\ge0}}\big((\omega,t),(\omega^\prime,t^\prime)\big)\defin d(\omega,\omega^\prime)+|t-t^\prime|$, then
$$
\tilde{T}_n^{\mathrm{CM}}\convl \int_{\Omega\times\R_{\ge0}} \tilde{\G}_0(\omega, t)^2 \, \rd F_{\omega}^{\mu_{\btheta_0}}(t) \, \rd \mu_{\btheta_0}(\omega)
$$
as $n \to \infty$. \label{thm:test_statistics_comp_ii}
\end{enumerate}
\end{theorem}
%---------------------------%

See Section \ref{ap:sec_proofs_null_asymptotics_thm_test_statistics_comp} in the Appendix for the proof of Theorem \ref{thm:test_statistics_comp}.
See Section~\ref{sec:numerical_experiments} for numerical illustrations of the convergence of $\tilde{T}_n^{\mathrm{KS}}$.

%----------------------------%
\begin{remark}\label{ap:rem_uniform_continuity_Fdot}
%---------------------------%
The uniform continuity assumption on $(\omega,t)\mapsto \dot F_{\omega}^{\btheta_0}(t)$ is automatically
satisfied in many parametric models. In particular, consider maximum likelihood, so that
$\mu_{\btheta}$ admits a density $p_{\btheta}$ with respect to a dominating measure and the score
function is given by $\psi_{\btheta}(x)=\frac{\partial}{\partial \btheta}\log p_{\btheta}(x)$.
Assume moreover that the regularity conditions allowing differentiation under the integral sign are satisfied.
Fix $(\omega,t)\in\Omega\times\R_{\ge0}$. Then
\begin{equation}\label{eq:rem_Fdot_as_expectation}
\begin{aligned}
\dot F_{\omega}^{\btheta_0}(t)&=
\left.\frac{\partial}{\partial\btheta}
\int y_{\omega,t}(x)\,p_{\btheta}(x)\,\rd x\right|_{\btheta=\btheta_0}
\\
&=
\int y_{\omega,t}(x)\,
p_{\btheta_0}(x)\,
\left.\frac{\partial}{\partial\btheta}\log p_{\btheta}(x)\right|_{\btheta=\btheta_0}
\,\rd x 
\\
&=
\mathsf{E}\left[y_{\omega,t}(X)\,\psi_{\btheta_0}(X)\right].
\end{aligned}
\end{equation}
Then, for any pair $(\omega,t),(\omega^\prime,t^\prime)\in\Omega\times\R_{\ge0}$, one has by Cauchy--Schwarz and
Lemma~\ref{ap:lem_relation_L2_prod} that
\begin{align*}
\big\|\dot F_{\omega}^{\btheta_0}(t)-\dot F_{\omega^\prime}^{\btheta_0}(t^\prime)\big\|
&\le
\Big(\mathsf{E}\|\psi_{\btheta_0}(X)\|^2\Big)^{1/2}\,
\Big(
\mathsf{E}\big(y_{\omega,t}(X)-y_{\omega^\prime,t^\prime}(X)\big)^2
\Big)^{1/2}
\\
&\le
\Big(\mathsf{E}\|\psi_{\btheta_0}(X)\|^2\Big)^{1/2}\,
C_1\,\sqrt{d_{\Omega\times\R_{\ge0}}\big((\omega,t),(\omega^\prime,t^\prime)\big)}
\end{align*}
for some constant $C_1>0$. Hence, $(\omega,t)\mapsto \dot F_{\omega}^{\btheta_0}(t)$ is uniformly
continuous with respect to $d_{\Omega\times\R_{\ge0}}$.
\end{remark}
%----------------------------%


Taking the supremum over all values of $\omega \in \Omega$ can be infeasible in practice, especially in high dimensions. 
For this reason, we provide conditions under which the supremum over the whole metric space can be replaced with maximization over the sample elements. 
Proposition \ref{prop:sup_max_ks_simple} deals precisely with this.

                          
%---------------------------%
\begin{proposition}\label{prop:sup_max_ks_simple}
%---------------------------%
    The following results hold:
    \begin{enumerate}[label=(\textit{\roman*}), ref=(\textit{\roman*})]
    \item Assume that $\Omega \subset \supp{\mu_0}$. Under the simple null hypothesis, it holds that \label{prop:sup_max_ks_simple_i}
    \begin{align}\label{eq:sup_max_ks_simple}
    \sup_{\omega\in \Omega, t\ge0} \left|\G_{n,\mu}^{\mu_0}(\omega, t)\right|
    =
    \max_{1\le i,j \le n}\left|\G_{n,\mu}^{\mu_0}(X_i,d(X_i,X_j))\right|
    + o_{\mathsf{P}^*}(1).
    \end{align}
\item Assume also that the conditions of Theorem~\ref{thm:convergence_G_n_composite} hold, $\Omega\subset\supp{\mu_{\btheta_0}}$, and that the map $(\omega,t)\mapsto  \dot{F}_{\omega}^{\btheta_0}(t)$ is uniformly continuous with respect to $d_{\Omega\times\R_{\ge0}}\big((\omega,t),(\omega^\prime,t^\prime)\big)\defin d(\omega,\omega^\prime)+|t-t^\prime|$. Then, under the composite null hypothesis, it holds that \label{prop:sup_max_ks_simple_ii}
\begin{align}\label{eq:sup_max_ks_comp}
    \sup_{\omega\in \Omega, t\ge0} \left|\G_{n,\mu}^{\mu_{\hat{\btheta}}}(\omega, t)\right| = \max_{1\le i,j \le n} \left| \G_{n,\mu}^{\mu_{\hat{\btheta}}}(X_i, d(X_i,X_j)) \right| + o_{\mathsf{P}^*}(1).
\end{align}
\end{enumerate}
\end{proposition}

See Section \ref{ap:sec_proofs_null_asymptotics_prop_sup_max_ks_simple} in the Appendix for the proof of Proposition \ref{prop:sup_max_ks_simple}.


%----------------------%
\section{Non-null asymptotics}\label{sec:non_null_asymptotics}
%----------------------%

%----------------------%
\subsection{Fixed alternatives}\label{sec:fixed_alternatives}
%----------------------%

Suppose that $X_1, \ldots, X_n$ are generated from a distribution $\mu$ that is different from $\mu_0$, the distribution under the null hypothesis.
In other words, we are interested in the behavior of the test statistics under alternatives of the form
$$
H_1: \mu \neq \mu_0 \text{ (simple null)},
$$
$$
H_1: \mu \notin \{\mu_{\btheta}:\btheta\in\bTheta\} \text{ (composite null)}.
$$

%-----------------------%
\subsubsection{Simple null hypothesis}\label{sec:fixed_alternatives_simple}
%-----------------------%

In terms of distance profiles, the alternative $H_1: \mu \neq \mu_0$ is equivalent to the existence of a pair $(\omega, t) \in \Omega \times \R_{\ge0}$ such that $F_{\omega}^{\mu}(t) \neq F_{\omega}^{\mu_0}(t)$.
Theorem \ref{thm:fixed_alt_divergence_ks} shows that the test based on $T_n^{\mathrm{KS}}$ is consistent against fixed alternatives.
%---------------------------%
\begin{theorem}[consistency of $T_n^{\mathrm{KS}}$ under a fixed alternative]\label{thm:fixed_alt_divergence_ks}
%---------------------------%
Under $H_1$, the Kolmogorov--Smirnov statistic $T_n^{\mathrm{KS}}$ diverges in outer probability.
\end{theorem}

A similar result holds for the Cramér--von Mises statistic $T_n^{\mathrm{CM}}$. 
Observe that under the alternative hypothesis, we only know that $F_{\omega}^{\mu}(t) \neq F_{\omega}^{\mu_0}(t)$ at a single point $(\omega,t) \in  \Omega \times \R_{\ge0}$. Theorem \ref{thm:fixed_alternative_divergence_cvm} is proved by extending $(\omega,t)$ to a neighborhood where the separation is uniform, in Lemma~\ref{lem:robust_separation}. Then, Lemma~\ref{lem:positive_mass_region} shows that the probability measure of this neighborhood is positive under the alternative, 
so that the Cramér von-Mises statistic can detect the separation between the distance profiles under the null and the alternative.
Finally, in Theorem \ref{thm:fixed_alternative_divergence_cvm}, a Glivenko--Cantelli argument allows us to replace the empirical ingredients in $T_n^{\mathrm{CM}}$ by their population counterparts, and the law of large numbers for $U$-statistics yields the desired result.

%---------------------------%
\begin{theorem}[divergence of $T_n^{\mathrm{CM}}$ under a fixed alternative]\label{thm:fixed_alternative_divergence_cvm}
%---------------------------%
% Assume the doubling condition \ref{a2}. This is done so that we work with the support of the distributions. The condition in \cite{DubeyChenMuller2024} did not talk about the support.
% Let $\mu_{0}$ be the distribution under $H_0$. Let $\mu$ ($\mu\neq\mu_{0})$ be the true distribution under $H_1$, and assume that \ref{b2} holds for $\mu$ and for $\mu_{0}$.
Assume that either \ref{a2} holds, or \ref{a1} holds and $\supp{\mu} = \Omega$. 
Then, the test statistic $T_n^{\mathrm{CM}}$ diverges in outer probability under $H_1$.
\end{theorem}
%---------------------------%
\begin{remark} 
%---------------------------%
Assumption~\ref{a2} guarantees the existence of a point $(\omega,t)$ at which
$F_{\omega}^{\mu}(t)\neq F_{\omega}^{\mu_0}(t)$, and such that $\omega\in\operatorname{supp}(\mu)$. In contrast, Assumption~\ref{a1} alone only ensures
the existence of such a point with $\omega\in\Omega$, which need not belong to
$\operatorname{supp}(\mu)$. The extra assumption $\supp{\mu} = \Omega$ guarantees $\omega\in\operatorname{supp}(\mu)$. 
This ensures that the differences between $F_\omega^{\mu}(t)$ and $F_\omega^{\mu_0}(t)$ are detected by the Cramér--von Mises statistic.
\end{remark}

See Section \ref{ap:sec_proofs_fixed_alternatives_simple} for the proof of Theorems \ref{thm:fixed_alt_divergence_ks} and \ref{thm:fixed_alternative_divergence_cvm}.

%-----------------------%
\subsubsection{Composite null hypothesis}\label{sec:fixed_alternatives_composite}
%-----------------------%

Focus now on the composite null hypothesis.
In terms of distance profiles, the alternative hypothesis means that for every $\btheta\in\bTheta$ there exists at least one pair
$(\omega,t)\in\Omega\times\R_{\ge 0}$ such that $F^\mu_{\omega}(t)\neq F^{\mu_{\btheta}}_{\omega}(t)$. 
Under $H_1$, the estimated parameter $\hat{\btheta}$ does not converge to a
``true'' parameter value, since the parametric model under which $\hat{\btheta}$ is computed is misspecified.
Hence, we work with a value $\btheta_1$, which is the limit to which $\hat{\btheta}$ converges in probability under $H_1$. 

\textcolor{blue}{Diego: The term ``true'' is written in quotation marks because, under the alternative, the data-generating distribution may lie outside the parametric family under the composite null, so a true parameter value may actually not exist.}

% \textcolor{blue}{\small\emph{Note:} An alternative route could be to assume compactness of $\bTheta$ and continuity of the map $\btheta\mapsto F^{\mu_{\btheta}}$ (as a map into $\ell^\infty(\Omega\times\R_{\ge0})$).}

%----------------------------------%
\begin{theorem}[divergence of $\tilde{T}_n^{\mathrm{KS}}$ under a fixed alternative]
%----------------------------------%
\label{thm:fixed_alt_divergence_ks_composite}
% Let $\mu$ be the true distribution under $H_1$. 
Consider an estimator
$\hat{\btheta}$, and assume that there exists $\btheta_1\in\bTheta$ such that $\hat{\btheta}\convp \btheta_1$.
Assume also that the map $\btheta\mapsto F^{\mu_{\btheta}}$ is continuous at $\btheta_1$ as a map from $\bTheta$ into $\ell^\infty(\Omega\times\R_{\ge0})$; that is,
\begin{equation}\label{eq:theta_star_profile_continuity}
\sup_{\omega\in\Omega, t\ge0}
\bigl|F^{\mu_{\btheta}}_\omega(t)-F^{\mu_{\btheta_1}}_\omega(t)\bigr|
\to 0
\quad\text{as }\btheta\to\btheta_1.
\end{equation}
Then, under $H_1$, the statistic $\tilde{T}_n^{\mathrm{KS}}$ diverges in outer probability, i.e., $\tilde{T}_n^{\mathrm{KS}}\xrightarrow{\mathsf{P}^*}\infty.$
\end{theorem}
Condition \eqref{eq:theta_star_profile_continuity} is a weaker analogue of Assumption~\ref{d} at $\btheta_1$, requiring continuity rather than differentiability.

% \begin{remark}
%     Observe that the differentiability of $\btheta\mapsto F^{\mu_{\btheta}}_{\omega}(t)$ and the boundedness condition on the derivative are analogous to the consequences of Assumption~\ref{d}, but replacing $\btheta_0$ by $\btheta_1$ and requiring the derivative bound to hold in a neighborhood of $\btheta_1$.
% \end{remark}
%---------------------------%
\begin{remark}[sufficient conditions for \eqref{eq:theta_star_profile_continuity}]\label{rem:examples_local_derivative_bound}
%---------------------------%

A sufficient condition for \eqref{eq:theta_star_profile_continuity} is that there exist $\rho>0$ and $L<\infty$ such that
\begin{equation}\label{eq:local_derivative_bound}
\sup_{\btheta\in B(\btheta_1,\rho)} \sup_{\omega\in\Omega, t\ge 0}
\bigl\|\dot{F}^{\btheta}_{\omega}(t)\bigr\|
\le L
\end{equation}
where the differentiability of $\btheta\mapsto F^{\mu_{\btheta}}_\omega(t)$ in $B(\btheta_1,\rho)$, for every $(\omega,t)\in\Omega\times\R_{\ge0}$, is implicitly assumed. Indeed, if \eqref{eq:local_derivative_bound} holds, then
 the mean value theorem gives, for every $\btheta\in B(\btheta_1,\rho)$, that
$
\sup_{\omega, t}
\bigl|F^{\mu_{\btheta}}_\omega(t)-F^{\mu_{\btheta_1}}_\omega(t)\bigr|
\le L\|\btheta-\btheta_1\|,
$
which implies \eqref{eq:theta_star_profile_continuity}. 

Recall the score representation in \eqref{eq:rem_Fdot_as_expectation} (in the setting of Remark \ref{ap:rem_uniform_continuity_Fdot}). It follows that
$
\bigl\|\dot F_{\omega}^{\btheta}(t)\bigr\|
=
\bigl\|\mathsf{E}_{\btheta}\bigl(y_{\omega,t}(X)\,\psi_{\btheta}(X)\bigr)\bigr\|
\le
\mathsf{E}_{\btheta}\bigl\|\psi_{\btheta}(X)\bigr\|.
$
Consequently, \eqref{eq:local_derivative_bound} follows if there exists $\rho>0$ such that
\begin{equation}\label{eq:local_derivative_bound_score}
\sup_{\btheta\in B(\btheta_1,\rho)}
\mathsf{E}_{\btheta}\bigl\|\psi_{\btheta}(X)\bigr\|
<\infty,
\end{equation}
which is independent of $(\omega,t)$. 

Condition \eqref{eq:local_derivative_bound_score} is satisfied, for instance, by the vMF distribution on $\mathbb{S}^q$ with canonical parameter $\boldsymbol{\xi}=\kappa\boldsymbol{\mu}$, since by \eqref{eq:vMF_score},
$
\|\psi_{\boldsymbol{\xi}}(x)\|
\le
\|x\|+A_q(\|\boldsymbol{\xi}\|)
\le 2
$
for all $x\in\mathbb{S}^q$. For $X\sim\mathcal{N}(\mu_0,\sigma^2)$ with $\mu_0$ fixed and $\sigma>0$ unknown, it holds by \eqref{eq:dot_F_normal_sigma} that $\sup_{\omega,t}\bigl|\dot F_{\omega}^{\sigma}(t)\bigr|\le \frac{2}{\sigma}\sup_{u\in\R}|u|\varphi(u)=\frac{2}{\sigma\sqrt{2\pi e}}$. Thus, \eqref{eq:local_derivative_bound} holds for any $\rho\in(0,\sigma_1)$. For $X\sim\mathcal{N}(\mu_0,\sigma^2)$ with $\mu$ unknown and $\sigma_0>0$ fixed, by \eqref{eq:dot_F_normal_mu} and the bound $\phi(u)\le \frac{1}{\sqrt{2\pi}}$ for all $u\in\R$, we have $\sup_{\omega,t}\bigl|\dot F^\mu_\omega(t)\bigr|\le \frac{\sqrt{2}}{\sigma_0\sqrt{\pi}}<\infty$. Hence, \eqref{eq:local_derivative_bound} holds for every $\rho>0$.
The case when both $\mu$ and $\sigma$ are unknown follows by combining the two previous cases.
\end{remark}


Theorem \ref{thm:fixed_alt_divergence_cvm_composite} is the analog of Theorem \ref{thm:fixed_alternative_divergence_cvm} for the composite null hypothesis.

%---------------------------%
\begin{theorem}[divergence of $\tilde{T}_n^{\mathrm{CM}}$ under a fixed alternative, composite null]
%---------------------------%
\label{thm:fixed_alt_divergence_cvm_composite}
Consider an estimator
$\hat{\btheta}$, and assume that there exists $\btheta_1\in\bTheta$ such that $\hat{\btheta}\xrightarrow{\mathsf{P}_\mu}\btheta_1$.
Assume also that the continuity condition \eqref{eq:theta_star_profile_continuity} holds.
Assume that either \ref{a2} holds, or \ref{a1} holds and $\supp{\mu} = \Omega$.
It holds that $\tilde{T}_n^{\mathrm{CM}}\xrightarrow{\mathsf{P}_{\mu}^{\star}}\infty$ under $H_1$.
\end{theorem}

See Section \ref{ap:sec_proofs_fixed_alternatives_composite} for the proof of the results in this section.

%----------------------%
\subsection{Local alternatives}\label{sec:local_alternatives}
%----------------------%

We study the asymptotic behavior of the Kolmogorov--Smirnov and Cramér--von Mises statistics under local alternatives, for simple and composite null hypotheses.
Fix a probability measure $G$ on $(\Omega,\mathcal{B}(\Omega))$, \textcolor{blue}{which we assume that satisfies \ref{b2}, so that Theorem 5.1 in \cite{DubeyChenMuller2024} holds (Diego: maybe not necessary to say it, since we stated that \ref{b2} is assumed globally in the paper).}

%----------------------%
\subsubsection{Simple null hypothesis}\label{sec:local_alternatives_simple}
%----------------------%

Define the sequence of local alternatives converging to $H_0$ at rate $n^{-1/2}$ by
\begin{equation}\label{eq:local_alt_Qn_def}
Q_n \defin (1-n^{-\frac{1}{2}})\,\mu_{0} + n^{-\frac{1}{2}}\,G,
\qquad n\in\mathbb{N}.
\end{equation}
For $(\omega,t)\in\Omega\times \R_{\ge0}$, this implies
\begin{align*}
F^{Q_n}_\omega(t)
\defin
(1-n^{-\frac{1}{2}})F^{\mu_0}_\omega(t)+n^{-\frac{1}{2}}F^G_\omega(t)
=
F^{\mu_0}_\omega(t)+n^{-\frac{1}{2}}h(\omega,t),
\end{align*}
where 
$h(\omega,t)\defin F^G_\omega(t)-F^{\mu_0}_\omega(t).$
Hence,
\begin{equation}\label{eq:local_alt_process_decomp}
\sqrt{n}\Big(F_{n,\omega}^{Q_n}(t)-F^{\mu_0}_\omega(t)\Big)
=
\sqrt{n}\Big(F_{n,\omega}^{Q_n}(t)-F^{Q_n}_\omega(t)\Big) + h(\omega,t)
=
\G_{n,Q_n}^{Q_n}(\omega,t) + h(\omega,t),
\end{equation}
where 
$\G_{n,Q_n}^{Q_n}(\omega,t)
\defin
\sqrt{n}\Big(F_{n,\omega}^{Q_n}(t)-F^{Q_n}_\omega(t)\Big)
=
\sqrt{n}(P_n-Q_n)y_{\omega,t}.$
Observe that $\G_{n,Q_n}^{Q_n}$ is a triangular-array empirical process, since the underlying law $Q_n$ varies with $n$. 

The following result shows that $\G_{n,Q_n}^{Q_n}$ 
admits a coupling with $\G_{n,\mu_0}^{\mu_0}$ such that their difference converges in probability to zero in $\ell^\infty(\cF)$.


%----------------------------%
\begin{proposition}\label{prop:bernoulli}
%----------------------------%
Let $\cF$ be a $\mu_{0}$- and $G$-Donsker class of measurable real-valued functions. Assume that
$\sup_{f\in\cF}\big|(G-\mu_{0})f\big|<\infty$.
Then, under the local alternative \eqref{eq:local_alt_Qn_def}, $\G_{n,Q_n}^{Q_n}$ admits a coupling with $\G_{n,\mu_{0}}^{\mu_{0}}$ such that
\begin{equation}\label{eq:sup_bernoulli_prob}
\sup_{f \in \cF}
\Big|
\G_{n,Q_n}^{Q_n}(f)
-
\G_{n,\mu_{0}}^{\mu_{0}}(f)\Big|
\convp 0. 
\end{equation}
\end{proposition}
%----------------------------%

%----------------------------%
\begin{corollary}[asymptotic distribution of $\G_{n,Q_n}^{Q_n}$ under the local alternative]\label{cor:bernoulli}
%----------------------------%
Under the assumptions of Proposition \ref{prop:bernoulli}, one has 
\begin{equation}\label{eq:local_alt_goal_GnQn}
\G_{n,Q_n}^{Q_n} \convl  \G_{0}
\quad\text{in }\ell^\infty(\cF),
\end{equation}
under the local alternative \eqref{eq:local_alt_Qn_def}.
\end{corollary}
\begin{proof}[Proof of Corollary \ref{cor:bernoulli}]
Using the coupled versions provided by Proposition~\ref{prop:bernoulli}, one has $\sup_{f\in\cF}\big|\G_{n,Q_n}^{Q_n}(f)-\G_{n,\mu_{0}}^{\mu_{0}}(f)\big|\convp0$. Since $\G_{n,\mu_{0}}^{\mu_{0}}\convl \G_{0}$ in $\ell^\infty(\cF)$ (Donsker), Slutsky's theorem (recall that $\G_{0}$ is separable) yields \eqref{eq:local_alt_goal_GnQn}.
\end{proof}

The following theorem gives the weak limit of the Kolmogorov--Smirnov statistic under the local alternative.
%----------------------------%
\begin{theorem}[asymptotic distribution of $T_n^{\mathrm{KS}}$ under the local alternative]\label{thm:local_alt_ks_simple}
%----------------------------%
Under the local alternative \eqref{eq:local_alt_Qn_def}, it holds that
\begin{equation*}
T_n^{\mathrm{KS}}
\convl
\sup_{(\omega,t)\in\Omega\times\R_{\ge0}}
\left|\G_{0}(\omega,t)+h(\omega,t)\right|.
\end{equation*}
\end{theorem}

Theorem \ref{thm:local_alt_CvM_simple} is the analog of Theorem \ref{thm:local_alt_ks_simple} for the Cramér--von Mises statistic.
%----------------------------%
\begin{theorem}[asymptotic distribution of $T_n^{\mathrm{CM}}$ under the local alternative]\label{thm:local_alt_CvM_simple}
%----------------------------%
Under the local alternative \eqref{eq:local_alt_Qn_def}, it holds that
$$
T_n^{\mathrm{CM}}
 \convl 
\int_{\Omega\times\mathbb R_{\ge0}}
\big(\G_0(\omega,t)+h(\omega,t)\big)^2\,
\rd F^{\mu_0}_\omega(t)\,\rd\mu_{0}(\omega).
$$
\end{theorem}
See Section \ref{ap:sec_local_alternatives} in the Appendix for the proofs of Proposition \ref{prop:bernoulli} and Theorems \ref{thm:local_alt_ks_simple} and \ref{thm:local_alt_CvM_simple}.

%----------------------------%

%----------------------------%
\subsubsection{Composite null hypothesis}\label{sec:local_alternatives_composite}
%----------------------------%

For the composite case, define the sequence of local alternatives converging to $H_0$ at rate $n^{-1/2}$ by
\begin{equation}\label{eq:local_alt_Qn_def_composite}
Q_n \defin (1-n^{-\frac{1}{2}})\,\mu_{\btheta_0} + n^{-\frac{1}{2}}\,G,
\qquad n\in\mathbb{N}.
\end{equation}
Accordingly, $h(\omega,t)\defin F^G_\omega(t)-F^{\mu_{\btheta_0}}_\omega(t)$.


Under the local alternative, there are two relevant parameters: the estimator $\hat{\btheta}$ based on $P_n$, and the true population parameter under the local alternative, $\tilde{\btheta}_n$, defined as any solution to $\mathsf{E}_{Q_n}\big(\psi_{\btheta}(X)\big) = 0$. We first give sufficient local conditions under which $\hat{\btheta}$ and $\tilde{\btheta}_n$ both converge ($\hat{\btheta}$ in probability) to $\btheta_0$ under the local alternative \eqref{eq:local_alt_Qn_def_composite}. 


\textcolor{blue}{Diego: The minimal assumptions I am able to come up with are below. It all boils down to finding a good neighborhood of $\btheta_0$ in which the score functions are ``sufficiently well-behaved''. The only complaint I have about this is that we assume that the estimated parameters $\hat{\btheta}$ and $\tilde{\btheta}_n$ belong to that neighborhood, which is more than just saying that there is \textbf{some} neighborhood that contains them. Still, this is less restrictive in general than assuming the consistency of $\hat{\btheta}$ and $\tilde{\btheta}_n$ directly, and all the conditions are local around $\btheta_0$. It would be interesting to check how restrictive our conditions are in examples, but right now I prioritize making progress.}

%----------------------------%
\begin{proposition}[sufficient conditions for the convergence of $\hat{\btheta}$ and $\tilde{\btheta}_n$ to $\btheta_0$]\label{prop:local_conv_roots_Qn}
%----------------------------%
Assume conditions \ref{cp2} and \ref{cp3} and $\mathsf E_G\|\psi_{\btheta_0}(X)\|^2<\infty$.
Consider a sufficiently small bounded neighborhood $U$ of $\btheta_0$. Assume that there exists a measurable function $L$ such that $\mathsf E_{\mu_{\btheta_0}}\big(L(X)^2\big)<\infty,$ $\mathsf E_G\big(L(X)^2\big)<\infty,$ and
\begin{equation}\label{eq:local_cp1_prop_local_conv}
\begin{aligned}
\|\psi_{\btheta_1}(x)-\psi_{\btheta_2}(x)\|
\le
L(x)\|\btheta_1-\btheta_2\|
\quad\text{for all }\btheta_1,\btheta_2\in U\text{ and all }x.
\end{aligned}
\end{equation}
The following results hold under the local alternative \eqref{eq:local_alt_Qn_def_composite}:
\begin{enumerate}[label=(\textit{\roman*}), ref=(\textit{\roman*})]
\item Let $\tilde\btheta_n\in\Theta$ be such that $\mathsf{E}_{Q_n}\big(\psi_{\tilde\btheta_n}(X)\big)=0$ and $\tilde\btheta_n\in U$ for all $n$ sufficiently large. Then, $\tilde\btheta_n\to\btheta_0$ as $n\to\infty$. \label{prop:local_conv_roots_Qn_i}
\item Let $\hat\btheta\in\Theta$ be such that $\mathsf{E}_{P_n}\big(\psi_{\hat\btheta}(X)\big)=0$ and $\mathsf P(\hat\btheta\in U)\to1$. Then, $\hat\btheta\convp\btheta_0$ as $n\to\infty$. \label{prop:local_conv_roots_Qn_ii}
\end{enumerate}
\end{proposition}
Observe that condition \eqref{eq:local_cp1_prop_local_conv} is a local version of \ref{cp1}.
%-----------------------------%

Obtaining the asymptotic distribution of the test statistics under the local alternative requires a Bahadur-type asymptotic linear expansion of $\hat{\btheta}$ and $\tilde{\btheta}_n$ as in \eqref{eq:bahadur}, under the local alternative. Since the parametric model is misspecified under $Q_n$, it is important to provide sufficient conditions under which such expansions hold. This is given in Theorem \ref{thm:local_alt_conv_process_correction}. To make the following result more general, we only require that $\tilde\btheta_n$ and $\hat\btheta$ converge to $\btheta_0$, instead of assuming the conditions of Proposition \ref{prop:local_conv_roots_Qn} directly.


%----------------------%
\begin{theorem}[a Bahadur expansion of $\hat{\btheta}$ and $\tilde{\btheta}_n$ under the local alternative]\label{thm:local_alt_conv_process_correction}
%----------------------%
Assume conditions \ref{cp1}--\ref{cp3}, $\mathsf E_G\big(L(X)^2\big)<\infty$, and $\mathsf E_G\|\psi_{\btheta_0}(X)\|^2<\infty$.
Under the local alternative, the following results hold:
\begin{enumerate}[label=(\textit{\roman*}), ref=(\textit{\roman*})]
\item Let $\tilde\btheta_n\in\Theta$ be such that $\mathsf E_{Q_n}\big(\psi_{\tilde\btheta_n}(X)\big)=0$ 
and $\tilde\btheta_n\to\btheta_0$. Then, \label{thm:local_alt_conv_process_correction_i}
\begin{align}\label{eq:thm_local_alt_conv_process_correction_i}
\tilde{\btheta}_n - \btheta_0 = - \frac{1}{\sqrt{n}} \bV_{\btheta_0}^{-1} \mathsf E_G\big(\psi_{\btheta_0}(X)\big) + o(n^{-\frac{1}{2}}),
\end{align}
\item Let $\hat\btheta\in\Theta$ be such that $\mathsf E_{P_n}\big(\psi_{\hat\btheta}(X)\big)=o_P(n^{-1/2})$ and $\hat\btheta\convp\btheta_0$. Then,\label{thm:local_alt_conv_process_correction_ii} 
\begin{align}\label{eq:thm_local_alt_conv_process_correction_ii}
\hat{\btheta} - \btheta_0 = - \frac{1}{n} \bV_{\btheta_0}^{-1} \sum_{i=1}^{n} \psi_{\tilde{\btheta}_n}(X_i) -  \frac{1}{\sqrt{n}} \bV_{\btheta_0}^{-1} \mathsf E_G\big(\psi_{\btheta_0}(X)\big) + o_P(n^{-\frac{1}{2}}),
\end{align}
\end{enumerate}
\end{theorem}


%----------------------%

The following result gives the asymptotic distribution of $\G_{n,Q_n}^{\mu_{\hat{\btheta}}} \defin \sqrt{n}\bigl(F_{n,\omega}^{Q_n}(t)-F^{\mu_{\hat{\btheta}}}_\omega(t)\bigr)$ under the local alternative \eqref{eq:local_alt_Qn_def_composite} in the composite null case.

%----------------------%
\begin{theorem}[asymptotic distribution of $\G_{n,Q_n}^{\mu_{\hat{\btheta}}}$ under the local alternative]\label{thm:local_alt_Gn_composite}
%----------------------%
Assume conditions \ref{cp1}--\ref{cp3} and \ref{d}, $\mathsf E_G\big(L(X)^2\big)<\infty$, and $\mathsf E_G\|\psi_{\btheta_0}(X)\|^2<\infty$. Let $\tilde{\btheta}_n\in\Theta$ be such that $\mathsf E_{Q_n}\big(\psi_{\tilde{\btheta}_n}(X)\big)=0$ and $\tilde{\btheta}_n\to\btheta_0$. Let also $\hat{\btheta}\in\Theta$ be such that $\mathsf E_{P_n}\big(\psi_{\hat{\btheta}}(X)\big)=o_P(n^{-1/2})$ and $\hat{\btheta}\convp\btheta_0$ under the local alternative \eqref{eq:local_alt_Qn_def_composite}. Then, it holds that
\begin{equation}\label{eq:local_alt_composite_consistency}
\G_{n,Q_n}^{\mu_{\hat{\btheta}}}
\convl
\tilde{\G}_{0}(\omega,t)
+
h(\omega,t)
+
\dot F^{\btheta_0}_\omega(t)^\top \bV_{\btheta_0}^{-1}\mathsf E_G\big(\psi_{\btheta_0}(X)\big)
\quad\text{in }\ell^\infty(\cF)
\end{equation}
under the local alternative.
% where $ \G_{0}$ is the centered Gaussian process from Theorem~\ref{thm:convergence_G_n_composite} (with covariance given in \eqref{eq:cov_comp}).
\end{theorem}
%----------------------%
\begin{remark}
%----------------------%
In Theorem \ref{thm:local_alt_Gn_composite}, we have not assumed directly the asymptotic linear expansions for $\hat{\btheta}$ and $\tilde{\btheta}_n$ given in \eqref{eq:thm_local_alt_conv_process_correction_i} and \eqref{eq:thm_local_alt_conv_process_correction_ii}. This is because the conditions of Theorem \ref{thm:local_alt_conv_process_correction} are also used in the proof of Theorem \ref{thm:local_alt_Gn_composite} beyond the derivation of these expansions themselves.
\end{remark}

The next theorem uses this convergence result to derive the asymptotic behavior of the Kolmogorov--Smirnov and Cramér--von Mises statistics under local alternatives for the composite null hypothesis.

%----------------------%
\begin{theorem}[asymptotic distribution of $\tilde{T}_n^{\mathrm{KS}}$ and $\tilde{T}_n^{\mathrm{CM}}$ under the local alternative]\label{thm:consistency_local_composite}
%----------------------%
Assuming the conditions of Theorem~\ref{thm:local_alt_Gn_composite}, it holds under the local alternative that
$$
\tilde{T}_n^{\mathrm{KS}}
=
\sup_{\omega\in\Omega, t \ge0}
\left|
\G_{n,Q_n}^{\mu_{\hat{\btheta}}}(\omega,t)
\right|
 \convl 
\sup_{\omega\in\Omega, t \ge0}
\left|
\tilde{\G}_{0}(\omega,t)
+
h(\omega,t)
+
\dot F^{\btheta_0}_\omega(t)^\top \bV_{\btheta_0}^{-1}\mathsf{E}_G(\psi_{ \btheta_0}(X))
\right|.
$$
Assume also that $(\omega,t)\mapsto \dot F^{\btheta_0}_\omega(t)$ is uniformly continuous with respect to $d_{\Omega\times\mathbb R_{\ge0}}$.
Then, under the local alternative, it holds that
$$
\tilde{T}_n^{\mathrm{CM}}
\convl
\int_{\Omega\times\mathbb R_{\ge0}}
\Big(
\tilde{\G}_{0}(\omega,t)
+
h(\omega,t)
+
\dot F^{\btheta_0}_\omega(t)^\top \bV_{\btheta_0}^{-1}\mathsf{E}_G(\psi_{\btheta_0}(X))
\Big)^2\,
dF^{\mu_{\btheta_0}}_\omega(t)\,d\mu_{\btheta_0}(\omega).
$$
\end{theorem}

See Section \ref{ap:sec_local_alternatives_composite} for the proof of the results from this section.




%-----------------------------%
\section{Multiplier bootstrap}\label{sec:bootstrap}
%-----------------------------%

%One way to estimate $P(T_n >t |H_0)$ is with bootstrap. 
Calculations with the theoretical distribution of the test statistics, such as critical values or p-values, can be complicated in practice. For this reason, we propose a multiplier bootstrap approach and prove its consistency under some general conditions.

Some basic concepts and notation are required before introducing the multiplier bootstrap. Let $\mathcal H=\{h_{\omega,t}:\omega\in\Omega,\ t\ge0\}$ denote the indexing class, with $\mathcal H=\cY$ (simple null) or $\mathcal H=\cF$ (composite null), with $\cY$ and $\cF$ as defined in \eqref{eq:def_Y} and \eqref{eq:def_F}, respectively.
Consider an iid sequence $X_1,X_2,\ldots$ of random objects in $(\Omega,d)$ with common law $P$, where $P=\mu_0$ (simple null), or $P=\mu_{\btheta_0}$ (composite null). Let $\bxi = (\xi_1, \xi_2, \ldots)$ be a sequence of nonnegative iid random variables, independent of $X_1,X_2,\ldots$, with mean $0<m<\infty$ and variance $0<\tau^2<\infty$. Assume that $\|\xi\|_{2,1} < \infty$ where $\|\xi\|_{2,1} = \int_0^\infty \sqrt{\mathsf{P}(|\xi| > x)} \, dx$. Let $\mathsf{E}_{*}$ and $\mathsf{P}_{*}$ denote conditional expectation and conditional probability with respect to the multipliers $\bxi$, given $X_1,\ldots,X_n$, i.e., we condition on the sample, and the remaining randomness is in the multipliers.

A key result \citep{Kosorok2008} is that, for a sequence of random elements $(\mathbb{Z}_n)_{n\ge1}$ and a random element $\mathbb{Z}$ in $\ell^\infty(\mathcal H)$, endowed with the supremum norm, $\mathbb{Z}_n \convl \mathbb{Z}$ if and only if
\begin{align}\label{eq:BL_convergence}
\sup_{\varphi\in \mathrm{BL}_1(\ell^\infty(\mathcal H))}
\left| \mathsf{E}^*\varphi(\mathbb{Z}_n) - \mathsf{E}\varphi(\mathbb{Z}) \right| \to 0,
\end{align}
where $\mathsf{E}^*$ denotes outer expectation. Here, $\mathrm{BL}_1(\ell^\infty(\mathcal H))$ is the class of bounded Lipschitz functions $\varphi: \ell^\infty(\mathcal H) \mapsto \R$ satisfying $\|\varphi\|_{\infty} \le 1$ and $|\varphi(z)-\varphi(z^\prime)| \le \|z-z^\prime\|_{\mathcal H}$ for all $z,z^\prime\in\ell^\infty(\mathcal H)$. This result is used to define the convergence of the bootstrap distributions.



For a tight random element $\mathbb{Z}$ in $\ell^\infty(\mathcal H)$, the convergence $\mathbb{Z}_n^{*}\stackrel[*]{P}{\to}\mathbb{Z}$ \citep{Kosorok2008} means that, conditionally on $X_1,\ldots,X_n$,
\begin{align}\label{eq:BL_convergence_bootstrap}
\sup_{\varphi\in \mathrm{BL}_1(\ell^\infty(\mathcal H))}
\left| \mathsf{E}_{*} \varphi(\mathbb{Z}_n^{*}) - \E \varphi(\mathbb{Z}) \right| \convp 0,
\end{align}
and for all $\varphi\in \mathrm{BL}_1(\ell^\infty(\mathcal H))$, $\mathsf{E}_{*} \overline{\varphi(\mathbb{Z}_n^{*})} - \mathsf{E}_{*} \underline{\varphi(\mathbb{Z}_n^{*})} \convp 0$. The quantities $\overline{\varphi(\mathbb{Z}_n^{*})}$ and $\underline{\varphi(\mathbb{Z}_n^{*})}$ are measurable majorants and minorants, respectively, with respect to the data, including the weights $\bxi$. Contrary to \eqref{eq:BL_convergence}, the convergence in \eqref{eq:BL_convergence_bootstrap} is in probability because the conditional expectations depend on the sample. 

Let $P_n^{*}$ denote the (multiplier) bootstrapped empirical measure, given by
$$
P_n^{*} f = \frac{1}{n} \sum_{i=1}^n \frac{\xi_i}{\bar{\xi}_n} f(X_i),
$$
where $\bar{\xi}_n = \frac{1}{n} \sum_{i=1}^n \xi_i$ and $P_n^{*} \equiv 0$ if $\bar{\xi}_n=0$.
Define 
$$
F_{n,\omega}^{*}(t)
\defin
P_n^{*}y_{\omega,t}
=
\frac{1}{n}\sum_{i=1}^n\frac{\xi_i}{\bar{\xi}_n} \allowbreak y_{\omega,t}(X_i).
$$
We set $F_{n,\omega}^{*}(t)=0$ if $\bar{\xi}_n=0$.

Theorem 2.6 in \cite{Kosorok2008} gives the asymptotic distribution of $\frac{m}{\tau}\sqrt{n}(P_n^{*}-P_n)$. This motivates the following definition of the multiplier bootstrap processes $\G_n^{*}$ and $\tilde{\G}_n^{*}$, which are the bootstrap analogues of $\G_{n,\mu}^{\mu_0}$ and $\G_{n,\mu}^{\mu_{\hat{\btheta}}}$, respectively.
Under a simple null hypothesis, set
$$
\G_n^{*}(\omega,t)
\defin
\frac{m}{\tau}\sqrt{n}
\left[
\left\{F_{n,\omega}^{*}(t)-F_{\omega}^{\mu_0}(t)\right\}
-
\left\{F_{n,\omega}^{\mu}(t)-F_{\omega}^{\mu_0}(t)\right\}
\right]
=
\frac{m}{\tau}\sqrt{n}\left(F_{n,\omega}^{*}(t)-F_{n,\omega}^{\mu}(t)\right).
$$
This definition connects with the bootstrapped empirical process, since $\G_n^{*}(\omega,t)=\frac{m}{\tau}\sqrt{n}(P_n^{*}-P_n)y_{\omega,t}.$

For a composite null hypothesis, let $\hat{\btheta}^{*}$ be the (multiplier) bootstrap version of $\hat{\btheta}$, obtained by replacing $P_n$ with $P_n^{*}$ in the estimating procedure. We assume that $\hat{\btheta}^{*}$ is conditionally measurable given the sample.
Define
$$
\tilde{\G}_n^{*}(\omega,t)
\defin
\frac{m}{\tau}\sqrt{n}
\left[
\left\{F_{n,\omega}^{*}(t)-F_{\omega}^{\mu_{\hat{\btheta}^{*}}}(t)\right\}
-
\left\{F_{n,\omega}^{\mu}(t)-F_{\omega}^{\mu_{\hat{\btheta}}}(t)\right\}
\right].
$$
For the composite case, the connection between $\tilde{\G}_n^{*}$ and the bootstrapped empirical process is given in Theorem~\ref{thm:bootstrap_processes}. This is the essential ingredient to obtain the asymptotic distribution of $\tilde{\G}_n^{*}$ (similarly for $\G_n^{*}$).

%---------------------------%
\begin{theorem}[asymptotic distribution of the multiplier bootstrap processes]\label{thm:bootstrap_processes}
%---------------------------%
The following statements are true:
\begin{enumerate}[label=(\textit{\roman*}), ref=(\textit{\roman*})]
\item Under the simple null hypothesis, it holds that \label{thm:bootstrap_processes_i}
$$
\G_n^{*}\stackrel[*]{P}{\to}\G_0
\quad \text{in } \ell^\infty(\cY).
$$

\item Assume the conditions of Theorem~\ref{thm:convergence_G_n_composite} and 
the bootstrap analogue of Bahadur's representation \eqref{eq:bahadur}:
\begin{align}\label{eq:weighted_bahadur_bootstrap}
\frac{m}{\tau}\sqrt{n}\left(\hat{\btheta}^{*}-\hat{\btheta}\right)
=
-\bV_{\btheta_0}^{-1}
\frac{m}{\tau}\frac{1}{\sqrt n}\sum_{i=1}^n
\left(\frac{\xi_i}{\bar\xi_n}-1\right)\bpsi_{\btheta_0}(X_i)
+o_{P_{*}}(1).
\end{align}
Then, it holds under the composite null hypothesis that \label{thm:bootstrap_processes_ii}
\begin{align}\label{eq:bootstrap_comp_process_remainder_statement}
\sup_{\omega\in\Omega, t\ge0}
\left|
\tilde{\G}_n^{*}(\omega,t)
-
\frac{m}{\tau}\sqrt n(P_n^{*}-P_n)f_{\omega,t}
\right|
=o_{P_{*}}(1).
\end{align}
This yields
$$
\tilde{\G}_n^{*}\stackrel[*]{P}{\to}\tilde{\G}_0
\quad \text{in } \ell^\infty(\cF).
$$
\end{enumerate}
\end{theorem}


The bootstrap versions of the test statistics under the simple null hypothesis are then
$$
\begin{aligned}
T_n^{*\mathrm{KS}}\defin \sup_{\omega \in \Omega, t \ge 0}\Big|\G_n^{*}(\omega, t)\Big|, 
\quad 
T_n^{*\mathrm{CM}} \defin \int_{\Omega\times\R_{\ge0}} \left( \G_n^{*}(\omega, t) \right)^2 \,\rd F_{n,\omega}^{\mu}(t)\,\rd P_n(\omega),
\end{aligned}
$$
and for the composite null hypothesis,
$$
\begin{aligned}
\tilde{T}_n^{*\mathrm{KS}}\defin \sup_{\omega \in \Omega, t \ge 0}\Big|\tilde{\G}_n^{*}(\omega, t)\Big|, 
\quad 
\tilde{T}_n^{*\mathrm{CM}} \defin \int_{\Omega\times\R_{\ge0}} \left( \tilde{\G}_n^{*}(\omega, t) \right)^2 \,\rd F_{n,\omega}^{\mu}(t)\,\rd P_n(\omega).
\end{aligned}
$$

%---------------------------%
\begin{theorem}[asymptotic distribution of the bootstrap test statistics]\label{thm:bootstrap_test_statistics}
%---------------------------%
The following statements are true:
\begin{enumerate}[label=(\textit{\roman*}), ref=(\textit{\roman*})]
\item Under the simple null hypothesis, it holds that \label{thm:bootstrap_test_statistics_i}
$$
T_n^{*\mathrm{KS}}\stackrel[*]{P}{\to}\sup_{\omega\in\Omega, t\ge0}|\G_0(\omega,t)|,
\quad
T_n^{*\mathrm{CM}}\stackrel[*]{P}{\to}
\int_{\Omega\times\R_{\ge0}}\G_0(\omega,t)^2\,\rd F_{\omega}^{\mu_0}(t)\,\rd\mu_0(\omega).
$$

\item Under the composite null hypothesis, assuming the conditions of Theorem~\ref{thm:bootstrap_processes}, one has that \label{thm:bootstrap_test_statistics_ii}
$$
\tilde{T}_n^{*\mathrm{KS}}\stackrel[*]{P}{\to}\sup_{\omega\in\Omega, t\ge0}|\tilde{\G}_0(\omega,t)|.
$$
If, in addition, the map $(\omega,t)\mapsto\dot F_{\omega}^{\btheta_0}(t)$ is uniformly continuous with respect to $d_{\Omega\times\R_{\ge0}}$ and, for every $(\omega,t)\in\Omega\times\R_{\ge0}$, the map $\btheta\mapsto F_\omega^{\mu_{\btheta}}(t)$ is Borel measurable, then
$$
\tilde{T}_n^{*\mathrm{CM}}\stackrel[*]{P}{\to}
\int_{\Omega\times\R_{\ge0}}\tilde{\G}_0(\omega,t)^2\,\rd F_{\omega}^{\mu_{\btheta_0}}(t)\,\rd\mu_{\btheta_0}(\omega).
$$
\end{enumerate}
\end{theorem}
See Section~\ref{ap:sec_proofs_multiplier_bootstrap} in the Appendix for the proof of Theorems~\ref{thm:bootstrap_processes} and \ref{thm:bootstrap_test_statistics}.


%---------------------------%
\section{Numerical experiments}\label{sec:numerical_experiments}
%---------------------------%
We present numerical results to validate some theoretical results derived in the previous sections. Figures~\ref{fig:mvnorm_distance_profiles}--\ref{fig:lg_aitchison_distance_profiles} illustrate the theoretical distance profiles obtained in Proposition~\ref{prop:dp_examples} by comparing them with empirical distance profiles from simulated samples. In each plot, the theoretical distance profile is compared with $10$ empirical distance profiles obtained from $10$ different samples, and pointwise $95$\% confidence bands (based on the normal distribution) for the distance profiles are also included. 

\ifincludeimages
\begin{figure}[htbp]
\centering
\begin{subfigure}[htbp]{0.32\textwidth}
    \centering
    \includegraphics[width=\textwidth]{img/mvnorm_dp_omega_1_n50.png}
    \caption{$\bomega = (3, -2, 1)$, $n = 50$}
    \label{fig:mvnorm_omega1_n50}
\end{subfigure}
\hfill
\begin{subfigure}[htbp]{0.32\textwidth}
    \centering
    \includegraphics[width=\textwidth]{img/mvnorm_dp_omega_2_n50.png}
    \caption{$\bomega = (1, 1, 1)$, $n = 50$}
    \label{fig:mvnorm_omega2_n50}
\end{subfigure}
\hfill
\begin{subfigure}[htbp]{0.32\textwidth}
    \centering
    \includegraphics[width=\textwidth]{img/mvnorm_dp_omega_3_n50.png}
    \caption{$\bomega = (-1, 0.5, 0.5)$, $n = 50$}
    \label{fig:mvnorm_omega3_n50}
\end{subfigure}

\begin{subfigure}[htbp]{0.32\textwidth}
    \centering
    \includegraphics[width=\textwidth]{img/mvnorm_dp_omega_1_n200.png}
    \caption{$\bomega = (3, -2, 1)$, $n = 200$}
    \label{fig:mvnorm_omega1_n200}
\end{subfigure}
\hfill
\begin{subfigure}[htbp]{0.32\textwidth}
    \centering
    \includegraphics[width=\textwidth]{img/mvnorm_dp_omega_2_n200.png}
    \caption{$\bomega = (1, 1, 1)$, $n = 200$}
    \label{fig:mvnorm_omega2_n200}
\end{subfigure}
\hfill
\begin{subfigure}[htbp]{0.32\textwidth}
    \centering
    \includegraphics[width=\textwidth]{img/mvnorm_dp_omega_3_n200.png}
    \caption{$\bomega = (-1, 0.5, 0.5)$, $n = 200$}
    \label{fig:mvnorm_omega3_n200}
\end{subfigure}
\caption{Distance profiles for multivariate normal distribution $\mathcal{N}_3(\bmu, \boldsymbol{\Sigma})$ with $\bmu = (3, -2, 1)$ and $\boldsymbol{\Sigma} = \left(\begin{smallmatrix} 1.18 & 0.82 & -3.07 \\ 0.82 & 6.96 & -10.10 \\ -3.07 & -10.10 & 18.02 \end{smallmatrix}\right)$. Each plot shows the theoretical distance profile, $10$ empirical distance profiles from $10$ different samples, and 95\% confidence bands (gray ribbons) for different values of $\bomega$ and sample sizes $n=50$, $200$.}
\label{fig:mvnorm_distance_profiles}
\end{figure}
\fi

\ifincludeimages
\begin{figure}[hb]
\centering
\begin{subfigure}[htbp]{0.32\textwidth}
    \centering
    \includegraphics[width=\textwidth]{img/vmf_chordal_dp_omega_1_n50.png}
    \caption{$\bomega = (0, 0, 1)$, $n = 50$}
    \label{fig:vmf_chordal_omega1_n50}
\end{subfigure}
\hfill
\begin{subfigure}[hb]{0.32\textwidth}
    \centering
    \includegraphics[width=\textwidth]{img/vmf_chordal_dp_omega_2_n50.png}
    \caption{$\bomega = (1, 0, 0)$, $n = 50$}
    \label{fig:vmf_chordal_omega2_n50}
\end{subfigure}
\hfill
\begin{subfigure}[htbp]{0.32\textwidth}
    \centering
    \includegraphics[width=\textwidth]{img/vmf_chordal_dp_omega_3_n50.png}
    \caption{$\bomega = (0, 0, -1)$, $n = 50$}
    \label{fig:vmf_chordal_omega3_n50}
\end{subfigure}


\begin{subfigure}[htbp]{0.32\textwidth}
    \centering
    \includegraphics[width=\textwidth]{img/vmf_chordal_dp_omega_1_n200.png}
    \caption{$\bomega = (0, 0, 1)$, $n = 200$}
    \label{fig:vmf_chordal_omega1_n200}
\end{subfigure}
\hfill
\begin{subfigure}[htbp]{0.32\textwidth}
    \centering
    \includegraphics[width=\textwidth]{img/vmf_chordal_dp_omega_2_n200.png}
    \caption{$\bomega = (1, 0, 0)$, $n = 200$}
    \label{fig:vmf_chordal_omega2_n200}
\end{subfigure}
\hfill
\begin{subfigure}[htbp]{0.32\textwidth}
    \centering
    \includegraphics[width=\textwidth]{img/vmf_chordal_dp_omega_3_n200.png}
    \caption{$\bomega = (0, 0, -1)$, $n = 200$}
    \label{fig:vmf_chordal_omega3_n200}
\end{subfigure}
\caption{Distance profiles for a $\mathrm{vMF}(\bmu, \kappa)$ distribution on $\mathbb{S}^2$ endowed with the chordal distance, with $\bmu = (0, 0, 1)$ and $\kappa = 2$. }\label{fig:dp_vmf_chordal}
\end{figure}

\begin{figure}[htbp]
\centering
\begin{subfigure}[htbp]{0.32\textwidth}
    \centering
    \includegraphics[width=\textwidth]{img/vmf_geodesic_dp_omega_1_n50.png}
    \caption{$\bomega = (0, 0, 1)$, $n = 50$}
    \label{fig:vmf_geodesic_omega1_n50}
\end{subfigure}
\hfill
\begin{subfigure}[htbp]{0.32\textwidth}
    \centering
    \includegraphics[width=\textwidth]{img/vmf_geodesic_dp_omega_2_n50.png}
    \caption{$\bomega = (1, 0, 0)$, $n = 50$}
    \label{fig:vmf_geodesic_omega2_n50}
\end{subfigure}
\hfill
\begin{subfigure}[htbp]{0.32\textwidth}
    \centering
    \includegraphics[width=\textwidth]{img/vmf_geodesic_dp_omega_3_n50.png}
    \caption{$\bomega = (0, 0, -1)$, $n = 50$}
    \label{fig:vmf_geodesic_omega3_n50}
\end{subfigure}


\begin{subfigure}[htbp]{0.32\textwidth}
    \centering
    \includegraphics[width=\textwidth]{img/vmf_geodesic_dp_omega_1_n200.png}
    \caption{$\bomega = (0, 0, 1)$, $n = 200$}
    \label{fig:vmf_geodesic_omega1_n200}
\end{subfigure}
\hfill
\begin{subfigure}[htbp]{0.32\textwidth}
    \centering
    \includegraphics[width=\textwidth]{img/vmf_geodesic_dp_omega_2_n200.png}
    \caption{$\bomega = (1, 0, 0)$, $n = 200$}
    \label{fig:vmf_geodesic_omega2_n200}
\end{subfigure}
\hfill
\begin{subfigure}[htbp]{0.32\textwidth}
    \centering
    \includegraphics[width=\textwidth]{img/vmf_geodesic_dp_omega_3_n200.png}
    \caption{$\bomega = (0, 0, -1)$, $n = 200$}
    \label{fig:vmf_geodesic_omega3_n200}
\end{subfigure}
\caption{Same description as Figure \ref{fig:dp_vmf_chordal}, but for the geodesic distance on $\mathbb{S}^2$.}
\label{fig:vmf_geodesic_distance_profiles}
\end{figure}
\fi

\ifincludeimages
\begin{figure}[htbp]
\centering
\begin{subfigure}[htbp]{0.32\textwidth}
    \centering
    \includegraphics[width=\textwidth]{img/hvmf_geodesic_dp_omega_1_n50.png}
    \caption{$\bomega = \big(\sqrt{2}, \frac{1}{\sqrt{2}}, \frac{1}{\sqrt{2}}\big)$, $n=50$}
    \label{fig:hvmf_geodesic_omega1_n50}
\end{subfigure}
\hfill
\begin{subfigure}[htbp]{0.32\textwidth}
    \centering
    \includegraphics[width=\textwidth]{img/hvmf_geodesic_dp_omega_2_n50.png}
    \caption{$\bomega = \big(\sqrt{2}, -\frac{1}{\sqrt{2}}, \frac{1}{\sqrt{2}}\big)$, $n=50$}
    \label{fig:hvmf_geodesic_omega2_n50}
\end{subfigure}
\hfill
\begin{subfigure}[htbp]{0.32\textwidth}
    \centering
    \includegraphics[width=\textwidth]{img/hvmf_geodesic_dp_omega_3_n50.png}
    \caption{$\bomega = \big(\sqrt{2}, -\frac{1}{\sqrt{2}}, -\frac{1}{\sqrt{2}}\big)$, $n=50$}
    \label{fig:hvmf_geodesic_omega3_n50}
\end{subfigure}

\begin{subfigure}[htbp]{0.32\textwidth}
    \centering
    \includegraphics[width=\textwidth]{img/hvmf_geodesic_dp_omega_1_n200.png}
    \caption{$\bomega = \big(\sqrt{2}, \frac{1}{\sqrt{2}}, \frac{1}{\sqrt{2}}\big)$, $n=200$}
    \label{fig:hvmf_geodesic_omega1_n200}
\end{subfigure}
\hfill
\begin{subfigure}[htbp]{0.32\textwidth}
    \centering
    \includegraphics[width=\textwidth]{img/hvmf_geodesic_dp_omega_2_n200.png}
    \caption{$\bomega = \big(\sqrt{2}, -\frac{1}{\sqrt{2}}, \frac{1}{\sqrt{2}}\big)$, \\ $n=200$}
    \label{fig:hvmf_geodesic_omega2_n200}
\end{subfigure}
\hfill
\begin{subfigure}[htbp]{0.32\textwidth}
    \centering
    \includegraphics[width=\textwidth]{img/hvmf_geodesic_dp_omega_3_n200.png}
    \caption{$\bomega = \big(\sqrt{2}, -\frac{1}{\sqrt{2}}, -\frac{1}{\sqrt{2}}\big)$, $n=200$}
    \label{fig:hvmf_geodesic_omega3_n200}
\end{subfigure}
\caption{Distance profiles for a $\mathrm{HvMF}(\bmu, \kappa)$ distribution on $\mathbb{H}^2$ endowed with the geodesic distance, with $\bmu = \big(\sqrt{2}, \frac{1}{\sqrt{2}}, \frac{1}{\sqrt{2}}\big)$ and $\kappa = 50$. }
\label{fig:hvmf_geodesic_distance_profiles}
\end{figure}
\fi

\ifincludeimages
\begin{figure}[htbp]
\centering
\begin{subfigure}[htbp]{0.32\textwidth}
    \centering
    \includegraphics[width=\textwidth]{img/lg_aitchison_dp_omega_1_n50.png}
    \caption{$\bomega = \left(\frac{1}{3}, \frac{1}{3}, \frac{1}{3}\right)$, $n=50$}
    \label{fig:lg_aitchison_omega1_n50}
\end{subfigure}
\hfill
\begin{subfigure}[htbp]{0.32\textwidth}
    \centering
    \includegraphics[width=\textwidth]{img/lg_aitchison_dp_omega_2_n50.png}
    \caption{$\bomega = (0.6, 0.3, 0.1)$, $n=50$}
    \label{fig:lg_aitchison_omega2_n50}
\end{subfigure}
\hfill
\begin{subfigure}[htbp]{0.32\textwidth}
    \centering
    \includegraphics[width=\textwidth]{img/lg_aitchison_dp_omega_3_n50.png}
    \caption{$\bomega = (0.2, 0.7, 0.1)$, $n=50$}
    \label{fig:lg_aitchison_omega3_n50}
\end{subfigure}


\begin{subfigure}[htbp]{0.32\textwidth}
    \centering
    \includegraphics[width=\textwidth]{img/lg_aitchison_dp_omega_1_n200.png}
    \caption{$\bomega = \left(\frac{1}{3}, \frac{1}{3}, \frac{1}{3}\right)$, $n=200$}
    \label{fig:lg_aitchison_omega1_n200}
\end{subfigure}
\hfill
\begin{subfigure}[htbp]{0.32\textwidth}
    \centering
    \includegraphics[width=\textwidth]{img/lg_aitchison_dp_omega_2_n200.png}
    \caption{$\bomega = (0.6, 0.3, 0.1)$, $n=200$}
    \label{fig:lg_aitchison_omega2_n200}
\end{subfigure}
\hfill
\begin{subfigure}[htbp]{0.32\textwidth}
    \centering
    \includegraphics[width=\textwidth]{img/lg_aitchison_dp_omega_3_n200.png}
    \caption{$\bomega = (0.2, 0.7, 0.1)$, $n=200$}
    \label{fig:lg_aitchison_omega3_n200}
\end{subfigure}
\caption{Distance profiles for the logistic Gaussian distribution on the simplex $\Delta^{2}$ endowed with the Aitchison distance, with $\bX = \bsigma(\bY)$ where $\bY \sim \mathcal{N}(\bmu, \bSigma)$, $\bmu = (1, 0.5, -0.5)$, and $\bSigma = \left(\begin{smallmatrix} 1.0 & 0.3 & 0.1 \\ 0.3 & 0.8 & 0.2 \\ 0.1 & 0.2 & 0.6 \end{smallmatrix}\right)$.}
\label{fig:lg_aitchison_distance_profiles}
\end{figure}
\fi



Figures~\ref{fig:convergence_process_normal_simple_sigma_5}--\ref{fig:convergence_process_vmf_comp_mu100} show the convergence of $T_n^{\mathrm{KS}}$ and $\tilde{T}_n^{\mathrm{KS}}$ to their respective weak limits under the null hypothesis, for the multivariate normal and vMF distributions. The suprema defining the Kolmogorov--Smirnov statistics are approximated on finite $10\times10$ grids in $\Omega\times\R_{\ge0}$. This grid varies depending on the example considered. %In the normal experiments, the grid is given by $t_j=j t_{\max}/9$ and $\omega_j=-3-t_{\max}+j(6+2t_{\max})/9$, $j=0,\ldots,9$, with $t_{\max}=5\Phi^{-1}(0.995)$. In the vMF experiments, the $\omega$-grid is a $10$-point canonical lattice on $\mathbb{S}^2$, and $t_j=10^{-8}+j(\pi-2\cdot10^{-8})/9$, $j=0,\ldots,9$. 
For each setting, we simulate $M=10000$ realizations of $\sup_{\omega, t}\big|\G_{0}(\omega, t)\big|$ (or $\sup_{\omega, t}\big|\tilde{\G}_{0}(\omega, t)\big|$) and of $T_n^{\mathrm{KS}}$ (or $\tilde{T}_n^{\mathrm{KS}}$) for each sample size $n=50,100,500$. In composite settings, the unknown parameters were estimated by maximum likelihood in each simulated sample. The first row of each figure displays kernel density estimates of the simulated distributions, while the second row gives the corresponding quantile-quantile plots of the approximation errors with respect to the weak limits. For the kernel density estimation, the Gaussian kernel was employed, and the bandwidth was selected by biased cross-validation.



\begin{comment}
\ifincludeimages
\begin{figure}[!htbp]
\begin{subfigure}[htbp]{0.32\textwidth}
    \centering
    \includegraphics[width=\textwidth]{img/simple_mu0_sigma1_M10000_grid10x10.png}
    \caption{$\mathcal{N}(0,1)$}
\end{subfigure}
\hfill
\begin{subfigure}[htbp]{0.32\textwidth}
    \centering
    \includegraphics[width=\textwidth]{img/simple_mu3_sigma1_M10000_grid10x10.png}
    \caption{$\mathcal{N}(3,1)$}
\end{subfigure}
\hfill
\begin{subfigure}[htbp]{0.32\textwidth}
    \centering
    \includegraphics[width=\textwidth]{img/simple_mu-3_sigma1_M10000_grid10x10.png}
    \caption{$\mathcal{N}(-3,1)$}
\end{subfigure}
\par\medskip
\begin{subfigure}[htbp]{0.32\textwidth}
    \centering
    \includegraphics[width=\textwidth]{img/qq_simple_mu0_sigma1_M10000_grid10x10.png}
    \caption{QQ plot for $\mathcal{N}(0,1)$}
\end{subfigure}
\hfill
\begin{subfigure}[htbp]{0.32\textwidth}
    \centering
    \includegraphics[width=\textwidth]{img/qq_simple_mu3_sigma1_M10000_grid10x10.png}
    \caption{QQ plot for $\mathcal{N}(3,1)$}
\end{subfigure}
\hfill
\begin{subfigure}[htbp]{0.32\textwidth}
    \centering
    \includegraphics[width=\textwidth]{img/qq_simple_mu-3_sigma1_M10000_grid10x10.png}
    \caption{QQ plot for $\mathcal{N}(-3,1)$}
\end{subfigure}
\caption{Illustration of the convergence result given in \eqref{eq:thm_5.1_chendubeymuller} (simple $H_0$). The figure shows how the test statistic $T_n = \sup_{\Omega \times \R_{\geq 0}}\left|\G_{n,\mu}^{\mu_0}(\omega, t)\right|$ approaches $\sup_{\Omega \times \R_{\geq 0}}\left|\G_{0}(\omega, t)\right|$ as the sample size $n$ grows. Specifically, the test contrasts whether the data comes from a $\mathcal{N}(0,1)$ (left), $\mathcal{N}(3,1)$ (center), and $\mathcal{N}(-3,1)$ (right) distributions. The first row shows the density plots (using kernel density estimation) and the second row, the corresponding quantile-quantile plots (QQ plots) for the same parameter configurations.}
\label{fig:convergence_process_normal_simple_mu_sigma_1}
\end{figure}
\fi
\end{comment}

\ifincludeimages
\begin{figure}[!htbp]
\begin{subfigure}[htbp]{0.32\textwidth}
    \centering
    \includegraphics[width=\textwidth]{img/simple_mu0_sigma5_M10000_grid10x10.png}
    \caption{$\mathcal{N}(0,25)$}
\end{subfigure}
\hfill
\begin{subfigure}[htbp]{0.32\textwidth}
    \centering
    \includegraphics[width=\textwidth]{img/simple_mu3_sigma5_M10000_grid10x10.png}
    \caption{$\mathcal{N}(3,25)$}
\end{subfigure}
\hfill
\begin{subfigure}[htbp]{0.32\textwidth}
    \centering
    \includegraphics[width=\textwidth]{img/simple_mu-3_sigma5_M10000_grid10x10.png}
    \caption{$\mathcal{N}(-3,25)$}
\end{subfigure}
\par\medskip
\begin{subfigure}[htbp]{0.32\textwidth}
    \centering
    \includegraphics[width=\textwidth]{img/qq_simple_mu0_sigma5_M10000_grid10x10.png}
    \caption{QQ plot for $\mathcal{N}(0,25)$}
\end{subfigure}
\hfill
\begin{subfigure}[htbp]{0.32\textwidth}
    \centering
    \includegraphics[width=\textwidth]{img/qq_simple_mu3_sigma5_M10000_grid10x10.png}
    \caption{QQ plot for $\mathcal{N}(3,25)$}
\end{subfigure}
\hfill
\begin{subfigure}[htbp]{0.32\textwidth}
    \centering
    \includegraphics[width=\textwidth]{img/qq_simple_mu-3_sigma5_M10000_grid10x10.png}
    \caption{QQ plot for $\mathcal{N}(-3,25)$}
\end{subfigure}
\caption{Illustration of the convergence result for $T_n^{\mathrm{KS}}$ given in Theorem~\ref{thm:test_statistics_simp}\ref{thm:test_statistics_simp_i}. Specifically, the test contrasts whether the data comes from a $\mathcal{N}(0,25)$ (left), $\mathcal{N}(3,25)$ (center), and $\mathcal{N}(-3,25)$ (right) distributions.}
\label{fig:convergence_process_normal_simple_sigma_5}
\end{figure}
\fi

\begin{comment}
\ifincludeimages
\begin{figure}[!htbp]
\begin{subfigure}[htbp]{0.32\textwidth}
    \centering
    \includegraphics[width=\textwidth]{img/comp_unk_mu_mu_0_sigma_1_M10000_grid10x10.png}
    \caption{$\mathcal{N}(0,1)$}
\end{subfigure}
\hfill
\begin{subfigure}[htbp]{0.32\textwidth}
    \centering
    \includegraphics[width=\textwidth]{img/comp_unk_mu_mu_3_sigma_1_M10000_grid10x10.png}
    \caption{$\mathcal{N}(3,1)$}
\end{subfigure}
\hfill
\begin{subfigure}[htbp]{0.32\textwidth}
    \centering
    \includegraphics[width=\textwidth]{img/comp_unk_mu_mu_-3_sigma_1_M10000_grid10x10.png}
    \caption{$\mathcal{N}(-3,1)$}
\end{subfigure}
\par\medskip
\begin{subfigure}[htbp]{0.32\textwidth}
    \centering
    \includegraphics[width=\textwidth]{img/qq_comp_unk_mu_mu_0_sigma_1_M10000_grid10x10.png}
    \caption{QQ plot for $\mathcal{N}(0,1)$}
\end{subfigure}
\hfill
\begin{subfigure}[htbp]{0.32\textwidth}
    \centering
    \includegraphics[width=\textwidth]{img/qq_comp_unk_mu_mu_3_sigma_1_M10000_grid10x10.png}
    \caption{QQ plot for $\mathcal{N}(3,1)$}
\end{subfigure}
\hfill
\begin{subfigure}[htbp]{0.32\textwidth}
    \centering
    \includegraphics[width=\textwidth]{img/qq_comp_unk_mu_mu_-3_sigma_1_M10000_grid10x10.png}
    \caption{QQ plot for $\mathcal{N}(-3,1)$}
\end{subfigure}
\caption{Illustration of the convergence of the test statistic $T_n \defin \sup_{\omega \in \Omega, t \ge 0} \left| \G_{n,\mu}^{\mu_{\hat{\btheta}}}(\omega, t) \right|$ to  $\sup_{\Omega \times \R_{\geq 0}}\left|\tilde{\G}_{0}(\omega, t)\right|$ as $n$ grows (composite null). The test contrasts whether the data comes from a normal distribution with a given (known) standard deviation, and some unknown $\mu$.  The estimate $\hat{\mu}$ of $\mu$ was obtained through maximum likelihood. In the three scenarios, the data follow the distributions $\mathcal{N}(0,1)$ (left), $\mathcal{N}(3,1)$ (center), and $\mathcal{N}(-3,1)$ (right).}
\label{fig:convergence_process_normal_comp_unk_mu_sigma_1}
\end{figure}
\fi
\end{comment}

\ifincludeimages
\begin{figure}[!htbp]
\begin{subfigure}[htbp]{0.32\textwidth}
    \centering
    \includegraphics[width=\textwidth]{img/comp_unk_mu_mu_0_sigma_5_M10000_grid10x10.png}
    \caption{$\mathcal{N}(0,25)$}
\end{subfigure}
\hfill
\begin{subfigure}[htbp]{0.32\textwidth}
    \centering
    \includegraphics[width=\textwidth]{img/comp_unk_mu_mu_3_sigma_5_M10000_grid10x10.png}
    \caption{$\mathcal{N}(3,25)$}
\end{subfigure}
\hfill
\begin{subfigure}[htbp]{0.32\textwidth}
    \centering
    \includegraphics[width=\textwidth]{img/comp_unk_mu_mu_-3_sigma_5_M10000_grid10x10.png}
    \caption{$\mathcal{N}(-3,25)$}
\end{subfigure}
\par\medskip
\begin{subfigure}[htbp]{0.32\textwidth}
    \centering
    \includegraphics[width=\textwidth]{img/qq_comp_unk_mu_mu_0_sigma_5_M10000_grid10x10.png}
    \caption{QQ plot for $\mathcal{N}(0,25)$}
\end{subfigure}
\hfill
\begin{subfigure}[htbp]{0.32\textwidth}
    \centering
    \includegraphics[width=\textwidth]{img/qq_comp_unk_mu_mu_3_sigma_5_M10000_grid10x10.png}
    \caption{QQ plot for $\mathcal{N}(3,25)$}
\end{subfigure}
\hfill
\begin{subfigure}[htbp]{0.32\textwidth}
    \centering
    \includegraphics[width=\textwidth]{img/qq_comp_unk_mu_mu_-3_sigma_5_M10000_grid10x10.png}
    \caption{QQ plot for $\mathcal{N}(-3,25)$}
\end{subfigure}
\caption{Illustration of the convergence result for $\tilde{T}_n^{\mathrm{KS}}$ given in Theorem~\ref{thm:test_statistics_comp}\ref{thm:test_statistics_comp_i}. The test contrasts whether the data comes from a normal distribution with a given (known) standard deviation, and some unknown $\mu$.  The estimate $\hat{\mu}$ of $\mu$ was obtained through maximum likelihood. In the three scenarios, the data follow the distributions $\mathcal{N}(0,25)$ (left), $\mathcal{N}(3,25)$ (center), and $\mathcal{N}(-3,25)$ (right).}
\label{fig:convergence_process_normal_comp_unk_mu_sigma_5}
\end{figure}
\fi

\ifincludeimages
\begin{figure}[!htbp]
\begin{subfigure}[htbp]{0.32\textwidth}
    \centering
    \includegraphics[width=\textwidth]{img/comp_unk_sigma_mu_0_sigma_1_M10000_grid10x10.png}
    \caption{$\mathcal{N}(0,1)$}
\end{subfigure}
\hfill
\begin{subfigure}[htbp]{0.32\textwidth}
    \centering
    \includegraphics[width=\textwidth]{img/comp_unk_sigma_mu_0_sigma_2_M10000_grid10x10.png}
    \caption{$\mathcal{N}(0,4)$}
\end{subfigure}
\hfill
\begin{subfigure}[htbp]{0.32\textwidth}
    \centering
    \includegraphics[width=\textwidth]{img/comp_unk_sigma_mu_0_sigma_5_M10000_grid10x10.png}
    \caption{$\mathcal{N}(0,25)$}
\end{subfigure}
\par\medskip
\begin{subfigure}[htbp]{0.32\textwidth}
    \centering
    \includegraphics[width=\textwidth]{img/qq_comp_unk_sigma_mu_0_sigma_1_M10000_grid10x10.png}
    \caption{QQ plot for $\mathcal{N}(0,1)$}
\end{subfigure}
\hfill
\begin{subfigure}[htbp]{0.32\textwidth}
    \centering
    \includegraphics[width=\textwidth]{img/qq_comp_unk_sigma_mu_0_sigma_2_M10000_grid10x10.png}
    \caption{QQ plot for $\mathcal{N}(0,4)$}
\end{subfigure}
\hfill
\begin{subfigure}[htbp]{0.32\textwidth}
    \centering
    \includegraphics[width=\textwidth]{img/qq_comp_unk_sigma_mu_0_sigma_5_M10000_grid10x10.png}
    \caption{QQ plot for $\mathcal{N}(0,25)$}
\end{subfigure}
\caption{Same description as Figure \ref{fig:convergence_process_normal_comp_unk_mu_sigma_5}, but $\mu$ is assumed to be known and $\sigma$ is unknown. The test contrasts whether the data comes from a normal distribution with fixed $\mu$ and unknown $\sigma$, which was estimated through maximum likelihood.}
\label{fig:convergence_process_normal_comp_unk_sigma_mu_0}
\end{figure}
\fi

\begin{comment}
\ifincludeimages
\begin{figure}[!htbp]
\begin{subfigure}[htbp]{0.32\textwidth}
    \centering
    \includegraphics[width=\textwidth]{img/comp_unk_sigma_mu_3_sigma_1_M10000_grid10x10.png}
    \caption{$\mathcal{N}(3,1)$}
\end{subfigure}
\hfill
\begin{subfigure}[htbp]{0.32\textwidth}
    \centering
    \includegraphics[width=\textwidth]{img/comp_unk_sigma_mu_3_sigma_2_M10000_grid10x10.png}
    \caption{$\mathcal{N}(3,4)$}
\end{subfigure}
\hfill
\begin{subfigure}[htbp]{0.32\textwidth}
    \centering
    \includegraphics[width=\textwidth]{img/comp_unk_sigma_mu_3_sigma_5_M10000_grid10x10.png}
    \caption{$\mathcal{N}(3,25)$}
\end{subfigure}
\par\medskip
\begin{subfigure}[htbp]{0.32\textwidth}
    \centering
    \includegraphics[width=\textwidth]{img/qq_comp_unk_sigma_mu_3_sigma_1_M10000_grid10x10.png}
    \caption{QQ plot for $\mathcal{N}(3,1)$}
\end{subfigure}
\hfill
\begin{subfigure}[htbp]{0.32\textwidth}
    \centering
    \includegraphics[width=\textwidth]{img/qq_comp_unk_sigma_mu_3_sigma_2_M10000_grid10x10.png}
    \caption{QQ plot for $\mathcal{N}(3,4)$}
\end{subfigure}
\hfill
\begin{subfigure}[htbp]{0.32\textwidth}
    \centering
    \includegraphics[width=\textwidth]{img/qq_comp_unk_sigma_mu_3_sigma_5_M10000_grid10x10.png}
    \caption{QQ plot for $\mathcal{N}(3,25)$}
\end{subfigure}
\caption{Same description as Figure \ref{fig:convergence_process_normal_comp_unk_sigma_mu_0}, but with $\mu=3$.}
\label{fig:convergence_process_normal_comp_unk_sigma_mu_3}
\end{figure}
\fi
\end{comment}

\begin{comment}
\ifincludeimages
\begin{figure}[!htbp]
\begin{subfigure}[htbp]{0.32\textwidth}
    \centering
    \includegraphics[width=\textwidth]{img/comp_unk_both_mu_0_sigma_1_M10000_grid10x10.png}
    \caption{$\mathcal{N}(0,1)$}
\end{subfigure}
\hfill
\begin{subfigure}[htbp]{0.32\textwidth}
    \centering
    \includegraphics[width=\textwidth]{img/comp_unk_both_mu_0_sigma_2_M10000_grid10x10.png}
    \caption{$\mathcal{N}(0,4)$}
\end{subfigure}
\hfill
\begin{subfigure}[htbp]{0.32\textwidth}
    \centering
    \includegraphics[width=\textwidth]{img/comp_unk_both_mu_0_sigma_5_M10000_grid10x10.png}
    \caption{$\mathcal{N}(0,25)$}
\end{subfigure}
\par\medskip
\begin{subfigure}[htbp]{0.32\textwidth}
    \centering
    \includegraphics[width=\textwidth]{img/qq_comp_unk_both_mu_0_sigma_1_M10000_grid10x10.png}
    \caption{QQ plot for $\mathcal{N}(0,1)$}
\end{subfigure}
\hfill
\begin{subfigure}[htbp]{0.32\textwidth}
    \centering
    \includegraphics[width=\textwidth]{img/qq_comp_unk_both_mu_0_sigma_2_M10000_grid10x10.png}
    \caption{QQ plot for $\mathcal{N}(0,4)$}
\end{subfigure}
\hfill
\begin{subfigure}[htbp]{0.32\textwidth}
    \centering
    \includegraphics[width=\textwidth]{img/qq_comp_unk_both_mu_0_sigma_5_M10000_grid10x10.png}
    \caption{QQ plot for $\mathcal{N}(0,25)$}
\end{subfigure}
\caption{The test contrasts whether the data comes from a normal distribution with unknown parameters $(\mu, \sigma)$, which are estimated through maximum likelihood. % The true distribution of the data was $\mathcal{N}(0,1)$ (left), $\mathcal{N}(0,4)$ (center), and $\mathcal{N}(0,25)$ (right).
}
\label{fig:convergence_process_normal_comp_unk_both}
\end{figure}
\fi
\end{comment}

\ifincludeimages
\begin{figure}[!htbp]
\begin{subfigure}[htbp]{0.32\textwidth}
    \centering
    \includegraphics[width=\textwidth]{img/comp_unk_both_mu_0_sigma_1_M10000_grid10x10.png}
    \caption{$\mathcal{N}(0,1)$}
\end{subfigure}
\hfill
\begin{subfigure}[htbp]{0.32\textwidth}
    \centering
    \includegraphics[width=\textwidth]{img/comp_unk_both_mu_3_sigma_1_M10000_grid10x10.png}
    \caption{$\mathcal{N}(3,1)$}
\end{subfigure}
\hfill
\begin{subfigure}[htbp]{0.32\textwidth}
    \centering
    \includegraphics[width=\textwidth]{img/comp_unk_both_mu_-3_sigma_1_M10000_grid10x10.png}
    \caption{$\mathcal{N}(-3,1)$}
\end{subfigure}
\par\medskip
\begin{subfigure}[htbp]{0.32\textwidth}
    \centering
    \includegraphics[width=\textwidth]{img/qq_comp_unk_both_mu_0_sigma_1_M10000_grid10x10.png}
    \caption{QQ plot for $\mathcal{N}(0,1)$}
\end{subfigure}
\hfill
\begin{subfigure}[htbp]{0.32\textwidth}
    \centering
    \includegraphics[width=\textwidth]{img/qq_comp_unk_both_mu_3_sigma_1_M10000_grid10x10.png}
    \caption{QQ plot for $\mathcal{N}(3,1)$}
\end{subfigure}
\hfill
\begin{subfigure}[htbp]{0.32\textwidth}
    \centering
    \includegraphics[width=\textwidth]{img/qq_comp_unk_both_mu_-3_sigma_1_M10000_grid10x10.png}
    \caption{QQ plot for $\mathcal{N}(-3,1)$}
\end{subfigure}
\caption{The test contrasts whether the data comes from a normal distribution with unknown parameters $(\mu, \sigma)$, which are estimated through maximum likelihood. % The true distribution of the data was $\mathcal{N}(0,1)$ (left), $\mathcal{N}(3,1)$ (center), and $\mathcal{N}(-3,1)$ (right).
}
\label{fig:convergence_process_normal_comp_unk_both_sigma_1}
\end{figure}
\fi


\ifincludeimages
\begin{figure}[!htbp]
\begin{subfigure}[htbp]{0.32\textwidth}
    \centering
    \includegraphics[width=\textwidth]{img/simple_vmf_kappa0.5_mu1p0_0p0_0p0_geodesic_M10000_grid10x10.png}
    \caption{$\mathrm{vMF}(\bmu,\frac{1}{2})$}
\end{subfigure}
\hfill
\begin{subfigure}[htbp]{0.32\textwidth}
    \centering
    \includegraphics[width=\textwidth]{img/simple_vmf_kappa1_mu1p0_0p0_0p0_geodesic_M10000_grid10x10.png}
    \caption{$\mathrm{vMF}(\bmu,1)$}
\end{subfigure}
\hfill
\begin{subfigure}[htbp]{0.32\textwidth}
    \centering
    \includegraphics[width=\textwidth]{img/simple_vmf_kappa5_mu1p0_0p0_0p0_geodesic_M10000_grid10x10.png}
    \caption{$\mathrm{vMF}(\bmu,5)$}
\end{subfigure}
\par\medskip
\begin{subfigure}[htbp]{0.32\textwidth}
    \centering
    \includegraphics[width=\textwidth]{img/qq_simple_vmf_kappa0.5_mu1p0_0p0_0p0_geodesic_M10000_grid10x10.png}
    \caption{QQ plot for $\mathrm{vMF}(\bmu,\frac{1}{2})$}
\end{subfigure}
\hfill
\begin{subfigure}[htbp]{0.32\textwidth}
    \centering
    \includegraphics[width=\textwidth]{img/qq_simple_vmf_kappa1_mu1p0_0p0_0p0_geodesic_M10000_grid10x10.png}
    \caption{QQ plot for $\mathrm{vMF}(\bmu,1)$}
\end{subfigure}
\hfill
\begin{subfigure}[htbp]{0.32\textwidth}
    \centering
    \includegraphics[width=\textwidth]{img/qq_simple_vmf_kappa5_mu1p0_0p0_0p0_geodesic_M10000_grid10x10.png}
    \caption{QQ plot for $\mathrm{vMF}(\bmu,5)$}
\end{subfigure}
\caption{Same description as Figure \ref{fig:convergence_process_normal_simple_sigma_5}, but for the $\mathrm{vMF}$ distribution with canonical parameter $\bxi$ (geodesic distance). Specifically, the test contrasts whether the data comes from a $\mathrm{vMF}$ with canonical parameter $\bxi = \frac{1}{2}\bmu$ (left), $\bxi = \bmu$ (center), and $\bxi = 5\bmu$ (right), with $\bmu = (1,0,0)^\top$.}
\label{fig:convergence_process_vmf_simple_mu_100}
\end{figure}
\fi

\begin{comment}
\ifincludeimages
\begin{figure}[!htbp]
\begin{subfigure}[htbp]{0.32\textwidth}
    \centering
    \includegraphics[width=\textwidth]{img/simple_vmf_kappa0.5_mu0p6_0p6_0p6_geodesic_M10000_grid10x10.png}
    \caption{$\mathrm{vMF}(\bmu, \frac{1}{2})$}
\end{subfigure}
\hfill
\begin{subfigure}[htbp]{0.32\textwidth}
    \centering
    \includegraphics[width=\textwidth]{img/simple_vmf_kappa1_mu0p6_0p6_0p6_geodesic_M10000_grid10x10.png}
    \caption{$\mathrm{vMF}(\bmu, 1)$}
\end{subfigure}
\hfill
\begin{subfigure}[htbp]{0.32\textwidth}
    \centering
    \includegraphics[width=\textwidth]{img/simple_vmf_kappa5_mu0p6_0p6_0p6_geodesic_M10000_grid10x10.png}
    \caption{$\mathrm{vMF}(\bmu, 5)$}
\end{subfigure}
\par\medskip
\begin{subfigure}[htbp]{0.32\textwidth}
    \centering
    \includegraphics[width=\textwidth]{img/qq_simple_vmf_kappa0.5_mu0p6_0p6_0p6_geodesic_M10000_grid10x10.png}
    \caption{QQ plot for $\mathrm{vMF}(\bmu, \frac{1}{2})$}
\end{subfigure}
\hfill
\begin{subfigure}[htbp]{0.32\textwidth}
    \centering
    \includegraphics[width=\textwidth]{img/qq_simple_vmf_kappa1_mu0p6_0p6_0p6_geodesic_M10000_grid10x10.png}
    \caption{QQ plot for $\mathrm{vMF}(\bmu, 1)$}
\end{subfigure}
\hfill
\begin{subfigure}[htbp]{0.32\textwidth}
    \centering
    \includegraphics[width=\textwidth]{img/qq_simple_vmf_kappa5_mu0p6_0p6_0p6_geodesic_M10000_grid10x10.png}
    \caption{QQ plot for $\mathrm{vMF}(\bmu, 5)$}
\end{subfigure}
\caption{Same description as Figure \ref{fig:convergence_process_vmf_simple_mu_100}, but for $\bmu= \big(\frac{1}{\sqrt{3}}, \frac{1}{\sqrt{3}}, \frac{1}{\sqrt{3}}\big)^\top$.}
\label{fig:convergence_process_vMF_simple_mu_sqrt_3}
\end{figure}
\fi

\ifincludeimages
\begin{figure}[!htbp]
\begin{subfigure}[htbp]{0.32\textwidth}
    \centering
    \includegraphics[width=\textwidth]{img/simple_vmf_kappa0.5_mu-0p7_0p0_-0p7_geodesic_M10000_grid10x10.png}
    \caption{$\mathrm{vMF}(\bmu, \frac{1}{2})$}
\end{subfigure}
\hfill
\begin{subfigure}[htbp]{0.32\textwidth}
    \centering
    \includegraphics[width=\textwidth]{img/simple_vmf_kappa1_mu-0p7_0p0_-0p7_geodesic_M10000_grid10x10.png}
    \caption{$\mathrm{vMF}(\bmu, 1)$}
\end{subfigure}
\hfill
\begin{subfigure}[htbp]{0.32\textwidth}
    \centering
    \includegraphics[width=\textwidth]{img/simple_vmf_kappa5_mu-0p7_0p0_-0p7_geodesic_M10000_grid10x10.png}
    \caption{$\mathrm{vMF}(\bmu, 5)$}
\end{subfigure}
\par\medskip
\begin{subfigure}[htbp]{0.32\textwidth}
    \centering
    \includegraphics[width=\textwidth]{img/qq_simple_vmf_kappa0.5_mu-0p7_0p0_-0p7_geodesic_M10000_grid10x10.png}
    \caption{QQ plot for $\mathrm{vMF}(\bmu, \frac{1}{2})$}
\end{subfigure}
\hfill
\begin{subfigure}[htbp]{0.32\textwidth}
    \centering
    \includegraphics[width=\textwidth]{img/qq_simple_vmf_kappa1_mu-0p7_0p0_-0p7_geodesic_M10000_grid10x10.png}
    \caption{QQ plot for $\mathrm{vMF}(\bmu, 1)$}
\end{subfigure}
\hfill
\begin{subfigure}[htbp]{0.32\textwidth}
    \centering
    \includegraphics[width=\textwidth]{img/qq_simple_vmf_kappa5_mu-0p7_0p0_-0p7_geodesic_M10000_grid10x10.png}
    \caption{QQ plot for $\mathrm{vMF}(\bmu, 5)$}
\end{subfigure}
\caption{Same description as Figure \ref{fig:convergence_process_vmf_simple_mu_100}, but for $\bmu=\big(-\frac{1}{\sqrt{2}}, 0, -\frac{1}{\sqrt{2}}\big)^\top$.}
\label{fig:convergence_process_vMF_simple_mu_sqrt_2}
\end{figure}
\fi
\end{comment}





\ifincludeimages
\begin{figure}[!htbp]
\begin{subfigure}[htbp]{0.32\textwidth}
    \centering
    \includegraphics[width=\textwidth]{img/comp_vmf_xi_kappa0.5_mu1p0_0p0_0p0_geodesic_M10000_grid10x10.png}
    \caption{$\mathrm{vMF}(\bmu, \frac{1}{2})$}
\end{subfigure}
\hfill
\begin{subfigure}[htbp]{0.32\textwidth}
    \centering
    \includegraphics[width=\textwidth]{img/comp_vmf_xi_kappa1_mu1p0_0p0_0p0_geodesic_M10000_grid10x10.png}
    \caption{$\mathrm{vMF}(\bmu, 1)$}
\end{subfigure}
\hfill
\begin{subfigure}[htbp]{0.32\textwidth}
    \centering
    \includegraphics[width=\textwidth]{img/comp_vmf_xi_kappa5_mu1p0_0p0_0p0_geodesic_M10000_grid10x10.png}
    \caption{$\mathrm{vMF}(\bmu, 5)$}
\end{subfigure}
\par\medskip
\begin{subfigure}[htbp]{0.32\textwidth}
    \centering
    \includegraphics[width=\textwidth]{img/qq_comp_vmf_xi_kappa0.5_mu1p0_0p0_0p0_geodesic_M10000_grid10x10.png}
    \caption{QQ plot for $\mathrm{vMF}(\bmu, \frac{1}{2})$}
\end{subfigure}
\hfill
\begin{subfigure}[htbp]{0.32\textwidth}
    \centering
    \includegraphics[width=\textwidth]{img/qq_comp_vmf_xi_kappa1_mu1p0_0p0_0p0_geodesic_M10000_grid10x10.png}
    \caption{QQ plot for $\mathrm{vMF}(\bmu, 1)$}
\end{subfigure}
\hfill
\begin{subfigure}[htbp]{0.32\textwidth}
    \centering
    \includegraphics[width=\textwidth]{img/qq_comp_vmf_xi_kappa5_mu1p0_0p0_0p0_geodesic_M10000_grid10x10.png}
    \caption{QQ plot for $\mathrm{vMF}(\bmu, 5)$}
\end{subfigure}
\caption{Same description as Figure \ref{fig:convergence_process_normal_comp_unk_mu_sigma_5}, for the $\mathrm{vMF}$ distribution with canonical parameter $\bxi=\kappa\bmu$ (geodesic distance). The test contrasts whether the data come from a $\mathrm{vMF}$ distribution with some unknown canonical parameter $\bxi = \kappa \bmu$. The estimate $\hat{\bxi}$ of $\bxi$ was obtained through maximum likelihood. In each scenario, the data follow either a $\mathrm{vMF}(\bmu,\frac{1}{2})$ (left), a $\mathrm{vMF}(\bmu,1)$ (center), and a $\mathrm{vMF}(\bmu,5)$ (right) distributions, with $\bmu= (1,0,0)^\top$.}
\label{fig:convergence_process_vmf_comp_mu100}
\end{figure}
\fi



Figures~\ref{fig:vmf_bahadur} and \ref{fig:vmf_bahadur_loglog} provide numerical evidence consistent with the Bahadur representation established in Proposition~\ref{prop:bahadur_vmf} for the MLE of the canonical parameter $\bxi=\kappa \bmu$ in the $\mathrm{vMF}(\bmu,1)$ model on $\mathbb{S}^2$. For several mean directions $\bmu$, and for increasing sample sizes $n$, we examine the magnitude of the difference between $\sqrt{n}(\hat{\bxi}-\bxi)$ and the linear approximation given in \eqref{eq:bahadur_vMF}. %If the expansion holds, this discrepancy should tend to $0$ as $n\to\infty$. 
For each value of $n$, the figures display the pointwise mean over $M=500$ simulated samples together with pointwise $95\%$ confidence intervals (based on the normal distribution), and include the reference curve $\frac{1}{\sqrt{n}}$ for comparison. The results confirm that the convergence rate is of order $\frac{1}{\sqrt{n}}$, as expected. 

\ifincludeimages
\begin{figure}[htbp]
    \centering
    {\color{red}
    \begin{subfigure}[htbp]{0.32\textwidth}
        \centering
        \includegraphics[width=\textwidth]{img/vmf_bahadur_linear_mu_1_0_0.png}
        \caption{$\bmu = (1,0,0)$}
    \end{subfigure}
    \hfill
    \begin{subfigure}[htbp]{0.32\textwidth}
        \centering
        \includegraphics[width=\textwidth]{img/vmf_bahadur_linear_mu_inv_sqrt_3_inv_sqrt_3_inv_sqrt_3.png}
        \caption{$\bmu = \left(\frac{1}{\sqrt{3}},\frac{1}{\sqrt{3}},\frac{1}{\sqrt{3}}\right)$}
    \end{subfigure}
    \hfill
    \begin{subfigure}[htbp]{0.32\textwidth}
        \centering
        \includegraphics[width=\textwidth]{img/vmf_bahadur_linear_mu_minus_inv_sqrt_3_inv_sqrt_3_inv_sqrt_3.png}
        \caption{$\bmu = \left(-\frac{1}{\sqrt{3}},\frac{1}{\sqrt{3}},\frac{1}{\sqrt{3}}\right)$}
    \end{subfigure}
    }
\caption{Validation of the Bahadur representation for the MLE of a $\mathrm{vMF}(\bmu, 1)$ on $\mathbb{S}^2$. The mean and (red dashed line) and the confidence intervals ($95\%$, in a lighter red) were calculated pointwise, based on $500$ realizations (on $500$ different samples) of $\|\mathrm{LHS} - \mathrm{RHS}\|_2$ for a different mean direction $\bmu$, as sample size $n$ increases, where $\mathrm{LHS} = \sqrt{n}(\hat{\bxi}_n - \bxi)$, and  $\mathrm{RHS} = -\big(\dot{\psi}_{\bxi}\big)^{-1} \frac{1}{\sqrt{n}} \sum_{i=1}^n \psi_{\bxi}(\bX_i)$. The convergence rate corresponding to $\frac{1}{\sqrt{n}}$ (dashed black line) is plotted for comparison.}
\label{fig:vmf_bahadur}
\end{figure}
\fi


\ifincludeimages
\begin{figure}[htbp]
    \centering
    {\color{red}
    \begin{subfigure}[htbp]{0.32\textwidth}
        \centering
        \includegraphics[width=\textwidth]{img/vmf_bahadur_loglog_mu_1_0_0.png}
        \caption{$\bmu = (1,0,0)$}
    \end{subfigure}
    \hfill
    \begin{subfigure}[htbp]{0.32\textwidth}
        \centering
        \includegraphics[width=\textwidth]{img/vmf_bahadur_loglog_mu_inv_sqrt_3_inv_sqrt_3_inv_sqrt_3.png}
        \caption{$\bmu = \left(\frac{1}{\sqrt{3}},\frac{1}{\sqrt{3}},\frac{1}{\sqrt{3}}\right)$}
    \end{subfigure}
    \hfill
    \begin{subfigure}[htbp]{0.32\textwidth}
        \centering
        \includegraphics[width=\textwidth]{img/vmf_bahadur_loglog_mu_minus_inv_sqrt_3_inv_sqrt_3_inv_sqrt_3.png}
        \caption{$\bmu = \left(-\frac{1}{\sqrt{3}},\frac{1}{\sqrt{3}},\frac{1}{\sqrt{3}}\right)$}
    \end{subfigure}
    }
\caption{Same description as Figure \ref{fig:vmf_bahadur}, but in log-log scale.}
\label{fig:vmf_bahadur_loglog}
\end{figure}
\fi


\begin{comment}
\ifincludeimages
\begin{figure}[!htbp]
\begin{subfigure}[htbp]{0.32\textwidth}
    \centering
    \includegraphics[width=\textwidth]{img/comp_vmf_xi_kappa0.5_mu0p6_0p6_0p6_geodesic_M10000_grid10x10.png}
    \caption{$\mathrm{vMF}(\bmu, \frac{1}{2})$}
\end{subfigure}
\hfill
\begin{subfigure}[htbp]{0.32\textwidth}
    \centering
    \includegraphics[width=\textwidth]{img/comp_vmf_xi_kappa1_mu0p6_0p6_0p6_geodesic_M10000_grid10x10.png}
    \caption{$\mathrm{vMF}(\bmu, 1)$}
\end{subfigure}
\hfill
\begin{subfigure}[htbp]{0.32\textwidth}
    \centering
    \includegraphics[width=\textwidth]{img/comp_vmf_xi_kappa5_mu0p6_0p6_0p6_geodesic_M10000_grid10x10.png}
    \caption{$\mathrm{vMF}(\bmu, 5)$}
\end{subfigure}
\par\medskip
\begin{subfigure}[htbp]{0.32\textwidth}
    \centering
    \includegraphics[width=\textwidth]{img/qq_comp_vmf_xi_kappa0.5_mu0p6_0p6_0p6_geodesic_M10000_grid10x10.png}
    \caption{QQ plot for $\mathrm{vMF}(\bmu, \frac{1}{2})$}
\end{subfigure}
\hfill
\begin{subfigure}[htbp]{0.32\textwidth}
    \centering
    \includegraphics[width=\textwidth]{img/qq_comp_vmf_xi_kappa1_mu0p6_0p6_0p6_geodesic_M10000_grid10x10.png}
    \caption{QQ plot for $\mathrm{vMF}(\bmu, 1)$}
\end{subfigure}
\hfill
\begin{subfigure}[htbp]{0.32\textwidth}
    \centering
    \includegraphics[width=\textwidth]{img/qq_comp_vmf_xi_kappa5_mu0p6_0p6_0p6_geodesic_M10000_grid10x10.png}
    \caption{QQ plot for $\mathrm{vMF}(\bmu, 5)$}
\end{subfigure}
\caption{Same description as Figure \ref{fig:convergence_process_vmf_comp_mu100}, but for $\bmu=(\frac{1}{\sqrt{3}}, \frac{1}{\sqrt{3}}, \frac{1}{\sqrt{3}})^\top$.}
\label{fig:convergence_process_vmf_comp_mu_sqrt_3}
\end{figure}
\fi

\ifincludeimages
\begin{figure}[!htbp]
\begin{subfigure}[htbp]{0.32\textwidth}
    \centering
    \includegraphics[width=\textwidth]{img/comp_vmf_xi_kappa0.5_mu-0p7_0p0_-0p7_geodesic_M10000_grid10x10.png}
    \caption{$\mathrm{vMF}(\bmu, \frac{1}{2})$}
\end{subfigure}
\hfill
\begin{subfigure}[htbp]{0.32\textwidth}
    \centering
    \includegraphics[width=\textwidth]{img/comp_vmf_xi_kappa1_mu-0p7_0p0_-0p7_geodesic_M10000_grid10x10.png}
    \caption{$\mathrm{vMF}(\bmu, 1)$}
\end{subfigure}
\hfill
\begin{subfigure}[htbp]{0.32\textwidth}
    \centering
    \includegraphics[width=\textwidth]{img/comp_vmf_xi_kappa5_mu-0p7_0p0_-0p7_geodesic_M10000_grid10x10.png}
    \caption{$\mathrm{vMF}(\bmu, 5)$}
\end{subfigure}
\par\medskip
\begin{subfigure}[htbp]{0.32\textwidth}
    \centering
    \includegraphics[width=\textwidth]{img/qq_comp_vmf_xi_kappa0.5_mu-0p7_0p0_-0p7_geodesic_M10000_grid10x10.png}
    \caption{QQ plot for $\mathrm{vMF}(\bmu, \frac{1}{2})$}
\end{subfigure}
\hfill
\begin{subfigure}[htbp]{0.32\textwidth}
    \centering
    \includegraphics[width=\textwidth]{img/qq_comp_vmf_xi_kappa1_mu-0p7_0p0_-0p7_geodesic_M10000_grid10x10.png}
    \caption{QQ plot for $\mathrm{vMF}(\bmu, 1)$}
\end{subfigure}
\hfill
\begin{subfigure}[htbp]{0.32\textwidth}
    \centering
    \includegraphics[width=\textwidth]{img/qq_comp_vmf_xi_kappa5_mu-0p7_0p0_-0p7_geodesic_M10000_grid10x10.png}
    \caption{QQ plot for $\mathrm{vMF}(\bmu, 5)$}
\end{subfigure}
\caption{Same description as Figure \ref{fig:convergence_process_vmf_comp_mu100}, but for $\bmu=(-\frac{1}{\sqrt{2}}, 0, -\frac{1}{\sqrt{2}})^\top$.}
\label{fig:convergence_process_vmf_comp_mu_sqrt_2}
\end{figure}
\fi
\end{comment}

%----------------------%
%---------------------------%
\section{Real data application}
%---------------------------%

Some ideas:
\begin{enumerate}
    \item Test the goodness-of-fit of sunspot positions with a mixture of two small circle distributions to model the positions in both hemispheres. Recall $f_{\mathrm{SC}}(\bx;\bmu,\tau,\mu)\defin c(\tau,\mu) \exp(-\tau(\bx^\top\bmu - \mu)^2)$. It is natural that the SCs for both hemispheres share the location parameter, have the same weights, and have symmetric effect of the parameters. So we can actually merge the mixtures under a rotationally symmetric model:
    $$
    f(\bx;\bmu,\tau,\mu)=\frac{c(\tau,\mu)}{2}\lrb{\exp(-\tau(\bx^\top\bmu - \mu)^2)+\exp(-\tau(\bx^\top\bmu + \mu)^2)}.
    $$
    We can estimate by MLE the axis $\bmu$ (very close to $(0,0,1)^\top$), concentration $\tau$, and location $\mu$.
    \item \cite{Jensen1981} fitted von Mises distributions to wind direction data to test independence between speed and direction.
    % \item Early works on the hyperbolic von Mises--Fisher family report real datasets allegedly fitting a ``hyperbolic distribution'' extremely well, typically assessed via graphical checks \citep[see, e.g.,][]{Barndorff-Nielsen1978}. It would be of interest to revisit those analyses using our goodness-of-fit tests to provide proper statistical guarantees. The datasets could be reconstructed via reverse engineering from the published figures. Update: these works are for Euclidean data; perhaps it is not of interest from a metric-space perspective.
    \item Comets:
    \begin{itemize}
        \item Puedes replicar el test de los long-period comets con el abordaje de los distance profiles.
        \item El resultado del análisis de los short-period comets te lo deja a huevo (historia + justificación) para continuar el análisis de datos reales en tu artículo: el modelo adecuado para estos cometas es uno rotacionalmente simétrico y unimodal (más fácil) pero que se concentra rápidamente en su moda. Creo que una distribución esférica Cauchy (hay dos tipos) o la distribución de Jones and Pewsey podrían ser buenas candidatas. La vMF creo que no va a dar un buen ajuste porque no se concentra lo suficientemente rápido en la moda.
        \item Bonus: debido a la equivalencia entre la cardioide y Watson para concentraciones bajas, podría ser que los long-period comets también se pudiesen modelar adecuadamente con una distribución de Watson. Esto sería doblemente interesante, porque si consideras una distribución small circle (generaliza a la Watson) y calculas su maquinaria de proyecciones, esta podría ser doblemente útil: para modelar los long-period comets y para modelar las manchas solares (en este caso, con una mixtura de dos o más small circles con el mismo eje).
        \item Sobre las distribuciones:
        \begin{itemize}
            \item Ver Jones and Pewsey (2005), Sección 5. Tendrás que ver cómo simularla. Inversión de la cdf proyectada es lo más directo: mira \verb|r_sph_car()| en el código que te mando como plantilla.
            \item Mira la prueba del Teorema 9, página 47, para ver cómo calcular distribuciones proyectadas de modelos rotacionalmente simétricos.
            \item En el artículo de Alberto y mío en TEST puedes ver una variante de la spherical Cauchy, Sección B.1. Puedes simularla con \verb|?sphunif::r_alt(alt = "C")|.
        \end{itemize}
    \end{itemize}


 Para ti, el artículo es relevante por lo siguiente:
 Usa el abordaje projected-ecdf, muy relacionado con los distance profiles. La diferencia es que en las proyecciones el signo se tiene en cuenta.
 La distribución de las proyecciones en la cardioide es más sencilla que la de la vMF, por lo que será más fácil de usar con la filosofía de los distance profiles. En la prueba del Teorema 8 tienes la "representación de Bahadur" del estimador MLE.
Se discuten varios pesos en el estadístico de CvM. No encontré grandes diferencias, por lo que podría estar justificado tomar el más sencillo (Pn). Si bien los estudios numéricos no son exhaustivos y por tanto tampoco concluyentes.
 Puedes replicar el test de los long-period comets con el abordaje de los distance profiles.
 El regalito: el resultado del análisis de los short-period comets te lo deja a huevo (historia + justificación) para continuar el análisis de datos reales en tu artículo: el modelo adecuado para estos cometas es uno rotacionalmente simétrico y unimodal (más fácil) pero que se concentra rápidamente en su moda. Creo que una distribución esférica Cauchy (hay dos tipos) o la distribución de Jones and Pewsey podrían ser buenas candidatas. La vMF creo que no va a dar un buen ajuste porque no se concentra lo suficientemente rápido en la moda.
 Bonus: debido a la equivalencia entre la cardioide y Watson para concentraciones bajas, podría ser que los long-period comets también se pudiesen modelar adecuadamente con una distribución de Watson. Esto sería doblemente interesante, porque si consideras una distribución small circle (generaliza a la Watson) y calculas su maquinaria de proyecciones, esta podría ser doblemente útil: para modelar los long-period comets y para modelar las manchas solares (en este caso, con una mixtura de dos o más small circles con el mismo eje).
 Sobre las distribuciones:
 - Ver Jones and Pewsey (2005), Sección 5. Tendrás que ver cómo simularla. Inversión de la cdf proyectada es lo más directo: mira $r_sph_car$() en el código que te mando como plantilla.
 - Mira la prueba del Teorema 9, página 47, para ver cómo calcular distribuciones proyectadas de modelos rotacionalmente simétricos.
 - En el artículo de Alberto y mío en TEST puedes ver una variante de la spherical Cauchy, Sección B.1. Puedes simumarla con $?sphunif::r_alt(alt = "C")$.
\end{enumerate}

\bibliographystyle{imsart-nameyear}

\bibliography{sample,bibliography}

\newpage

\appendix

%---------------------------%
\section{\texorpdfstring{Proof of Proposition \ref{prop:dp_examples}}{Proof of Proposition}}\label{ap:proof_prop_dp_examples}
%---------------------------%

\begin{proof}[Proof of Proposition~\ref{prop:dp_examples}]
The proof is done following the same order as Proposition~\ref{prop:dp_examples}.

\begin{itemize}
\item Example~\ref{ex:dp_mvn}.
The distance profile at $\bomega \in \R^q$ is given by
$$
F_{\bomega}^{\mu}(t) = \Prob{d(\bomega, \bX) \le t} = \Prob{\| \bomega - \bX \|_2 \le t} = \Prob{\sqrt{\bZ^\top \bZ} \le t},
$$
where $\bZ = \bX - \bomega$ and $\|\cdot\|_2$ is the Euclidean norm in $\R^q$. Thus, $\bZ \sim \mathcal{N}\left(\bmu - \bomega, \bSigma \right)$.
We are interested in the distribution of $\bZ^\top \bZ$. Consider $\boldsymbol{\eta} = \bU^\top \bZ$, where $\bU$ is an orthogonal matrix such that $\bSigma = \bU \bLambda \bU^\top$ (spectral decomposition). Thus, $\boldsymbol{\eta} \sim \mathcal{N}_q(\bU^\top (\bmu - \bomega), \bLambda)$. We conclude that
$$\bZ^\top \bZ =
%\boldsymbol{\eta}^\top \bU^\top \bU \boldsymbol{\eta} =
\boldsymbol{\eta}^\top \boldsymbol{\eta} = \sum_{i=1}^q \eta_i^2 = \sum_{i=1}^q \lambda_i P_i,
$$
with $\eta_i$ independent, $\eta_i \sim \mathcal{N}(\nu_i, \lambda_i)$, and thus $P_i$ are also independent and follow a noncentral chi-squared distribution, $P_i \sim \chi_1^2\left(\frac{\nu_i^2}{\lambda_i}\right)$.
Here, $\boldsymbol{\nu} = \bU^\top (\bmu - \bomega)$ and $\blambda$ is the vector of eigenvalues of $\bSigma$. In the special case when $\bomega = \bmu$, we have that
\begin{align*}
\bZ^\top \bZ = \sum_{i=1}^q \lambda_i P_i,
\end{align*}
where $P_i$ are independent, $P_i \sim \chi_1^2$. If additionally, $\bSigma=\bI_q$, where $\bI_q$ is the $q\times q$ identity matrix, then $\sqrt{\bZ^\top \bZ}$ follows a $\chi_q$ distribution.

\item Example~\ref{ex:dp_brownian}.
The distance profile of $B$ at a point $\omega \in L^2[0,1]$ is given by
$$
F_{\omega}(t) = \Prob{ d_{L^2}(B, \omega) \le t } = \Prob{ \|B - \omega\|_{L^2} \le t}.
$$
In order to analyze the distribution of $ \|B - \omega\|_{L^2}$, express $B$ and $\omega$ in the Karhunen-Loève expansion. Let $\{e_k\}_{k=1}^\infty$ denote the eigenfunctions of the covariance kernel of $B$, given by $K_B(s,t) = \Cov{B(s)}{B(t)} = \min\{s, t\}$. Then, $B$ and $\omega$ can be expressed as
$$
B(t) = \sum_{i=1}^\infty Z_i e_i(t), \text{ and } \omega(t) = \sum_{i=1}^\infty c_i e_i(t),
$$
where $Z_i$ are independent random variables with $Z_i \sim \mathcal{N}(0, \lambda_i)$ and the coefficients $c_i$ are given by $c_i = \langle \omega, e_i \rangle = \int_{0}^1\omega(t)e_i(t) \, \rd t$. Hence,
$
B(t) - \omega(t) = \sum_{i=1}^\infty ( Z_i- c_i )e_i(t),
$
where $\lambda_k=\frac{1}{\left(k-\frac{1}{2}\right)^2 \pi^2}$.
Thus,
$$
\|B - \omega\|_{L^2}^2 = \sum_{i=1}^\infty ( Z_i- c_i )^2 = \sum_{i=1}^\infty X_i^2,
$$
where $X_i\defin c_i-Z_i$; with $X_1, \ldots$ a sequence of independent random variables, $X_i \sim \mathcal{N}(c_i, \lambda_i)$.
As a special case, if $\omega(t)=0$ for all $t \in [0,1]$, then $\|B - \omega\|_{L^2}^2 =  \sum_{i=1}^\infty \lambda_i Q_i$, where $Q_i$ are iid and $Q_i \sim \chi_1^2$. Observe the similarity to the distance profile for the multivariate normal when $\bomega = \bmu$.

\item Example~\ref{ex:dp_vmf}.
Assume first that $|\bmu^\top\bomega|<1$. To characterize the distance profile, the distribution of $T = \bomega^\top \bX$ plays a central role. To simplify notation, let $\bmu_\perp \defin \bmu - (\bmu^\top \bomega)\bomega$, where $\| \bmu_\perp\| = \sqrt{1 - (\bmu^\top\bomega)^2}$. Define also $\Tilde{\bmu} = \frac{\bmu_\perp}{\| \bmu_\perp \|}$.
The tangent-normal decomposition allows expressing any $\bx \in \mathbb{S}^{q}$ as $\bx = t\bomega + \sqrt{1-t^2}\bu$, where $\bu \in \mathbb{S}^{q-1}$ and $\bu^\top \bomega = 0$. We have that
$$
\bmu^\top \bx = t\bmu^\top \bomega + \sqrt{1-t^2} \bmu^\top \bu = t \bmu^\top\bomega +  \sqrt{1-t^2} \bmu_\perp^\top \bu = t \bmu^\top\bomega + \sqrt{(1-t^2)(1-(\bmu^\top\bomega)^2)}\Tilde{\bmu}^\top \bu.
$$
Hence,
$$
f_\mathrm{vMF}(\bx;\bmu,\kappa) = c_{q}^\mathrm{vMF}(\kappa)e^{ \kappa\bmu^\top\bomega t} e^{\kappa \sqrt{(1-t^2)(1-(\bmu^\top\bomega)^2)} \Tilde{\bmu}^\top \bu }.
$$
We know that $d\sigma_{\mathbb{S}^{q}}(\bx) = (1-t^2)^{\frac{q-2}{2}} d\sigma_{\mathbb{S}^{q-1}}(\bu)$, where $d\sigma_{\mathbb{S}^{q}}$ is the Lebesgue measure on the $q$-dimensional unit sphere. So
\begin{align}\label{integral_form_density}
  f_{T}(t) = c_{q}^\mathrm{vMF}(\kappa)e^{ \kappa\bmu^\top\bomega t} (1-t^2)^{\frac{q-2}{2}}
\int_{\bu \in \mathbb{S}^{q-1}} e^{\kappa \sqrt{(1-t^2)(1-(\bmu^\top\bomega)^2)} \Tilde{\bmu}^\top \bu} \, \rd \sigma_{\mathbb{S}^{q-1}}(\bu).
\end{align}
The integrand in \eqref{integral_form_density} corresponds to the density of a $\mathrm{vMF}\big(\Tilde{\bmu}, \kappa\sqrt{(1-t^2)(1-(\bmu^\top\bomega)^2)}\big)$ (up to the normalization constant). We conclude that \eqref{eq:density_t_vmf} holds. If $\bomega=\pm\bmu$, then the tangent-normal decomposition gives $\bmu^\top \bx=t\,\bmu^\top\omega$. Integrating over $\mathbb{S}^{q-1}$ yields \eqref{eq:density_t_vmf}. Hence, \eqref{eq:density_t_vmf} holds for all $\bomega\in\mathbb{S}^q$.
Obtaining $F_T$, the CDF of $T$, from \eqref{eq:density_t_vmf} is not direct.

For the chordal distance, if $0\le t\le 2$, then
$$
F_{\bomega}^{\mu}(t) = \Prob{\sqrt{2(1-\bX^\top \bomega)}\le t} = 1 - F_T\left(1-\frac{t^2}{2}\right),
$$
whereas $F_{\bomega}^{\mu}(t)=1$ for $t>2$. Under the geodesic distance, if $0\le t\le \pi$, then
$$
F_{\bomega}^{\mu}(t) = \Prob{\cos^{-1}(\bX^\top \bomega)\le t} = 1 - F_T\left(\cos(t)\right),
$$
whereas $F_{\bomega}^{\mu}(t)=1$ for $t>\pi$.

\item Example~\ref{ex:dp_hvmf}.
Following \cite{Jensen1981}, for $\bx \in \mathbb{H}^q$, define the hyperbolic-spherical coordinates:
\begin{align*}
\begin{cases}
    x_1=\cosh(x_r)\\
    x_2=\sinh(x_r)
\end{cases}
\end{align*}
when $q=1$, with $x_r\ge0$, and
$$
\begin{cases}
\begin{aligned}
    x_1      =&\; \cosh(x_r) \\
    x_2      =&\; \sinh(x_r)\cos(\theta_1) =\sinh(x_r)(x_{s})_1\\
    x_3      =&\; \sinh(x_r)\sin(\theta_1)\cos(\theta_2) =\sinh(x_r)(x_{s})_2\\
    &\;\;\vdots \\
    x_q      =&\; \sinh(x_r)\sin(\theta_1)\cdots\sin(\theta_{q-2})\cos(\theta_{q-1})=\sinh(x_r)(x_{s})_{q-1} \\
    x_{q+1}  =&\; \sinh(x_r)\sin(\theta_1)\cdots\sin(\theta_{q-2})\sin(\theta_{q-1})=\sinh(x_r)(x_{s})_{q}
\end{aligned}
\end{cases}
$$
when $q>1$, with $(\theta_1,\ldots,\theta_{q-1})^\top\in (0,\pi)^{q-2}\times[0,2\pi)$ being the angular hyperspherical coordinates of $\bx_s\in\mathbb{S}^{q-1}$.
Denoting by $\mu_q$ the Lebesgue measure on $\mathbb{H}^q$ and $\sigma_{q-1}$ the Lebesgue measure on $\mathbb{S}^{q-1}$, it holds that $\mu_q(\rd \bx)= \sinh^{q-1}(x_r)\,\rd x_r \,\sigma_{q-1}(\rd \bx_s)$.

Define the group of hyperbolic transformations on $\R^{q+1}$ as:
$$
\mathrm{SH}(q) \defin \{
\bA \in \R^{(q+1) \times (q+1)}: \bA^\top \tilde{\bI}_{q+1} \bA = \tilde{\bI}_{q+1}, \, \det(\bA)=1, \, A_{11}>0
\},
$$
where
$$
\tilde{\bI}_{q+1}
=
\begin{bmatrix}
-1 & \boldsymbol{0}_q \\
\boldsymbol{0}_q & \bI_q \\
\end{bmatrix}.
$$
The group of hyperbolic transformations acts transitively on $\mathbb{H}^q$ \citep[see Lemma 2(ii) of][]{Jensen1981}, i.e., for any pair of mean vectors $\bx, \by \in \mathbb{H}^q$, there exists an action $\bA\in \mathrm{SH}(q)$ such that $\bx = \bA \by$.
The Minkowski pseudo-inner product is invariant under hyperbolic transformations, i.e., $(\bx,\by) = (\bA\bx,\bA\by)$ for every $\bx, \by \in \mathbb{H}^q$ and every $\bA \in \mathrm{SH}(q)$.
The invariance follows from the identity $(\bx, \by) = \bx^\top \tilde{\bI}_{q+1} \by$ for any $\bx, \by \in \R^{q+1}$.
Take $A \in \mathrm{SH}(q)$ such that $\bA\bomega = \be_1 = (1,0,\ldots,0)^\top$.

We are interested in the distance profile
\begin{align*}
F_{\bomega}^{\mu}(t) &= \Prob{d_{\mathbb{H}^q}(\bX, \bomega) \le t} = \Prob{-(\bX, \bomega) \le \cosh(t)} = \Prob{-(\bA \bX, \bA\bomega) \le \cosh(t)} \\
 & = \Prob{-(\bA\bX, \be_1) \le \cosh(t)} = \Prob{(\bA\bX)_1 \le \cosh(t)}
\end{align*}
for $t\ge0$. Therefore,
$$
F_{\bomega}^{\mu}(t) = \int_{\mathbb{H}^q} 1_{\{(\bA\bx)_1 \le \cosh(t)\}} f_\mathrm{HvMF}(\bx; \bmu, \kappa) \, \mu_q(\rd \bx).
$$
%where $\mu_q$ is the Riemannian volume measure on $\mathbb{H}^q$ \citep[see Section 2.C of][]{Jensen1981}.
Write $\bz \defin \bA\bx = (\cosh(z_r), \sinh(z_r) \bz_s)^\top$, where $(z_r, \bz_s)^\top \in \R_{\ge0} \times \mathbb{S}^{q-1}$. We thus have that
\begin{align}
F_{\bomega}^{\mu}(t) &= \int_{\mathbb{H}^q} 1_{\{z_r \le t\}} f_\mathrm{HvMF}(\bA^{-1}\bz; \bmu, \kappa) \, \mu_q(\rd (\bA^{-1}\bz)) \label{eq:example_Hq_change_variables} \\
&= \int_{\mathbb{H}^q} 1_{\{z_r \le t\}} f_\mathrm{HvMF}(\bz; \bA\bmu, \kappa) \, \mu_q(\rd (\bA^{-1}\bz)) \label{eq:example_Hq_densities_equality}  \\
&= \int_{\mathbb{H}^q} 1_{\{z_r \le t\}} f_\mathrm{HvMF}(\bz; \bA\bmu, \kappa) \, \mu_q(\rd \bz)  \label{eq:example_Hq_measure_invariance} \\
&= \int_{0}^t \int_{\mathbb{S}^{q-1}} f_\mathrm{HvMF}\left((\cosh(z_r), \sinh(z_r) \bz_s); \bA\bmu, \kappa\right) \sinh^{q-1}(z_r) \sigma_{q-1}(\rd \bz_s)  \rd z_r \notag
\end{align}
Equality \eqref{eq:example_Hq_change_variables} follows from the substitution $\bx = \bA^{-1}\bz$.
Observe that the domain of integration remains invariant since the mapping $\bx \mapsto \bA^{-1}\bx$ is a bijection from $\mathbb{H}^q$ to itself.
Surjectivity holds because for any $\bz \in \mathbb{H}^q$, $\bz = \bA^{-1}(\bA\bx)$, and $\bA\bx \in \mathbb{H}^q$ because the unit hyperboloid is invariant under hyperbolic transformations \citep[Lemma $1$(ii) of][]{Jensen1981}.
The equality \eqref{eq:example_Hq_densities_equality}, i.e., $f_\mathrm{HvMF}(\bA^{-1}\bz; \bmu, \kappa) = f_\mathrm{HvMF}(\bz; \bA\bmu, \kappa)$, is verified because $(\bA^{-1}\bz, \bmu) = (\bz, \bA\bmu)$, and the normalization constant of the HvMF distribution, $c_q^{\mathrm{HvMF}}(\kappa)$, does not depend on the mean vector \citep[see the proof of Proposition 4.1 in][]{Serrano2025}.
Equality \eqref{eq:example_Hq_measure_invariance} holds since $\mu_q$ is invariant under hyperbolic transformations \citep[see Section $2$C of][]{Jensen1981}. Note that $\det(\bA^{-1}) = 1$.

Assume first that $\rho>0$. Writing $\bA\bmu = (\cosh(\rho), \sinh(\rho) \bxi)^\top$, with $(\rho, \bxi) \in \R_{\ge0} \times \mathbb{S}^{q-1}$, we have that
$$
(\bomega,\bmu)=(\bA\bomega,\bA\bmu) = (\be_1,\bA\bmu) = -\cosh(\rho).
$$
Hence,
\begin{equation}
\begin{aligned}
(\bA\bmu, \bz)
&= - \cosh(\rho) \cosh(z_r)
 + \sinh(\rho) \sinh(z_r) \langle \bz_s, \bxi \rangle
\\
&= (\bmu, \bomega) \cosh(z_r)
 + \sinh(\rho) \sinh(z_r) \langle \bz_s, \bxi \rangle.
\end{aligned}
\end{equation}
Therefore,
\begin{align*}
F_{\bomega}^{\mu}(t) = c_q^{\mathrm{HvMF}}(\kappa) \int_{0}^t e^{\kappa(\bmu, \bomega)\cosh(z_r)} \sinh^{q-1}(z_r) \int_{\mathbb{S}^{q-1}} e^{\kappa\sinh(z_r) \sinh(\rho) \langle \bxi, \bz_s \rangle} \sigma_{q-1}(\rd \bz_s)  \rd z_r
\end{align*}
Thus, $F_{\bomega}^{\mu}(t)$ is the CDF of the real random variable $Z_r \defin d_{\mathbb{H}^q}(\bX,\bomega)= \arccosh((\bA\bX)_1)$, with density given by
$$
\begin{aligned}
f_{Z_r}(z_r) &= c_{q}^\mathrm{HvMF}(\kappa) (\sinh(z_r))^{q-1} e^{\kappa(\bmu, \bomega) \cosh(z_r)} \\
&\quad \times \int_{\mathbb{S}^{q-1}} e^{\kappa \sinh(z_r) \sinh(\rho) \langle \bxi, \bz_s \rangle} \, \rd \sigma_{\mathbb{S}^{q-1}}(\bz_s).
\end{aligned}$$
This integral is proportional to the normalization constant of a von Mises--Fisher distribution on $\mathbb{S}^{q-1}$:
$$
\int_{\mathbb{S}^{q-1}} e^{\kappa \sinh(z_r) \sinh(\rho) \langle\bxi, \bz_s\rangle} \, \rd \sigma_{\mathbb{S}^{q-1}}(\bz_s) = \frac{1}{c_{q-1}^\mathrm{vMF}(\kappa \sinh(z_r) \sinh(\rho))}.
$$
Thus, the density of $Z_r$ is
\begin{align*}
f_{Z_r}(z_r) = \frac{c_{q}^\mathrm{HvMF}(\kappa)}{c_{q-1}^\mathrm{vMF}(\kappa\sinh(z_r)\sinh(\rho))} (\sinh(z_r))^{q-1} e^{\kappa(\bmu, \bomega) \cosh(z_r)}, \quad z_r \ge 0,
\end{align*}
for $\rho>0$.

If $\rho=0$, then $\bomega=\bmu$, and hence $\bA\bmu=\be_1$. Therefore, $(\bA\bmu,\bz)=-\cosh(z_r)$. Integrating over $\mathbb{S}^{q-1}$ yields \eqref{eq:density_zr_hvmf}. Hence, \eqref{eq:density_zr_hvmf} holds for all $\rho\ge0$.
The distance profile is then given by
$$
F_{\bomega}^{\mu}(t) = \Prob{Z_r \le t}.
$$

\begin{comment}
\item \textcolor{blue}{Wishart distribution on the space of symmetric positive definite matrices.}
\textcolor{blue}{In this setting, the distance profile of $\bS$ at $\bW \in \mathcal{S}_q^+$ is given by}
{\color{blue}
$$
F_{\bW}^{\mu}(t) = \Prob{d_{\mathrm{AI}}(\bW, \bS) \le t} = \mathsf{P}\Big\{\sum_{i=1}^q \log(\lambda_i)^2 \le t ^2\Big\}
$$
}
\textcolor{blue}{where $\{\lambda_i\}_{i=1}^q$ are the eigenvalues of $\bW^{-1}\bS$. \cite{Zanella2009} gives the joint density of the ordered eigenvalues of a Wishart-distributed random matrix, which characterizes the density of $\sum_{i=1}^q \log(\lambda_i)^2$.}
\textcolor{blue}{However, obtaining an explicit form for $F_{\bW}^{\mu}(t)$ is challenging.}
\end{comment}
\begin{comment}
\item Example~\ref{ex:dp_dirichlet}.
Using the Gamma representation of the Dirichlet distribution, we have $\log(X_i) = \log(G_i)- \log(S)$, where $G_i$ are independent Gamma variables, $G_i \sim \Gamma(\alpha_i, 1)$, and $S = \sum_{i=1}^{q+1} G_i$. Hence,
$$
\mathrm{clr}(\bX) = \big(\log(G_1) - \frac{1}{q+1}\sum_{i=1}^{q+1} \log(G_i), \ldots, \log(G_{q+1}) - \frac{1}{q+1}\sum_{i=1}^{q+1} \log(G_i)\big).
$$
Let $\bL = (\log(G_1), \dots, \log(G_{q+1}))^\top$
and denote by $\bH = \bI_{q+1} - \frac{1}{q+1} \boldsymbol{1}_{q+1}\boldsymbol{1}_{q+1}^\top$
the centering matrix projecting onto $H$, with $\boldsymbol{1}_{q+1}$ a $(q+1)$-dimensional vector of ones. Then}
$$
\mathrm{clr}(\bX) = \bH \bL.
$$
\textcolor{blue}{\textit{Shall I include the expected value and the variance of $\mathrm{clr}(\bX)$?}}.

The components $L_k = \log(G_k)$ are independent and have characteristic functions
$$
\varphi_{L_k}(t) = \E{e^{i t L_k}}
= \frac{\Gamma(\alpha_k + i t)}{\Gamma(\alpha_k)}, \quad t \in \R.}
$$
Therefore, for any $\bt \in \R^{q+1}$,
$$
\varphi_{\mathrm{clr}(\bX)}(\bt)
= \prod_{k=1}^{q+1} \frac{\Gamma\big(\alpha_k + i (\bH \bt)_k\big)}{\Gamma(\alpha_k)}.}
$$
For the Euclidean distance, the corresponding CDF can be expressed as the push-forward of the Dirichlet density:
$$
F_{\bomega}^{\mu}(t)
= \int_{\{\bx \in \Delta^{q} : \|\bx - \bomega\|_2 \le t\}}
    \frac{1}{B(\boldsymbol{\alpha})} \prod_{i=1}^{q+1} x_i^{\alpha_i - 1} \, \rd x.
$$
For the Euclidean distance, if $q=1$, then $\bX = (X_1, 1 - X_1)$ with $X_1 \sim \mathrm{Beta}(\alpha_1, \alpha_2)$.
Since $\|\bX - \bomega\|_2 = \sqrt{2}  |X_1 - \omega_1|$, the distance profile is given by
$$
F_{\bomega}^{\mu}(t)
= \Probbig{ |X_1 - \omega_1| \le t / \sqrt{2} }
= F_{X_1}\left(\omega_1 + \frac{t}{\sqrt{2}}\right)
- F_{X_1}\left(\omega_1 - \frac{t}{\sqrt{2}}\right),
$$
where $F_{X_1}$ denotes the CDF of $X_1$.

For the Aitchison distance, if $q=1$, we have $X_1 \sim \mathrm{Beta}(\alpha_1, \alpha_2)$ and
$$
\mathrm{clr}(\bX)
= \left( \frac{1}{2} \log \left(\frac{X_1}{1-X_1}\right), -\frac{1}{2} \log\left(\frac{X_1}{1-X_1}\right) \right).
$$
Hence,
$$
d_A(\bX, \bomega)
= \frac{\sqrt{2}}{2} \left| \log \left( \frac{X_1/(1-X_1)}{\omega_1/(1-\omega_1)} \right) \right|.
$$
The corresponding CDF is
$$
F_{\bomega}^{\mu}(t)
= \Prob{
\left| \log\left(\frac{X_1/(1-X_1)}{\omega_1/(1-\omega_1)}\right) \right| \le \sqrt{2} t
}
= F_{X_1}\big(\mathrm{logit}^{-1}\big(\mathrm{logit}(\omega_1) + \sqrt{2}t\big)\big) - F_{X_1}\big(\mathrm{logit}^{-1}\big(\mathrm{logit}(\omega_1) - \sqrt{2}t\big)\big),
$$
where $\mathrm{logit}(x) = \log\big(\frac{x}{1-x}\big)$, and $\mathrm{logit}^{-1}(x) = \frac{e^x}{1 + e^x}$ is the logistic function.
\end{comment}

\item Example~\ref{ex:dp_logistic_gaussian}.
With the notation of Example~\ref{ex:dp_logistic_gaussian}, let
$$
\bZ \defin \bH \bY - \mathrm{clr}(\bomega), \qquad \bW \defin \bU^\top \bZ, \qquad \bv \defin \bU^\top \bm.
$$
Since $\mathrm{clr}(\bX) = \mathrm{clr}(\bsigma(\bY)) = \bH\bY$, one has $\mathrm{clr}(\bX)\sim\mathcal{N}(\bH\bmu,\bH\bSigma\bH^\top) = \mathcal{N}(\bH\bmu,\boldsymbol{C})$. Hence, $\bZ\sim \mathcal{N}(\bm,\boldsymbol{C})$ and $\bW\sim \mathcal{N}(\bv,\bA)$. Since $\boldsymbol{\Sigma}$ is positive definite, $\operatorname{rank}(\boldsymbol{C})=q$, so $\bA=\mathrm{diag}(\alpha_1,\ldots,\alpha_q,0)$. Choosing $\bu_{q+1}$ proportional to $\boldsymbol{1}_{q+1}$, one has $W_{q+1}=0$ almost surely because $\bZ\in H$ almost surely. Since $\bU$ is orthogonal,
$$
d_A(\bX, \bomega)^2 = \bZ^\top \bZ = \bW^\top \bW = \sum_{i=1}^{q+1} W_i^2 = \sum_{i=1}^{q} W_i^2,
$$
where
$W_i \sim \mathcal{N}(v_i,\alpha_i),$
for $i=1,\dots,q$. Hence, we have that $W_i = v_i + \sqrt{\alpha_i}\,\varepsilon_i$ where $\varepsilon_i$ are independent and identically distributed, $\varepsilon_i \sim \mathcal{N}(0,1)$. Hence,
$W_i^2 = \alpha_i\left(\varepsilon_i + \frac{v_i}{\sqrt{\alpha_i}}\right)^2.$
By definition, $\left(\varepsilon_i + \frac{v_i}{\sqrt{\alpha_i}}\right)^2$ is a noncentral $\chi^2_1$ random variable with one degree of freedom and noncentrality parameter
$\delta_i = \frac{v_i^2}{\alpha_i} = \frac{(\bu_i^\top\bm)^2}{\alpha_i}$. Therefore, setting $Q_i \defin W_i^2/\alpha_i$, one has $Q_i \sim \chi_1^2(\delta_i)$. This completes the characterization of the distance profile of the logistic Gaussian distribution.

Finally, we prove the equality $\|\mathrm{clr}(\bx) - \mathrm{clr}(\bomega)\|_2 = \|\mathrm{ilr}(\bx) - \mathrm{ilr}(\bomega)\|_2$, shown in Remark \ref{rmk:ilr_clr_equivalence}.
\citet[][Equation (28)]{Egozcue2003} show that $\mathrm{ilr}(\bx) = \mathbf{B}^\top \mathrm{clr}(\bx)$, where $\mathbf{B}$ is the $(q+1)\times q$ matrix from Equation (17) of \citet{Egozcue2003}. The columns $\mathbf{b}_i \in \R^{q+1}$ of $\mathbf{B}$ form a basis of $H = \{\bz \in \R^{q+1} : \sum_{i=1}^{q+1} z_i = 0\}$, and they are orthonormal with respect to the Euclidean inner product in $\R^{q+1}$. Hence,
$$
\begin{aligned}
\| \mathrm{ilr}(\bx) - \mathrm{ilr}(\bomega) \|_2^2
&= \| \mathbf{B}^\top \mathrm{clr}(\bx) - \mathbf{B}^\top \mathrm{clr}(\bomega) \|_2^2 \\
&= \left(\mathrm{clr}(\bx) - \mathrm{clr}(\bomega)\right)^\top \mathbf{B} \mathbf{B}^\top \left(\mathrm{clr}(\bx) - \mathrm{clr}(\bomega)\right) \\
&= \| \mathrm{clr}(\bx) - \mathrm{clr}(\bomega) \|_2^2
= d_A(\bx, \bomega)^2.
\end{aligned}
$$
The third equality follows because $\mathbf{B} \mathbf{B}^\top$ is the projection matrix onto $H$ and $(\mathrm{clr}(\bx) - \mathrm{clr}(\bomega)) \in H$.
\end{itemize}
\end{proof}

\begin{comment}
%---------------------------%
\section{Proof of the continuity of the supremum}\label{ap:continuitysup}
%---------------------------%
\textcolor{blue}{This is a sanity check from my baby-steps.}

We now show that $\sup_{\omega,t}(\cdot)$ is $1$-Lipschitz. To simplify notation, let $S = \Omega \times \R_{\geq 0}$ and  
$$
A\defin\|f\|_{\infty}=\sup _{s \in S}|f(s)|, \quad B\defin\|g\|_{\infty}=\sup _{s \in S}|g(s)|.
$$ 
First, observe that 
$$
\big||f(s)|-|g(s)|\big| \leq|f(s)-g(s)|, \quad \text{for every } s \in S.
$$
Taking the supremum over $s$ yields
$$
\sup _{s \in S}| | f(s)|-|g(s)|| \leq \sup _{s \in S}|f(s)-g(s)|=\|f-g\|_{\infty}.
$$

We now prove that $|A-B|  \le \sup _{s \in S}| | f(s)|-|g(s)||$. For every $\varepsilon>0$, there exists $s_{\varepsilon} \in S$ such that $\left|f\left(s_{\varepsilon}\right)\right|>A-\varepsilon$.
Hence,
$$
A-B < \left|f\left(s_{\varepsilon}\right)\right|-\left|g\left(s_{\varepsilon}\right)\right|+\varepsilon \leq \sup _{s \in S}(|f(s)|-|g(s)|)+\varepsilon.
$$
Since this holds for every $\varepsilon>0$ we conclude
$$
A-B \leq \sup _{s \in S}(|f(s)|-|g(s)|).
$$
A similar reasoning ensures
$$
|A-B| \leq \sup _{s \in S}| | f(s)|-|g(s)|| .
$$
Finally, we can conclude that
$$
\Big| \sup _{s \in S} f - \sup _{s \in S} g \Big|=|A-B| \leq \sup _{s \in S} \big| | f(s)|-| g(s)| \big| \leq\| f-g \|_{\infty},
$$
as desired.
%-----------------------%
\end{comment}


%-----------------------%
\section{Proofs of Section \ref{sec:null_asymptotics}}\label{ap:sec_proofs_null_asymptotics}
%-----------------------%

%------------------------%
\subsection{Proof of Theorem \ref{thm:convergence_G_n_composite}}\label{ap:sec_proofs_null_asymptotics_thm_convergence_G_n_composite}
%------------------------%
\begin{proof}[Proof of Theorem \ref{thm:convergence_G_n_composite}]
First, observe that
$$
\G_{n,\mu}^{\mu_{\hat{\btheta}}}(\omega, t) =
%\sqrt{n}(F_{n,\omega}^{\mu}(t) - F_{\omega}^{\mu_{\hat{\btheta}}}(t)) = 
\G_{n,\mu}^{\mu_{\btheta_0}}(\omega, t) 
- \sqrt{n} \left( 
F_{\omega}^{\mu_{\hat{\btheta}}}(t) - F_{\omega}^{\mu_{\btheta_0}}(t)
\right).
$$
We know that $\G_{n,\mu}^{\mu_{\btheta_0}}(\omega, t)$ converges weakly to a Gaussian process by Theorem 5.1 of \cite{DubeyChenMuller2024}:
\begin{align}\label{eq:thm_5.1_chendubeymuller_theta_0}
\G_{n,\mu}^{\mu_{\btheta_0}}(\omega, t)  \convl \G_{\btheta_0}(\omega, t) \text { in } \ell^{\infty}(\mathcal{Y}),
\end{align}
where $\G_{\btheta_0}(\omega, t)$ is a zero-mean Gaussian process with covariance given by 
\begin{align}\label{eq:cov_G_theta_0}
\Cov{y_{(\omega_1, t_1)}}{y_{(\omega_2, t_2)}} = \Prob{d(\omega_1, X) \le t_1, d(\omega_2, X) \le t_2} - F_{\omega_1}^{\mu_{\btheta_0}}(t_1)F_{\omega_2}^{\mu_{\btheta_0}}(t_2),
\end{align}
for $\omega_1, \omega_2 \in \Omega$ and $t_1, t_2 \in \R_{\ge0}$.

Now, focus on $F_{\omega}^{\mu_{\hat{\btheta}}}(t) - F_{\omega}^{\mu_{\btheta_0}}(t)$. By differentiability of $\btheta \mapsto F_{\omega}^{\btheta}(t)$ at $\btheta_0$, it holds that
\begin{align}\label{eq:taylor_F_omega_t}
F_{\omega}^{\mu_{\hat{\btheta}}}(t) - F_{\omega}^{\mu_{\btheta_0}}(t) = 
\dot{F}_{\omega}^{\btheta_0}(t)^\top \left(\hat{\btheta} - \btheta_0\right) + o_P\left(\|\hat{\btheta} - \btheta_0 \| \right).   
\end{align}
%Note that for Equation \eqref{eq:bahadur} to hold, we are implicitly assuming that $\mathsf{E}(\|\bpsi_{{\btheta}_0}(X)\|^2) < \infty $. 
From \eqref{eq:taylor_F_omega_t} and Bahadur's representation \eqref{eq:bahadur}, one obtains that
\begin{align}\label{eq:taylor_bahadur}
\sqrt{n} \left(F_{\omega}^{\mu_{\hat{\btheta}}}(t) - F_{\omega}^{\mu_{\btheta_0}}(t)\right) =
-\dot{F}_{\omega}^{\btheta_0}(t)^\top \bV_{\btheta_0}^{-1} \frac{1}{\sqrt{n}} \sum_{i=1}^n \bpsi_{\btheta_0}(X_i)+o_{P}(1) + \sqrt{n} o_{P}\left(\|\hat{\btheta} - \btheta_0 \| \right).
\end{align}
Observe that $\dot{F}_{\omega}^{\btheta_0}(t)^\top o_P(1) = o_P(1)$ by \ref{d}.
By Bahadur's representation \eqref{eq:bahadur}, $\sqrt{n}\|\hat{\btheta} - \btheta_0\| = O_P(1)$, so $\sqrt{n} o_P\big(\|\hat{\btheta} - \btheta_0 \|\big) = o_P(1)$. 
Hence, by \ref{d}, the remainders in \eqref{eq:taylor_bahadur} are $o_P(1)$ uniformly over $(\omega,t)\in\Omega\times\R_{\ge0}$.
%Here we are implicitly using boundedness of the derivative of F, because it goes inside the o_P term.
\begin{comment}
Since
$$
\frac{1}{\sqrt{n}}\sum_{i=1}^n \bpsi_{\btheta_0}(X_i) \convl Z \sim \mathcal{N}(0,\boldsymbol{\Sigma}), \quad \text{where } \boldsymbol{\Sigma} \defin\E{\bpsi_{\btheta_0}(X) \bpsi_{\btheta_0}(X)^{\top}},
$$
we conclude that
\begin{align}\label{eq:bias_convergence}
    \sqrt{n} \left(F_{\omega}^{\mu_{\hat{\btheta}}}(t) - F_{\omega}^{\mu_{\btheta_0}}(t)\right)
\convl \dot{F}_{\omega}^{\btheta_0}(t)^\top Z.
\end{align}

From \eqref{eq:thm_5.1_chendubeymuller_theta_0} and \eqref{eq:bias_convergence}, one cannot conclude directly that $T_n$ converges weakly to the sum of marginal limits of the two processes. Idea: the sum of two Donsker classes is Donsker. See Section 2.10.3 of \cite{VanderVaartWellner2023} for results about transformations of Donsker classes.
\end{comment}
Therefore, \eqref{eq:thm_5.1_chendubeymuller_theta_0} and \eqref{eq:taylor_bahadur} yield that
$$
G_{n,\mu}^{\mu_{\hat{\btheta}}}(\omega, t)
=
\frac{1}{\sqrt{n}}\sum_{i=1}^n
\Big( y_{\omega,t}(X_i) - F_{\omega}^{\mu_{\btheta_0}}(t) + g_{\omega,t}(X_i) \Big)
+ o_{P}(1)
=
\frac{1}{\sqrt{n}} \sum_{i=1}^n f_{\omega, t}(X_i) + o_P(1),
$$
where 
\begin{align*}
g_{\omega,t}(x) \defin \dot{F}_{\omega}^{\btheta_0}(t)^\top \bV_{\btheta_0}^{-1} \bpsi_{\btheta_0}(x), 
\end{align*}
and
\begin{align*}
f_{\omega,t}(x) \defin y_{\omega, t}(x) - F_{\omega}^{\mu_{\btheta_0}}(t) + g_{\omega, t}(x).
\end{align*}
Observe that $\mathsf{E}(f_{\omega, t}(X)) = 0$, so $G_{n,\mu}^{\mu_{\hat{\btheta}}}$
is, up to an $o_P(1)$ remainder, the empirical process indexed by the class 
\begin{align*}
\cF \defin \{f_{\omega, t}: \ \omega \in \Omega, \ t \ge0\}.
\end{align*}
Thus, proving that $\cF$ is Donsker concludes the proof.
To prove the Donsker property, write $\cF \subset \cY + \mathcal{G}$, where $\cY \defin \{y_{\omega, t} - F_{\omega}^{\mu_{\btheta_0}}(t): \ \omega \in \Omega, \ t \ge 0\}$ and $\mathcal{G} \defin \{g_{\omega, t}: \ \omega\in \Omega, \ t \ge 0\}$. 
% Theorem 8 in Ingrid's notes shows that, 
If $\cY$ and $\mathcal{G}$ are Donsker and $\sup_{h\in\cY \cup \mathcal{G}} {\mathsf{E}|h(X)|} < \infty$, then  $\cY + \mathcal{G}$ is also Donsker \citep[see, e.g., Example 2.10.9 in][]{vandervaart1998}, and since $\cF$ is a subset of a Donsker class, it is Donsker.

The class $\cY$ is Donsker by \eqref{eq:thm_5.1_chendubeymuller}. Corollary 16.1 in \cite{DasGupta2011} shows that a sufficient condition for $\mathcal{G}$ to be Donsker is that $\mathcal{G}$ is finite-dimensional, and
its envelope function is in $L_2(\mathsf{P})$. The fact that $\mathcal{G}$ is finite-dimensional can be seen by writing
$$
g_{\omega,t}(x) = \dot{F}_{\omega}^{\btheta_0}(t)^\top \bV_{\btheta_0}^{-1}\bpsi_{\btheta_0}(x) = \sum_{i=1}^p \alpha_i(\omega, t) \psi_{\btheta_0, i}(x),
$$
where $\boldsymbol{\alpha}(\omega, t)^\top \defin \dot{F}_{\omega}^{\btheta_0}(t)\bV_{\btheta_0}^{-1} \in \R^p$ (coefficients), and $\psi_{\btheta_0,i}$ represents the $i$-th component of $\bpsi_{\btheta_0}$, $i=1, \ldots,p$. Therefore, one has $\mathcal{G} \subset \mathrm{span}\{\bpsi_{\btheta_0, 1}, \ldots, \bpsi_{\btheta_0, p}\}$, a linear subspace of dimension at most $p$. 
By Assumption \ref{d}, it holds that $\sup_{\omega, t} \|\dot{F}_{\omega}^{\btheta_0}(t) \| < \infty $, hence Assumption \ref{cp2} guarantees that the envelope function of $\mathcal{G}$ is in $L_2(\mathsf{P})$. 
Therefore, $\mathcal{G}$ is Donsker.

It is only left to show that $\sup_{h\in\cY \cup \mathcal{G}} {\mathsf{E}(h(X))} < \infty$. 
This is trivial for class $\cY$. For $\mathcal{G}$, it follows from Assumption \ref{cp3}, since $\mathsf{E}(g_{\omega,t}(X)) =0$ for every $(\omega, t)\in \Omega \times \R_{\ge 0}$.
Therefore, $\cF$ is Donsker. 
This implies that $G_{n,\mu}^{\mu_{\hat{\btheta}}}$ converges weakly to a Gaussian process $\tilde{\G}_0$, with mean zero and covariance $\tilde{\mathcal{C}}_{(\omega_1,t_1), (\omega_2, t_2)} \defin %\Cov{\tilde{\G}_0(f_{\omega_1,t_1})}{\tilde{\G}_0(f_{\omega_2,t_2})} =
\mathsf{E}(f_{\omega_1, t_1}(X) f_{\omega_2, t_2}(X))$, which is given in \eqref{eq:cov_comp}.
%($\mathsf{E}\left(\bpsi_{\theta_0}(X)\right)=\boldsymbol{0}$) 
%$\mathsf{E}(\|\bpsi_{{\btheta}_0}(X)\|^2) < \infty $
\end{proof}
%-----------------------%

%------------------------%
\subsection{Proof of Theorem \ref{thm:test_statistics_simp}}\label{ap:sec_proofs_null_asymptotics_thm_test_statistics_simp}
%------------------------%

This section contains the proof of Theorem \ref{thm:test_statistics_simp}, together with auxiliary results used in the proof.
In what follows, we denote by $d_{\Omega \times \R_{\ge 0}}$ the product distance on $(\Omega, d) \times (\R_{\ge 0}, |\cdot|)$, that is,
\begin{align*}
d_{\Omega \times \R_{\ge 0}}\big((\omega,t), (\omega^\prime, t^\prime)\big) \defin d(\omega, \omega^\prime) + |t - t^\prime|.
\end{align*}


\begin{comment}
%---------------------------%
\begin{remark}
%---------------------------%
Although our application of the Helly--Bray theorem requires only that the map $t \mapsto G_{\btheta_0}(\omega,t)$ be continuous for each fixed $\omega$, we have proved continuity jointly in
$(\omega,t)$. This is justified because the proof reuses directly the entropy condition \eqref{eq:entropy_dubeychenmuller} of
\cite{DubeyChenMuller2024} for $\cF$, and the set $\Omega \times \R_{\ge 0}$ can be trivially identified with $\cF$. Using \eqref{eq:entropy_dubeychenmuller}, the additional technical work to prove joint continuity is straightforward.

An alternative approach is to fix $\omega\in \Omega$ and prove continuity in $t\in \R_{\ge0}$ by considering the subclass
$\cF_{\omega} = \{ f_{\omega,t} : t \ge 0 \}$, since \eqref{eq:entropy_dubeychenmuller} also holds for $\cF_\omega$. 
However, the structure of the proof remains the same, and the same assumptions are required. Thus, joint continuity is stronger and comes at no significant additional cost.
\end{remark}
\end{comment}


\begin{comment}
%----------------------------%
\begin{lemma}\label{ap:lem_DCT_G}
%----------------------------%
    If $\int_{\R_{\ge 0}} \G_{0}(\omega, t)^2 \, \rd (F_{n,\omega} - F_{n,\omega}^{\btheta_0})(t) \to 0$ as $n \to \infty$, then also $\int_{\Omega \times\R_{\ge 0}} \G_{0}(\omega,t)^2 \, \rd (F_{n,\omega} - F_{n,\omega}^{\btheta_0})(t) \, \allowbreak \rd \mu(\omega) \to 0$ as $n \to \infty$.
\end{lemma}
\begin{proof}
The result follows by an application of the DCT. Let 
$$
I_n(\omega) \defin \int_{\R_{\ge 0}} \G_{0}(\omega, t)^2 \, \rd (F_{n,\omega} - F_{n,\omega}^{\btheta_0})(t).
$$
By assumption, $I_n(\omega) \to 0$ as $n \to \infty$ for every fixed $\omega \in \Omega$. To obtain a bound for $|I_n(\omega)|$, we use a result from Page 88 in Folland's Real Analysis (Exercise 3, \url{https://apachepersonal.miun.se/~andrli/Bok.pdf}). For $\mu$ a signed measure and $f\in \rho_{1,y}$,
\begin{align}\label{eq:aux_folland}
\left|\int f(x) \,\rd \mu(x)\right|\leq \int |f| \, \rd |\mu|(x),
\end{align}
where $|\mu|$ is the total variation of $\mu$. The total variation of a measurable subset $E$ is, in particular,
$$
|\mu|(E)=\sup_{\text{Partitions $E_1,\ldots,E_m$ of $E$}}\lrb{\sum_{i=1}^m|\mu(E_j)|}
$$
(Exercise 8).

To apply this result, one needs to show that $\G_{0}(\omega, t)^2$ is in $\rho_{1,y}$, where $\mu = \hat{\mu}_{\omega} - \mu_{\omega}^{\btheta_0}$ is induced by $F_{n,\omega}(t)=P_n((-\infty,t])$ and $F_{\omega}^{\mu_0}(t)=\mu_0((-\infty,t])$. First, observe that $\mu$ is a finite signed measure, because $|\mu|(\R) \le 2$, where $|\mu|$ is the total variation of $\mu$. This follows easily from the calculation
\begin{align*}
    |\mu|(\R)&=\sup \lrb{\sum_{i=1}^m|\mu(E_j)|}=\sup \lrb{\sum_{i=1}^m|P_n(E_j)-\mu_\omega^{\theta_0}(E_j)|}\\
    &\leq \sup \lrb{\sum_{i=1}^m|P_n(E_j)+\mu_\omega^{\theta_0}(E_j)|}= \sup \lrb{\sum_{i=1}^mP_n(E_j)+\sum_{i=1}^m\mu_\omega^{\theta_0}(E_j)}\\
    &= \sup \lrb{1+1}=2.
\end{align*}

Since $\cF$ is a Donsker class, the sample paths of $\G_{0}(\omega, t)$ are bounded almost surely, and by the finiteness of $\mu$, $\G_{0}(\omega, t)^2 \in \rho_{1,y}$  for every $\omega \in \Omega$. Thus, one can apply \eqref{eq:aux_folland} to show that
\begin{align*}
|I_n(\omega)| \le \int_{\R_{\ge 0}} \G_{0}(\omega, t)^2 \, d \big|F_{n,\omega} - F_{n,\omega}^{\btheta_0}\big|(t) \le \sup_{t\in \R_{\ge0}} \G_{0}(\omega, t)^2 \big| F_{n,\omega} - F_{n,\omega}^{\btheta_0} \big| (\R) \le 2 \sup_{t \in \R_{\ge 0}} \G_{0}(\omega, t)^2 \defin M(\omega).
\end{align*}
Above, we have denoted by $|F_{n,\omega}-F_{\omega}^{\mu_0}|$ the total variation $|\mu|$ of $\mu$. Since $\cF$ is Donsker, $\sup_{\omega \in \Omega} M(\omega) < \infty$ almost surely, so $M(\omega) \in \ell^{\infty}(\Omega)$ is bounded almost surely, hence integrable with respect to $\mu$. Therefore, by the DCT, it holds that
$$
\int_{\Omega \times \R_{\ge 0}} \G_{0}(\omega,t)^2 \, \rd (F_{n,\omega} - F_{n,\omega}^{\btheta_0})(t) \, \rd \mu(\omega) \convas 0 \text{ as } n \to \infty,
$$
as desired.
\end{proof}
\end{comment}


Lemma \ref{ap:lem_bound_int_g_epsilon} is an auxiliary result in the proof of Lemma \ref{ap:lem_sup_r_n_r}. 

%------------------------%
\begin{lemma}\label{ap:lem_bound_int_g_epsilon}
%----------------------------%
    Fix $\varepsilon > 0$. Let $g: \Omega \times \R_{\ge0} \to \R$ be a bounded, uniformly continuous function on $\Omega \times \R_{\ge0}$. For each $\omega \in \Omega$, define the following regularization of $g$:
\begin{equation}\label{ap:eq_g_epsilon}
        g_\varepsilon(\omega, t) \defin \inf_{s \in [0, M]} \left\{ g(\omega, s) + \frac{2\|g\|_\infty}{\delta} |t - s| \right\}.
\end{equation}    
    Here, $M \defin \mathrm{diam}(\Omega) = \sup_{\omega, \omega^\prime \in \Omega} d(\omega, \omega^\prime)$, and $\delta(\varepsilon) > 0$ is such that for every $(\omega, t), (\omega^\prime, t^\prime) \in \Omega \times [0, M]$ with $d_{\Omega \times \R_{\ge 0}}\big((\omega, t), (\omega^\prime, t^\prime)\big) < \delta$, it holds that $|g(\omega, t) - g(\omega^\prime, t^\prime)| < \varepsilon$. 
    For every fixed $\omega\in\Omega$, the map $t\mapsto g_\varepsilon(\omega,t)$ is Lipschitz on $[0,M]$, with Lipschitz constant $\mathrm{Lip}(g_{\varepsilon}) = \frac{2\|g\|_\infty}{\delta}$, 
    and for every $\omega \in \Omega$ and $t \in [0, M]$, it holds a.s. that
    $$
    \left| \int_{0}^{M} g_\varepsilon(\omega, t) \, \rd (F_{n,\omega}^{\mu} - F_{\omega}^{\mu_0})(t) \right|
    \le
    \frac{2 M \|g\|_\infty}{\delta} \sup_{t \in [0, M]} |F_{n,\omega}^{\mu}(t) - F_{\omega}^{\mu_0}(t)|.
    $$
\end{lemma}

\begin{proof}
    Observe that $M < \infty$, since $\Omega$ is totally bounded, which is implied by \ref{b1}.
    Thus, $F_{\omega}([0, M]) = F_{n,\omega}^{\mu}([0, M])=1$ and for any measurable function $\phi$,
    $\int_{\R_{\ge0}} \phi(t) \, \rd F_{\omega}(t) = \int_{[0, M]} \phi(t) \, \rd F_{\omega}(t)$ (similarly for $F_{n,\omega}$).
    Since $g$ is uniformly continuous on $\Omega \times \R_{\ge0}$, it is also uniformly continuous on $\Omega\times [0, M]$.

    We first show that $g_\varepsilon(\omega, \cdot)$ is Lipschitz with Lipschitz constant
    $\mathrm{Lip}(g_{\varepsilon}) = \frac{2\|g\|_\infty}{\delta}$. Indeed, let $t,t^\prime\in[0,M]$ with $t<t^\prime$. For every $s\in[0,M]$, one has
    $$
    g(\omega,s)+\frac{2\|g\|_\infty}{\delta}|t-s|
    \le
    g(\omega,s)+\frac{2\|g\|_\infty}{\delta}|t^\prime-s|
    +\frac{2\|g\|_\infty}{\delta}|t-t^\prime|,
    $$
    and taking the infimum over $s$ yields
    $
    g_\varepsilon(\omega,t)
    \le
    g_\varepsilon(\omega,t^\prime)
    +\frac{2\|g\|_\infty}{\delta}|t-t^\prime|.
    $
    By symmetry, one also has $g_\varepsilon(\omega, t^\prime) \le g_\varepsilon(\omega, t) + \frac{2\|g\|_\infty}{\delta} |t - t^\prime|$, so that $|g_\varepsilon(\omega, t) - g_\varepsilon(\omega, t^\prime)| \le \frac{2\|g\|_\infty}{\delta} |t - t^\prime|$. 

    Fix $\omega\in\Omega$. Since $g_\varepsilon(\omega,\cdot)$ is Lipschitz on $[0,M]$, it is absolutely continuous and differentiable almost everywhere, with $|g_\varepsilon'(\omega,t)|\le\mathrm{Lip}(g_\varepsilon)$. An integration by parts argument gives
    \begin{align}\label{ap:eq_bound_d_BL_3}
    \left| \int_{0}^{M} g_\varepsilon(\omega,t)\, \rd(F_{n,\omega}^{\mu}-F_{\omega}^{\mu_0})(t) \right| \notag
    &=
    \left| -\int_{0}^{M} g_\varepsilon'(\omega,t)\,
    \big(F_{n,\omega}^{\mu}(t)-F_{\omega}^{\mu_0}(t)\big)\, \rd t \right| \\ \notag
    &\le
    \int_{0}^{M} |g_\varepsilon'(\omega,t)|\,
    |F_{n,\omega}^{\mu}(t)-F_{\omega}^{\mu_0}(t)|\, \rd t \\
    &\le
     \frac{2 M \|g\|_\infty}{\delta}
    \sup_{t\in[0,M]}|F_{n,\omega}^{\mu}(t)-F_{\omega}^{\mu_0}(t)|.
    \end{align}
    This holds almost surely and concludes the proof.
    % We use \ref{b2} to ensure that F(0)=0 and (a.s.) F_n(0)=1 
\end{proof}


Lemma \ref{ap:lem_sup_r_n_r} is used to prove convergence to zero of a term in Lemma \ref{ap:lem_DCT_P_hat}.
%----------------------------%
\begin{lemma}\label{ap:lem_sup_r_n_r}
%----------------------------%
    Let $g: \Omega \times \R_{\ge0} \to \R$ be a bounded, uniformly continuous function on $\Omega \times \R_{\ge0}$.
    For each $\omega \in \Omega$, define
    \begin{align}\label{eq:def_r_r_n}
    r(\omega) \defin \int_{\R_{\ge0}} g(\omega, t) \, \rd F_{\omega}^{\mu_0}(t), \quad r_n(\omega) \defin  \int_{\R_{\ge0}} g(\omega, t) \, \rd F_{n,\omega}^{\mu}(t).
    \end{align}
    It holds that $\sup_{\omega \in \Omega} |r_n(\omega) - r(\omega)| \convas 0$ as $n \to \infty$.
\end{lemma}
\begin{proof}
    Observe that
    $$
    r(\omega) = \int_{0}^{M} g(\omega, t) \, \rd F_{\omega}^{\mu_0}(t), \quad r_n(\omega) = \int_{0}^{M} g(\omega, t) \, \rd F_{n,\omega}^{\mu}(t).
    $$
Also, one has
\begin{equation}\label{ap:eq_h_n}
    |r_n(\omega) - r(\omega)| \le \left| \int_{0}^{M} (g(\omega, t) - g_\varepsilon(\omega, t)) \, \rd (F_{n,\omega}^{\mu} - F_{\omega}^{\mu_0})(t) \right| + \left| \int_{0}^{M} g_\varepsilon(\omega, t) \, \rd (F_{n,\omega}^{\mu} - F_{\omega}^{\mu_0})(t) \right|.
\end{equation}
Begin with the first term on the RHS of \eqref{ap:eq_h_n}.    
    Fix $\varepsilon > 0$, and define $g_\varepsilon$ as in \eqref{ap:eq_g_epsilon}. 
    If one takes $s = t$ in the definition of $g_\varepsilon(\omega, t)$, one has $g_\varepsilon(\omega, t) \le g(\omega, t)$ for every $(\omega, t) \in \Omega \times [0, M]$. 
    To obtain a lower bound for $g_\varepsilon(\omega, t)$, we consider two cases. If $|t-s|<\delta$, by uniform continuity of $g$, one has
    $g(\omega,s)\ge g(\omega,t)-\varepsilon$, so,
    $$
    g(\omega,s)+\frac{2\|g\|_\infty}{\delta}|t-s|\ge g(\omega,t)-\varepsilon.
    $$
    If
    $|t-s|\ge\delta$, then
    $$
    g(\omega,s)+\frac{2\|g\|_\infty}{\delta}|t-s|
    \ge
    -\|g\|_\infty+2\|g\|_\infty
    =
    \|g\|_\infty
    \ge
    g(\omega,t)
    \ge
    g(\omega,t)-\varepsilon. 
    $$
Taking infima over $s$ yields $g_\varepsilon(\omega, t) \ge g(\omega, t) - \varepsilon$ for every $(\omega, t) \in \Omega \times [0, M]$.    
    Therefore, one has $|g_\varepsilon(\omega, t) - g(\omega, t)| \le \varepsilon \text{ for every } (\omega, t) \in \Omega \times [0, M]$. This provides an upper bound for the second term in the RHS of \eqref{ap:eq_h_n}:
    $$
    \left| \int_{0}^{M} (g_\varepsilon(\omega, t) - g(\omega, t)) \, \rd (F_{n,\omega}^{\mu} - F_{\omega}^{\mu_0})(t) \right| \le \varepsilon \left|\mu_{\omega,n} - \mu_{\omega,0}\right|([0,M]) 
    \le \varepsilon \left|\gamma\right|(\R) \le 2\varepsilon
    $$
    where $|\gamma|$ denotes the total variation of a finite signed measure $\gamma = \mu_{\omega,n} - \mu_{\omega,0}$. Here, $\mu_{\omega,n}$ and $\mu_{\omega,0}$ are the measures induced by $F_{n,\omega}$ and $F_{\omega}^{\mu_0}$, respectively.
     The total variation of a measurable subset $E$ is, in particular, 
    $$
    |\gamma|(E)=\sup_{\text{Partitions $E_1,\ldots,E_m$ of $E$}}\lrb{\sum_{i=1}^m|\gamma(E_j)|}
    $$
    \citep[Exercise 8 in][]{Folland1999}. That $|\gamma|(\R) \le 2$ follows easily from the calculation
\begin{align*}
    |\gamma|(\R)&=\sup \lrb{\sum_{i=1}^m|\gamma(E_j)|}=\sup \lrb{\sum_{i=1}^m|\mu_{\omega,n}(E_j)-\mu_{\omega,0}(E_j)|}\\
    &\leq \sup \lrb{\sum_{i=1}^m|\mu_{\omega,n}(E_j)+\mu_{\omega,0}(E_j)|}= \sup \lrb{\sum_{i=1}^m\mu_{\omega,n}(E_j)+\sum_{i=1}^m\mu_{\omega,0}(E_j)}\\
    &= \sup \lrb{1+1}=2.
\end{align*}
    Now, focus on the second term in the RHS of \eqref{ap:eq_h_n}. Since the class $\cY = \{ y_{\omega,t} : (\omega,t)\in\Omega\times\mathbb R_{\ge0} \}$ is Donsker and $\G_0$ is separable, %we need separability to apply Slutsky, see van der Vaart and Wellner
we obtain the Glivenko--Cantelli property:
$$
\sup_{(\omega,t)\in\Omega\times\R_{\ge0}} \big|F_{n,\omega}^{\mu}(t)-F_\omega^{\mu_0}(t)\big|
=
\sup_{y\in\cY} \big|(P_n - \mu_0)y\big|
\convas 0,
$$
see page 130 in \cite{VanderVaartWellner2023}.
%  \textcolor{blue}{see Lemma \ref{ap:lem_donsker_gc_slutsky}} and 
Thus, applying Lemma \ref{ap:lem_bound_int_g_epsilon} and taking the supremum over $\omega \in \Omega$ in \eqref{ap:eq_bound_d_BL_3} shows that the second term in the RHS of \eqref{ap:eq_h_n} converges to zero as $n \to \infty$. Therefore, we can conclude that
    $$
    \sup_{\omega \in \Omega} |r_n(\omega) - r(\omega)| \le 3\varepsilon,
    $$ 
    for $n$ sufficiently large such that $\frac{2 M \|g\|_\infty}{\delta} \sup_{\omega\in\Omega}\sup_{t\in[0,M]} |F_{n,\omega}^\mu(t) - F_{\omega}^{\mu_0}(t)| \le \varepsilon$.  Since $\varepsilon > 0$ is arbitrary, this shows that $\sup_{\omega \in \Omega} |r_n(\omega) - r(\omega)| \to 0$ as $n \to \infty$.
\end{proof}
    
Lemma \ref{ap:lem_continuity_r} is used in the proof of Lemma \ref{ap:lem_DCT_P_hat}.
%----------------------------%
\begin{lemma}\label{ap:lem_continuity_r}
%----------------------------%
    The function $r$ defined in \eqref{eq:def_r_r_n} is continuous on $\Omega$.
\end{lemma}
\begin{proof}
    By the definition of expectation (law of the unconscious statistician), one has 
    $$
    r(\omega) = \int_{0}^{M} g(\omega,t) \, \rd F_{\omega}^{\mu_0}(t) = \E{g(\omega, d(\omega, X))}.
    $$

    Consider a sequence $(\omega_k)_{k\ge 1}$ with $d(\omega_k, \omega_0) \to 0$ as $k \to \infty$. We want to prove $\lim_{k \to \infty} \E{g(\omega_k, d(\omega_k, X))} =  \E{g(\omega_0, d(\omega_0, X))}$. Since $d(\omega_k, x) \to d(\omega_0, x)$ for every $x \in \Omega$ and $g$ is continuous on $\Omega \times [0, M]$, it holds that $g(\omega_k, d(\omega_k, X)) \convas g(\omega_0, d(\omega_0, X))$ as $k\to \infty$. Now, since $g$ is bounded, an application of the DCT gives
    $$
    \lim_{k \to \infty} r(\omega_k) = \lim_{k \to \infty} \E{g(\omega_k, d(\omega_k, X))} = \E{g(\omega_0, d(\omega_0, X))} = r(\omega_0).
    $$
    Thus, $r$ is continuous on $\Omega$. 
\end{proof}


Lemma \ref{ap:lem_DCT_P_hat} is used to prove convergence of a term in the proof of Theorem \ref{thm:test_statistics_simp}.
%----------------------------%
\begin{lemma}\label{ap:lem_DCT_P_hat}
%----------------------------%
    Let $g: \Omega \times \R_{\ge0} \to \R$ be a bounded, uniformly continuous function on $\Omega \times \R_{\ge0}$. It holds under $H_0$ that
    $$
    \int_{\Omega \times \R_{\ge0}} g(\omega, t) \, \rd F_{n,\omega}^{\mu}(t) \, \rd P_n(\omega) \convas \int_{\Omega \times \R_{\ge0}} g(\omega, t) \, \rd F_{\omega}^{\mu_0}(t) \, \rd \mu_0(\omega) \text{ as } n \to \infty.
    $$
\end{lemma}
\begin{proof}
    One has
    $$
    \begin{aligned}
    \left|\int_{\Omega} r_n(\omega) \, \rd P_n(\omega) - \int_{\Omega} r(\omega) \, \rd \mu_0(\omega) \right|
    &\le
    \left| \int_{\Omega} r_n(\omega) - r(\omega) \, \rd P_n(\omega) \right| +
    \left| \int_{\Omega} r(\omega) \, \rd \left( P_n(\omega) -  \mu_0(\omega)\right) \right| \\
    &\le
    \sup_{\omega \in \Omega} |r_n(\omega) - r(\omega)| +
    \left| \int_{\Omega} r(\omega) \, \rd \left( P_n(\omega) -  \mu_0(\omega)\right) \right| \to 0
    \end{aligned}
    $$
    as $n \to \infty$.  The convergence follows by Lemma \ref{ap:lem_sup_r_n_r} for the first term. For the second term, convergence follows by the weak convergence of $P_n$ to $\mu_0$ and the continuity of $r$ by Lemma \ref{ap:lem_continuity_r}. For the weak convergence of $P_n$ to $\mu_0$, see Theorems 1 and 3 of \cite{Varadarajan1958}.
\end{proof}

%----------------------------%
\begin{comment}
%----------------------------%
\begin{lemma}\label{ap:lem_dist_prod_0_dist_L1_0_old_L2}
%----------------------------%
   Let $d_{\Omega \times \R_{\ge 0}}$  denote the product distance on $(\Omega, d) \times (\R_{\ge 0}, |\cdot|)$, and let $d_{L^2(P)}$ be the $L^2(P)$ pseudometric induced by a random object $X$ taking values in $\Omega$. That is, for $p\ge1$, we define
   \begin{align}\label{eq:ap_distances_old_L2}
    d_{\Omega \times \R_{\ge 0}}\big((\omega,t), (\omega^\prime, t^\prime)\big) \defin d(\omega, \omega^\prime) + |t - t^\prime|, \quad d_{L^p(P)}\big((\omega,t), (\omega^\prime, t^\prime)\big) \defin \left(\mathsf{E}\big|y_{\omega,t}(X) - y_{\omega^\prime, t^\prime}(X)\big|^p\right)^{\frac{1}{p}}.
   \end{align}
Assume that the distribution of $d(\omega,X)$ is continuous for every $\omega\in \Omega$. If $d_{\Omega \times \R_{\ge 0}}\big((\omega_k,t_k), (\omega, t)\big) \to 0$ as $k\to\infty$, then also $d_{L^2(P)}\big((\omega_k,t_k), (\omega, t)\big) \to 0$ as $k\to\infty$.
\end{lemma}
%----------------------------%
\begin{proof}
    \textit{Note: The proof is very similar to that of Lemma \ref{lem:continuity_omega}}.

Consider a sequence $\{(\omega_k, t_k)\}_{k\ge 1}$ with
$d_{\Omega \times \R_{\ge 0}}((\omega_k,t_k),(\omega,t))\to0$. In particular $d(\omega_k,\omega)\to0$ and $t_k\to t$.

Fix $x \in \Omega$ such that $d(\omega,x) \neq t$. One has
    \begin{align}\label{eq:conv_indicators_old_L2}
    1_{\{d(\omega_k,x)  \le t_k\}} \to 1_{\{d(\omega,x) \le t\}} \ \text{ as } k \to \infty.
    \end{align}   
To see this, first assume $d(\omega,x)<t$, and take $\varepsilon= \frac{t - d(\omega,x)}{2} >0$. Then, we have
$$
d(\omega_k,x) - t_k = (d(\omega,x) - d(\omega_k, \omega)) + (d(\omega,x) - t) + (t-t_k) < 0,
$$
for sufficiently large $k$, so that $|d(\omega_k,x) - d(\omega,x)| < \varepsilon$ and $|t_k - t| < \varepsilon$. Analogously, if $d(\omega,x)>t$, take $\varepsilon= \frac{d(\omega,x) - t}{2}$, for sufficiently large $k$ it holds that
$$
t_k - d(\omega_k,x) = (t_k - t) + (t - d(\omega,x)) + (d(\omega,x) - d(\omega_k,x)) < 0,
$$
Thus \eqref{eq:conv_indicators_old_L2} holds for every $x$ with $d(\omega,x)\neq t$.  Observe that since $\Prob{d(\omega_k,x) = t_k} = 0$ (by the continuity of the distribution of $d(\omega,X)$ for every $\omega \in \Omega$), the convergence in \eqref{eq:conv_indicators_old_L2} holds $P$-a.s. 

Since $1_{\{d(\omega_k,x) \le t_k\}}$ is bounded for every $k\ge0$, an application of the DCT gives
$$
\begin{aligned}
d_{L^2(P)}\big((\omega_k,t_k),(\omega,t)\big)^2
= \E{ ( 1_{\{d(\omega_k,X)\le t_k\}} - 1_{\{d(\omega,X)\le t\}} )^2 }
= \int_\Omega \big( 1_{\{d(\omega_k,x)\le t_k\}} - 1_{\{d(\omega,x)\le t\}} \big)^2 \, \rd \mu(x)
\to 0
\end{aligned}
$$
as $k\to\infty$. This proves the result.
\end{proof}
\end{comment}

\begin{comment}
Lemma \ref{ap:lem_aux_omega_0_delta_0} provides an entropy bound needed to obtain uniform continuity in Lemma \ref{ap:lem_continuity_sample_paths_G}.
%----------------------------%
\begin{lemma}\label{ap:lem_aux_omega_0_delta_0}[Replaced by Lemma \ref{ap:lem_continuity_sample_paths_G_rho2}]
%----------------------------%
    \textcolor{blue}{ Let $\omega(\delta) \defin  \int_0^\delta \sqrt{\log N(\varepsilon, \Omega \times \R_{\ge 0}, \rho_{2,y})} \rd \varepsilon$. 
    Assuming that $\log(N(\varepsilon, \Omega, d)) = O(\varepsilon^{-\alpha})$ for some $\alpha<1$, it holds that $\omega(\delta) \to 0$ as $\delta \to 0$.}
\end{lemma}
\begin{proof}
The first step is to prove 
\begin{align}\label{ap:eq_entropy_bound_dubey_entropy}
    \int_0^\infty \sqrt{\log N_{[]}(\varepsilon, \Omega \times \R_{\ge 0}, \rho_{2,y})} \rd \varepsilon < \infty.
\end{align}

\cite{DubeyChenMuller2024} show that
\begin{align}\label{eq:entropy_dubeychenmuller}
\int_0^\infty \sqrt{\log N_{[]}(\varepsilon,\cY, \rho_{1,y})} d \varepsilon < \infty
\end{align}
where $\cY$ is the class of indicator functions $\cY = \{y_{\omega, t}: \, \omega \in \Omega, \, t\in\R_{\ge 0}\}$. %Theorem 5.32 in \cite{vanHandel} requires subGaussianity (and we also have to check separability). SubGaussianity is a property that depends on the norm. SubGaussianity is implied by Gaussianity under the $L^2(\mu)$ norm, not under $\rho_{1,y}$. We must either prove subGaussianity for $\rho_{1,y}$, or prove the other conditions for $L^2(\mu)$. I try the second approach. 
Observe that this result is for $\rho_{1,y}$, not $L^2(\mu)$. Since for any pair $(\omega,t), (\omega^\prime, t^\prime) \in \Omega \times \R_{\ge 0}$, it holds that 
\begin{align}\label{eq:L_1_L_2_relation}
\rho_{1,y}((\omega,t), (\omega^\prime, t^\prime)) = \rho_{2,y}((\omega,t), (\omega^\prime, t^\prime))^2,
\end{align}
 then $N_{[]}(\varepsilon, \Omega \times \R_{\ge 0}, \rho_{2,y}) = N_{[]}(\varepsilon^2, \Omega \times \R_{\ge 0}, \rho_{1,y})$. Therefore,
$$
\int_0^\infty \sqrt{\log N_{[]}(\varepsilon, \cY, \rho_{2,y})} \rd \varepsilon = \frac{1}{2} \int_0^\infty \sqrt{\log N_{[]}(\varepsilon, \cY, \rho_{1,y})} \frac{d\varepsilon}{\sqrt{\varepsilon}}.
$$
So if $\int_0^\infty \sqrt{\log N_{[]}(\varepsilon, \cY, \rho_{1,y})} \frac{d\varepsilon}{\sqrt{\varepsilon}} < \infty$, then \eqref{ap:eq_entropy_bound_dubey_entropy} holds.

If $\varepsilon>1$, then $N_{[]}(\varepsilon, \cY, \rho_{1,y}) = 1$ because $\cY$ is a class of indicator functions, so %the distance is in fact a probability, which is bounded by $1$. Hence,
$$
\int_1^\infty \sqrt{\log N_{[]}(\varepsilon, \cY, \rho_{1,y})} \frac{d\varepsilon}{\sqrt{\varepsilon}} = 0
$$
Now, focus on $0<\varepsilon\le1$. \cite[][page 5 of supplementary material]{DubeyChenMuller2024} show that
$$
\sqrt{\log{N_{[]}(\varepsilon, \cY, \rho_{1,y})}} \le  \sqrt{\log\left(\frac{4}{\varepsilon}\right) + \log N\left(\frac{\varepsilon}{4\bar{\Delta}}, \Omega, d\right)}  \le  \sqrt{\log \left(\frac{4}{\varepsilon}\right)} + \sqrt{\log N\left(\frac{\varepsilon}{4\bar{\Delta}}, \Omega, d\right)} ,
$$
For the first term on the RHS of the inequality, one has 
$$
\int_{0}^{1} \sqrt{\log \left(\frac{4}{\varepsilon}\right)} \frac{d\varepsilon}{\sqrt{\varepsilon}} = \int_{0}^{1} e^{-\frac{u}{2}} \sqrt{\log(4) + u} \, \rd u < \infty,
$$
with $u=-\log(\varepsilon)$.

Let $\bar{\Delta}$ as defined in Assumption \ref{b2}, and assume that $\bar{\Delta} < \frac{1}{4}$. The proof also follows (in fact, more simply) if $\bar{\Delta} \ge \frac{1}{4}$. It holds that 
$$
\begin{aligned}
\int_{0}^{1} \sqrt{\log N\left(\frac{\varepsilon}{4\bar{\Delta}}, \Omega, d\right)} \frac{d\varepsilon}{\sqrt{\varepsilon}} &= 2\sqrt{\bar{\Delta}} \int_0^{\frac{1}{4\bar{\Delta}}} \sqrt{\log N(v, \Omega, d)} \; \frac{\rd v}{\sqrt{v}} \\ 
&= 2\sqrt{\bar{\Delta}} \left( \int_0^{1} \sqrt{\log N(v, \Omega, d)} \; \frac{\rd v}{\sqrt{v}} + \int_1^{ \frac{1}{4\bar{\Delta}}} \sqrt{\log N(v, \Omega, d)} \; \frac{\rd v}{\sqrt{v}} \right)
\end{aligned}
$$
with $v=\frac{\varepsilon}{4\bar{\Delta}}$. Under Assumptions 1 and 2 in \cite{DubeyChenMuller2024}, it is proved in ibid that $\int_{0}^{1} \sqrt{\log N(\frac{\varepsilon}{4\bar{\Delta}}, \Omega, d)} \, \rd \varepsilon < \infty$.
Hence, if $v \in \left[1, \frac{1}{4\bar{\Delta}}\right]$, then $\frac{1}{\sqrt{v}} \le 1$, and one has 
$$
\int_{1}^{\frac{1}{4\bar{\Delta}}} \sqrt{\log N(v, \Omega, d)} \; \frac{\rd v}{\sqrt{v}} \le \int_{1}^{\frac{1}{4\bar{\Delta}}} \sqrt{\log N(v, \Omega, d)} \, \rd v < \infty.
$$  
It is only left to check that 
$$
\int_{0}^{1} \sqrt{\log\left(N(v, \Omega, d)\right)} \frac{\rd v}{\sqrt{v}} < \infty.
$$
But this holds due to the assumption that $\log(N(\varepsilon, \Omega, d)) = O(\varepsilon^{-\alpha})$ for some $\alpha<1$. 

% Assumption 1 in \cite{DubeyChenMuller2024} ensures $\varepsilon \log(N(\varepsilon)) \to 0$ as $\varepsilon \to 0$, but this is not enough; take, e.g., $\log(N(\varepsilon)) = \frac{1}{\varepsilon\log(1/\varepsilon)}$. So we need a restriction on how fast $\varepsilon \log(N(\varepsilon))$ converges to $0$.


Finally, since $N(\varepsilon, \cF, L^2(\mu)) \le N_{[]}(2\varepsilon, \cF, L^2(\mu))$, then $\int_0^{\infty} \sqrt{\log N_{[]}(\varepsilon, \cF, L^2(\mu))} \, d\varepsilon < \infty$ implies 
$$
\int_0^{\infty} \sqrt{\log N(\varepsilon, \Omega\times \R_{\ge 0}, L^2(\mu))} \,\allowbreak d\varepsilon < \infty
$$
(see Remark \ref{ap:rem_F_Omega_R}). Because $\sqrt{\log N(\varepsilon, \Omega, \rho_{2,y})}$ is integrable and the Lebesgue integral is a continuous functional, it follows that $\omega(\delta) \to 0$ as $\delta \to 0$.
\end{proof}
%----------------------------%

%----------------------------%
\begin{remark}\label{ap:rem_F_Omega_R}
%----------------------------%
We are interested in the covering number $N(\varepsilon, \Omega \times \mathbb{R}_{\ge 0}, \rho_{p,y})$, whereas \citet{DubeyChenMuller2024} work with $N(\varepsilon, \cY, \rho_{p,y})$. 
These two covering numbers can be identified via the map $(\omega,t)\mapsto y_{\omega,t}$. 
With a slight abuse of notation, we denote 
$\rho_{p,y}((\omega,t),(\omega',t')):=\rho_{p,y}(y_{\omega,t},y_{\omega',t'})$.
Although the map $(\omega,t)\mapsto y_{\omega,t}$ need not be injective, it induces an isometric identification between $(\cY,\rho_{p,y})$ and the quotient of $\Omega\times\mathbb{R}_{\ge 0}$ obtained by identifying pairs that generate the same function. 
Consequently, the covering numbers agree:
$N(\varepsilon,\cY,\rho_{p,y})=N(\varepsilon,\Omega\times\mathbb{R}_{\ge 0},\rho_{p,y})$
(and similarly for any $p\ge 1$). We will use this identification without further mention. 
Note that this argument is specific to covering numbers; it does not directly apply to bracketing numbers, which rely on an order structure.
\end{remark}
%----------------------------%
\begin{remark}\label{rem:alternative_donsker_proof}
%----------------------------%
By Theorem~2.3 in \cite{Kosorok2008}, the condition \eqref{ap:eq_entropy_bound_dubey_entropy} implies that $\cF$ is
$\mu$-Donsker. Therefore, Lemma \ref{ap:lem_aux_omega_0_delta_0} gives an alternative way to prove the Donsker property of $\cF$ established in \cite{DubeyChenMuller2024}. 
\end{remark}
\end{comment}


\begin{comment}
%----------------------------%
\begin{lemma}\label{ap:lem_dist_prod_0_dist_L1_0_old_rho1}
%----------------------------%
   Let $d_{\Omega \times \R_{\ge 0}}$  denote the product distance on $(\Omega, d) \times (\R_{\ge 0}, |\cdot|)$, and let $\rho_{1,y}$ be the $\rho_{1,y}$ distance induced by a random object $X$ taking values in $\Omega$. That is,
   $$
    d_{\Omega \times \R_{\ge 0}}\big((\omega,t), (\omega^\prime, t^\prime)\big) \defin d(\omega, \omega^\prime) + |t - t^\prime|, \quad \rho_{1,y}\big((\omega,t), (\omega^\prime, t^\prime)\big) \defin \mathsf{E}\big|y_{\omega,t}(X) - y_{\omega^\prime, t^\prime}(X)\big|.
   $$
Assume that the distribution of $d(\omega,X)$ is continuous for every $\omega\in \Omega$. If $d_{\Omega \times \R_{\ge 0}}\big((\omega_k,t_k), (\omega, t)\big) \to 0$ as $k\to\infty$, then also $\rho_{1,y}\big((\omega_k,t_k), (\omega, t)\big) \to 0$ as $k\to\infty$.
\end{lemma}
%----------------------------%
\begin{proof}
    \textit{Note: The proof is very similar to that of Lemma \ref{lem:continuity_omega}}.
    Consider a sequence $\{(\omega_k, t_k)\}_{k\ge 1}$ in $\Omega$ such that $d_{\Omega \times \R_{\ge 0}}((\omega_k, t_k), (\omega, t)) \to 0$. Obviously, $|t-t_k| \to 0$. Also, since $d(\cdot, x)$ is continuous for every $x \in \Omega$, then $d(\omega_k, x) \to d(\omega,x)$ as $k\to \infty$ for every $x \in \Omega$. 

    Fix $\delta \ge 0$, and consider $x \in \Omega$ such that $|d(\omega,x) - t| \neq \delta$. One has
    \begin{align}\label{eq:conv_indicators_old_rho1}
    1_{\{|d(\omega_k,x) - t_k| \le \delta\}} \to 1_{\{|d(\omega,x) - t| \le \delta\}} \ \text{ as } k \to \infty.
    \end{align}
    To see this, first assume that $|d(\omega,x) - t| < \delta$, and let $\varepsilon = \frac{\delta - |d(\omega,x) - t|}{2} >0$. One has
    $$
    |d(\omega_k,x) - t_k| \le |d(\omega_k,x) - d(\omega, x)| +  |d(\omega,x) - t| + |t - t_k| < \delta
    $$
    for $k$ sufficiently large, so that $|d(\omega_k,x) - d(\omega, x)| \le \varepsilon$ and $|t - t_k| \le \varepsilon$. If instead, $|d(\omega,x) - t| > \delta$, let $\varepsilon = \frac{|d(\omega,x) - t| - \delta}{2} > 0$. Then, by applying the reverse triangle inequality,
    $$
    |d(\omega_k,x) - t_k| \ge |d(\omega, x) - t| - |d(\omega_k,x) - d(\omega,x)| - |t - t_k| \ge \delta
    $$    
    for $k$ sufficiently large. This proves \eqref{eq:conv_indicators_old_rho1}.
    
    Observe that since $\Prob{|d(\omega_k,x) - t_k| = \delta} = 0$ (by the continuity of the distribution of $d(\omega,X)$ for every $\omega \in \Omega$), the convergence in \eqref{eq:conv_indicators_old_rho1} holds $P$-a.s. Now, since $1_{\{|d(\omega_k,x) - t_k| \le \delta\}}$ is bounded for every $k\ge0$, an application of the DCT gives
    $$
    d_{L^1(P)}\big((\omega_k,t_k), (\omega, t)\big) = \int_{\Omega} \left|1_{\{|d(\omega_k,x) - t_k| \le \delta\}} - 1_{\{|d(\omega,x) - t| \le \delta\}} \right| \, dP(x) \to 0 \text{ as } k\to\infty.
    $$
    %That is, $h(\omega_k, \delta) \to h(\omega, \delta)$ when $\omega_k \to \omega$, i.e., $h(\cdot, \delta)$ is continuous on $\Omega$ for every $\delta\ge0$.
    This concludes the proof.
\end{proof}
\end{comment}

\begin{comment}
NOT NECESSARY ANYMORE
%----------------------------%
\begin{lemma}\label{ap:lem_cov_number_greater}
%----------------------------%
    There exist constants $c, q>0$ such that $N(\Omega \times \R_{\ge 0}, d_{L^2(\Omega)}, \varepsilon) \geq(c / \varepsilon)^q$ for all $\varepsilon>0$.
\end{lemma}
%----------------------------%
\begin{proof}
    Maybe this is not even necessary!
\end{proof}
\end{comment}





Lemma \ref{ap:lem_relation_L2_prod} connects $d_{\Omega \times \R_{\ge 0}}$ with the $L^2(\mu)$ pseudometric induced by $\G_0$, which we denote by $\rho_{2,\G_0}$.
This allows obtaining uniform continuity with respect to $\rho_{2,\G_0}$ from uniform continuity with respect to $d_{\Omega \times \R_{\ge 0}}$, which is required in Lemma \ref{ap:lem_continuity_sample_paths_G_rho2}.
%----------------------------%
\begin{lemma}\label{ap:lem_relation_L2_prod}
%----------------------------%
% Let $d_{\Omega \times \R_{\ge 0}}$  denote the product distance on $(\Omega, d) \times (\R_{\ge 0}, |\cdot|)$. 
Let $\rho_{2,y}$ and $\rho_{2,\G_0}$ be the $L^2(\mu)$ pseudometrics induced by the class of functions $\cY$ and by the Gaussian process $\G_0$, respectively. 
That is, for $p\ge1$ and $(\omega,t), (\omega^\prime, t^\prime) \in \Omega \times \R_{\ge 0}$, we define
\begin{equation}\label{eq:ap_distances}
\begin{aligned}
\rho_{p,y}\big((\omega,t), (\omega^\prime, t^\prime)\big)
&\defin
\Big(\mathsf{E}\big|y_{\omega,t}(X) - y_{\omega^\prime, t^\prime}(X)\big|^p\Big)^{\frac{1}{p}},
\\
\rho_{p,\G_0}\big((\omega,t),(\omega^\prime,t^\prime)\big)
&\defin
\Big(
\mathsf E\big|
\G_0(\omega,t)-\G_0(\omega^\prime,t^\prime)
\big|^p
\Big)^{1/p}.
\end{aligned}
\end{equation}
%Under Assumption \ref{b2},  For all $(\omega,t), (\omega^\prime,t^\prime) \in \Omega \times \R_{\ge 0}$. 
It holds that 
\begin{align}\label{eq:rho2_product_bound_simple}
\rho_{2,\G_0}\big((\omega,t),(\omega^\prime,t^\prime)\big)
\le
\rho_{2,y}\big((\omega,t),(\omega^\prime,t^\prime)\big)
\le
2\sqrt{\bar\Delta 
d_{\Omega\times\R_{\ge0}}
\big((\omega,t),(\omega^\prime,t^\prime)\big)
},
\end{align}
In particular, any function that is uniformly continuous with respect to $\rho_{2,\G_0}$ is also uniformly continuous with respect to $d_{\Omega\times\R_{\ge0}}$.
\end{lemma}
%----------------------------%


\begin{proof}


To prove the first inequality in \eqref{eq:rho2_product_bound_simple}, recall the covariance structure \eqref{eq:cov_G_theta_0}
$$
\Cov{\G_0(\omega,t)}{\G_0(\omega^\prime,t^\prime)} 
=
\mu(y_{\omega,t}y_{\omega^\prime,t^\prime})-\mu y_{\omega,t}\,\mu y_{\omega^\prime,t^\prime}.
$$
One has
\begin{align*}
\rho_{2,\G_0}\big((\omega,t),(\omega^\prime,t^\prime)\big)^2
&=
\Ebig{
\G_0(\omega,t)-\G_0(\omega^\prime,t^\prime)
}^2
\\
&=
\mu\Big(
\big(y_{\omega,t}-y_{\omega^\prime,t^\prime}\big)^2
\Big)
-
\Big(
\mu\big(y_{\omega,t}-y_{\omega^\prime,t^\prime}\big)
\Big)^2
\\
&\le
\mu\Big(
\big(y_{\omega,t}-y_{\omega^\prime,t^\prime}\big)^2
\Big)
=
\rho_{2,y}\big((\omega,t),(\omega^\prime,t^\prime)\big)^2,
\end{align*}
which yields the desired inequality.


Now, prove the second inequality in \eqref{eq:rho2_product_bound_simple}. 
Observe that since $y_{\omega, t}$ and $y_{\omega^\prime, t^\prime}$ are indicator functions, one has $(y_{\omega, t} - y_{\omega^\prime, t^\prime})^2 = 1_{B(\omega,t)\triangle B(\omega^\prime,t^\prime)}$, where $A\triangle B \defin (A\setminus B)\cup(B\setminus A)$ denotes the symmetric difference between two sets $A$ and $B$. 
Therefore, it is verified that
$$
\rho_{2,y}\big((\omega,t), (\omega^\prime,t^\prime)\big)^2 
= \E{(y_{\omega, t} - y_{\omega^\prime, t^\prime})^2} 
= \Probbig{ B(\omega,t) \triangle B(\omega^\prime,t^\prime) }.
$$

Define $r \defin d_{\Omega \times \R_{\ge 0}}((\omega,t),(\omega^\prime,t^\prime))$. Then,
\begin{align}\label{ap:aux_lem_containment}
    B(\omega,t) \triangle B(\omega^\prime,t^\prime) \subset \{ x : |d(\omega,x)-t| \le r \} \cup \{ x : |d(\omega^\prime,x)-t^\prime| \le r \}.
\end{align}
To see this, consider any $x \in B(\omega,t) \triangle B(\omega^\prime,t^\prime)$. Without loss of generality, assume $x \in B(\omega,t) \setminus B(\omega^\prime,t^\prime)$ (the other case is symmetric). Then, we have $d(\omega,x) \le t$ and $d(\omega^\prime,x) > t^\prime$

From the reverse triangle inequality, one obtains $d(\omega,x) \ge d(\omega^\prime,x) - d(\omega,\omega^\prime) > t^\prime - d(\omega,\omega^\prime)$. This provides a bound for $|d(\omega,x) - t| = t - d(\omega,x)$:
$$
t - d(\omega,x) < t - \big(t^\prime - d(\omega,\omega^\prime)\big) = (t - t^\prime) + d(\omega,\omega^\prime) \le |t-t^\prime| + d(\omega,\omega^\prime) = r.
$$
Thus, $|d(\omega,x) - t| \le r$. Analogously, if $x \in B(\omega^\prime,t^\prime) \setminus B(\omega,t)$, one obtains $|d(\omega^\prime,x) - t^\prime| \le r$. Therefore,
$$
x \in \big\{ x : |d(\omega,x)-t| \le r \big\} \cup \big\{ x : |d(\omega^\prime,x)-t^\prime| \le r \big\},
$$
which proves \eqref{ap:aux_lem_containment}.


From \eqref{ap:aux_lem_containment}, it readily follows that 
$$
\Probbig{ B(\omega,t) \triangle B(\omega^\prime,t^\prime) } \le \Probbig{ |d(\omega,X)-t| \le r } + \Probbig{ |d(\omega^\prime,X)-t^\prime| \le r }.
$$
By Assumption \ref{b2}, for every $(\omega,t) \in \Omega \times \R_{\ge 0}$ it holds that $\sup_{\omega,t} f_\omega(t) \le \bar\Delta$. Thus, we obtain:
$$
\Probbig{ |d(\omega,X)-t| \le r } = \int_{t-r}^{t+r} f_\omega(s) \,\rd s \le \int_{t-r}^{t+r} \bar\Delta \,\rd s  = 2\bar\Delta r.
$$
The same estimate applies to $\Probbig{ |d(\omega^\prime,X)-t^\prime| \le r }$. Hence,
\begin{equation}\label{eq:rho2y_prod_bound}
\rho_{2,y}\big((\omega,t), (\omega^\prime,t^\prime)\big)^2 \le 4 \bar\Delta \, r = 4 \bar\Delta \, d_{\Omega \times \R_{\ge 0}}((\omega,t),(\omega^\prime,t^\prime))
\end{equation}
which gives the second inequality in \eqref{eq:rho2_product_bound_simple}.
\end{proof}
%----------------------------%

Lemma \ref{ap:lem_continuity_sample_paths_G_rho2} shows that the Gaussian process $\G_{0}(\omega, t)$ has almost surely uniformly continuous sample paths on $(\Omega \times \R_{\ge 0}, d_{\Omega \times \R_{\ge 0}})$. This result is used, combined with Lemma \ref{ap:lem_DCT_P_hat}, in the proof of Theorem \ref{thm:test_statistics_simp}.


%----------------------------%
\begin{lemma}\label{ap:lem_continuity_sample_paths_G_rho2}
%----------------------------%
The process $\G_0(\omega,t)$ has almost surely uniformly continuous sample paths on the metric space $(\Omega\times\R_{\ge0}, d_{\Omega\times\R_{\ge0}})$, considering the sample paths as functions from $(\Omega\times\R_{\ge0}, d_{\Omega\times\R_{\ge0}})$ to $(\R,|\cdot|)$.
\end{lemma}

\begin{proof}
Under the null hypothesis, the class $\cY$ is $\mu$-Donsker. Hence, the limiting process $\G_0$ is a tight centered Gaussian process in $\ell^\infty(\cY)$.

By Lemma 1.5.9 in \citet{VanderVaartWellner2023}, there exists a semimetric $\rho$ on $\cY$ such that $\cY$ is totally bounded with respect to $\rho$, and almost all sample paths of $\G_0$ are uniformly $\rho$-continuous. 
If, in addition, the map $\cY\mapsto\R$ defined by $y\mapsto \G_0(y)$ is uniformly $\rho$-continuous in second mean, then Lemma 1.5.9 in \citet{VanderVaartWellner2023} (equivalence between (i) and (iii)) can be applied (with $p=2$). Therefore, $\cY$ is totally bounded with respect to $\rho_{2,\G_0}$, and almost all sample paths of $\G_0$ are uniformly $\rho_{2,\G_0}$-continuous.
Here, $\rho_{2,\G_0}(y,y^\prime) = \big(\mathsf E\big|\G_0(y)-\G_0(y^\prime)\big|^2\big)^{1/2}$, which is the same as \eqref{eq:ap_distances}, but as a function of $\cY\times\cY$, rather than $(\Omega\times\R_{\ge0})\times(\Omega\times\R_{\ge0})$.
However, for every $y,y^\prime \in \cY$, there exist $(\omega,t), (\omega^\prime,t^\prime) \in \Omega\times\R_{\ge0}$ such that $y=y_{\omega,t}$ and $y^\prime = y_{\omega^\prime,t^\prime}$. Hence, all sample paths of $\G_0(\omega,t)$ are uniformly $\rho_{2,\G_0}$-continuous as functions from $(\Omega\times\R_{\ge0}, \rho_{2,\G_0})$ to $(\R,|\cdot|)$.
It follows from Lemma~\ref{ap:lem_relation_L2_prod} that any function that is uniformly continuous with respect to $\rho_{2,\G_0}$ is also uniformly continuous with respect to $d_{\Omega\times\R_{\ge0}}$. Therefore, $\G_0$ has almost surely uniformly continuous sample paths as a function from $(\Omega\times\R_{\ge0}, d_{\Omega\times\R_{\ge0}})$ to $(\R,|\cdot|)$.

It remains to show that the map $y\mapsto \G_0(y)$ is uniformly $\rho$-continuous in second mean. Let $(y_n)_{n\ge0}$ and $(y_n^\prime)_{n\ge0}$ be sequences in $\cY$ such that $\rho(y_n,y_n^\prime)\to0$, as $n\to\infty$.
Since almost all sample paths of $\G_0$ are uniformly $\rho$-continuous, it follows that $\G_0(y_n)-\G_0(y_n^\prime) \convas 0$, and so
$\G_0(y_n)-\G_0(y_n^\prime) \convl 0$.

For every $n$, define
$$
Z_n:=\G_0(y_n)-\G_0(y_n^\prime).
$$  
Since $\G_0$ is a centered Gaussian process, one has that $Z_n\sim N(0,\sigma_n^2)$, with $\sigma_n^2=\mathsf E(Z_n^2).$
Since $Z_n \convl 0$, it holds that $\exp(-\sigma_n^2u^2/2)\to1$ for every $u\in\R,$
which implies that $\sigma_n^2\to0$. Therefore,
$$
\mathsf E\big|
\G_0(y_n)-\G_0(y_n^\prime)
\big|^2
\to0.
$$
This proves that $y\mapsto \G_0(y)$ is uniformly $\rho$-continuous in second mean.
\end{proof}
%----------------------------%


\begin{comment}
%----------------------------%
\begin{lemma}[replaced by Lemma \ref{ap:lem_continuity_sample_paths_G_rho2}]\label{ap:lem_continuity_sample_paths_G}
%----------------------------%
    \textcolor{blue}{Assume that $\log(N(\varepsilon, \Omega, d)) = O(\varepsilon^{-\alpha})$ for some $\alpha<1$. The process $\G_{0}(\omega, t)$ has almost surely uniformly continuous sample paths on the metric space $(\Omega \times \R_{\ge 0}, d_{\Omega \times \R_{\ge 0}})$, considering the sample paths as functions from $(\Omega \times \R_{\ge 0}, d_{\Omega \times \R_{\ge 0}})$ to $(\R, |\cdot|)$.}
\end{lemma}
\begin{proof}
\textcolor{blue}{Note: For a Gaussian process, continuity in probability is equivalent to mean-square continuity, and continuity with probability one (a.s.) is equivalent to sample continuity}. %But the definition of sample continuity is very confusing to me, because in many places it is defined as ``there exists a version of the process with a.s. continuous sample paths''. I think this would be in conflict with what we are doing, because we are already using a specific version of the process (Skorohod's construction for a.s. convergence)! There is a result by Kolmogorov which guarantees sample continuity but it also talks about ``a version of the process''. After a very exhaustive search, 
    %\item   \textit{It is important that the paths are a.s. continuous with respect to the usual metric on $\Omega\times \R_{\ge 0}$ (denote it by $d$), because this is what we need in the (extended) Helly-Bray Theorem. We need the usual distance both in the domain and the codomain.}

Since $\G_{0}(\omega, t)$ is a separable $L^2(\mu)$-subGaussian process on $(\Omega \times \R_{\ge 0}, \rho_{2,y})$ (it is a Gaussian process), then the function
$$
\omega(\delta)=\int_0^\delta \sqrt{\log N(\varepsilon, \Omega \times \R_{\ge 0}, \rho_{2,y})} \rd \varepsilon
$$
is a modulus of continuity for $\left\{X_t\right\}_{t \in T}$. See, e.g., Theorem 5.32 in \cite{vanHandel}, page 140 in \cite{VanderVaartWellner2023} or \cite{Talagrand2021} for this result.
%In particular, we have
%$$
%\mathbf{E}\left[\sup _{t, s \in T} \frac{X_t-X_s}{\omega(d(t, s))}\right]<\infty.
%$$
This means that there is a random variable $K$ such that
$$
y_{\omega,t} - y_{\omega^\prime, t^\prime} \le K w\left(d_{L^2(P)}((\omega,t),(\omega^\prime,t^\prime))\right) \ \text{for all } t,s \in T.
$$
Hence, under the assumption that $\log(N(\varepsilon, \Omega, d)) = O(\varepsilon^{-\alpha})$ for some $\alpha<1$ and condition \ref{b2}, it follows from Lemmas \ref{ap:lem_relation_L2_prod} and \ref{ap:lem_aux_omega_0_delta_0} that $y_{\omega,t}$ is uniformly almost surely continuous as a function from $(\Omega \times \R_{\ge0}, d_{\Omega \times \R_{\ge 0}})$ to $(\R, |\cdot|)$. 
See page 769 in \cite{DubeyChenMuller2024} for an analysis and examples of metric spaces satisfying $\log(N(\varepsilon, \Omega, d)) = O(\varepsilon^{-\alpha})$ for some $\alpha<1$, and see also page 70 in \cite{Dudley1973} for a discussion of $\alpha$ under the name \emph{exponent of entropy}. 
\end{comment}
\begin{comment}
Notice that we are working with functions of the type $y_{\omega, t}$. We are not interested in $|y_{\omega,t} - y_{\omega_k,t_k}|$, but rather in $\G_{\theta_0}(\omega, t) - \G_{\theta_0}(\omega_k, t_k)$. But it can be seen that 
$
\E{
\bigl(\G_{0}(y_{\omega_1,t_1})-\G_{0}(y_{\omega_2,t_2})\bigr)^2
} \le \|y_{\omega_1,t_1}-y_{\omega_2,t_2}\|_{L^2(P)}.
$
%This is the L^2 metric based on the Gaussian process G. If the L^2(P) pseudo-metric (based on the indicators) goes to zero, then the metric based on the Gaussian process also goes to zero. This is useful because we have already shown that if the usual distance in $\Omega \times \R_{\ge0}$ goes to zero, then the L^2(P) pseudometric goes to zero and hence the L^2 metric based on G also goes to zero, so Theorem 5.31 provides continuity wrt the usual distance in $\Omega \times \R_{\ge0}$, which is what we want.
 It is only left to check that Theorem 5.32 applies for G, and not only for the indicators (which is what we have proved so far). The entropy condition will be trivially satisfied for the L^2 distance induced by G.
The covariance structure \eqref{eq:cov_simp} of $\G_{0}$ gives
$$
\E\left[
\bigl(\G_{0}(y_{\omega_1, t_1})-\G_{0}(y_{\omega_2, t_2})\bigr)^2
\right]
=
\Var{y_{\omega_1, t_1}(X)-y_{\omega_2, t_2}(X)}.
$$
Expand the variance term:
$$
\Var{y_{\omega_1, t_1} - y_{\omega_2, t_2}}
%=
%\Ebig{(y_{\omega_1, t_1} - y_{\omega_2, t_2})^2}
%-
%\left(\E{y_{\omega_1, t_1} - y_{\omega_2, t_2}}\right)^2 
\le 
\Ebig{(y_{\omega_1, t_1} - y_{\omega_2, t_2})^2} 
= 
\|y_{\omega_1,t_1}-y_{\omega_2,t_2}\|_{L^2(P)}^2.
$$

Observation: Since $y_{\omega_1, t_1}$ and $y_{\omega_2, t_2}$ are indicator functions, $(y_{\omega_1, t_1} - y_{\omega_2, t_2})^2 = 1_{A\triangle B}$,
Therefore,
$$
\E[(y_{\omega_1, t_1} - y_{\omega_2, t_2})^2] 
= P(A \triangle B), \ \text{ and } \ \E[y_{\omega_1, t_1} - y_{\omega_2, t_2}] = P(A) - P(B),
$$
where $A\triangle B \defin (A\setminus B)\cup(B\setminus A)$ denotes the symmetric difference. Combining these expressions,
$$
\E{
\bigl(\G_{0}(y_{\omega_1,t_1})-\G_{0}(y_{\omega_2,t_2})\bigr)^2
\right]
=
P\left(\bar{B}(\omega_1,t_1)\triangle \bar{B}(\omega_2,t_2)\right)
-
\bigl(P(\bar{B}(\omega_1,t_1)) - P(\bar{B}(\omega_2,t_2))\bigr)^2.
$$
\end{proof}
\end{comment}

\begin{comment}
%----------------------------%
\begin{lemma}\label{ap:lem_separable_support_G_old_simple}
%----------------------------%
    The Gaussian process $\G_{0}(\omega, t)$ has a separable support in the space $\ell^{\infty}(\Omega \times \R_{\ge 0})$.
\end{lemma}
\begin{proof}
    The Gaussian process $\G_{0}(\omega, t)$ takes values in the space $\ell^{\infty}(\Omega \times \R_{\ge 0})$. By Lemma \ref{ap:lem_continuity_sample_paths_G}, the sample paths of $\G_{0}(\omega, t)$ are almost surely uniformly continuous on $(\Omega \times \R_{\ge 0}, d_{\Omega \times \R_{\ge 0}})$. Therefore, the support of $\G_{0}(\omega, t)$ is contained in the space $C_u(\Omega \times \R_{\ge 0})$ of bounded uniformly continuous functions from $\Omega \times \R_{\ge 0}$ to $\R$. The Banach space $C_u(\Omega \times \R_{\ge 0})$ is separable if (and only if)
$(\Omega \times \R_{\ge 0}, d_{\Omega \times \R_{\ge 0}})$ is totally bounded \citep[see, e.g., page 17 of][]{Gine2021}. Since $(\Omega \times \R_{\ge 0}, d_{\Omega \times \R_{\ge 0}})$ is a totally bounded metric space ($\Omega$ is totally bounded by assumption and $\R_{\ge 0}$ is totally bounded), then $C_u(\Omega \times \R_{\ge 0})$ is separable. Hence, the support of $\G_{0}(\omega, t)$ is also separable.
\end{proof}
\end{comment}

\begin{comment}
\textcolor{blue}{
Lemma \ref{ap:lem_donsker_gc_slutsky} is included to illustrate where separability is needed in Slutsky’s lemma to obtain Glivenko--Cantelli (in probability) from the Donsker property. This result and its almost sure version are classical \citep[see, e.g.,][]{VanderVaartWellner2023}. This lemma will be removed.
}
%----------------------------%
\begin{lemma}[a Donsker-to-Glivenko--Cantelli implication in probability]
%----------------------------%
\label{ap:lem_donsker_gc_slutsky}
Let $\cF$ be a $P$-Donsker class of measurable functions, that is,
$$
\sqrt{n}\,(P_n - \mathsf P) \convl \G
\quad\text{in } \ell^\infty(\cF),
$$
where the limiting process $\G$ is separable.
Then $\cF$ is Glivenko--Cantelli in probability, that is,
$$
\sup_{f\in\cF} \bigl|(P_n - \mathsf P)f\bigr|
\;\convp\; 0 .
$$
%Add P^* instead? We have a supremum over an uncountable set in our case.
\end{lemma}

\begin{proof}
Let $a_n = n^{-1/2}$, so that $a_n \to 0$ deterministically.
By Slutsky’s lemma for random elements in $\ell^\infty(\cF)$,
applied to the pair $(\sqrt n(P_n-P), a_n)$, we obtain
$$
(P_n - \mathsf P)
=
a_n \sqrt n(P_n - \mathsf P)
\;\convp\; 0
\quad\text{in } \ell^\infty(\cF).
$$
Equivalently,
$$
\sup_{f\in\cF} \bigl|(P_n - \mathsf P)f\bigr|
\;\convp\; 0 ,
$$
which proves the claim.
\end{proof}
%-----------------------%
\end{comment}

Using Lemmas \ref{ap:lem_bound_int_g_epsilon}--\ref{ap:lem_continuity_sample_paths_G_rho2}, Theorem \ref{thm:test_statistics_simp} is proved.

%-----------------------%
\begin{proof}[Proof of Theorem \ref{thm:test_statistics_simp}]
\begin{comment}
With a telescopic trick:
\begin{align*}
    T_n^{\mathrm{CM}} &\defin n \int_{\Omega\times\R_{\ge0}} \left(F_{n,\omega}(t) - F_{\omega}^{\mu_0}(t)\right)^2  \rd F{n,\omega}(t)\, \rd \mu(\omega)\\
    &=n \int_{\Omega\times\R_{\ge0}} \left(F_{n,\omega}(t) - F_{\omega}^{\mu_0}(t)\right)^2 \rd F_{\omega}^{\mu_0}(t)\, \rd \mu(\omega)+n \int_{\Omega\times\R_{\ge0}} \left(F_{n,\omega}(t) - F_{\omega}^{\mu_0}(t)\right)^2  d( F_{n,\omega}-F_{\omega}^{\mu_0})(t)\, \rd \mu(\omega).
\end{align*}
Now we do a simple bounding argument to show the second integral is $o_P(1)$. Page 88 in Folland's Real Analysis (Exercise 3, \url{https://apachepersonal.miun.se/~andrli/Bok.pdf}) states: For $\mu$ a signed measure and $f\in \rho_{1,y}$,
$$
\left|\int f(x) \,\rd \mu(x)\right|\leq \int |f|\, \rd |\mu|(x),
$$
where $|\mu|$ is the total variation of $\mu$. The total variation of a measurable subset $E$ is, in particular,
$$
|\mu|(E)=\sup_{\text{Partitions $E_1,\ldots,E_m$ of $E$}}\lrb{\sum_{i=1}^m|\mu(E_j)|}
$$
(Exercise 8).

We then have
\begin{align*}
     \left|\int \left(F_{n,\omega}(t) - F_{\omega}^{\mu_0}(t)\right)^2  \rd (F_{n,\omega}-F_{\omega}^{\mu_0})(t)\right|&\leq \int \left(F_{n,\omega}(t) - F_{\omega}^{\mu_0}(t)\right)^2  d|F_{n,\omega}-F_{\omega}^{\mu_0}|(t)\\
     &\leq \sup_t \big|F_{n,\omega}(t) - F_{\omega}^{\mu_0}(t)\big|^2 \int   d|F_{n,\omega}-F_{\omega}^{\mu_0}|(t)\\
     &=\sup_t \big|F_{n,\omega}(t) - F_{\omega}^{\mu_0}(t)\big|^2 |F_{n,\omega}-F_{\omega}^{\mu_0}|(\R)\\
     &\leq 2\sup_t \big|F_{n,\omega}(t) - F_{\omega}^{\mu_0}(t)\big|^2\\
     &=o_P(1).
\end{align*}
Above, we have denoted by $|F_{n,\omega}-F_{\omega}^{\mu_0}|$ to the total variation $|\mu|$ of the signed measure $\mu=P_n-\mu_\omega^{\theta_0}$ induced by $F_{n,\omega}(t)=P_n((-\infty,t])$ and $F_{\omega}^{\mu_0}(t)=\mu_\omega^{\theta_0}((-\infty,t])$. The bound on $|\mu|(\R)$ is easily seen:
\begin{align*}
    |\mu|(\R)&=\sup \lrb{\sum_{i=1}^m|\mu(E_j)|}=\sup \lrb{\sum_{i=1}^m|P_n(E_j)-\mu_\omega^{\theta_0}(E_j)|}\\
    &\leq \sup \lrb{\sum_{i=1}^m|P_n(E_j)+\mu_\omega^{\theta_0}(E_j)|}= \sup \lrb{\sum_{i=1}^mP_n(E_j)+\sum_{i=1}^m\mu_\omega^{\theta_0}(E_j)}\\
    &= \sup \lrb{1+1}=2.
\end{align*}
\end{comment}
We divide the proof into its two parts, \ref{thm:test_statistics_simp_i} and \ref{thm:test_statistics_simp_ii}.

\emph{Proof of \ref{thm:test_statistics_simp_i}.} The continuous mapping theorem gives the asymptotics of $T_n^{\mathrm{KS}}$. 
Since we can identify $\G_{n,\mu}^{\mu_0}(\omega,t) = \G_{n,\mu}^{\mu_0}(y_{\omega,t})$ (see the definition of $\G_{n,\mu}^{\mu_0}$ in Section \ref{sec:null_asymptotics_simple}), it holds that
% because the processes take the same values, the suprema are the same
$$
\sup_{(\omega,t)\in\Omega\times\R_{\ge0}} |\G_{n,\mu}^{\mu_0}(\omega,t)|=\sup_{y\in\cY}|\G_{n,\mu}^{\mu_0}(y)|.
$$
The mapping $\sup_{\cY}: \ell^\infty(\cY) \to \R$ defined by $f \mapsto \sup_{\cY} |f(y)|$ for every $f \in \ell^\infty(\cY)$ is in fact $1$-Lipschitz, hence it is continuous (with respect to the sup-norm). 


\emph{Proof of \ref{thm:test_statistics_simp_ii}.} Focus now on $T_n^{\mathrm{CM}}$. We know that $\G_0(\omega, t)$ is separable by Remark \ref{rem:separability_G_0}.
Therefore, Skorohod's construction \citep[see Theorem  1.10.3 in][page 58]{VanderVaartWellner2023} applies, which ensures that 
$\G_{n,\mu}^{\mu_0}(\omega,t)$ converges almost uniformly to $\G_{0}(\omega, t)$.
Egorov's theorem \citep[Lemma 1.9.2 in][]{VanderVaartWellner2023} states the equivalence of almost uniform and outer almost sure convergence for sequences of Borel measurable random variables.
Therefore, there exist random variables (stochastic processes) $(\G_{n,\mu}^{\mu_0})^\star(\omega, t)$ and $\G_0^\star(\omega, t)$ defined on the same probability space such that 
\begin{equation}\label{eq:aux_cvm_sup_convas}
 \sup_{(\omega,t)\in \Omega \times \R_{\ge 0}} \left| (\G_{n,\mu}^{\mu_0})^\star(\omega, t) - \G_0^\star(\omega, t) \right| \convas 0.   
\end{equation}
For simplicity, we will use the same notation for $(\G_{n,\mu}^{\mu_0})^\star$ and $\G_0^\star$ as the original processes.
See Definition 1.9.1 in \cite{VanderVaartWellner2023} for almost uniform and outer almost sure convergence. 
%---------------------------%
\begin{remark}
%---------------------------%
The almost sure convergence result \eqref{eq:aux_cvm_sup_convas} follows from the definition of outer almost sure convergence, using the supremum distance.
Observe that the metric space implicit in the weak convergence \eqref{eq:thm_5.1_chendubeymuller} is $\ell^\infty(\cY)$ endowed with the supremum norm (the metric is implicit in the concept of weak convergence through the boundedness and continuity of the functions).
\end{remark}
%Similarly, since $F_{n,\omega}(t) \convl F_{\omega}(t)$, by Skorohod's construction, there exist $\tilde{\hat{F}}_{\omega}(t)$ and $\tilde{F}_{\omega}(t)$ such that $\tilde{\hat{F}}_{\omega}(t) \convas \tilde{F}_{\omega}(t)$ a.s. $\ell^\infty(\Omega \times \R_{\ge 0})$.  
Now, observe that under $H_0$,
\begin{equation}\label{eq:cvm_split}
\begin{aligned}
\Bigg| \int_{\Omega \times \R_{\geq 0}} &(\G_{n,\mu}^{\mu_0}(\omega, t))^2 \, \rd F_{n,\omega}^{\mu_0}(t) \, \rd P_n(\omega) - \int_{\Omega \times \R_{\geq 0}} (\G_{0}(\omega, t))^2 \, \rd F_{\omega}^{\mu_0}(t) \, \rd \mu_0(\omega) \Bigg| \\
&\le
\bigg| \int_{\Omega \times \R_{\ge0}} \left((\G_{n,\mu}^{\mu_0}(\omega, t))^2 - (\G_{0}(\omega, t))^2  \right) \, \rd F_{n,\omega}^{\mu_0}(t) \, \rd P_n(\omega) \bigg| 
\\ &+
\left| \int_{\Omega \times \R_{\ge0}}  (\G_{0}(\omega, t))^2 \, \rd \left(F_{n,\omega}^{\mu_0}(t)\, P_n(\omega) - F_{\omega}^{\mu_0}(t) \, \mu_0 (\omega)\right) \right|. 
\end{aligned}
\end{equation}
For the first term of the RHS of \eqref{eq:cvm_split}, one has
$$
\begin{aligned}
\Bigl| \int_{\Omega \times \R_{\ge0}} \big(&(\G_{n,\mu}^{\mu_0}(\omega, t))^2 - (\G_{0}(\omega, t))^2 \big) \, \rd F_{n,\omega}^{\mu_0}(t) \, \rd P_n(\omega) \Bigr| \\
&\le 
\int_{\Omega \times \R_{\ge0}} \left| (\G_{n,\mu}^{\mu_0}(\omega, t))^2 - (\G_{0}(\omega, t))^2 \right| \, \rd F_{n,\omega}^{\mu_0}(t) \, \rd P_n(\omega)\\
&\le 
\sup_{(\omega, t) \in \Omega\times \R_{\ge0}} \left| (\G_{n,\mu}^{\mu_0}(\omega, t))^2 - (\G_{0}(\omega, t))^2 \right| \int_{\Omega \times \R_{\ge0}}  \, \rd F_{n,\omega}^{\mu_0}(t) \, \rd P_n(\omega)\\
&=
%d\hatF is a CDF
\sup_{(\omega, t) \in \Omega\times \R_{\ge0}} \left| (\G_{n,\mu}^{\mu_0}(\omega, t))^2 - (\G_{0}(\omega, t))^2 \right|.
\end{aligned}
$$
Now, observe that
\begin{align*}
\left\|
(\G_{n,\mu}^{\mu_0})^2-\G_0^2
\right\|_\infty
&\le
\left\|
\G_{n,\mu}^{\mu_0}-\G_0
\right\|_\infty^2
+
2\left\|\G_0\right\|_\infty
\left\|
\G_{n,\mu}^{\mu_0}-\G_0
\right\|_\infty
\convas 0,
\end{align*}
by \eqref{eq:aux_cvm_sup_convas} and the fact that $\G_{0}$ has almost surely bounded sample paths. 
Boundedness holds because $\G_0$ is a map into $\ell^\infty(\cY)$.
\begin{comment}
\textcolor{forestgreen}{
\begin{equation}
\begin{aligned}\label{eq:aux_cvm_first}
\int_{\Omega \times \R_{\ge0}} \left(\G_{n,\mu}^{\mu_0}(\omega, t)^2 - \G_{0}(\omega, t)^2 \right) \, \rd F_{n,\omega}^{\mu_0}(t) \, \rd P_n(\omega) 
& =
\int_{\Omega \times \R_{\ge0}} \left(\G_{n,\mu}^{\mu_0}(\omega, t) - \G_{0}(\omega, t) \right)^2 \, \rd F_{n,\omega}^{\mu_0}(t) \, \rd P_n(\omega) \\
&+ 2\int_{\Omega \times \R_{\ge0}} \G_{0}(\omega, t) \left(\G_{n,\mu}^{\mu_0}(\omega, t) - \G_{0}(\omega, t) \right) \, \rd F_{n,\omega}^{\mu_0}(t) \, \rd P_n(\omega)
\end{aligned}
\end{equation}
For the first term in the RHS of \eqref{eq:aux_cvm_first}, one has
$$
\begin{aligned}
 \int_{\Omega \times \R_{\ge0}} \left(\G_{n,\mu}^{\mu_0}(\omega, t) - \G_{0}(\omega, t) \right)^2 \, \rd F_{n,\omega}^{\mu_0}(t) \, \rd P_n(\omega) &\le
     \sup_{(\omega, t) \in \Omega\times \R_{\ge0}} \left( \G_{n,\mu}^{\mu_0}(\omega, t) - \G_{0}(\omega, t) \right)^2 \int_{\Omega \times \R_{\ge0}}  \, \rd F_{n,\omega}^{\mu_0}(t) \, \rd P_n(\omega) \\
    & = \sup_{(\omega, t) \in \Omega\times \R_{\ge0}} \left( \G_{n,\mu}^{\mu_0}(\omega, t) - \G_{0}(\omega, t) \right)^2 \convd 0
\end{aligned}
$$
by \eqref{eq:thm_5.1_chendubeymuller} and a simple application of the continuous mapping theorem. For the second term in the RHS of \eqref{eq:aux_cvm_first}, one has
$$
\begin{aligned}
\int_{\Omega \times \R_{\ge0}} &\left|\G_{0}(\omega, t)\right| \left|\G_{n,\mu}^{\mu_0}(\omega, t) - \G_{0}(\omega, t) \right| \, \rd F_{n,\omega}^{\mu_0}(t) \, \rd P_n(\omega) \\
&\le
\sup_{(\omega, t) \in \Omega\times \R_{\ge0}} \left|\G_{0}(\omega, t)\right| \sup_{(\omega, t) \in \Omega\times \R_{\ge0}} \left|\G_n(\omega, t) - \G_{0}(\omega, t) \right| \\
& \quad \times \int_{\Omega \times \R_{\ge0}} \, \rd F_{n,\omega}(t) \, \rd \mu(\omega) \\
& \le \sup_{(\omega, t) \in \Omega\times \R_{\ge0}} \left|\G_{0}(\omega, t)\right| \sup_{(\omega, t) \in \Omega\times \R_{\ge0}} \left|\G_n(\omega, t) - \G_{0}(\omega, t) \right| \\
&\quad \convd 0
\end{aligned}
$$
since by \eqref{eq:thm_5.1_chendubeymuller}, the first term is bounded in probability (Donsker) and the second term converges to zero. Both terms in the RHS of \eqref{eq:aux_cvm_first} converge to zero in probability, so
$$
\int_{\Omega \times \R_{\ge0}} \left(\G_n(\omega, t)^2 - \G_{0}(\omega, t)^2 \right) \, \rd F_{n,\omega}(t) \, \rd \mu(\omega) \convd 0,
$$
which proves the claim.
}
\end{comment}

Now, focus on the second term of the RHS of \eqref{eq:cvm_split}. The sample paths of the Gaussian process $\G_{0}(\omega, t)$ are bounded and uniformly continuous almost surely. Uniform continuity holds by Lemma \ref{ap:lem_continuity_sample_paths_G_rho2}. 
\begin{comment}
For every sample path and every $\omega\in\Omega$, an application of the Helly-Bray Theorem \citep[see, e.g., page 117 in][]{Rao1973} shows that
\begin{align}\label{eq:aux_rhs_cvm_conv_R}
\left| \int_{\R_{\ge0}}  \G_{0}(\omega, t)^2 \,\rd \left(F_{n,\omega}(t) - F_{\omega}^{\mu_0}(t)\right) \right| \convas 0 \text{ as } n\to \infty.
\end{align}
\end{comment}
An application of Lemma \ref{ap:lem_DCT_P_hat} to every sample path gives
\begin{align*}
    \int_{\Omega \times \R_{\ge0}}  \G_{0}(\omega, t)^2 \, \rd \left(F_{n,\omega}^{\mu_0}(t)\, P_n(\omega) - F_{\omega}^{\mu_0}(t) \, \mu_0(\omega)\right) \convas 0 \text{ as } n\to \infty,
\end{align*}
which is the second term of the RHS of \eqref{eq:aux_cvm_sup_convas}. Thus, almost surely, 
$$
T_n^{\mathrm{CM}} \convl \int_{\Omega \times \R_{\geq 0}} \G_{0}(\omega, t)^2 \,\rd F_{\omega}^{\mu_0}(t) \, \rd \mu_0(\omega) \text{ as } n \to \infty.
$$
This concludes the proof.
\end{proof}
%-----------------------%

%-----------------------%
\subsection{Proof of Theorem \ref{thm:test_statistics_comp}}\label{ap:sec_proofs_null_asymptotics_thm_test_statistics_comp}
%-----------------------%

This section contains the proof of Theorem \ref{thm:test_statistics_comp}, together with the Lemmas needed in the proof. 
Observe that Lemmas \ref{ap:lem_bound_int_g_epsilon}--\ref{ap:lem_DCT_P_hat} also hold when $\mu_{0}$ is replaced by $\mu_{\btheta_0}$.
% In what follows, we define $d_{\Omega \times \R_{\ge 0}}$ as in \eqref{ap:eq_distances}, 
% and redefine $\rho_{2,y}$ in terms of the class
% $\cF= \{f_{\omega,t} : (\omega,t) \in \Omega \times \R_{\ge 0}\}$.
% That is, for $p\ge1$ and any pair $((\omega_1,t_1), (\omega_2, t_2))$, let
% \begin{align}\label{eq:ap_distance_Lp_composite_old_sketch}
% \rho_{p,y}\big((\omega_1,t_1), (\omega_2, t_2)\big)
% \defin
% \left(\mathsf{E}\big|f_{\omega_1,t_1}(X) - f_{\omega_2, t_2}(X)\big|^p\right)^{\frac{1}{p}} .
% \end{align}



\begin{comment}
%----------------------------%
\begin{lemma}\label{ap:lem_dist_prod_0_dist_L1_0_comp}
%----------------------------%
   Define $d_{\Omega \times \R_{\ge 0}}$ as in \eqref{ap:eq_distances}, and redefine $\rho_{2,y}$ in terms of the (Donsker) class $\cF= \{f_{\omega,t} : (\omega,t) \in \Omega \times \R_{\ge 0}\}$. That is, for $p\ge1$ and any pair $((\omega_1,t_1), (\omega_2, t_2))$, let
   \begin{align}\label{eq:ap_distance_Lp_composite_old_cont}
   \rho_{p,y}\big((\omega_1,t_1), (\omega_2, t_2)\big) \defin \left(\mathsf{E}\big|f_{\omega_1,t_1}(X) - f_{\omega_2, t_2}(X)\big|^p\right)^{\frac{1}{p}}.
   \end{align}
Assume that the map $(\omega,t) \mapsto \dot{F}_{\omega}^{\btheta_0}(t)$ is continuous with respect to $d_{\Omega\times\R_{\ge0}}$ (\textcolor{blue}{note: we assumed uniform continuity with respect to $\omega$ for Proposition \ref{prop:sup_max_ks_composite}}). Under Assumptions \ref{b2}, \ref{c3} and \ref{d}, if $d_{\Omega \times \R_{\ge 0}}\big((\omega_k,t_k), (\omega, t)\big)\allowbreak \to 0$ as $k\to\infty$, then also $\rho_{2,y}\big((\omega_k,t_k), \allowbreak (\omega, t)\big) \to 0$ as $k\to\infty$.
\end{lemma}
%----------------------------%
\begin{proof}
Recall that
$$
f_{\omega,t}(x)
=
y_{\omega,t}(x)
-
F_{\omega}^{\mu_0}(t)
-
g_{\omega,t}(x),
\qquad
g_{\omega,t}(x)
=
- \dot{F}_{\omega}^{\btheta_0}(t)^\top
\bV_{\btheta_0}^{-1}
\bpsi_{\btheta_0}(x).
$$
Note that $f_{\omega, t}(x)$ depends on $x$ only through $y_{\omega,t}(x)$ and $\bpsi_{\btheta_0}(x)$.

Let $\{(\omega_k,t_k)\}_{k\ge1}$ be a sequence such that
$$
d_{\Omega \times \R_{\ge 0}}\big((\omega_k,t_k),(\omega,t)\big) \to 0,
$$
that is, $d(\omega_k,\omega)\to0$ and $t_k\to t$, and write
\begin{align*}
f_{\omega_k,t_k}(X)-f_{\omega,t}(X)
=
\big(y_{\omega_k,t_k}(X)-y_{\omega,t}(X)\big)
-
\big(F_{\omega_k}^{\btheta_0}(t_k)-F_{\omega}^{\mu_0}(t)\big) 
-
\big(g_{\omega_k,t_k}(X)-g_{\omega,t}(X)\big).
\end{align*}
To show $\rho_{2,y}\big((\omega_k,t_k),(\omega,t)\big) \to 0$, by the inequality $(a+b+c)^2 \le 3(a^2+b^2+c^2)$ (Cauchy-Schwarz), it suffices to show that each of the three squared terms converges to zero in mean. By %Lemma~\ref{ap:lem_dist_prod_0_dist_L1_0}, 
Lemma old continuity BUT WE WANT UNIFORM!
it follows that $\Ebig{(y_{\omega_k,t_k}(X)-y_{\omega,t}(X))^2} \to 0$.

For the second term, observe that
$$
\Ebig{\big(F_{\omega_k}^{\btheta_0}(t_k)-F_{\omega}^{\mu_0}(t)\big)^2}
=
\big(F_{\omega_k}^{\btheta_0}(t_k)-F_{\omega}^{\mu_0}(t)\big)^2,
$$
so it is enough to note that $(\omega,t)\mapsto F_{\omega}^{\mu_0}(t)$ is continuous. Observe that $\big|F_{\omega_k}^{\btheta_0}(t_k)-F_{\omega}^{\mu_0}(t)\big| \le \mathsf{E}\big|y_{\omega_k,t_k}(X)-y_{\omega,t}(X)\big|$. The result then follows by Lemma \ref{ap:lem_relation_L2_prod} and the relation \eqref{eq:L_1_L_2_relation}. 

The third term is
\begin{align*}
g_{\omega_k,t_k}(X)-g_{\omega,t}(X)
=
-
\Big(
\dot{F}_{\omega_k}^{\btheta_0}(t_k)
-
\dot{F}_{\omega}^{\btheta_0}(t)
\Big)^\top
\bV_{\btheta_0}^{-1}
\bpsi_{\btheta_0}(X).
\end{align*}
Hence,
\begin{align}\label{eq:g_diff_bound}
\Ebig{\big(g_{\omega_k,t_k}(X)-g_{\omega,t}(X)\big)^2}
\le
\Big\|
\dot{F}_{\omega_k}^{\btheta_0}(t_k)
-
\dot{F}_{\omega}^{\btheta_0}(t)
\Big\|^2
\,
\Big\|
\bV_{\btheta_0}^{-1}
\Big\|^2
\,
\Ebig{\|\bpsi_{\btheta_0}(X)\|^2}.
\end{align}

The first term in the RHS of \eqref{eq:g_diff_bound} converges to zero when $(\omega_k,t_k) \to (\omega,t)$ due to the assumption of continuity of $\dot{F}_{\omega}^{\btheta_0}(t)$, and by Assumption~\ref{c4} the matrix $\bV_{\btheta_0}$ is non-singular. By Assumption~\ref{c3}, the last expectation in \eqref{eq:g_diff_bound} is finite. 

We conclude that $\Ebig{\big(f_{\omega_k,t_k}(X)-f_{\omega,t}(X)\big)^2} \to 0$, that is, $\rho_{2,y}\big((\omega_k,t_k),(\omega,t)\big)\to0$.
\end{proof}
\end{comment}



\begin{comment}
Lemma \ref{ap:lem_aux_omega_0_delta_0_composite} extends Lemma \ref{ap:lem_aux_omega_0_delta_0} for the composite null hypothesis.
%----------------------------%
\begin{lemma}[replaced by Lemma \ref{ap:lem_continuity_sample_paths_G_rho2_comp}]\label{ap:lem_aux_omega_0_delta_0_composite}
%----------------------------%
\textcolor{blue}{    Let $\cF = \{f_{\omega,t}: (\omega,t) \in \Omega \times \R_{\ge 0}\}$, with $f_{\omega,t}$ defined in \eqref{eq:f_omega_t}. 
    Define $d_{\Omega \times \R_{\ge 0}}$ as in \eqref{eq:ap_distance_product}, and redefine $\rho_{2,y}$ in terms of the (Donsker) class $\cF= \{f_{\omega,t} : (\omega,t) \in \Omega \times \R_{\ge 0}\}$. 
    That is, for $p\ge1$ and any pair $((\omega_1,t_1), (\omega_2, t_2))$, let
    \begin{align}\label{eq:ap_distance_Lp_composite_old_entropy}
    \rho_{p,y}\big((\omega_1,t_1), (\omega_2, t_2)\big)
    \defin
    \left(\mathsf{E}\big|f_{\omega_1,t_1}(X) - f_{\omega_2, t_2}(X)\big|^p\right)^{\frac{1}{p}}.
    \end{align}
    Define $\omega(\delta) \defin \int_0^\delta \sqrt{\log N(\varepsilon, \cF, \rho_{2,y})} \allowbreak \, \rd \varepsilon,$ with $\rho_{2,y}$ as defined in \eqref{eq:ap_distance_Lp_composite_old_entropy}.
    Assume that $\log(N(\varepsilon, \Omega, d)) = O(\varepsilon^{-\alpha})$ for some $\alpha<1$. Assume also the conditions \ref{cp2}, \ref{c4}, and \ref{d}. Then, $\omega(\delta) \to 0$ as $\delta \to 0$.}
\end{lemma}

\begin{proof}
    For each function $f_{\omega,t}\in\cF$, one has
    $
    f_{\omega,t} = y_{\omega,t} - F_{\omega}^{\mu_0}(t) - g_{\omega,t},
    $
    with $g_{\omega,t}$ as defined in \eqref{eq:g_omega_t}. If $[l,u]$ is an $\varepsilon$-bracket for $\cY = \{y_{\omega,t}: (\omega,t) \in \Omega \times \R_{\ge 0}\}$ and $[l',u']$ is an $\varepsilon$-bracket for $\mathcal{G} = \{g_{\omega,t}: (\omega,t) \in \Omega \times \R_{\ge 0}\}$, then $[l-l', u-u']$ is a $2\varepsilon$-bracket for $\cF$. Therefore,
    \begin{equation}\label{eq:bracket_max}
    N_{[]}(\varepsilon, \cF, L^2(\mu)) 
    \le
    N_{[]}\left(\frac{\varepsilon}{2}, \cY, L^2(\mu)\right)
    N_{[]}\left(\frac{\varepsilon}{2}, \mathcal{G}, L^2(\mu)\right).
    \end{equation}
    Hence, it suffices to control the entropy integrals of $\cY$ and $\mathcal{G}$ separately. Lemma~\ref{ap:lem_aux_omega_0_delta_0} deals with $N_{[]}(\varepsilon, \cY, L^2(\mu))$. 
    Thus, we focus only on $\mathcal{G}$.

    Write
    $
    g_{\omega,t}(x)=\sum_{j=1}^{p}\alpha_j(\omega,t)\,\psi_j(x),
    $
    where $\psi_j$ are the components of $\psi_{\btheta_0}$ and
    $
    \boldsymbol{\alpha}(\omega,t)
    =
    -\dot{F}_{\omega}^{\btheta_0}(t)^{\top}
    \bigl(\bV_{\btheta_0}\bigr)^{-1}.
    $
    By hypotheses~\ref{c4} and~\ref{d}, the set
    $\{\boldsymbol{\alpha}(\omega,t):(\omega,t)\in \Omega \times \R_{\ge 0}\}$
    is bounded in $\R^{p}$. Moreover, by Assumption~\ref{cp2},
    $$
    \sup_{(\omega,t)\in\Omega\times\R_{\ge0}}
    \|g_{\omega,t}\|_{L^2(\mu)}
    \le
    \sup_{(\omega,t)\in\Omega\times\R_{\ge0}}\|\boldsymbol{\alpha}(\omega,t)\|
    \Big\|\big(\mathsf{E}\dot\psi_{\btheta_0}(X)\big)^{-1}\Big\|
    \|\psi_{\btheta_0}\|_{L^2(\mu)}
    <\infty.
    $$
    Hence, $\mathcal{G}$ is a bounded subset of a fixed $p$-dimensional linear subspace of $L^2(\mu)$.

    Let $R>0$ be such that $\|g\|_{L^2(\mu)}\le R$ for all $g\in\mathcal{G}$. For $0<\varepsilon\le R$, standard entropy bounds in finite-dimensional normed spaces yield
    $$
    \log N_{[]}(\varepsilon,\mathcal{G},L^2(\mu))
    \le
    Cp\log\left(\frac{R}{\varepsilon}\right) + Cp,
    $$
    where $C$ is a universal constant. Consequently,
    $$
    \int_0^{\infty}\sqrt{\log N_{[]}(\varepsilon,\mathcal{G},L^2(\mu))}\, \rd \varepsilon
    =
    \int_0^{2R}\sqrt{\log N_{[]}(\varepsilon,\mathcal{G},L^2(\mu))}\, \rd \varepsilon
    \le
    \int_0^{2R}\sqrt{Cp\log\left(\frac{2R}{\varepsilon}\right) + Cp}\, \rd \varepsilon
    <\infty.
    $$

    Since $N(\varepsilon, \cF, L^2(\mu)) \le N_{[]}(2\varepsilon, \cF, L^2(\mu))$, then $\int_0^{\infty} \sqrt{\log N_{[]}(\varepsilon, \cF, L^2(\mu))} \, d\varepsilon < \infty$ implies $\int_0^{\infty} \sqrt{\log N(\varepsilon, \Omega\times \R_{\ge 0}, L^2(\mu))} \,\allowbreak d\varepsilon < \infty$. 
    Because $\sqrt{\log N(\varepsilon, \Omega, \rho_{2,y})}$ is integrable and the Lebesgue integral is a continuous functional, it follows that $\omega(\delta) \to 0$ as $\delta \to 0$. 
    This concludes the result.
\end{proof}
\end{comment}

\begin{comment}
%----------------------------%
\begin{lemma}\label{ap:lem_dist_prod_0_dist_L1_0_old_comp}
%----------------------------%
   Let $d_{\Omega \times \R_{\ge 0}}$  denote the product distance on $(\Omega, d) \times (\R_{\ge 0}, |\cdot|)$, and let $d_{L^1(P)}$ be the $L^1(P)$ distance induced by a random object $X$ taking values in $\Omega$. That is,
   $$
    d_{\Omega \times \R_{\ge 0}}\big((\omega,t), (\omega^\prime, t^\prime)\big) \defin d(\omega, \omega^\prime) + |t - t^\prime|, \quad d_{L^1(P)}\big((\omega,t), (\omega^\prime, t^\prime)\big) \defin \mathsf{E}\big|y_{\omega,t}(X) - y_{\omega^\prime, t^\prime}(X)\big|.
   $$
Assume that the distribution of $d(\omega,X)$ is continuous for every $\omega\in \Omega$. If $d_{\Omega \times \R_{\ge 0}}\big((\omega_k,t_k), (\omega, t)\big) \to 0$ as $k\to\infty$, then also $d_{L^1(P)}\big((\omega_k,t_k), (\omega, t)\big) \to 0$ as $k\to\infty$.
\end{lemma}
%----------------------------%
\begin{proof}
    \textit{Note: The proof is very similar to that of Lemma \ref{lem:continuity_omega}}.
    Consider a sequence $\{(\omega_k, t_k)\}_{k\ge 1}$ in $\Omega$ such that $d_{\Omega \times \R_{\ge 0}}((\omega_k, t_k), (\omega, t)) \to 0$. Obviously, $|t-t_k| \to 0$. Also, since $d(\cdot, x)$ is continuous for every $x \in \Omega$, then $d(\omega_k, x) \to d(\omega,x)$ as $k\to \infty$ for every $x \in \Omega$. 

    Fix $\delta \ge 0$, and consider $x \in \Omega$ such that $|d(\omega,x) - t| \neq \delta$. One has
    \begin{align}\label{eq:conv_indicators_old_comp}
    1_{\{|d(\omega_k,x) - t_k| \le \delta\}} \to 1_{\{|d(\omega,x) - t| \le \delta\}} \ \text{ as } k \to \infty.
    \end{align}
    To see this, first assume that $|d(\omega,x) - t| < \delta$, and let $\varepsilon = \frac{\delta - |d(\omega,x) - t|}{2} >0$. One has
    $$
    |d(\omega_k,x) - t_k| \le |d(\omega_k,x) - d(\omega, x)| +  |d(\omega,x) - t| + |t - t_k| < \delta
    $$
    for $k$ sufficiently large, so that $|d(\omega_k,x) - d(\omega, x)| \le \varepsilon$ and $|t - t_k| \le \varepsilon$. If instead, $|d(\omega,x) - t| > \delta$, let $\varepsilon = \frac{|d(\omega,x) - t| - \delta}{2} > 0$. Then, by applying the reverse triangle inequality,
    $$
    |d(\omega_k,x) - t_k| \ge |d(\omega, x) - t| - |d(\omega_k,x) - d(\omega,x)| - |t - t_k| \ge \delta
    $$    
    for $k$ sufficiently large. This proves \eqref{eq:conv_indicators_old_comp}.
    
    Observe that since $\Prob{|d(\omega_k,x) - t_k| = \delta} = 0$ (by the continuity of the distribution of $d(\omega,X)$ for every $\omega \in \Omega$), the convergence in \eqref{eq:conv_indicators_old_comp} holds $P$-a.s. Now, since $1_{\{|d(\omega_k,x) - t_k| \le \delta\}}$ is bounded for every $k\ge0$, an application of the DCT gives
    $$
    d_{L^1(P)}\big((\omega_k,t_k), (\omega, t)\big) = \int_{\Omega} \left|1_{\{|d(\omega_k,x) - t_k| \le \delta\}} - 1_{\{|d(\omega,x) - t| \le \delta\}} \right| \, dP(x) \to 0 \text{ as } k\to\infty.
    $$
    %That is, $h(\omega_k, \delta) \to h(\omega, \delta)$ when $\omega_k \to \omega$, i.e., $h(\cdot, \delta)$ is continuous on $\Omega$ for every $\delta\ge0$.
    This concludes the proof.
\end{proof}
\end{comment}


Lemma \ref{ap:lem_relation_L2_prod_composite} extends Lemma \ref{ap:lem_relation_L2_prod} for the composite null hypothesis.
%----------------------------%
\begin{lemma}\label{ap:lem_relation_L2_prod_composite}
%----------------------------%
Let $\rho_{2,f}$ and $\rho_{2,\tilde{\G}_0}$ be the $L^2(P)$ pseudometrics induced by the class $\cF$ and by the Gaussian process $\tilde{\G}_0$, respectively. That is, for $p\ge1$ and $(\omega,t),(\omega^\prime,t^\prime)\in\Omega\times\R_{\ge0}$, define
\begin{align*}
\rho_{p,f}\big((\omega,t),(\omega^\prime,t^\prime)\big)
&\defin
\Big(
\Ebig{\big|f_{\omega,t}(X)-f_{\omega^\prime,t^\prime}(X)\big|^p}
\Big)^{1/p}
\\
\rho_{p,\tilde{\G}_0}\big((\omega,t),(\omega^\prime,t^\prime)\big)
&\defin
\Big(
\Ebig{\big|\tilde{\G}_0(\omega,t)-\tilde{\G}_0(\omega^\prime,t^\prime)\big|^p}
\Big)^{1/p}.
\end{align*}
Assume that the map $(\omega,t)\mapsto \dot F_{\omega}^{\btheta_0}(t)$ is uniformly continuous with respect to $d_{\Omega\times\R_{\ge0}}$. Assume also conditions~\ref{cp2}, \ref{cp3}, and \ref{d}. Then, there exists a nondecreasing function $\phi:[0,\infty)\to[0,\infty)$ with $\phi(r)\to0$ as $r\to0$ such that, for all $(\omega,t),(\omega^\prime,t^\prime)\in\Omega\times\R_{\ge0}$,
\begin{align}\label{eq:rho_2_tilde_G_0_rho_2_f_composite}
\rho_{2,\tilde{\G}_0}\big((\omega,t),(\omega^\prime,t^\prime)\big)
= 
\rho_{2,f}\big((\omega,t),(\omega^\prime,t^\prime)\big)
\le
\phi\!\left(
d_{\Omega\times\R_{\ge0}}
\big((\omega,t),(\omega^\prime,t^\prime)\big)
\right).
\end{align}
In particular, any function that is uniformly continuous with respect to $\rho_{2,\tilde{\G}_0}$ is also uniformly continuous with respect to $d_{\Omega\times\R_{\ge0}}$.
\end{lemma}

\begin{proof}
We first prove the equality
$$
\rho_{2,\tilde{\G}_0}\big((\omega,t),(\omega^\prime,t^\prime)\big)
=
\rho_{2,f}\big((\omega,t),(\omega^\prime,t^\prime)\big).
$$
Since $\tilde{\G}_0$ is the centered Gaussian process indexed by the class $\cF$, and $\Ebig{f_{\omega,t}(X)}=0$ for every $(\omega,t)\in\Omega\times\R_{\ge0}$, its covariance is given by
$$
\Cov{\tilde{\G}_0(\omega,t)}{\tilde{\G}_0(\omega^\prime,t^\prime)}
=
\Ebig{f_{\omega,t}(X)f_{\omega^\prime,t^\prime}(X)}.
$$
Hence,
\begin{align*}
\rho_{2,\tilde{\G}_0}\big((\omega,t),(\omega^\prime,t^\prime)\big)^2
&=
\Ebig{\big(\tilde{\G}_0(\omega,t)-\tilde{\G}_0(\omega^\prime,t^\prime)\big)^2}
\nonumber
\\
&=
\Ebig{\big(f_{\omega,t}(X)-f_{\omega^\prime,t^\prime}(X)\big)^2}
\nonumber
\\
&=
\rho_{2,f}\big((\omega,t),(\omega^\prime,t^\prime)\big)^2.
\end{align*}

Now, prove the inequality in \eqref{eq:rho_2_tilde_G_0_rho_2_f_composite}. Recall that
$$
f_{\omega,t}(x)
=
y_{\omega,t}(x)-F_{\omega}^{\mu_{\btheta_0}}(t)+g_{\omega,t}(x),
\qquad
g_{\omega,t}(x)
=
\dot F_{\omega}^{\btheta_0}(t)^\top \bV_{\btheta_0}^{-1}\bpsi_{\btheta_0}(x).
$$
Fix $(\omega,t),(\omega^\prime,t^\prime)\in\Omega\times\R_{\ge0}$, and set
$$
r\defin d_{\Omega\times\R_{\ge0}}\big((\omega,t),(\omega^\prime,t^\prime)\big).
$$
By the triangle inequality,
\begin{align*}
\rho_{2,f}\big((\omega,t),(\omega^\prime,t^\prime)\big)
&=
\Big(
\Ebig{\big(f_{\omega,t}(X)-f_{\omega^\prime,t^\prime}(X)\big)^2}
\Big)^{1/2}
\\
&\le
\Big(
\Ebig{\big(y_{\omega,t}(X)-y_{\omega^\prime,t^\prime}(X)\big)^2}
\Big)^{1/2}
+
\big|F_{\omega}^{\mu_{\btheta_0}}(t)-F_{\omega^\prime}^{\mu_{\btheta_0}}(t^\prime)\big|
\\
&\quad +
\Big(
\Ebig{\big(g_{\omega,t}(X)-g_{\omega^\prime,t^\prime}(X)\big)^2}
\Big)^{1/2}.
\end{align*}

For the first term, Lemma~\ref{ap:lem_relation_L2_prod} gives
$$
\Ebig{\big(y_{\omega,t}(X)-y_{\omega^\prime,t^\prime}(X)\big)^2}
\le
4\bar\Delta r.
$$

For the second term, since
$$
F_{\omega}^{\mu_{\btheta_0}}(t)=\Ebig{y_{\omega,t}(X)},
$$
one has
\begin{align*}
\big|F_{\omega}^{\mu_{\btheta_0}}(t)-F_{\omega^\prime}^{\mu_{\btheta_0}}(t^\prime)\big|
\le
\Ebig{\big|y_{\omega,t}(X)-y_{\omega^\prime,t^\prime}(X)\big|}
=
\Ebig{\big(y_{\omega,t}(X)-y_{\omega^\prime,t^\prime}(X)\big)^2}
\le
4\bar\Delta\, r.
\end{align*}

Finally, for the third term, observe that
$$
g_{\omega,t}(x)-g_{\omega^\prime,t^\prime}(x)
=
\Big(
\dot F_{\omega}^{\btheta_0}(t)-\dot F_{\omega^\prime}^{\btheta_0}(t^\prime)
\Big)^\top
\bV_{\btheta_0}^{-1}\bpsi_{\btheta_0}(x).
$$
Hence, by Cauchy--Schwarz,
$$
\Big(
\Ebig{\big(g_{\omega,t}(X)-g_{\omega^\prime,t^\prime}(X)\big)^2}
\Big)^{1/2}
\le
\big\|
\dot F_{\omega}^{\btheta_0}(t)-\dot F_{\omega^\prime}^{\btheta_0}(t^\prime)
\big\|
\,
\|\bV_{\btheta_0}^{-1}\|
\,
\Big(\Ebig{\|\bpsi_{\btheta_0}(X)\|^2}\Big)^{1/2}.
$$
Recall that by Assumption~\ref{cp2}, $\mathsf{E}\|\psi_{\btheta_0}(X)\|^2<\infty$ and by Assumption~\ref{cp3} the matrix $\bV_{\btheta_0}$ is non-singular. 
Define the following modulus of continuity of the map $(\omega,t)\mapsto \dot F_{\omega}^{\btheta_0}(t)$:
$$
\eta_{\dot F}(r)
\defin
\sup\Big\{
\big\|\dot F_{\omega}^{\btheta_0}(t)-\dot F_{\omega^\prime}^{\btheta_0}(t^\prime)\big\| :
(\omega,t),(\omega^\prime,t^\prime) \in \Omega \times \R_{\ge0}, \, 
d_{\Omega\times\R_{\ge0}}\big((\omega,t),(\omega^\prime,t^\prime)\big)\le r
\Big\}.
$$
Thus,
$$
\Big(
\Ebig{\big(g_{\omega,t}(X)-g_{\omega^\prime,t^\prime}(X)\big)^2}
\Big)^{1/2}
\le
\eta_{\dot F}(r)\big\|\bV_{\btheta_0}^{-1}\big\|\Big(\mathsf{E}\big\|\psi_{\btheta_0}(X)\big\|^2\Big)^{1/2}
$$
and by the assumed uniform continuity of $(\omega,t)\mapsto \dot F_{\omega}^{\btheta_0}(t)$,
$\eta_{\dot F}(r)\to0$ as $r\to0$.

Combining the three bounds, we obtain
$$
\rho_{2,f}\big((\omega,t),(\omega^\prime,t^\prime)\big)
\le
2\sqrt{\bar\Delta r}
+
4\bar\Delta\, r
+
\|\bV_{\btheta_0}^{-1}\|
\Big(\Ebig{\|\bpsi_{\btheta_0}(X)\|^2}\Big)^{1/2}
\eta_{\dot F}(r).
$$
Thus, the result holds with
$$
\phi(r)
\defin
2\sqrt{\bar\Delta r}
+
4\bar\Delta\, r
+
\|\bV_{\btheta_0}^{-1}\|
\Big(\Ebig{\|\bpsi_{\btheta_0}(X)\|^2}\Big)^{1/2}
\eta_{\dot F}(r).
$$
It holds that $\phi$ is nondecreasing and $\phi(r)\to0$ as $r\to0$, which concludes the proof.
\end{proof}
%----------------------------%

\begin{comment}
%----------------------------%
\begin{lemma}[replaced by Lemma \ref{ap:lem_continuity_sample_paths_G_rho2_comp}]\label{ap:lem_continuity_sample_paths_G_comp}
%----------------------------%
\textcolor{blue}{   Assume that $\log(N(\varepsilon, \Omega, d)) = O(\varepsilon^{-\alpha})$ for some $\alpha<1$, and that conditions~\ref{cp2}, \ref{c4}, and \ref{d} hold. Assume that the map $(\omega,t)\mapsto \dot F_{\omega}^{\btheta_0}(t)$ is uniformly continuous with respect to $d_{\Omega\times\R_{\ge0}}$. 
The process $\tilde{\G}_0(\omega, t)$ has almost surely uniformly continuous sample paths on the metric space $(\Omega \times \R_{\ge 0}, d_{\Omega \times \R_{\ge 0}})$, considering the sample paths as functions from $(\Omega \times \R_{\ge 0}, d_{\Omega \times \R_{\ge 0}})$ to $(\R, |\cdot|)$.}
\end{lemma}
\begin{proof}
Since $\tilde{\G}_0(\omega, t)$ is a separable $L^2(P)$-subGaussian process on $(\Omega \times \R_{\ge 0}, d_{L^2(P)})$ (it is a Gaussian process), the function
$$
\omega(\delta)=\int_0^\delta \sqrt{\log N(\varepsilon, \Omega \times \R_{\ge 0}, d_{L^2(P)})} \rd \varepsilon
$$
is a modulus of continuity for $\{\tilde{\G}_0(\omega,t)\}_{(\omega,t)\in \Omega\times\R_{\ge0}}$. See, e.g., Theorem 5.32 in \cite{vanHandel}, page 140 in \cite{VanderVaartWellner2023} or \cite{Talagrand2021} for this result. This means that there is a random variable $K$ such that
% In particular, we have
% $$
% \mathbf{E}\left[\sup _{t, s \in T} \frac{X_t-X_s}{\omega(d(t, s))}\right]<\infty.
% $$
\begin{align}\label{eq:modulus_cont_comp}
\big|\tilde{\G}_0(\omega, t) - \tilde{\G}_0(\omega^\prime, t^\prime)\big|
\le K\,\omega\left(d_{L^2(P)}\big((\omega,t),(\omega^\prime,t^\prime)\big)\right)
\quad \text{for all } (\omega,t),(\omega^\prime,t^\prime)\in \Omega\times\R_{\ge0}.
\end{align}
Hence, it follows from Lemma~\ref{ap:lem_aux_omega_0_delta_0_composite} that $\omega(\delta)\to0$ as $\delta\to0$. In addition, by Lemma~\ref{ap:lem_relation_L2_prod_composite}, the modulus of continuity \eqref{eq:modulus_cont_comp} yields that $\tilde{\G}_0$ has almost surely uniformly continuous sample paths on $(\Omega\times\R_{\ge0}, d_{\Omega\times\R_{\ge0}})$.
\end{proof}
%------------------------%
\end{comment}


Lemma \ref{ap:lem_continuity_sample_paths_G_rho2_comp} is the analogue of Lemma \ref{ap:lem_continuity_sample_paths_G_rho2} for the composite null hypothesis.

%------------------------%
\begin{lemma}\label{ap:lem_continuity_sample_paths_G_rho2_comp}
%------------------------%
    The process $\tilde{\G}_0(\omega,t)$ has almost surely uniformly continuous sample paths on the metric space $(\Omega\times\R_{\ge0}, d_{\Omega\times\R_{\ge0}})$, considering the sample paths as functions from $(\Omega\times\R_{\ge0}, d_{\Omega\times\R_{\ge0}})$ to $(\R,|\cdot|)$.
\end{lemma}
\begin{proof}
The proof is completely analogous to that of Lemma~\ref{ap:lem_continuity_sample_paths_G_rho2}, replacing the class $\cY$ by $\cF$, the Gaussian process $\G_0$ by $\tilde{\G}_0$, and Lemma~\ref{ap:lem_relation_L2_prod} by Lemma~\ref{ap:lem_relation_L2_prod_composite}.
\end{proof}


Using Lemmas \ref{ap:lem_relation_L2_prod_composite} and \ref{ap:lem_continuity_sample_paths_G_rho2_comp}, Theorem \ref{thm:test_statistics_comp} can be proved.

%-----------------------%
\begin{proof}[Proof of Theorem \ref{thm:test_statistics_comp}]
    
\begin{comment}
With a telescopic trick:
\begin{align*}
    T_n^{\mathrm{CM}} &\defin n \int_{\Omega\times\R_{\ge0}} \left(F_{n,\omega}(t) - F_{\omega}^{\mu_0}(t)\right)^2 \rd F{n,\omega}(t)\, \rd \mu(\omega)\\
    &=n \int_{\Omega\times\R_{\ge0}} \left(F_{n,\omega}(t) - F_{\omega}^{\mu_0}(t)\right)^2 \,\rd F_{\omega}^{\mu_0}(t)\, \rd \mu(\omega)+n \int_{\Omega\times\R_{\ge0}} \left(F_{n,\omega}(t) - F_{\omega}^{\mu_0}(t)\right)^2  d( F_{n,\omega}-F_{\omega}^{\mu_0})(t)\, \rd \mu(\omega).
\end{align*}
Now we do a simple bounding argument to show the second integral is $o_P(1)$. Page 88 in Folland's Real Analysis (Exercise 3, \url{https://apachepersonal.miun.se/~andrli/Bok.pdf}) states: For $\mu$ a signed measure and $f\in \rho_{1,y}$,
$$
\left|\int f(x) \,\rd \mu(x)\right|\leq \int |f|\, \rd |\mu|(x),
$$
where $|\mu|$ is the total variation of $\mu$. The total variation of a measurable subset $E$ is, in particular,
$$
|\mu|(E)=\sup_{\text{Partitions $E_1,\ldots,E_m$ of $E$}}\lrb{\sum_{i=1}^m|\mu(E_j)|}
$$
(Exercise 8).

We then have
\begin{align*}
     \left|\int \left(F_{n,\omega}(t) - F_{\omega}^{\mu_0}(t)\right)^2  \rd (F_{n,\omega}-F_{\omega}^{\mu_0})(t)\right|&\leq \int \left(F_{n,\omega}(t) - F_{\omega}^{\mu_0}(t)\right)^2  d|F_{n,\omega}-F_{\omega}^{\mu_0}|(t)\\
     &\leq \sup_t \big|F_{n,\omega}(t) - F_{\omega}^{\mu_0}(t)\big|^2 \int   d|F_{n,\omega}-F_{\omega}^{\mu_0}|(t)\\
     &=\sup_t \big|F_{n,\omega}(t) - F_{\omega}^{\mu_0}(t)\big|^2 |F_{n,\omega}-F_{\omega}^{\mu_0}|(\R)\\
     &\leq 2\sup_t \big|F_{n,\omega}(t) - F_{\omega}^{\mu_0}(t)\big|^2\\
     &=o_P(1).
\end{align*}
Above, we have denoted by $|F_{n,\omega}-F_{\omega}^{\mu_0}|$ to the total variation $|\mu|$ of the signed measure $\mu=P_n-\mu_\omega^{\theta_0}$ induced by $F_{n,\omega}(t)=P_n((-\infty,t])$ and $F_{\omega}^{\mu_0}(t)=\mu_\omega^{\theta_0}((-\infty,t])$. The bound on $|\mu|(\R)$ is easily seen:
\begin{align*}
    |\mu|(\R)&=\sup \lrb{\sum_{i=1}^m|\mu(E_j)|}=\sup \lrb{\sum_{i=1}^m|P_n(E_j)-\mu_\omega^{\theta_0}(E_j)|}\\
    &\leq \sup \lrb{\sum_{i=1}^m|P_n(E_j)+\mu_\omega^{\theta_0}(E_j)|}= \sup \lrb{\sum_{i=1}^mP_n(E_j)+\sum_{i=1}^m\mu_\omega^{\theta_0}(E_j)}\\
    &= \sup \lrb{1+1}=2.
\end{align*}
\end{comment}
We divide the proof into its two parts, \ref{thm:test_statistics_comp_i} and \ref{thm:test_statistics_comp_ii}.

\emph{Proof of \ref{thm:test_statistics_comp_i}.} As explained in the proof of Theorem \ref{thm:convergence_G_n_composite}, $G_{n,\mu}^{\mu_{\hat{\boldsymbol\theta}}}$ is, up to an $o_P(1)$ remainder (that converges uniformly in $\omega$ and $t$), the empirical process indexed by $\mathcal F$. Therefore, the continuous mapping theorem gives
$$
\tilde T_n^{\mathrm{KS}} \convl \sup_{\omega\in \Omega, t\ge 0}\big|\tilde{\G}_0(\omega,t)\big|.
$$

\emph{Proof of \ref{thm:test_statistics_comp_ii}.} The proof for $\tilde{T}_n^{\mathrm{CM}}$ is completely analogous to that of $T_n^{\mathrm{CM}}$ in Theorem~\ref{thm:test_statistics_simp}. 
It suffices to replace the class $\cY$ by $\cF$, the Gaussian process $\G_0$ by $\tilde{\G}_0$, and Lemmas~\ref{ap:lem_relation_L2_prod} and \ref{ap:lem_continuity_sample_paths_G_rho2} by Lemmas~\ref{ap:lem_relation_L2_prod_composite} and \ref{ap:lem_continuity_sample_paths_G_rho2_comp}, respectively. 
\begin{comment}
Focus now on $\tilde{T}_n^{\mathrm{CM}}$.
We know that the class $\cF$ defined in \eqref{eq:function_class_comp} is Donsker, thus $\tilde{\G}_{n,\mu}^{\mu_0}(\omega, t)$ converges weakly to a tight (thus separable) Gaussian process $\tilde{\G}_0(\omega, t)$ in $\ell^\infty(\cY)$. 
Skorohod representation \citep[see Theorem 1.10.3 of][page 58]{VanderVaartWellner2023} ensures that since $\tilde{\G}_0(\omega, t)$ is separable, there exist random variables (in this case, stochastic processes) $\tilde{\G}_n(\omega, t)$ and $\tilde{\G}_0(\omega, t)$ defined on the same probability space such that $\tilde{\G}_n(\omega, t) \convas \tilde{\G}_0(\omega, t)$ in $\ell^\infty(\cY)$. 
Note: see Egorov's theorem \citep[Lemma 1.9.2 in][]{VanderVaartWellner2023} for the equivalence of almost uniform and outer almost sure convergence (see Definition 1.9.1 on page 52 in ibid). 
For simplicity, we use the same notation as the original processes. We have 
\begin{equation}\label{eq:aux_cvm_sup_convas_comp}
 \sup_{(\omega,t)\in \Omega \times \R_{\ge 0}} \left| \tilde{\G}_{n,\mu}^{\mu_0}(\omega,t) - \tilde{\G}_0(\omega, t) \right| \convas 0.   
\end{equation}
Now, observe that
\begin{equation}\label{eq:cvm_split_comp}
\begin{aligned}
\Bigg| \int_{\Omega \times \R_{\geq 0}} &\tilde{\G}_{n,\mu}^{\mu_0}(\omega, t)^2 \, \rd F_{n,\omega}^{\mu}(t) \, \rd P_n(\omega) - \int_{\Omega \times \R_{\geq 0}} \tilde{\G}_0(\omega, t)^2 \, \rd F_{\omega}^{\mu_{\btheta_0}}(t) \, \rd \mu(\omega) \Bigg| \\
&\le
\bigg| \int_{\Omega \times \R_{\ge0}} \left(\tilde{\G}_{n,\mu}^{\mu_0}(\omega, t)^2 - \tilde{\G}_0(\omega, t)^2 \right)  \, \rd F_{n,\omega}^{\mu}(t) \, \rd P_n(\omega) \bigg| \\
&+
\left| \int_{\Omega \times \R_{\ge0}}  \tilde{\G}_0(\omega, t)^2 \, \rd \left(F_{n,\omega}^{\mu}(t)\, P_n(\omega) - F_{\omega}^{\mu_{\btheta_0}}(t) \, \mu(\omega)\right) \right| 
\end{aligned}
\end{equation}
For the first term of the RHS of the inequality, one has
$$
\begin{aligned}
\Bigl| \int_{\Omega \times \R_{\ge0}} \big(&\tilde{\G}_{n,\mu}^{\mu_0}(\omega, t)^2 - \tilde{\G}_0(\omega, t)^2 \big) \, \rd F_{n,\omega}^{\mu}(t) \, \rd P_n(\omega) \Bigr| \\
&\le
\sup_{(\omega, t) \in \Omega\times \R_{\ge0}} \left| \tilde{\G}_{n,\mu}^{\mu_0}(\omega, t)^2 - \tilde{\G}_0(\omega, t)^2 \right| \int_{\Omega \times \R_{\ge0}}  \, \rd F_{n,\omega}^{\mu}(t) \, \rd P_n(\omega) \\
&= \sup_{(\omega, t) \in \Omega\times \R_{\ge0}} \left| \tilde{\G}_{n,\mu}^{\mu_0}(\omega, t)^2 - \tilde{\G}_0(\omega, t)^2 \right|.
\end{aligned}
$$
Now, observe that
$$
\begin{aligned}
\sup_{(\omega, t) \in \Omega\times \R_{\ge0}} \left| \tilde{\G}_{n,\mu}^{\mu_0}(\omega, t)^2 - \tilde{\G}_0(\omega, t)^2 \right|
&\le 
\sup_{(\omega, t) \in \Omega\times \R_{\ge0}} \left| \tilde{\G}_{n,\mu}^{\mu_0}(\omega, t) - \tilde{\G}_0(\omega, t) \right|^2 \\ 
& + 2 \sup_{(\omega, t) \in \Omega\times \R_{\ge0}} \left|\tilde{\G}_0(\omega, t)\right|  \left| \tilde{\G}_{n,\mu}^{\mu_0}(\omega, t) - \tilde{\G}_0(\omega, t) \right|
\convas 0
\end{aligned}
$$
by \eqref{eq:aux_cvm_sup_convas_comp} and the fact that $\tilde{\G}_0(\omega, t)$ has bounded sample paths almost surely. Boundedness holds because $\cF$ is Donsker, thus $\tilde{\G}_0(\omega, t)$ is a map into $\ell^\infty(\cF)$ and $\Prob{\sup_{f\in\cF} |\tilde{\G}_0(f)| < \infty} = 1$.

Now, focus on the second term in the RHS of \eqref{eq:cvm_split_comp}. The sample paths of the Gaussian process $\tilde{\G}_0(\omega, t)$ are bounded and uniformly continuous almost surely. Uniform continuity holds by Lemma \ref{ap:lem_continuity_sample_paths_G_comp}. 
Observe that Lemmas \ref{ap:lem_bound_int_g_epsilon}--\ref{ap:lem_DCT_P_hat} can be applied using $\mu_{\btheta_0}$ instead of $\mu_0$. An application of Lemma \ref{ap:lem_DCT_P_hat} to every sample path gives
\begin{align}\label{eq:aux_rhs_cvm_conv_comp_measure}
    \left| \int_{\Omega \times \R_{\ge0}} \tilde{\G}_0(\omega, t)^2 \, \rd \left(F_{n,\omega}^\mu(t)\,P_n(\omega) - F_{\omega}^{\mu_{\btheta_0}}(t)\,\mu(\omega)\right) \right| \convas 0 \text{ as } n\to \infty,
\end{align}
which is the second term of the RHS of \eqref{eq:cvm_split_comp}. Thus, almost surely, 
$$
\tilde{T}_n^{\mathrm{CM}} \to \int_{\Omega \times \R_{\geq 0}} \tilde{\G}_0(\omega, t)^2 \,\rd F_{\omega}^{\mu_{\btheta_0}}(t) \, \rd \mu(\omega) \text{ as } n \to \infty.
$$
This concludes the proof.
\end{comment}
\end{proof}


\begin{comment}
Notice that we are working with functions of the type $y_{\omega, t}$. We are not interested in $|y_{\omega,t} - y_{\omega_k,t_k}|$, but rather in $\G_{\theta_0}(\omega, t) - \G_{\theta_0}(\omega_k, t_k)$. But it can be seen that 
$
\E\big[
\bigl(\G_{0}(y_{\omega_1,t_1})-\G_{0}(y_{\omega_2,t_2})\bigr)^2
\big] \le \|y_{\omega_1,t_1}-y_{\omega_2,t_2}\|_{L^2(P)}.
$
%This is the L^2 metric based on the Gaussian process G. If the L^2(P) pseudo-metric (based on the indicators) goes to zero, then the metric based on the Gaussian process also goes to zero. 
This is useful because we have already shown that if the usual distance in $\Omega \times \R_{\ge0}$ goes to zero, then the L^2(P) pseudometric goes to zero and hence the L^2 metric based on G also goes to zero, so Theorem 5.31 provides continuity wrt the usual distance in $\Omega \times \R_{\ge0}$, which is what we want. 
It is only left to check that Theorem 5.32 applies for G, and not only for the indicators (which is what we have proved so far). The entropy condition will be trivially satisfied for the L^2 distance induced by G.
The covariance structure \eqref{eq:cov_simp} of $\G_{0}$ gives
$$
\E\left[
\bigl(\G_{0}(y_{\omega_1, t_1})-\G_{0}(y_{\omega_2, t_2})\bigr)^2
\right]
=
\Var{y_{\omega_1, t_1}(X)-y_{\omega_2, t_2}(X)}.
$$
Expand the variance term:
$$
\Var{y_{\omega_1, t_1} - y_{\omega_2, t_2}}
%=
%\Ebig{(y_{\omega_1, t_1} - y_{\omega_2, t_2})^2}
%-
%\left(\E{y_{\omega_1, t_1} - y_{\omega_2, t_2}}\right)^2 
\le 
\Ebig{(y_{\omega_1, t_1} - y_{\omega_2, t_2})^2} 
= 
\|y_{\omega_1,t_1}-y_{\omega_2,t_2}\|_{L^2(P)}^2.
$$

Observation: Since $y_{\omega_1, t_1}$ and $y_{\omega_2, t_2}$ are indicator functions, $(y_{\omega_1, t_1} - y_{\omega_2, t_2})^2 = 1_{A\triangle B}$,
Therefore,
$$
\E[(y_{\omega_1, t_1} - y_{\omega_2, t_2})^2] 
= P(A \triangle B), \ \text{ and } \ \E[y_{\omega_1, t_1} - y_{\omega_2, t_2}] = P(A) - P(B),
$$
where $A\triangle B \defin (A\setminus B)\cup(B\setminus A)$ denotes the symmetric difference. Combining these expressions,
$$
\E{
\bigl(\G_{0}(y_{\omega_1,t_1})-\G_{0}(y_{\omega_2,t_2})\bigr)^2
}
=
P\left(\bar{B}(\omega_1,t_1)\triangle \bar{B}(\omega_2,t_2)\right)
-
\bigl(P(\bar{B}(\omega_1,t_1)) - P(\bar{B}(\omega_2,t_2))\bigr)^2.
$$
\end{comment}
\begin{comment}
%----------------------------%
\begin{lemma}\label{ap:lem_separable_support_G_old_comp}
%----------------------------%
    The Gaussian process $\G_{0}(\omega, t)$ has a separable support in the space $\ell^{\infty}(\Omega \times \R_{\ge 0})$.
\end{lemma}
\begin{proof}
    The Gaussian process $\G_{0}(\omega, t)$ takes values in the space $\ell^{\infty}(\Omega \times \R_{\ge 0})$. By Lemma \ref{ap:lem_continuity_sample_paths_G}, the sample paths of $\G_{0}(\omega, t)$ are almost surely uniformly continuous on $(\Omega \times \R_{\ge 0}, d_{\Omega \times \R_{\ge 0}})$. Therefore, the support of $\G_{0}(\omega, t)$ is contained in the space $C_u(\Omega \times \R_{\ge 0})$ of bounded uniformly continuous functions from $\Omega \times \R_{\ge 0}$ to $\R$. The Banach space $C_u(\Omega \times \R_{\ge 0})$ is separable if (and only if)
$(\Omega \times \R_{\ge 0}, d_{\Omega \times \R_{\ge 0}})$ is totally bounded \citep[see, e.g., page 17 of][]{Gine2021}. Since $(\Omega \times \R_{\ge 0}, d_{\Omega \times \R_{\ge 0}})$ is a totally bounded metric space ($\Omega$ is totally bounded by assumption and $\R_{\ge 0}$ is totally bounded), then $C_u(\Omega \times \R_{\ge 0})$ is separable. Hence, the support of $\G_{0}(\omega, t)$ is also separable.
\end{proof}
\end{comment}

%---------------------------%
%---------------------------%
\subsection{Proof of Proposition \ref{prop:sup_max_ks_simple}}\label{ap:sec_proofs_null_asymptotics_prop_sup_max_ks_simple}
%---------------------------%
This section contains the proof of Proposition \ref{prop:sup_max_ks_simple}, as well as technical lemmas required to prove this result.
% Lemmas

\begin{comment}
Lemma \ref{lem:continuity_omega} is an auxiliary result for the proof of Lemma \ref{lem:uniform-continuity-y}.
%---------------------------%
\begin{lemma}[continuity of $p_t(\cdot, \delta)$]\label{lem:continuity_omega}
%---------------------------%
Fix $t \in \R_{\ge 0}$. Define
$$
p_t(\omega,\delta) \defin \Prob{|d(\omega,X)-t| \le \delta}, \quad \omega \in \Omega, \ \delta \in \R_{\ge 0}.
$$
The map $\omega \mapsto p_t(\omega, \delta)$ is continuous on $\Omega$ for every $\delta \ge 0$.
\end{lemma}
\begin{proof}
    Fix $\delta \ge 0$, and consider a sequence $\{\omega_k\}_{k\ge 1}$ in $\Omega$ such that $\omega_k \to \omega$. Since $d(\cdot, x)$ is continuous for every $x \in \Omega$, then $d(\omega_k, x) \to d(\omega,x)$ as $k\to \infty$ for every $x \in \Omega$. 

    Now, consider $x \in \Omega$ such that $|d(\omega,x) - t| \neq \delta$. One has
    \begin{align}\label{eq:conv_indicators}
    1_{\{|d(\omega_k,x) - t| \le \delta\}} \to 1_{\{|d(\omega,x) - t| \le \delta\}} \ \text{ as } k \to \infty.
    \end{align}
    To see this, first assume that $|d(\omega,x) - t| < \delta$, and let $\varepsilon = \delta - |d(\omega,x) - t| >0$. One has
    $$
    |d(\omega_k,x) - t| \le |d(\omega_k,x) - d(\omega, x)| +  |d(\omega,x) - t| < \delta,
    $$
    for $k$ sufficiently large, so that $|d(\omega_k,x) - d(\omega, x)| \le \varepsilon$. 
    If instead, $|d(\omega,x) - t| > \delta$, let $\varepsilon = |d(\omega_k,x) - t| - \delta > 0$. 
    Then, applying the reverse triangle inequality,
    $$
    |d(\omega_k,x) - t| \ge |d(\omega, x) - t| - |d(\omega_k,x) - d(\omega,x)| \ge \delta,
    $$    
    for $k$ sufficiently large, so that $|d(\omega_k,x) - d(\omega,x)| \le \varepsilon$.
    
    Observe that since $\Prob{|d(\omega_k,X) - t| = \delta} = 0$ (by the continuity of the distribution of $d(\omega,X)$ for every $\omega \in \Omega$), the convergence in \eqref{eq:conv_indicators} holds $P$-a.s. Now, since $1_{\{|d(\omega_k,x) - t| \le \delta\}}$ is bounded for every $k\ge0$, an application of the DCT gives
    $$
    \Prob{|d(\omega_k,X) - t| < \delta} \to %\int_{\Omega} 1_{|d(\omega_k,x) - t| \le \delta} dP(x) = \int_{\Omega} 1_{|d(\omega,x) - t| \le \delta} dP(x)= 
    \Prob{|d(\omega,X) - t| < \delta} \quad \text{as } n \to \infty.
    $$
    That is, $p_t(\omega_k, \delta) \to p_t(\omega, \delta)$ when $\omega_k \to \omega$, i.e., $p_t(\cdot, \delta)$ is continuous on $\Omega$ for every $\delta\ge0$.
\end{proof}
%------------------------%

%------------------------%
\begin{lemma}[uniform continuity of $\omega \mapsto y_{\omega,t}$]
%------------------------%
\label{lem:uniform-continuity-y}
For every fixed $t\ge0$, the mapping $\omega \mapsto y_{\omega,t}$
is uniformly continuous from $(\Omega, d)$ into $(L_2(\mu), \rho_{\mu})$, where $\rho_{\mu}$ is the pseudometric induced by the $L_2$-norm: 
$$
\rho_{\mu}(f, g) \defin \left(\mu[(f-g)^2]\right)^{\frac{1}{2}} = \left(\mathsf{E}[(f(X)-g(X))^2]\right)^{\frac{1}{2}},
$$
for measurable functions $f$ and $g$.
\end{lemma}

\begin{proof}
Fix $t \in \R_{\ge 0}$ and $\omega, \omega^\prime \in \Omega$. We compute
$$
\rho_{\mu}(y_{\omega,t}, y_{\omega^\prime,t})^2
= \mathsf{E}\big( (1_{\{d(\omega,X)\le t\}} - 1_{\{d(\omega^\prime,X)\le t\}})^2 \big)
= \Prob{1_{\{ 1_{\{d(\omega,X)\le t\}} \neq 1_{\{d(\omega^\prime,X)\le t\}}\}}}.
$$
Hence,
$$
\rho_{\mu}^2(y_{\omega,t}, y_{\omega^\prime,t})
= \Prob{ \{d(\omega,X)\le t < d(\omega^\prime,X)\} \cup \{d(\omega^\prime,X)\le t < d(\omega,X)\}}.
$$
If we let $\delta = d(\omega,\omega^\prime)$, then
$$
\{d(\omega,X)\le t < d(\omega^\prime,X)\} \subseteq \{ |d(\omega,X)-t| \le \delta \},
$$
and similarly for the other term, giving
$$
\rho_{\mu}^2(y_{\omega,t}, y_{\omega^\prime,t})
\le 2\, \Prob{|d(\omega,X)-t| \le \delta}.
$$
% Uniform control without compactness (replaces the Dini argument)
Under Assumption~\ref{b2}, for every $\omega\in\Omega$ the random variable $d(\omega,X)$
admits a (continuous) density $f_\omega$, and $\sup_{\omega\in\Omega}\sup_{u\in\R} f_\omega(u)\le \bar\Delta<\infty.$
Therefore, for any fixed $t\in\R$ and $\delta>0$,
$$
p_t(\omega,\delta)
=\Prob{|d(\omega,X)-t|\le \delta}
=\int_{t-\delta}^{t+\delta} f_\omega(u)\,du
\le 2\delta\,\sup_{u\in\R} f_\omega(u)
\le 2\delta\,\bar\Delta.
$$
Taking the supremum over $\omega\in\Omega$ yields
$$
\sup_{\omega\in\Omega} p_t(\omega,\delta)\le 2\delta\,\bar\Delta \to 0 \quad \text{as } \delta\to 0^+,
$$
Finally, we can conclude that
$$
\sup_{d(\omega,\omega^\prime) < \delta}
\rho_{\mu}(y_{\omega,t}, y_{\omega^\prime,t})
\le \sqrt{ 2 \sup_{\omega \in \Omega} p_t(\omega, \delta) }
\to 0 \quad \text{as } \delta \to 0^+.
$$
This proves that $\omega \mapsto y_{\omega,t}$ is uniformly continuous from $(\Omega,d)$ to $(L_2(\mu),\rho_{\mu})$.
\end{proof}
%---------------------------%
\end{comment}




%---------------------------%
\begin{lemma}[sample forms an asymptotic $\delta$-net]
%---------------------------%
\label{lem:sample-delta-net}
Let $\mathcal{L}_n = \{X_1,X_2, \allowbreak \dots, X_n\}$ be iid draws from a   distribution $\mu$ such that $\Omega \subset \supp{\mu}$.
Then, the sample $\mathcal{L}_n$ is a.s. dense in $\Omega$ as $n$ diverges to infinity, i.e.,
$$
\sup_{\omega\in\Omega}\min_{1\le i\le n} d(\omega,X_i) \to 0 \ \text{ a.s. as } n\to\infty. 
$$
In particular, for every $\varepsilon>0$ the random set $\mathcal{L}_n$ is a.s. a $2\varepsilon$-net of $\Omega$, for $n$ sufficiently large.
\end{lemma}

\begin{proof}
Fix $\varepsilon>0$. By total boundedness of $\Omega$ (implied by Assumption \ref{b1}) there exists a finite $\varepsilon$-net
$$
\Omega \subset \bigcup_{j=1}^m B(\omega_j,\varepsilon)
$$
for some points $\omega_1,\dots,\omega_m\in\Omega$.

Define, for each $j=1,\dots,m$, the event
$$
A_{n,j} \defin \{B(\omega_j,\varepsilon)\cap \mathcal{L}_n = \emptyset\}
%= \big\{X_1\notin B(\omega_j,\varepsilon),\dots,X_n\notin B(\omega_j,\varepsilon)\big\}.
$$
%Because the $X_i$ are iid, 
For each fixed $j$ one has $\mathsf{P}\big(A_{n,j}\big) = \big(1 - \mathsf{P}\left(B(\omega_j,\varepsilon)\right)\big)^n$.
Set $c_\varepsilon \defin \min_{1\le j\le m} \mathsf{P}\big(B(\omega_j,\varepsilon)\big)$. Since $\omega_j\in\Omega\subset \mathrm{supp}(\mu)$, one has $c_\varepsilon>0$.

Consider the event 
$$
A_n \defin \Big\{\exists j\in\{1,\dots,m\}\text{ such that }B(\omega_j,\varepsilon)\cap \mathcal{L}_n=\emptyset\Big\}
= \bigcup_{j=1}^m A_{n,j}.
$$
It holds that
$$
\sum_{n=1}^\infty \mathsf{P}(A_n) \le \sum_{n=1}^\infty \sum_{j=1}^m \mathsf{P}(A_{n,j})
= \sum_{n=1}^\infty \sum_{j=1}^m \big(1 - \mathsf{P}\big(B(\omega_j,\varepsilon))\big)^n
\le \sum_{n=1}^\infty m \big(1 - c_\varepsilon\big)^n < \infty,
$$
because $1-c_\varepsilon<1$.
By the Borel--Cantelli lemma, with probability one, only finitely many of the events $A_n$ occur; equivalently, with probability one there exists a (random) integer $N = N(\varepsilon)$ such that for all $n\ge N(\varepsilon)$ no $A_{n,j}$ occurs, i.e. every ball $B(\omega_j,\varepsilon)$ contains at least one sample point. Since $\{\omega_1, \ldots, \omega_m\}$ forms an $\varepsilon$-net of $\Omega$, for every $\omega\in\Omega$ we can pick $1\le j\le m$ with $d(\omega,\omega_j)<\varepsilon$, and hence for any sample element $X_i\in B(\omega_j,\varepsilon)$, the triangle inequality gives $d(\omega,X_i) \le 2\varepsilon$ for all $\omega \in \Omega$ with probability one. 
Therefore,
$$
\Prob{\exists N(\varepsilon):\ \forall n\ge N(\varepsilon),\ 
\sup_{\omega\in\Omega}\min_{1\le i\le n} d(\omega,X_i)\le 2\varepsilon}=1.
$$
Equivalently,
$$
\Prob{\liminf_{n\to\infty}\Big\{\sup_{\omega\in\Omega}\min_{1\le i\le n} d(\omega,X_i)\le 2\varepsilon\Big\}}=1.
$$
Observe that $\sup_{\omega\in\Omega}\min_{1\le i\le n} d(\omega,X_i)$ is nonnegative and nonincreasing as a function of $n$. 
Therefore,
\begin{align}\label{eq:as_lem_dense}
\Prob{\lim_{n\to\infty}\sup_{\omega\in\Omega}\min_{1\le i\le n} d(\omega,X_i)\le 2\varepsilon}=1.
\end{align}
%i.e., $\sup_{\omega}\min_i d(\omega,X_i)\le 2\varepsilon$ almost surely when $n$ is sufficiently large.

Finally, let $\varepsilon_k = 1/k$ for $k\in\mathbb N$. The above argument can be repeated to obtain \eqref{eq:as_lem_dense} for every $\varepsilon_k$. Making $k\to \infty$, 
% Simply take the (countable) intersection of the event inside the probability in (11) over k\ge0. Since it is a countable intersection of probability 1 events, it has probability 1.
this implies
%$$
%\Prob{\lim_{n\to \infty} \sup_{\omega\in\Omega}\min_{1\le i\le n} d(\omega,X_i) = 0} = 1.
%$$
%Since $d(\omega,X_i) \ge 0$, 
%then the limsup is in fact the lim
that $\sup_{\omega\in\Omega}\min_{1\le i\le n} d(\omega,X_i) \to 0$ a.s. as $n\to\infty$.
\end{proof}

Using Lemmas \ref{ap:lem_relation_L2_prod} and \ref{lem:sample-delta-net}, Proposition \ref{prop:sup_max_ks_simple} can be proved.
\begin{proof}[Proof of Proposition \ref{prop:sup_max_ks_simple}]
We divide the proof into its two parts, \ref{prop:sup_max_ks_simple_i} and \ref{prop:sup_max_ks_simple_ii}.

\emph{Proof of \ref{prop:sup_max_ks_simple_i}.} Consider the class $\cY = \{y_{\omega,t} : \omega \in \Omega,\ t\ge0\}$. Under $H_0$, the class $\cY$ is Donsker.
Throughout the proof, we use the $L^2(\mu)$ pseudometric $\rho_{2,y}$ defined in Lemma~\ref{ap:lem_relation_L2_prod}, identifying $\cY$ with its index set $\Omega\times\R_{\ge0}$.

Since $\cY$ is Donsker, the empirical process is asymptotically equicontinuous with respect to $\rho_{2,y}$ \citep[page 139 in][]%[and Theorem 9 in Ingrid's notes]
{VanderVaartWellner2023}: for every $\eta >0$,
$$
\lim_{\delta \to 0^+} \limsup_{n\to \infty} \mathsf{P}^*\Bigg\{ \sup_{\substack{f,g\in\cY\\ \rho_{2,y}(f,g) < \delta}}  \left|\G_{n,\mu}^{\mu_0}(f-g)\right| > \eta \Bigg\} = 0.
$$

Fix $\tau>0$ and $\gamma>0$. By asymptotic equicontinuity, there exists $\delta = \delta(\tau,\gamma) > 0$ such that
$$
\limsup_{n\to\infty} \mathsf{P}^*\Bigg\{\sup_{\substack{f,g\in\cY\\ \rho_{2,y}(f,g)<\delta}} \left| \G_{n,\mu}^{\mu_0}(f) - \G_{n,\mu}^{\mu_0}(g)\right| > \tau \Bigg\} < \gamma.
$$
Hence, for $n$ sufficiently large, 
\begin{align}\label{eq:eta_equicont}
\mathsf{P}^*\Bigg\{\sup_{\substack{f,g\in\cY\\ \rho_{2,y}(f,g) < \delta}} \left| \G_{n,\mu}^{\mu_0}(f) - \G_{n,\mu}^{\mu_0}(g)\right| > \tau \Bigg\} < \gamma.
\end{align}

For any function $y_{\omega,t} \in \cY$, Lemma~\ref{ap:lem_relation_L2_prod} shows that the map $\omega \mapsto y_{\omega,t}$ is uniformly continuous with respect to $\rho_{2,y}$, uniformly in $t$. Indeed,
$$
\rho_{2,y}(y_{\omega,t}, y_{\omega^\prime,t}) \le 2\sqrt{\bar\Delta\, d(\omega,\omega^\prime)}
$$
for all $\omega, \omega^\prime \in \Omega$ and all $t\ge0$. Therefore, choosing $r > 0$ such that $r < \delta^2/(4\bar\Delta)$, it follows that for all $\omega, \omega^\prime \in \Omega$ and all $t\ge0$, $d(\omega, \omega^\prime) < r$ implies $\rho_{2,y}(y_{\omega,t}, y_{\omega^\prime,t}) < \delta$.

By Lemma~\ref{lem:sample-delta-net}, the sample $\mathcal{L}_n$ is dense in $\Omega$, so for $n$ sufficiently large:
$$
\Prob{\sup_{\omega \in \Omega} \min_{1 \le i \le n} d(\omega, X_i) < r} \geq 1 - \gamma.
$$
Thus, with probability at least $1 - \gamma$, for every $\omega \in \Omega$ there exists $X_i \in \mathcal{L}_n$ such that $d(\omega, X_i) < r$, and hence by uniform continuity $\rho_{2,y}(y_{\omega,t}, y_{X_i,t}) < \delta$ for all $t\ge0$. 
This means that $\cY_n \defin \{y_{X_i,t} : i=1,\ldots,n,\ t\ge0\}$ forms a $\delta$-net for $\cY$ with probability at least $1-\gamma$ when $n$ is sufficiently large, i.e., for every $y_{\omega,t} \in \cY$ there exists $y_{X_i,t} \in \cY_n$ such that $\rho_{2,y}(y_{\omega,t}, y_{X_i,t}) < \delta$.

By the triangle inequality, one has
$$
\left| \G_{n,\mu}^{\mu_0}(y_{\omega,t}) \right| \le \left|\G_{n,\mu}^{\mu_0}(y_{X_i,t})\right| + \left| \G_{n,\mu}^{\mu_0}(y_{\omega,t}) - \G_{n,\mu}^{\mu_0}(y_{X_i,t}) \right|.
$$
Therefore, if $\cY_n$ is a $\delta$-net of $\cY$, then
\begin{align*}
    \sup_{y\in \cY} \left| \G_{n,\mu}^{\mu_0}(y) \right| \le \max_{1 \le i \le n}\sup_{t\ge0} \left|\G_{n,\mu}^{\mu_0}(y_{X_i,t})\right| + \sup_{\substack{f,g\in\cY\\ \rho_{2,y}(f,g) < \delta}} \left| \G_{n,\mu}^{\mu_0}(f) - \G_{n,\mu}^{\mu_0}(g) \right|.
\end{align*}
The last term on the right is less than or equal to $\tau$ with (outer) probability at least $1-\gamma$ by \eqref{eq:eta_equicont}. Trivially, $\max_{1 \le i \le n}\sup_{t\ge0} \left|\G_{n,\mu}^{\mu_0}(y_{X_i,t})\right| \le \sup_{y\in \cY} \left| \G_{n,\mu}^{\mu_0}(y) \right|$. 
Hence, for every $\tau>0$ and $\gamma>0$, it holds that
$$
\mathsf{P}^*\Bigg\{\max_{1 \le i \le n}\sup_{t\ge0} \left|\G_{n,\mu}^{\mu_0}(y_{X_i,t})\right| \le \sup_{y\in \cY} \left| \G_{n,\mu}^{\mu_0}(y) \right| \le \max_{1 \le i \le n}\sup_{t\ge0} \left|\G_{n,\mu}^{\mu_0}(y_{X_i,t})\right| + \tau \Bigg\} \ge 1 - 2\gamma.
$$
This means that 
$$
\mathsf{P}^*\Bigg\{\Big|\sup_{y\in\cY} \left|\G_{n,\mu}^{\mu_0}(y)\right| - \max_{1\le i \le n}\sup_{t\ge0} \left| \G_{n,\mu}^{\mu_0}(y_{X_i,t}) \right| \Big| > \tau \Bigg\} < 2\gamma
$$
% use the probability of a union <= sum of probabilities
for $n$ sufficiently large. Therefore, 
$$
\sup_{y\in\cY} \left| \G_{n,\mu}^{\mu_0}(y)\right| - \max_{1\le i \le n}\sup_{t\ge0} \left| \G_{n,\mu}^{\mu_0}(y_{X_i,t}) \right| \stackrel{\mathsf{P}^*}{\rightarrow} 0 \text{ as } n \to \infty,
$$
that is,
$$
\sup_{\omega\in \Omega,\ t\ge0} \left|\G_{n,\mu}^{\mu_0}(\omega,t)\right| = \max_{1\le i \le n}\sup_{t\ge0} \left|\G_{n,\mu}^{\mu_0}(X_i,t)\right| + o_{\mathsf{P}^*}(1).
$$
Under Assumption~\ref{b2}, for each fixed $X_i$, the standard argument for the Kolmogorov--Smirnov statistic discretizes the supremum over $t$ into a maximum over the distance sample $d(X_i,X_1),\ldots,d(X_i,X_n)$.
This proves \eqref{eq:sup_max_ks_simple}.

\emph{Proof of \ref{prop:sup_max_ks_simple_ii}.} The proof of \eqref{eq:sup_max_ks_comp} is analogous to that of \eqref{eq:sup_max_ks_simple}. One replaces $\cY$ by $\cF$, and Lemma~\ref{ap:lem_relation_L2_prod} by Lemma~\ref{ap:lem_relation_L2_prod_composite}.
\end{proof}
%---------------------------%

%---------------------------%
\section{Proofs of Section \ref{sec:non_null_asymptotics}} \label{ap:sec_proofs_non_null_asymptotics}
%---------------------------%

%---------------------------%
\subsection{Fixed alternatives} \label{ap:sec_proofs_fixed_alternatives}
%---------------------------%

%---------------------------%
\subsubsection{Simple null hypothesis} \label{ap:sec_proofs_fixed_alternatives_simple}
%---------------------------%

Theorem \ref{thm:fixed_alt_divergence_ks} states the consistency of the Kolmogorov--Smirnov test statistic under fixed alternatives.

\begin{proof}[Proof of Theorem \ref{thm:fixed_alt_divergence_ks}]
For $(\omega,t)\in\Omega\times\R_{\ge 0}$, write
$$
F_{n,\omega}^\mu(t)-F^{\mu_0}_\omega(t)
=
\bigl(F_{n,\omega}^\mu(t)-F^\mu_\omega(t)\bigr)
+
\bigl(F^\mu_\omega(t)-F^{\mu_0}_\omega(t)\bigr),
$$
and define the empirical process under $\mu$ ($\mu \neq \mu_0$) by
\begin{align}\label{eq:G_n_nu}
\G_{n,\mu}^{\mu}(\omega,t)
\defin
\sqrt{n}\,\bigl(F_{n,\omega}^\mu(t)-F^\mu_\omega(t)\bigr).    
\end{align}
Next, define $\Delta(\omega,t)\defin F^\mu_\omega(t)-F^{\mu_0}_\omega(t)$.
Since $\mu\neq\mu_{0}$, there exists at least one pair $(\omega_0,t_0)\in\Omega\times\R_{\ge 0}$ such that $\Delta(\omega_0,t_0)\neq 0$. Therefore,
$$
\delta
\defin
\sup_{\omega\in\Omega, t\ge 0}
|\Delta(\omega,t)|
\ \ge\
|\Delta(\omega_0,t_0)|
\ >\ 0.
$$
\textcolor{blue}{Diego: I present two options to continue the proof from here. Option 1 is simpler. I let you choose.}

\noindent\emph{Option 1.}
Since $\cY$ is $\mu$-Donsker, it is also $\mu$-Glivenko--Cantelli. Therefore,
$$
\sup_{\omega\in\Omega, t\ge0}\bigl|F_{n,\omega}^\mu(t)-F_\omega^\mu(t)\bigr|
=
o_{\mathsf{P}^*}(1).
$$
Now, write
% Since $\cY$ is $\mu$-Donsker, there exists a tight Gaussian process $\G_{\mu}$ in
% $\cY$ such that
% $\G_{n,\mu}^{\mu}\convl \G_{\mu}$ in $\ell^\infty(\cY).$
% Since $\G_{\mu}$ is tight, by Lemma 1.3.8(ii) in \cite{VanderVaartWellner2023}, $\G_{n,\mu}^{\mu}$ is asymptotically tight (see Definition 1.3.7 and Problem 1.3.6 in ibid).
% \cite[see Definition 1.3.7 and Problem 1.3.6 in][]{VanderVaartWellner2023}. 
% It follows that $\|\G_{n,\mu}^{\mu}\|_\infty=O_{\mathsf{P}^*}(1)$. Moreover,
\begin{comment}
For every $\varepsilon>0$ there exists a compact set
$K\subset \ell^\infty(\Omega\times\R_{\ge 0})$ such that
$$
\liminf_{n\to\infty}\mathsf{P}^*\bigl(\G_{n,\mu}^{\mu}\in G\bigr)\ge 1-\varepsilon
$$
for every open set $G\supset K$. Since $K$ is compact, it is bounded, so there
exists $M<\infty$ such that
$$
K\subset \left\{f\in \ell^\infty(\Omega\times\R_{\ge0}): \|f\|_\infty<M\right\}.
$$
Applying the previous display to the open set
$G=\{f\in \ell^\infty(\Omega\times\R_{\ge0}): \|f\|_\infty<M\}$, we obtain
$$
\liminf_{n\to\infty}\mathsf{P}^*\left(\|\G_{n,\mu}^{\mu}\|_\infty<M\right)\ge 1-\varepsilon,
$$
that is,
$$
\limsup_{n\to\infty}\mathsf{P}^*\left(\|\G_{n,\mu}^{\mu}\|_\infty\ge M\right)\le \varepsilon.
$$
\end{comment}
\begin{align*}
\frac{T_n^{\mathrm{KS}}}{\sqrt n}
=
\sup_{\omega\in\Omega, t\ge0}\left|\Delta(\omega,t)+\bigl(F_{n,\omega}^\mu(t)-F_\omega^\mu(t)\bigr)\right|.
\end{align*}
Using the inequality $\bigl|\sup |a+b|-\sup |a|\bigr|\le \sup |b|$, we obtain
$$
\left|\frac{T_n^{\mathrm{KS}}}{\sqrt n}-\delta\right|
\le
\sup_{\omega\in\Omega, t\ge0}\bigl|F_{n,\omega}^\mu(t)-F_\omega^\mu(t)\bigr|
=
\frac{\|\G_{n,\mu}^{\mu}\|_\infty}{\sqrt n}
=
o_{\mathsf{P}^*}(1).
$$
Therefore,
$$
\frac{T_n^{\mathrm{KS}}}{\sqrt n}\xrightarrow{\mathsf{P}^*}\delta,
$$
and in particular,
$$
T_n^{\mathrm{KS}}\xrightarrow{\mathsf{P}^*}\infty.
$$


\noindent\textcolor{blue}{\emph{Option 2.}}
{\color{blue}
Using the reverse triangle inequality pointwise in $(\omega,t)$ and taking suprema, we obtain
\begin{align*}
T_n^{\mathrm{KS}} 
&\ge
\sup_{\omega\in\Omega, t\ge 0}\sqrt{n}\,|\Delta(\omega,t)|
-
\sup_{\omega\in\Omega, t\ge 0}\sqrt{n}\,\bigl|F_{n,\omega}^\mu(t)-F^\mu_\omega(t)\bigr|
=
\sqrt{n}\,\delta
-
\big\|\G_{n,\mu}^{\mu}\big\|_\infty.
\end{align*}

Since $\cY$ is Donsker under $\mu$, the empirical process $\G_{n,\mu}^{\mu}$ is asymptotically tight in
$\ell^\infty(\Omega\times\R_{\ge 0})$; \citep[see, e.g., page 139 in][]{VanderVaartWellner2023}.
By the continuous mapping theorem for asymptotic tightness
\citep[see, e.g., Problem~7(i) in Section~1.3 of][]{VanderVaartWellner2023}, the sequence $\|\G_{n,\mu}^{\mu}\|_\infty$ is
asymptotically tight in $\R$. In particular, for every $\varepsilon>0$ there exists $M<\infty$ such that
\begin{equation}\label{eq:def_asympt_tight}
    \limsup_{n\to\infty}\mathsf{P}_\mu^{*}\left(\|\G_{n,\mu}^{\mu}\|_\infty>M\right)\le \varepsilon.
\end{equation}

Take any $M_0>0$ and set $a_n \defin \sqrt{n}\,\delta-M_0$. Fix $\varepsilon>0$ and choose $M$ to satisfy \eqref{eq:def_asympt_tight}.
Then, there exists $n_0\in\mathbb{N}$ such that $a_n\ge M$ for all $n\ge n_0$, and hence, for all $n\ge n_0$, $\mathsf{P}_\mu^{*}\left(\|\G_{n,\mu}^{\mu}\|_\infty\ge a_n\right) \le \mathsf{P}_\mu^{*}\left(\|\G_{n,\mu}^{\mu}\|_\infty> M\right).$
Taking $\limsup_{n\to\infty}$ yields
$$
\limsup_{n\to\infty}\mathsf{P}_\mu^{*}\left(\|\G_{n,\mu}^{\mu}\|_\infty\ge a_n\right)\le \varepsilon.
$$
Since $\varepsilon>0$ is arbitrary, it follows that
$$
\mathsf{P}_\mu^{*}\left(\|\G_{n,\mu}^{\mu}\|_\infty\ge \sqrt{n}\,\delta-M_0\right)\to 0,
\qquad n\to\infty.
$$
Hence, $\mathsf{P}_\mu^{*}(T_n^{\mathrm{KS}}\le M_0)\to 0$ for every $M_0>0$. This proves
$$
T_n^{\mathrm{KS}}\xrightarrow{\mathsf{P}^*}\infty.
$$
}
\end{proof}
%---------------------------%


We now turn to the Cramér--von Mises test statistic. Theorem \ref{thm:fixed_alternative_divergence_cvm} states that the consistency of the Cramér--von Mises test statistic under fixed alternatives.
To prove this result, Lemmas \ref{lem:delta_regularity}, \ref{lem:robust_separation} and \ref{lem:positive_mass_region} are needed.


%---------------------------%
\begin{lemma}\label{lem:delta_regularity}
%---------------------------%
%Assume \ref{b2} holds for $\mu$ and for $\mu_{0}$. 
Let $M=\mathrm{diam}(\Omega)$ and define $\Delta(\omega,t)\defin F^\mu_\omega(t)-F^{\mu_0}_\omega(t)$ for $(\omega,t)\in\Omega\times[0,M]$.
There exist constants $\bar\Delta_\mu<\infty$ and $\bar\Delta_{\mu_0}<\infty$ (independent of $t$) such that, for any
$\omega,\omega^\prime\in\Omega$ and each $t\in[0,M]$,
$$
|\Delta(\omega,t)-\Delta(\omega^\prime,t)|
\le
(\bar\Delta_\mu+\bar\Delta_{\mu_0})\, d(\omega,\omega^\prime).
$$
That is, for each $t\in[0,M]$ the map $\omega\mapsto \Delta(\omega,t)$ is Lipschitz on $\Omega$ with Lipschitz constant
bounded uniformly over $t\in[0,M]$.
\end{lemma}

\begin{proof}
Fix $t\in[0,M]$ and $\omega,\omega^\prime\in\Omega$, and define $r\defin d(\omega,\omega^\prime)$. For every $x\in\Omega$, the triangle inequality yields
$$
d(\omega^\prime,x) \le d(\omega^\prime,\omega) + d(\omega,x) = r + d(\omega,x).
$$
Hence, $\{d(\omega,x)\le t\}\subset\{d(\omega^\prime,x)\le t+r\}$, and so
$F^\mu_\omega(t)\le F^\mu_{\omega^\prime}(t+r)$. Similarly, $F^\mu_{\omega^\prime}(t)\le F^\mu_{\omega}(t+r)$, so
$$
|F^\mu_\omega(t)-F^\mu_{\omega^\prime}(t)|
\le \max\{F^\mu_{\omega^\prime}(t+r)-F^\mu_{\omega^\prime}(t),\,F^\mu_{\omega}(t+r)-F^\mu_{\omega}(t)\}.
$$
By \ref{b2}, the densities $\{f^\mu_\omega\}$ are uniformly bounded above by some $\bar\Delta_\mu<\infty$, hence
$$
F^\mu_{\omega}(t+r)-F^\mu_{\omega}(t)
=\int_t^{t+r} f^\mu_{\omega}(u)\,\rd u
\le \bar\Delta_\mu \, r,
$$
and similarly for $\omega^\prime$. Therefore,
$$
|F^\mu_\omega(t)-F^\mu_{\omega^\prime}(t)|\le \bar\Delta_{\mu}\, d(\omega,\omega^\prime).
$$
The same argument applied to $\mu_{0}$ yields
$|F^{\mu_0}_\omega(t)-F^{\mu_0}_{\omega^\prime}(t)|\le \bar\Delta_{\mu_0}\, d(\omega,\omega^\prime)$ for a finite $\bar\Delta_{\mu_0}$.
Finally,
$$
|\Delta(\omega,t)-\Delta(\omega^\prime,t)|
\le |F^\mu_\omega(t)-F^\mu_{\omega^\prime}(t)| + |F^{\mu_0}_\omega(t)-F^{\mu_0}_{\omega^\prime}(t)|
\le (\bar\Delta_{\mu}+\bar\Delta_{\mu_0})\, d(\omega,\omega^\prime),
$$
which gives the result.
\end{proof}

%---------------------------%
\begin{lemma}\label{lem:robust_separation}
%---------------------------%
Assume that \ref{b2} holds. 
If there exist $\omega_0\in\Omega$ and $t_0\in[0,M]$ such that $\Delta(\omega_0,t_0)\neq 0$, then there exist $\varepsilon_\omega>0$, $\varepsilon_t>0$, and $c_0>0$ such that
$$
|\Delta(\omega,t)|\ge c_0
\qquad\forall\,\omega\in B(\omega_0,\varepsilon_\omega),\ \forall\,t\in (t_0-\varepsilon_t,t_0+\varepsilon_t)\cap[0,M].
$$
\end{lemma}

\begin{proof}
Let $\delta_0\defin|\Delta(\omega_0,t_0)|>0$. 
For each $\omega\in\Omega$, the map $t\mapsto \Delta(\omega,t)$ is continuous on $[0,M]$ by assumption \ref{b2}. 
Thus, there exists $\varepsilon_t>0$ such that
$$
|\Delta(\omega_0,t)-\Delta(\omega_0,t_0)|<\frac{\delta_0}{2}
\quad\text{for all $t\in[0,M]$ such that } |t-t_0|<\varepsilon_t,
$$
which implies $|\Delta(\omega_0,t)|\ge \frac{\delta_0}{2}$. Next, by Lemma~\ref{lem:delta_regularity}, for all $t\in[0,M]$ one has
$$
|\Delta(\omega,t)-\Delta(\omega_0,t)|
\le (\bar\Delta_{\mu}+\bar\Delta_{\mu_0})\, d(\omega,\omega_0).
$$
Observe that the right-hand side does not depend on $t$.

Choose $\varepsilon_\omega\defin\frac{\delta_0}{4(\bar\Delta_{\mu}+\bar\Delta_{\mu_0})}$ and set $c_0\defin\frac{\delta_0}{4}$. 
Then, for all $\omega\in B(\omega_0,\varepsilon_\omega)$ and all $t\in (t_0-\varepsilon_t,t_0+\varepsilon_t)\cap[0,M]$, it holds that
$$
|\Delta(\omega,t)|
\ge |\Delta(\omega_0,t)| - |\Delta(\omega,t)-\Delta(\omega_0,t)|
\ge \delta_0/2 - {(\bar\Delta_{\mu}+\bar\Delta_{\mu_0}) \varepsilon_\omega}
= \delta_0/4
\defin c_0 >0.
$$
This concludes the result.
\end{proof}

%----------------------------%
\begin{lemma}\label{lem:positive_mass_region}
%----------------------------%
% Assume \ref{b2} holds for $\mu$ and 
Let $\omega_0\in\supp{\mu}$. Let $B_0 \defin B(\omega_0,\varepsilon_\omega)$ for some $\varepsilon_\omega>0$, and let $I_0\subset[0, M]$ be (any) nonempty open interval (relative to $[0,M]$). If $X,Y$ are iid with law $\mu$, then
$$
\Probbig{X\in B_0,\ d(X,Y)\in I_0}>0.
$$
\end{lemma}
\begin{proof}
Fix $\omega\in\Omega$. Let $\mu_\omega$ denote the law of $d(\omega,Y)$. Under \ref{b2}, the random variable $d(\omega,Y)$ admits a continuous density $f^\mu_\omega$ on $[0,M]$ and satisfies $\underline{\Delta}_\omega\defin\inf_{t\in\supp{\mu_{\omega}}} f^\mu_\omega(t)>0$.
Set
$$
A_\omega\defin\{t\in[0,M]: f^\mu_\omega(t)>0\}.
$$
Since $f^\mu_\omega$ is continuous on $[0,M]$, the set $A_\omega$ is (relatively) open in $[0,M]$.

We first show that $A_\omega\subset\supp{\mu_\omega}$. Let $t\in A_\omega$, so that $f^\mu_\omega(t)>0$. By continuity of $f^\mu_\omega$ on $[0,M]$, there exists $\eta>0$ such that
$$
f^\mu_\omega(s)\ge \frac{f^\mu_\omega(t)}{2}>0
\qquad\text{for all }s\in (t-\eta,t+\eta)\cap[0,M].
$$
Let $U\subset[0,M]$ be a relatively open set with $t\in U$. Then there exists $\eta'>0$ such that
$$
(t-\eta',t+\eta')\cap[0,M]\subset U.
$$
Setting $\tilde\eta\defin\min\{\eta,\eta'\}$, we obtain
$$
\mu_\omega(U)=\int_U f^\mu_\omega(s)\,\rd s
\ge \int_{(t-\tilde\eta,t+\tilde\eta)\cap[0,M]} f^\mu_\omega(s)\,\rd s
>0.
$$
Hence, by the definition of support, $t\in\supp{\mu_\omega}$. 

Next, we show that $\supp{\mu_\omega}\subset A_\omega$. If $t\in\supp{\mu_\omega}$, then by Assumption \ref{b2} one has $f^\mu_\omega(t)\ge \underline{\Delta}_\omega>0,$
so $t\in A_\omega$.
Therefore, $A_\omega=\supp{\mu_\omega}$, which is relatively closed in $[0,M]$. Since $A_\omega$ is also relatively open in $[0,M]$, and $A_\omega\neq\varnothing$ (because $d(\omega,Y)$ admits a density on $[0,M]$), the connectedness of $[0,M]$ implies that $A_\omega=[0,M]$, and thus $f^\mu_\omega(t)>0$ for all $t\in[0,M]$.
Consequently, for any nonempty open interval $I_0\subset[0,M]$,

\begin{equation}\label{eq:positive_density_distance}
    \Probbig{d(\omega,Y)\in I_0}=\int_{I_0} f^\mu_\omega(t)\,dt>0.
\end{equation}
Finally, by the tower property,
$$
\Probbig{X\in B_0,\ d(X,Y)\in I_0}
=\mathsf{E}_\mu\left[1_{\{X\in B_0\}}1_{\{d(X,Y)\in I_0\}}\right]
=\mathsf{E}_\mu\left[1_{\{X\in B_0\}}\,\mathsf{P}_\mu\left(d(X,Y)\in I_0\,\middle|\,X\right)\right]
$$
and for every $\omega\in\Omega$, since $Y$ is independent of $X$,
$$
\mathsf{P}_\mu\left(d(X,Y)\in I_0\,\middle|\,X=\omega\right)
=\Probbig{d(\omega,Y)\in I_0}.
$$
Therefore,
$$
\Probbig{X\in B_0,\ d(X,Y)\in I_0}
=\int_{B_0}\Probbig{d(\omega,Y)\in I_0}\, \rd \mu(\omega).
$$
The integrand is strictly positive for every $\omega\in B_0$ by \eqref{eq:positive_density_distance}, and $\mu(B_0)>0$ because $\omega_0\in\supp{\mu}$. Hence, %the integral is strictly positive.
$\Prob{X\in B_0,\ D\in I_0}>0.$
\end{proof}

Using Lemmas \ref{lem:delta_regularity}, \ref{lem:robust_separation} and \ref{lem:positive_mass_region}, Theorem~\ref{thm:fixed_alternative_divergence_cvm} can be proved.
%----------------------------%
\begin{proof}[Proof of Theorem~\ref{thm:fixed_alternative_divergence_cvm}]
Since either \ref{a2} holds, or \ref{a1} holds and $\supp{\mu} = \Omega$, under $H_1$, there exist $\omega_0\in\supp{\mu}$ and $t_0\in(0,M]$ such that
$F^\mu_{\omega_0}(t_0)\neq F^{\mu_{0}}_{\omega_0}(t_0)$, i.e.  $\Delta(\omega_0,t_0)\neq 0$.
Apply Lemma~\ref{lem:robust_separation} to obtain $\varepsilon_\omega>0$, $\varepsilon_t>0$, and $c_0>0$ such that
$$
|\Delta(\omega,t)|\ge c_0
\quad\forall\,(\omega,t)\in B_0\times I_0,
$$
where $B_0=B(\omega_0,\varepsilon_\omega)$ and $I_0=(t_0-\varepsilon_t,t_0+\varepsilon_t)\cap[0,M]$.
Since $\omega_0\in\supp{\mu}$, Lemma~\ref{lem:positive_mass_region} yields
\begin{equation}\label{eq:positive_mass_event}
\Probbig{X\in B_0,\ d(X,Y)\in I_0}>0,
\end{equation}
for iid $X,Y\sim\mu$.

Define
$$
k(x,y)\defin\bigl(F^\mu_x(d(x,y))-F^{\mu_{0}}_x(d(x,y))\bigr)^2=\Delta\bigl(x,d(x,y)\bigr)^2,
\quad x,y\in\Omega,
$$
and let $J\defin\mathbb E_\mu\big[k(X,Y)\big]$, for a pair of iid random objects $X,Y\sim\mu$.


By Lemma~\ref{lem:robust_separation}, one has
$$
k(X,Y)\, 1_{\{X\in B_0,\ d(X,Y)\in I_0\}}
\;\ge\;
c_0^2\, 1_{\{X\in B_0,\ d(X,Y)\in I_0\}}.
$$
% The indicator is needed to restrict to B_0 x I_0, where Lemma~\ref{lem:robust_separation} applies.
Taking expectations and using \eqref{eq:positive_mass_event} yields
\begin{equation}\label{eq:J_positive}
J=\mathsf{E}_\mu[k(X,Y)]
\ge
\mathsf{E}_\mu\left[k(X,Y)\, 1_{\{X\in B_0,\ d(X,Y)\in I_0\}}\right]
\ge
c_0^2\,\Probbig{X\in B_0,\ d(X,Y)\in I_0}
>0.
\end{equation}

Next, define
$$
e_n\defin\sup_{(\omega,t)\in\Omega\times\R_{\ge0}}\big|F_{n,\omega}^{\mu}(t)-F^\mu_\omega(t)\big|.
$$
By Glivenko--Cantelli, one has $e_n\convp0$. Observe that one can write
$$
\frac{1}{n} T_n^{\mathrm{CM}}
=\frac{1}{n^2}\sum_{i=1}^n\sum_{j=1}^n
\bigl(F_{n,X_i}^{\mu}(d(X_i,X_j))-F^{\mu_{0}}_{X_i}(d(X_i,X_j))\bigr)^2.
$$
Define $k_n(x,y)\defin\bigl(F_{n,x}^{\mu}(d(x,y))-F^{\mu_{0}}_x(d(x,y))\bigr)^2,$
so that $\frac{1}{n}T_n^{\mathrm{CM}} = \frac{1}{n^2}\sum_{i,j} k_n(X_i,X_j)$.
For any $x,y\in\Omega$ write
$$
F_{n,x}^{\mu}(d(x,y))-F^{\mu_{0}}_x(d(x,y))
=
\bigl(F^\mu_x(d(x,y))-F^{\mu_{0}}_x(d(x,y))\bigr)
+\bigl(F_{n,x}^{\mu}(d(x,y))-F^\mu_x(d(x,y))\bigr).
$$
The second term is bounded in absolute value by $e_n$, uniformly in $x,y$. The first term is bounded by $1$ in absolute value.
Therefore, using $|(a+b)^2-a^2|\le 2|a||b|+b^2$ for any $a,b\in\R$, it holds that
$$
\sup_{x,y\in\Omega}\,|k_n(x,y)-k(x,y)|\le 2e_n+e_n^2.
$$
Consequently,
\begin{equation}\label{eq:T_n_minus_V}
\left|\frac{T_n^{\mathrm{CM}}}{n}-V_n\right|
\le 2e_n+e_n^2
\xrightarrow{\mathsf{P}_\mu}0,
\end{equation} 
where
$$
V_n\defin\frac1{n^2}\sum_{i=1}^n\sum_{j=1}^n k(X_i,X_j).
$$

It remains to show $V_n\xrightarrow{\mathsf{P}_\mu}J$. Since $k$ is a positive function uniformly bounded by $1$, it is integrable.
Decompose
$$
V_n=\frac1{n^2}\sum_{i\neq j}k(X_i,X_j)+\frac1{n^2}\sum_{i=1}^n k(X_i,X_i),
$$
For the second term, one has $0\le n^{-2}\sum_{i}k(X_i,X_i)\le n^{-1} \to 0$ a.s.
Define the symmetrized kernel $\tilde{k}(x,y)\defin\frac12(k(x,y)+k(y,x))$, so $\tilde{k}(x,y)=\tilde{k}(y,x)$ and $0\le \tilde{k}\le 1$.
Then
$$
\sum_{i\neq j}k(X_i,X_j)
=2\sum_{1\le i<j\le n} \tilde{k}(X_i,X_j),
$$
and hence
$$
\frac1{n^2}\sum_{i\neq j}k(X_i,X_j)
=\left(1-\frac1n\right)U_n, \quad \text{where } \
U_n
\defin\frac{1}{\binom n2}\sum_{1\le i<j\le n} \tilde{k}(X_i,X_j).
$$
By the strong law of large numbers for U-statistics of order two applied to $\tilde{k}$ \citep[see, e.g., Theorem 5.4A of][]{Serfling1980}, one has $U_n\xrightarrow{\mathrm{a.s.}} \mathbb E_\mu\left(\tilde{k}(X,Y)\right).$
Since $X,Y$ are iid, $\mathbb E_\mu\left(k(X,Y)\right)=\mathbb E_\mu\left(k(Y,X)\right)$, thus from the definition of $\tilde{k}$ it follows that
$\mathbb E_\mu\left(\tilde{k}(X,Y)\right)=\mathbb E_\mu \left(k(X,Y)\right)=J$.
Therefore, $V_n\xrightarrow{\mathrm{a.s.}}J$, and in particular $V_n\xrightarrow{\mathsf{P}}J$. Combining this with \eqref{eq:T_n_minus_V} yields
$\frac{1}{n}T_n^{\mathrm{CM}}\xrightarrow{\mathsf{P}_\mu}J$, and hence $T_n^{\mathrm{CM}}\xrightarrow{\mathsf{P}_\mu}\infty$.
\begin{comment}
We show that this implies the divergence of $\sqrt{n} T_n^{\mathrm{CM}}$ (in probability). By \eqref{eq:J_positive}, one has $J>0$. We know that, for every $\varepsilon>0$,
$$
\mathsf{P}_\mu\left(\left|\frac{T_n^{\mathrm{CM}}}{n}-J\right|>\varepsilon\right)\to 0
$$
as $n\to\infty$. Now, set $\varepsilon=\frac{J}{2}>0$. If $\frac{T_n^{\mathrm{CM}}}{n}\le \varepsilon$, then
$$
\left|\frac{T_n^{\mathrm{CM}}}{n}-J\right|
=
J-\frac{T_n^{\mathrm{CM}}}{n}
\ge
\varepsilon,
$$
and thus
$$
\left\{\frac{T_n^{\mathrm{CM}}}{n}\le \varepsilon\right\}
\subset
\left\{\left|\frac{T_n^{\mathrm{CM}}}{n}-J\right|\ge \varepsilon\right\}.
$$
Consequently,
$$
\mathsf{P}_\mu\left(\frac{T_n^{\mathrm{CM}}}{n}\le \varepsilon\right)
\le
\mathsf{P}_\mu\left(\left|\frac{T_n^{\mathrm{CM}}}{n}-J\right|\ge \varepsilon\right)
\le
\mathsf{P}_\mu\left(\left|\frac{T_n^{\mathrm{CM}}}{n}-J\right|>\varepsilon/2\right)
\to 0.
$$
Take any $M_0>0$. Since $M_0/n^{3/2}\to 0$, there exists $n_0\in\mathbb{N}$ such that $M_0/n^{3/2}<\varepsilon$ for all $n\ge n_0$. Hence, for all $n\ge n_0$,
$$
\mathsf{P}_\mu\left(\sqrt{n}\,T_n^{\mathrm{CM}} \le M_0\right)
 = \mathsf{P}_\mu\left(\frac{T_n^{\mathrm{CM}}}{n}\le \frac{M_0}{n^{3/2}}\right)
\le \mathsf{P}_\mu\left(\frac{T_n^{\mathrm{CM}}}{n}\le \varepsilon\right)
\to 0,
$$
which proves $\sqrt{n}T_n^{\mathrm{CM}}\xrightarrow{\mathsf{P}_\mu}\infty$.
\end{comment}
\end{proof}

%----------------------------------%
% Fixed alternatives: composite null
%----------------------------------%
%----------------------------------%
\subsubsection{Composite null} \label{ap:sec_proofs_fixed_alternatives_composite}
%----------------------------------%

Theorem \ref{thm:fixed_alt_divergence_ks_composite} states the consistency of the Kolmogorov--Smirnov test statistic under fixed alternatives, when the null hypothesis is composite and a parameter is estimated.

%---------------------------%
\begin{proof}[Proof of Theorem~\ref{thm:fixed_alt_divergence_ks_composite}]
For $(\omega,t)\in\Omega\times\R_{\ge 0}$, write
$$
\sqrt{n}\,\bigl(F_{n,\omega}^\mu(t)-F^{\mu_{\hat{\btheta}}}_\omega(t)\bigr)
=
\G_{n,\mu}^{\mu}(\omega,t)
+
\sqrt{n}\,\bigl(F^\mu_\omega(t)-F^{\mu_{\hat{\btheta}}}_\omega(t)\bigr),
$$
with $\G_{n,\mu}^{\mu}$ as in \eqref{eq:G_n_nu}. 
Let $\mu_{\btheta_1}$ be the probability measure corresponding to $\btheta_1$.
Define
$$
\Delta^\star(\omega,t)\defin F^\mu_\omega(t)-F^{\mu_{\btheta_1}}_\omega(t),
\qquad
\delta^\star\defin
\sup_{(\omega,t)\in\Omega\times\R_{\ge 0}}
\bigl|\Delta^\star(\omega,t)\bigr|.
$$
Since $\mu\notin\{\mu_{\btheta}:\btheta\in\bTheta\}$, in particular $\mu\neq\mu_{\btheta_1}$, and hence there exists
$(\omega_0,t_0)\in\Omega\times\R_{\ge 0}$ such that $\Delta^\star(\omega_0,t_0)\neq 0$. Therefore, $\delta^\star \ge |\Delta^\star(\omega_0,t_0)|>0$.

Define
$$
r_n \defin \sup_{\omega\in\Omega, t\ge 0}
\bigl|F^{\mu_{\hat{\btheta}}}_\omega(t)-F^{\mu_{\btheta_1}}_\omega(t)\bigr|.
$$
By \eqref{eq:theta_star_profile_continuity}, for every $\varepsilon>0$ there exists $\eta>0$ such that
$$
\|\btheta-\btheta_1\|<\eta
\quad\Longrightarrow\quad
\sup_{\omega,t}\bigl|F^{\mu_{\btheta}}_\omega(t)-F^{\mu_{\btheta_1}}_\omega(t)\bigr|<\varepsilon.
$$
Since $\hat{\btheta}\convp\btheta_1$, it follows that
$\mathsf{P}^*(r_n>\varepsilon)\le
\mathsf{P}(\|\hat{\btheta}-\btheta_1\|\ge\eta)\to0$.
Thus, $r_n\xrightarrow{\mathsf{P}^*}0$.

Using the triangle inequality pointwise in $(\omega,t)$ and taking suprema,
\begin{align*}
\tilde{T}_n^{\mathrm{KS}}
&=
\sup_{\omega,t}\sqrt{n}\,\bigl|F_{n,\omega}^\mu(t)-F^{\mu_{\hat{\btheta}}}_\omega(t)\bigr|
\\
&\ge
\sup_{\omega,t}\sqrt{n}\,\bigl|F^\mu_\omega(t)-F^{\mu_{\hat{\btheta}}}_\omega(t)\bigr|
-
\sup_{\omega,t}\sqrt{n}\,\bigl|F_{n,\omega}^\mu(t)-F^\mu_\omega(t)\bigr|
\\
&=
\sup_{\omega,t}\sqrt{n}\,\bigl|F^\mu_\omega(t)-F^{\mu_{\hat{\btheta}}}_\omega(t)\bigr|
-
\|\G_{n,\mu}^{\mu}\|_\infty
\\
&\ge
\sup_{\omega,t}\sqrt{n}\,\bigl|F^\mu_\omega(t)-F^{\mu_{\btheta_1}}_\omega(t)\bigr|
-
\sup_{\omega,t}\sqrt{n}\,\bigl|F^{\mu_{\hat{\btheta}}}_\omega(t)-F^{\mu_{\btheta_1}}_\omega(t)\bigr|
-
\|\G_{n,\mu}^{\mu}\|_\infty
\\
&\ge
\sqrt{n}\,\delta^\star - \sqrt{n}\,r_n - \|\G_{n,\mu}^{\mu}\|_\infty.
\end{align*}


Fix any $M_0>0$ and define $a_n \defin \sqrt{n}\,\delta^\star - M_0.$
Then,
$$
\mathsf{P}^*(\tilde{T}_n^{\mathrm{KS}}\le M_0)
\le
\mathsf{P}^*\left(\|\G_{n,\mu}^{\mu}\|_\infty + \sqrt{n}\,r_n \ge a_n\right).
$$
Since the class $\cY$ is $\mu$-Donsker, it is also $\mu$-Glivenko--Cantelli. Hence,
$$
\sup_{\omega\in\Omega, t\ge0}
\left|
F^\mu_{n,\omega}(t)-F^\mu_\omega(t)
\right|
=
o_{\mathsf{P}^*}(1).
$$
Equivalently,
$$
\frac{1}{\sqrt n}\|\G_{n,\mu}^{\mu}\|_\infty
=
o_{\mathsf{P}^*}(1).
$$
Therefore,
$$
\mathsf{P}^*\left(\|\G_{n,\mu}^{\mu}\|_\infty \ge \frac{a_n}{2}\right)
=
\mathsf{P}^*\left(
\frac{1}{\sqrt n}\|\G_{n,\mu}^{\mu}\|_\infty
\ge
\frac{\delta^\star}{2}-\frac{M_0}{2\sqrt n}
\right)\to0.
$$

Also, since $r_n\to 0$ in $\mathsf{P}^*$, we have $\mathsf{P}^*(\sqrt{n}\,r_n\ge a_n/2)\to 0$. To see this, note that
$$
\mathsf{P}^*\left(\sqrt{n}\,r_n\ge \frac{a_n}{2}\right)
=
\mathsf{P}^*\left(r_n\ge \frac{a_n}{2\sqrt{n}}\right)
=
\mathsf{P}^*\left(r_n\ge \frac{\delta^\star}{2}-\frac{M_0}{2\sqrt{n}}\right).
$$
Since $\frac{M_0}{2\sqrt{n}}\to 0$, there exists $n_0$ such that $\frac{\delta^\star}{2}-\frac{M_0}{2\sqrt{n}}\ge \frac{\delta^\star}{4}$ for all $n\ge n_0$, and hence, for $n\ge n_0$,
$$
\mathsf{P}^*\left(\sqrt{n}\,r_n\ge \frac{a_n}{2}\right)
\le
\mathsf{P}^*\left(r_n\ge \frac{\delta^\star}{4}\right)
\to 0.
$$
Therefore,
$$
\mathsf{P}^*\left(\|\G_{n,\mu}^{\mu}\|_\infty + \sqrt{n}\,r_n \ge a_n\right)
\le
\mathsf{P}^*\left(\|\G_{n,\mu}^{\mu}\|_\infty \ge \frac{a_n}{2}\right)
+
\mathsf{P}^*\left(\sqrt{n}\,r_n \ge \frac{a_n}{2}\right)
\to 0.
$$
Hence, $\mathsf{P}^*(\tilde{T}_n^{\mathrm{KS}}\le M_0)\to 0$ for every $M_0>0$, which proves $\tilde{T}_n^{\mathrm{KS}}\xrightarrow{\mathsf{P}^*}\infty$.
\end{proof}
%----------------------------%

Theorem \ref{thm:fixed_alt_divergence_cvm_composite} proves the consistency of the Cramér--von Mises test statistic under fixed alternatives, when the null hypothesis is composite and a parameter is estimated.

%----------------------------%
\begin{proof}[Proof of Theorem~\ref{thm:fixed_alt_divergence_cvm_composite}]
Let $M=\mathrm{diam}(\Omega)$ and define $\Delta^\star(\omega,t):=F^\mu_\omega(t)-F^{\mu_{\btheta_1}}_\omega(t)$ for $(\omega,t)\in\Omega\times[0,M]$.
Since either \ref{a2} holds, or \ref{a1} holds and $\supp{\mu} = \Omega$, and since $\mu\neq\mu_{\btheta_1}$, there exist $\omega_0\in\supp{\mu}$ and $t_0\in(0,M]$
such that $\Delta^\star(\omega_0,t_0)\neq 0$.
Applying Lemma~\ref{lem:robust_separation} (with $\btheta_0$ replaced by $\btheta_1$),
we obtain $\varepsilon_\omega>0$, $\varepsilon_t>0$, and $c_0>0$ such that
\begin{equation}\label{eq:uniform_lower_bound_Delta_cvm}
|\Delta^\star(\omega,t)|\ge c_0
\quad\text{for all }(\omega,t)\in B_0\times I_0,
\end{equation}
where $B_0=B(\omega_0,\varepsilon_\omega)$ and $I_0=(t_0-\varepsilon_t,t_0+\varepsilon_t)\cap[0,M]$.
Since $\omega_0\in\supp{\mu}$, Lemma~\ref{lem:positive_mass_region} yields
\begin{equation}\label{eq:positive_mass_event_cvm_composite}
\mathsf{P}\bigl(X\in B_0,\ d(X,Y)\in I_0\bigr)>0
\end{equation}
for iid $X,Y\sim\mu$.

Define the kernel
$$
k^\star(x,y):=\bigl(F^\mu_x(d(x,y))-F^{\mu_{\btheta_1}}_x(d(x,y))\bigr)^2
=\Delta^\star(x,d(x,y))^2,\qquad x,y\in\Omega,
$$
and let $J^\star:=\mathsf{E}_\mu[k^\star(X,Y)]$ for iid $X,Y\sim\mu$.
From \eqref{eq:uniform_lower_bound_Delta_cvm} and \eqref{eq:positive_mass_event_cvm_composite},
one gets $J^\star>0$ in the same way \eqref{eq:J_positive} was obtained in the proof of Theorem \ref{thm:fixed_alternative_divergence_cvm}.

Next, define the stochastic errors
$$
e_n:=\sup_{(\omega,t)\in\Omega\times\R_{\ge0}}\bigl|F_{n,\omega}^{\mu}(t)-F^{\mu}_\omega(t)\bigr|,
\qquad
r_n:=\sup_{(\omega,t)\in\Omega\times\R_{\ge0}}\bigl|F^{\mu_{\hat{\btheta}}}_\omega(t)-F^{\mu_{\btheta_1}}_\omega(t)\bigr|.
$$
By Glivenko--Cantelli, $e_n\convp0$. %P^*? We are taking the supremum.
Moreover, by \eqref{eq:theta_star_profile_continuity} and
$\hat{\btheta}\convp\btheta_1$, the same argument as in the proof of Theorem~\ref{thm:fixed_alt_divergence_ks_composite} gives
$r_n\xrightarrow{\mathsf{P}^{\star}}0$.

One can express the Cramér--von Mises statistic as
$$
\frac{1}{n}\tilde{T}_n^{\mathrm{CM}}
=\frac{1}{n^2}\sum_{i=1}^n\sum_{j=1}^n
\bigl(F_{n,X_i}^{\mu}(d(X_i,X_j))-F^{\mu_{\hat{\btheta}}}_{X_i}(d(X_i,X_j))\bigr)^2,
$$
write, for any $x,y\in\Omega$,
\begin{align*}
F_{n,x}^{\mu}(d(x,y))-F^{\mu_{\hat{\btheta}}}_x(d(x,y))
=&
\bigl(F^\mu_x(d(x,y))-F^{\mu_{\btheta_1}}_x(d(x,y))\bigr)\\
&+\bigl(F_{n,x}^{\mu}(d(x,y))-F^\mu_x(d(x,y))\bigr)\\
&-\bigl(F^{\mu_{\hat{\btheta}}}_x(d(x,y))-F^{\mu_{\btheta_1}}_x(d(x,y))\bigr).
\end{align*}
The second term is bounded by $e_n$ uniformly in $x,y$, and the third term is bounded by $r_n$ uniformly in $x,y$.
Let
$$
k_n^\star(x,y):=\bigl(F_{n,x}^{\mu}(d(x,y))-F^{\mu_{\hat{\btheta}}}_x(d(x,y))\bigr)^2.
$$
Now, use $|(a+b)^2-a^2|\le 2|a||b|+b^2$ for any $a,b\in \R$, taking $a \defin \bigl(F^\mu_x(d(x,y))-F^{\mu_{\btheta_1}}_x(d(x,y))\bigr)$ and 
$b \defin \bigl(F_{n,x}^{\mu}(d(x,y))-F^\mu_x(d(x,y))\bigr)
-\bigl(F^{\mu_{\hat{\btheta}}}_x(d(x,y))-F^{\mu_{\btheta_1}}_x(d(x,y))\bigr)$. Since $|a|\le 1$, we obtain
$$
\sup_{x,y\in\Omega}\bigl|k_n^\star(x,y)-k^\star(x,y)\bigr|
\le 2(e_n+r_n)+(e_n+r_n)^2
\xrightarrow{\mathsf{P}^{\star}}0.
$$
Therefore, with
$$
V_n^\star:=\frac{1}{n^2}\sum_{i=1}^n\sum_{j=1}^n k^\star(X_i,X_j),
\qquad
\frac{1}{n}\tilde{T}_n^{\mathrm{CM}}=\frac{1}{n^2}\sum_{i,j}k_n^\star(X_i,X_j),
$$
we have $\bigl|\frac{1}{n}\tilde{T}_n^{\mathrm{CM}}-V_n^\star\bigr|\xrightarrow{\mathsf{P}^{\star}}0$.

Finally, $k^\star$ is bounded and nonnegative, hence integrable, and the same U-statistic LLN argument as in the proof of
Theorem~\ref{thm:fixed_alternative_divergence_cvm} yields $V_n^\star\xrightarrow{\mathsf{P}^{\star}}J^\star$.
Thus, $\frac{1}{n}\tilde{T}_n^{\mathrm{CM}}\xrightarrow{\mathsf{P}^{\star}}J^\star>0$, and consequently $\tilde{T}_n^{\mathrm{CM}}\xrightarrow{\mathsf{P}^{\star}}\infty$.
\end{proof}


%----------------------------%
\subsection{Local alternatives}\label{ap:sec_local_alternatives}
%----------------------------%
\subsubsection{Simple null hypothesis}\label{ap:sec_local_alternatives_simple}
%----------------------------%

This section contains the proofs of Proposition \ref{prop:bernoulli} and Theorems \ref{thm:local_alt_ks_simple} and \ref{thm:local_alt_CvM_simple}, which establish the behavior of the Kolmogorov--Smirnov and Cramér--von Mises test statistics under local alternatives.

\begin{proof}[Proof of Proposition~\ref{prop:bernoulli}]
Define the empirical process under $Q_n$ by $\G_{n,Q_n}^{Q_n}(f) \defin \sqrt{n}(P_n-Q_n)f$, $f\in \cF$. Consider two mutually independent iid sequences $U_i$ and $V_i$, $i\ge1$, where $U_i\sim \mu_{0}$ and $V_i\sim G$. For each $n\in\mathbb{N}$, consider an iid sequence $B_{n,i}$, $i\ge1$, independent of $(U_i)_{i\ge1}$ and $(V_i)_{i\ge1}$, where $B_{n,i}\sim\mathrm{Bernoulli}(n^{-1/2})$. Define $X_{n,i} \defin U_i$ if $B_{n,i}=0$ and $X_{n,i} \defin V_i$ if $B_{n,i}=1$. 
Then $X_{n,1},\ldots,X_{n,n}$ are iid with common distribution $Q_n=\big(1-\frac{1}{\sqrt{n}}\big)\mu_{0}+\frac{1}{\sqrt{n}} G$. Observe that for any $f \in \cF$, it holds that $f(X_{n,i})=(1-B_{n,i})f(U_i)+B_{n,i} f(V_i)$. Write $N_{n,1}\defin\sum_{i=1}^n B_{n,i}$ and $N_{n,0}\defin n-N_{n,1}$. One has that
\begin{align}
\G_{n,Q_n}^{Q_n}(f)
&=\frac{1}{\sqrt{n}}\sum_{i=1}^n\Bigl((1-B_{n,i})f(U_i)+B_{n,i} f(V_i)-Q_n f\Bigr)\nonumber\\
&=\frac{1}{\sqrt{n}}\sum_{i=1}^n(1-B_{n,i})\Bigl(f(U_i)-\mu_{0}f\Bigr)\nonumber
+\frac{1}{\sqrt{n}}\sum_{i=1}^n B_{n,i}\Bigl(f(V_i)-Gf\Bigr)
\\
& \quad +\frac{1}{\sqrt{n}}\Bigl(N_{n,0}\,\mu_{0}f+N_{n,1}\,Gf-nQ_n f\Bigr)\nonumber\\
&=\frac{1}{\sqrt{n}}\sum_{i=1}^n(1-B_{n,i})\Bigl(f(U_i)-\mu_{0}f\Bigr) \label{eq:local_alt_split_1}
+\frac{1}{\sqrt{n}}\sum_{i=1}^n B_{n,i}\Bigl(f(V_i)-Gf\Bigr)\\ 
&\quad
+\sqrt{n}\left(\frac{N_{n,0}}{n}-\Big(1-\frac{1}{\sqrt{n}}\Big)\right)\mu_{0}f\label{eq:local_alt_split_2}
+\sqrt{n}\left(\frac{N_{n,1}}{n}-\frac{1}{\sqrt{n}}\right)Gf,
\end{align}
%we introduce \frac{1}{\sqrt{n}} for centering
Since $\frac{N_{n,0}}{n}-\big(1-\frac{1}{\sqrt{n}}\big)=-\big(\frac{N_{n,1}}{n}-\frac{1}{\sqrt{n}}\big)$, \eqref{eq:local_alt_split_2} combines into
$$
R_n(f)
\defin
\frac{N_{n,1} - \sqrt{n}}{\sqrt{n}}\big(G-\mu_{0}\big)f.
$$
Since $N_{n,1}\sim\mathrm{Bin}(n,n^{-1/2})$, it follows that 
$\frac{\E{(N_{n,1}-\sqrt{n})^2}}{n} = \frac{\mathrm{Var}(N_{n,1})}{n}=\frac{1}{\sqrt{n}}\left(1-\frac{1}{\sqrt{n}}\right)\to 0$. Thus, since $\sup_{f\in\cF}\big|(G-\mu_{0})f\big|<\infty$ by assumption, it holds that
\begin{equation}\label{eq:local_alt_Rn_vanish}
\sup_{f\in\cF}|R_n(f)|\convp0.
\end{equation}

Next, consider the second sum in \eqref{eq:local_alt_split_1}. Define, for $n\ge0$, the empirical process based on $V_i\sim G$ by
$$
\G_{n,G}^{G}(f)\defin
\begin{cases}
\frac{1}{\sqrt{n}}\sum_{i=1}^{n}\Big(f(V_i)-Gf\Big), & n\ge1,\\
0, & n=0.
\end{cases}
$$
Let $I_{n,1}\defin\{i\in\{1,\ldots,n\}:B_{n,i}=1\}$ and, on the event $\{N_{n,1}\ge 1\}$, write $I_{n,1}=\{j_{n,1}<\cdots<j_{n,N_{n,1}}\}$. Then,
$$
\frac{1}{\sqrt{n}}\sum_{i=1}^n B_{n,i}\Big(f(V_i)-Gf\Big)
=
\frac{1}{\sqrt{n}}\sum_{\ell=1}^{N_{n,1}}\Big(f(V_{j_{n,\ell}})-Gf\Big).
$$
Since $(V_i)_{i\ge1}$ are iid and independent of $(B_{n,i})_{i=1}^n$ for every $n$, then conditional on $B_{n,1},\ldots,B_{n,n}$, if $N_{n,1}=m$, then $(V_{j_{n,\ell}})_{\ell=1}^{m}$ has the same law as $(V_\ell)_{\ell=1}^{m}$. Hence, the following equality in distribution also holds unconditionally:
$$
\frac{1}{\sqrt{n}}\sum_{i=1}^n B_{n,i}\Big(f(V_i)-Gf\Big)
\stackrel{d}{=}
\sqrt{\frac{N_{n,1}}{n}}\ \G_{N_{n,1},G}^{G}(f),
$$
as random elements in $\ell^\infty(\cF)$.
By the Donsker property, $\G_{n,G}^{G}\convl \G_{G}$ in $\ell^\infty(\cF)$ for some tight limit $\G_{G}$. Since $\cF$ is Donsker under $G$, the empirical process $\G_{n,G}^{G}$ is asymptotically tight in
$\ell^\infty(\cF)$; \citep[see, e.g., page 139 in][]{VanderVaartWellner2023}.
By the continuous mapping theorem for asymptotic tightness
\citep[see, e.g., Problem~7(i) in Section~1.3 of][]{VanderVaartWellner2023}, the sequence $\sup_{f \in \cF}|\G_{n,G}^{G}(f)|$ is
asymptotically tight in $\R$. In particular, for every $\varepsilon>0$ there exists $M_\varepsilon<\infty$ such that
\begin{equation}\label{eq:local_alt_asympt_tight_G}
\limsup_{n\to\infty}\mathsf{P}\left(\sup_{f \in \cF}|\G_{n,G}^{G}(f)|>M_\varepsilon\right)\le \frac{\varepsilon}{2}.
\end{equation}

For every fixed $m>0$, one has
\begin{equation}\label{eq:local_alt_random_index_bound}
\mathsf{P}\left(\sup_{f \in \cF}|\G_{N_{n,1},G}^{G}(f)|>M_\varepsilon\right)
\le
\mathsf{P}(N_{n,1}\le m)+
\sup_{k\ge m}\mathsf{P}\left(\sup_{f \in \cF}|\G_{k,G}^{G}(f)|>M_\varepsilon\right)
\end{equation}
% For each k\ge 1 define
% $$
% A_k \defin \left\{\sup_{f \in \cF}|\G_{k,G}^{G}(f)|>M_\varepsilon\right\}.
% $$
% Then, for any fixed $m\in\mathbb{Z}_{>0}$,
% $$x
% \mathsf{P}(A_{N_{n,1}})
% =\mathsf{P}(A_{N_{n,1}}\cap\{N_{n,1}\le m\})+\mathsf{P}(A_{N_{n,1}}\cap\{N_{n,1}> m\})
% \le \mathsf{P}(N_{n,1}\le m)+\mathsf{P}(A_{N_{n,1}}\cap\{N_{n,1}> m\}).
% $$
% Moreover,
% $$
% A_{N_{n,1}}\cap\{N_{n,1}> m\}\subseteq \bigcup_{k\ge m+1}\bigl(A_k\cap\{N_{n,1}=k\}\bigr),
% $$
% hence
% $$
% \mathsf{P}(A_{N_{n,1}}\cap\{N_{n,1}> m\})
% \le \sum_{k\ge m+1}\mathsf{P}(A_k\cap\{N_{n,1}=k\})
% \le \sum_{k\ge m+1}\mathsf{P}(A_k)\,\mathsf{P}(N_{n,1}=k)
% \le \sup_{k\ge m+1}\mathsf{P}(A_k)
% \le \sup_{k\ge m}\mathsf{P}(A_k).
% $$
% Therefore,
% $$
% \mathsf{P}(A_{N_{n,1}}) \le \mathsf{P}(N_{n,1}\le m)+\sup_{k\ge m}\mathsf{P}(A_k).
% $$
By \eqref{eq:local_alt_asympt_tight_G}, there exists $m_0\in\mathbb{N}$ such that
$$
\sup_{k\ge m_0}\mathsf{P}\left(\sup_{f \in \cF}|\G_{k,G}^{G}(f)|>M_\varepsilon\right)\le \varepsilon.
$$
Applying \eqref{eq:local_alt_random_index_bound} with $m=m_0$, and since $N_{n,1}\xrightarrow{\mathsf{P}}\infty$, we obtain
$$
\limsup_{n\to\infty}\mathsf{P}\left(\sup_{f \in \cF}|\G_{N_{n,1},G}^{G}(f)|>M_\varepsilon\right)
\le
\limsup_{n\to\infty}\mathsf{P}(N_{n,1}\le m_0)+\varepsilon
=
\varepsilon.
$$
Therefore,
$$
\sup_{f \in \cF}|\G_{N_{n,1},G}^{G}(f)|=O_{\mathsf{P}}(1).
$$
Moreover, $\frac{N_{n,1}}{n} \convp 0$ because $\Ebig{\frac{N_{n,1}}{n}} = n^{-1/2}$. Consequently,
\begin{align}\label{eq:local_alt_G_part_vanish}
\sup_{f \in \cF}
\bigg|
\frac{1}{\sqrt{n}}\sum_{i=1}^n B_{n,i}\Big(f(V_i)-Gf\Big)
\bigg|
\stackrel{d}{=}
\sqrt{\frac{N_{n,1}}{n}} \sup_{f \in \cF}|\G_{N_{n,1},G}^{G}(f)|
\convp 0.
\end{align}

It remains to handle the first term in \eqref{eq:local_alt_split_1}. Define, for $n\ge0$, the empirical process based on $U_i\sim\mu_0$ by
$$
\G_{n,\mu_0}^{\mu_0}(f)\defin
\begin{cases}
\frac{1}{\sqrt{n}}\sum_{i=1}^{n}\big(f(U_i)-\mu_{0}f\big), & n\ge1,\\
0, & n=0.
\end{cases}
$$
Then, one has
$$
\frac{1}{\sqrt{n}}\sum_{i=1}^n (1-B_{n,i})\Big(f(U_i)-\mu_{0}f\Big)
=
\G_{n,\mu_0}^{\mu_0}(f)
-
\frac{1}{\sqrt{n}}\sum_{i=1}^n B_{n,i}\Big(f(U_i)-\mu_{0}f\Big).
$$
For the second term on the RHS of the equality, by the definition of $j_{n,\ell}$, one has
$$
\frac{1}{\sqrt{n}}\sum_{i=1}^n B_{n,i}\Big(f(U_i)-\mu_{0}f\Big)
=
\frac{1}{\sqrt{n}}\sum_{\ell=1}^{N_{n,1}}\Big(f(U_{j_{n,\ell}})-\mu_{0}f\Big).
$$
Since $(U_i)_{i\ge1}$ are iid and independent of $(B_{n,i})_{i=1}^n$, then conditional on $B_{n,1},\ldots,B_{n,n}$, $(U_{j_{n,\ell}})_{\ell=1}^{m}$ has the same law as $(U_\ell)_{\ell=1}^{m}$, if $N_{n,1}=m$. Hence, the following equality in distribution also holds unconditionally:
$$
\frac{1}{\sqrt{n}}\sum_{i=1}^n B_{n,i}\Big(f(U_i)-\mu_{0}f\Big)
\stackrel{d}{=}
\sqrt{\frac{N_{n,1}}{n}}\ \G_{N_{n,1}, \mu_0}^{\mu_0}(f)
$$
as random elements in $\ell^\infty(\cF)$.
Since $\cF$ is $\mu_{0}$-Donsker, the same argument as above yields
$\sup_{f \in \cF}|\G_{N_{n,1},\mu_0}^{\mu_0}(f)|=O_{\mathsf{P}}(1)$, and hence
\begin{equation}\label{eq:local_alt_P_part_correction_vanish}
\sup_{f \in \cF}
\Big|
\frac{1}{\sqrt{n}}\sum_{i=1}^n B_{n,i}\Big(f(U_i)-\mu_{0}f\Big)
\Big|
\stackrel{d}{=}
\sqrt{\frac{N_{n,1}}{n}} \sup_{f \in \cF}|\G_{N_{n,1},\mu_0}^{\mu_0}(f)|
\convp 0.
\end{equation}

Finally, observe that applying the triangle inequality and taking suprema, it follows from \eqref{eq:local_alt_split_1}--\eqref{eq:local_alt_split_2}  that 
\begin{align*}
\sup_{f \in \cF}\big|\G_{n,Q_n}^{Q_n}(f)-\G_{n,\mu_0}^{\mu_0}(f)\big|
&\le
\sup_{f \in \cF}
\Big|\frac{1}{\sqrt n}\sum_{i=1}^n B_{n,i}\Big(f(U_i)-\mu_{0}f\Big)\Big|
\\
&\quad +
\sup_{f \in \cF}
\Big|\frac{1}{\sqrt n}\sum_{i=1}^n B_{n,i}\Big(f(V_i)-Gf\Big)\Big|
\\
&\quad + \sup_{f \in \cF}|R_n(f)|.
\end{align*}
Combining \eqref{eq:local_alt_Rn_vanish}, \eqref{eq:local_alt_G_part_vanish}, and \eqref{eq:local_alt_P_part_correction_vanish} gives \eqref{eq:sup_bernoulli_prob}.
\end{proof}
%----------------------------%

Theorem \ref{thm:local_alt_ks_simple} gives the weak limit of the Kolmogorov--Smirnov statistic under local alternatives, when the null hypothesis is simple.

%----------------------------%
\begin{proof}[Proof of Theorem \ref{thm:local_alt_ks_simple}]
Observe that since $0\le y_{\omega,t}\le 1$ for all $(\omega,t)\in\Omega\times\R_{\ge0}$, it holds that $\sup_{\omega,t}\big|(G-\mu_{0})y_{\omega,t}\big|\le 2$, so the conditions of Proposition \ref{prop:bernoulli} are satisfied for $\cY$. Under the local alternative defined in \eqref{eq:local_alt_Qn_def}, an application of Corollary \ref{cor:bernoulli} to the class of functions $\cY$ yields $\G_{n,Q_n}^{Q_n} \convl  \G_{0}$ in $\ell^\infty(\cY)$.
Combined with \eqref{eq:local_alt_process_decomp}, this gives
\begin{equation}\label{eq:local_alt_weak_limit_process}
\sqrt{n}\Big(F_{n,\omega}^{Q_n}(t)-F^{\mu_0}_\omega(t)\Big)
\convl
\G_{0}(\omega,t)+h(\omega,t)
\qquad\text{in }\ell^\infty(\cY).
\end{equation}
Finally, by an application of the continuous mapping theorem, we obtain
\begin{equation*}
T_n^{\mathrm{KS}}
\convl
\sup_{\omega\in\Omega, t \ge 0}
\left|\G_{0}(\omega,t)+h(\omega,t)\right|.
\end{equation*}
This concludes the result.
\end{proof}
%----------------------------%

%----------------------------%

Lemma \ref{ap:lem_DCT_P_hat_local} is the analogue of Lemma~\ref{ap:lem_DCT_P_hat} for the local alternative.
Let $M=\operatorname{diam}(\Omega)$. For $0<B<\infty$ and a modulus of continuity $w:[0,\infty)\to[0,\infty)$, let $\mathcal G_{B,w}$ denote the class of functions on $\Omega\times[0,M]$ that are uniformly bounded by $B$ and have modulus of continuity $w$, that is,
$$
\mathcal G_{B,w}
:=
\left\{
g:\Omega\times[0,M]\to\mathbb R:
\|g\|_\infty\le B,\ 
|g(z)-g(z')|\le w\big(d_{\Omega\times\mathbb R_{\ge0}}(z,z')\big)
\text{ for all } z,z'\in\Omega\times[0,M]
\right\}.
$$

Here we only obtain convergence in probability, whereas Lemma~\ref{ap:lem_DCT_P_hat} yields almost sure convergence. This is why we prove below a uniform version, which is the extra ingredient needed in the proof of Theorem~\ref{thm:local_alt_CvM_simple}.
%----------------------------%
\begin{lemma}\label{ap:lem_DCT_P_hat_local}
%----------------------------%
Under the local alternative \eqref{eq:local_alt_Qn_def}, it holds that
$$
\sup_{g\in\mathcal G_{B,w}}
\left|
\int_{\Omega\times[0,M]} g(\omega,t)\,\rd F_{n,\omega}^{Q_n}(t)\,\rd P_n(\omega)
-
\int_{\Omega\times[0,M]} g(\omega,t)\,\rd F_{\omega}^{\mu_0}(t)\,\rd \mu_0(\omega)
\right|
=o_P(1).
$$
\end{lemma}
\begin{proof}
For $g\in\mathcal G_{B,w}$, define
$$
r_{n,g}(\omega) \defin \int_{[0,M]} g(\omega, t) \, \rd F_{n,\omega}^{Q_n}(t),
\qquad
r_g(\omega) \defin \int_{[0,M]} g(\omega, t) \, \rd F_{\omega}^{\mu_0}(t).
$$
Then,
\begin{equation}\label{eq:local_alt_DCT_split}
\sup_{g\in\mathcal G_{B,w}}
\left|\int_{\Omega} r_{n,g}(\omega) \, \rd P_n(\omega) - \int_{\Omega} r_g(\omega) \, \rd \mu_0(\omega)\right|
\le
A_n+B_n,
\end{equation}
where
$$
A_n \defin \sup_{g\in\mathcal G_{B,w}}\sup_{\omega\in\Omega}|r_{n,g}(\omega)-r_g(\omega)|
$$
and
$$
B_n \defin
\sup_{g\in\mathcal G_{B,w}}
\left|
\int_{\Omega} r_g(\omega) \, \rd(P_n-\mu_0)(\omega)
\right|.
$$

We first control $A_n$. Fix $\varepsilon>0$, and choose $\delta>0$ such that $w(\delta)\le \varepsilon$. Repeating the regularization argument in the proof of Lemma~\ref{ap:lem_sup_r_n_r}, with
$$
g_\varepsilon(\omega,t)
:=
\inf_{s\in[0,M]}
\left\{
g(\omega,s)+\frac{2B}{\delta}|t-s|
\right\},
$$
gives, uniformly over $g\in\mathcal G_{B,w}$,
$$
A_n
\le
C\varepsilon
+
C_{B,\delta,M}
\sup_{\omega\in\Omega, t\in[0,M]}
\left|
F_{n,\omega}^{Q_n}(t)-F_\omega^{\mu_0}(t)
\right|,
$$
for constants $C$ and $C_{B,\delta,M}$ independent of $n$. It remains to show that
\begin{equation}\label{eq:local_alt_DCT_uniform_profiles}
\sup_{\omega \in\Omega, t\in[0,M]}
\left|F_{n,\omega}^{Q_n}(t)-F_\omega^{\mu_0}(t)\right|
\convp 0.
\end{equation}
The decomposition
$$
F_{n,\omega}^{Q_n}(t)-F_\omega^{\mu_0}(t)
=
(P_n-Q_n)y_{\omega,t}+(Q_n-\mu_0)y_{\omega,t}
$$
gives
$$
\sup_{\omega \in\Omega, t\in[0,M]}
\left|F_{n,\omega}^{Q_n}(t)-F_\omega^{\mu_0}(t)\right|
\le
\frac{1}{\sqrt{n}}
\sup_{\omega \in\Omega, t\in[0,M]}
\left|\G_{n,Q_n}^{Q_n}(\omega,t)\right|
+
\frac{1}{\sqrt{n}}
\sup_{\omega \in\Omega, t\in[0,M]}
\left|(G-\mu_0)y_{\omega,t}\right|.
$$
By Corollary~\ref{cor:bernoulli} applied to $\cY$, one has $\G_{n,Q_n}^{Q_n}\convl \G_0$ in $\ell^\infty(\cY)$, and therefore
$$
\sup_{\omega \in\Omega, t\in[0,M]}
\left|\G_{n,Q_n}^{Q_n}(\omega,t)\right|
=
O_P(1).
$$
Moreover, one has $\sup_{\omega,t}|(G-\mu_0)y_{\omega,t}|\le 2$. Thus, \eqref{eq:local_alt_DCT_uniform_profiles} holds and hence $A_n=o_P(1)$.

We next control $B_n$. Write
\begin{align}\label{eq:local_alt_DCT_Bn_split}
B_n
\le
\sup_{g\in\mathcal G_{B,w}} |(P_n-Q_n)r_g|
+
\sup_{g\in\mathcal G_{B,w}} |(Q_n-\mu_0)r_g|.
\end{align}
Since $|r_g(\omega)|\le B$ for every $g\in\mathcal G_{B,w}$ and $\omega\in\Omega$, one has
$$
\sup_{g\in\mathcal G_{B,w}} |(Q_n-\mu_0)r_g|
\le
\frac{2B}{\sqrt{n}}.
$$

For the first term in \eqref{eq:local_alt_DCT_Bn_split}, note that
$$
r_g(\omega)=\int_\Omega g(\omega,d(\omega,x))\,\rd\mu_0(x).
$$
Hence, for $\omega,\omega^\prime\in\Omega$,
$$
|r_g(\omega)-r_g(\omega^\prime)|
\le
\int_\Omega
\left|
g(\omega,d(\omega,x))-g(\omega^\prime,d(\omega^\prime,x))
\right|\,\rd\mu_0(x).
$$
The triangle inequality gives
$$
d_{\Omega\times\mathbb R_{\ge0}}
\big((\omega,d(\omega,x)),(\omega^\prime,d(\omega^\prime,x))\big)
\le 2d(\omega,\omega^\prime),
$$
and therefore
$$
|r_g(\omega)-r_g(\omega^\prime)|
\le
w\big(2d(\omega,\omega^\prime)\big),
$$
uniformly in $g\in\mathcal G_{B,w}$. Thus, the family $\{r_g:g\in\mathcal G_{B,w}\}$ is uniformly bounded by $B$ and equicontinuous with common modulus $u\mapsto w(2u)$.

Fix $\eta>0$. Choose $\delta>0$ such that $w(2\delta)\le \eta/3$, and let $\omega_1,\dots,\omega_m$ be a finite $\delta$-net of $\Omega$. Consider the set
$$
S:=\left\{
\big(r_g(\omega_1),\dots,r_g(\omega_m)\big):g\in\mathcal G_{B,w}
\right\}\subset[-B,B]^m.
$$
Since $S$ is totally bounded, it admits a finite $\eta/3$-net with centers belonging to $S$. Hence, there exist $g_1,\dots,g_N\in\mathcal G_{B,w}$ such that, for every $g\in\mathcal G_{B,w}$, there exists $k\in\{1,\dots,N\}$ satisfying
$$
\max_{1\le j\le m}|r_g(\omega_j)-r_{g_k}(\omega_j)|\le \eta/3.
$$
If $d(\omega,\omega_j)<\delta$, then
$$
|r_g(\omega)-r_{g_k}(\omega)|
\le
|r_g(\omega)-r_g(\omega_j)|
+
|r_g(\omega_j)-r_{g_k}(\omega_j)|
+
|r_{g_k}(\omega_j)-r_{g_k}(\omega)|
\le \eta.
$$
Thus, $\|r_g-r_{g_k}\|_\infty\le \eta$, and so
$$
\sup_{g\in\mathcal G_{B,w}} |(P_n-Q_n)r_g|
\le
\max_{1\le k\le N}|(P_n-Q_n)r_{g_k}|+2\eta.
$$

For each fixed $k$,
$$
\Var{(P_n-Q_n)r_{g_k}}
=
\frac{\Var{r_{g_k}(X_1)}}{n}
\le
\frac{B^2}{n},
$$
so $(P_n-Q_n)r_{g_k}=o_P(1)$ by Chebyshev's inequality. Since $N$ is finite,
$$
\max_{1\le k\le N}|(P_n-Q_n)r_{g_k}|=o_P(1).
$$
Letting $\eta\downarrow0$, we obtain
$$
\sup_{g\in\mathcal G_{B,w}} |(P_n-Q_n)r_g|=o_P(1).
$$
Therefore, $B_n=o_P(1)$, and the result follows.
\end{proof}

%----------------------------%
\begin{proof}[Proof of Theorem~\ref{thm:local_alt_CvM_simple}]
The proof is analog to that of Theorem~\ref{thm:test_statistics_simp}.

From \eqref{eq:local_alt_process_decomp}, we have
$$
T_n^{\mathrm{CM}}
=
\int_{\Omega\times\mathbb R_{\ge0}}
\big(\G_{n,Q_n}^{Q_n}(\omega,t)+h(\omega,t)\big)^2\,
\rd F_{n,\omega}^{Q_n}(t)\,\rd P_n(\omega).
$$
By Corollary \ref{cor:bernoulli},
$$
\G_{n,Q_n}^{Q_n}+h \convl  \G_0+h
\quad\text{in }\ell^\infty(\cY).
$$
In the same way as in the proof of Theorem \ref{thm:test_statistics_simp}, it can be shown that $\G_0+h$ is separable, and therefore the Skorohod representation theorem yields a probability space
on which it holds that
\begin{align}\label{eq:aux_cvm_sup_convas_2}
\sup_{\omega\in\Omega, t \ge 0}
\big|\G_{n,Q_n}^{Q_n}(\omega,t)-\G_0(\omega,t)\big|
\convas 0.
\end{align}
Now, write 
\begin{align}
&\left| \notag
\int \big(\G_{n,Q_n}^{Q_n}+h\big)^2\,\rd F_{n,\omega}^{Q_n}\,\rd P_n
-
\int \big(\G_0+h\big)^2\,\rd F^{\mu_0}_\omega\,\rd\mu_{0}
\right|
\\[0.15cm]
&\le \label{eq:local_cvm_aux_ineq_1}
\left|
\int \big(\G_{n,Q_n}^{Q_n}+h\big)^2-\big(\G_0+h\big)^2 \,
\rd F_{n,\omega}^{Q_n}(t)\,\rd P_n(\omega)
\right| 
\\
&\quad+ \label{eq:local_cvm_aux_ineq_2}
\left|
\int \big(\G_0(\omega,t)+h(\omega,t)\big)^2\,
\rd F_{n,\omega}^{Q_n}(t)\,\rd P_n(\omega)
-
\int \big(\G_0(\omega,t)+h(\omega,t)\big)^2\,
\rd F^{\mu_0}_\omega(t)\,\rd\mu_{0}(\omega)
\right|.
\end{align}
For \eqref{eq:local_cvm_aux_ineq_1}, one has:
\begin{align*}
\left|
\int
\Big(
\big(\G_{n,Q_n}^{Q_n}+h\big)^2-\big(\G_0+h\big)^2
\Big)
\,
\rd F_{n,\omega}^{Q_n}(t)\,\rd P_n(\omega)
\right|
&\le
\sup_{\omega,t}
\bigg|
\big(\G_{n,Q_n}^{Q_n}(\omega,t)+h(\omega,t)\big)^2
\\
&\quad -
\big(\G_0(\omega,t)+h(\omega,t)\big)^2
\bigg| \convas 0.
\end{align*}
The convergence to $0$ follows from the identity $(a+h)^2-(b+h)^2=(a-b)(a+b+2h)$, the almost sure convergence
\eqref{eq:aux_cvm_sup_convas_2}, and the facts that $|h|\le 1$ and
$\sup_{\omega,t}|\G_0(\omega,t)|<\infty$ almost surely. %this was shown in the proof of Theorem~\ref{thm:test_statistics_simp}. 
%Boundedness of $G_{n,Q_n}$ follows from boundedness of $G_{\btheta_0}$ and the almost sure convergence \eqref{eq:aux_cvm_sup_convas_2}.

To control \eqref{eq:local_cvm_aux_ineq_2}, observe first that $\G_0$ has bounded and uniformly continuous sample paths almost surely on $\Omega\times[0,M]$, by Lemma~\ref{ap:lem_continuity_sample_paths_G_rho2}. Moreover, $h$ is deterministic and Lipschitz on $\Omega\times[0,M]$.

Fix $\eta>0$. Since $\|\G_0\|_\infty<\infty$ almost surely, there exists $L<\infty$ such that
$$
\mathsf{P}\big(\|\G_0\|_\infty\le L\big)\ge 1-\eta/2.
$$
For $\delta>0$, define
$$
W(\delta)
:=
\sup_{\substack{
d_{\Omega\times\R_{\ge0}}((\omega,t),(\omega^\prime,t^\prime))\le\delta}}
|\G_0(\omega,t)-\G_0(\omega^\prime,t^\prime)|.
$$
Since $\G_0$ has uniformly continuous sample paths almost surely, $W(\delta)\to0$ in probability as $\delta\downarrow0$. Hence, one may choose a decreasing sequence $(\delta_m)_{m\ge1}$ such that
$$
\mathsf{P}\big(W(\delta_m)>2^{-m}\big)\le \eta\,2^{-m-1},
\qquad m\ge1.
$$
Define $w_0:[0,\infty)\to[0,\infty)$ by
$$
w_0(u):=
\begin{cases}
0, & u=0,\\
2^{-m}, & \delta_{m+1}<u\le\delta_m,\\
\max\{2L,2^{-1}\}, & u>\delta_1.
\end{cases}
$$
Then $w_0$ is nondecreasing and $w_0(u)\to0$ as $u\downarrow0$. Let
$$
A_\eta
:=
\left\{\|\G_0\|_\infty\le L\right\}
\cap
\bigcap_{m\ge1}\left\{W(\delta_m)\le 2^{-m}\right\},
$$
which satisfies $\mathsf{P}(A_\eta)\ge 1-\eta$. On $A_\eta$, one has
$$
|\G_0(\omega,t)-\G_0(\omega^\prime,t^\prime)|
\le
w_0\!\left(
d_{\Omega\times\mathbb R_{\ge0}}((\omega,t),(\omega^\prime,t^\prime))
\right)
$$
for all $(\omega,t),(\omega^\prime,t^\prime)\in\Omega\times[0,M]$.

By \eqref{eq:rho2y_prod_bound} applied with the measure $\nu\in\{\mu_0,G\}$, there exists $C_\nu<\infty$ such that
$$
\left|F^\nu_\omega(t)-F^\nu_{\omega^\prime}(t^\prime)\right|
=
\left|\nu\big(y_{\omega,t}-y_{\omega^\prime,t^\prime}\big)\right|
\le
\nu\big(\big(y_{\omega,t}-y_{\omega^\prime,t^\prime}\big)^2\big)
\le
C_\nu\, d_{\Omega\times\mathbb R_{\ge0}}\big((\omega,t),(\omega^\prime,t^\prime)\big).
$$
Hence, $(\omega,t)\mapsto F^\nu_\omega(t)$ is Lipschitz for each $\nu\in\{\mu_0,G\}$, and therefore $h(\omega,t)=F_{\omega}^G(t)-F_{\omega}^{\mu_0}(t)$ is Lipschitz. Let $L_h$ be a Lipschitz constant of $h$. On $A_\eta$, $\|(\G_0+h)^2\|_\infty\le (L+\|h\|_\infty)^2,$
and, using $|a^2-b^2|\le(|a|+|b|)|a-b|$, it holds that
\begin{align*}
|(\G_0+h)^2(\omega,t)-(\G_0+h)^2(\omega^\prime,t^\prime)|
&\le
2(L+\|h\|_\infty)\Big(
w_0\!\left(d_{\Omega\times\mathbb R_{\ge0}}((\omega,t),(\omega^\prime,t^\prime))\right)
\nonumber
\\
&\quad +
L_h\,d_{\Omega\times\mathbb R_{\ge0}}((\omega,t),(\omega^\prime,t^\prime))
\Big)
\end{align*}
for every $(\omega,t), (\omega',t')\in\Omega\times[0,M].$
Therefore, defining
$$
B:=(L+\|h\|_\infty)^2,
\qquad
w(u):=2(L+\|h\|_\infty)\big(w_0(u)+L_hu\big),
$$
then
$$
\mathsf{P}\big((\G_0+h)^2\in\mathcal G_{B,w}\big)\ge1-\eta.
$$
Therefore, for every $\varepsilon>0$,
\begin{align*}
\mathsf{P}&\Bigg(
\left|
\int (\G_0+h)^2\,\rd F_{n,\omega}^{Q_n}(t)\,\rd P_n(\omega)
-
\int (\G_0+h)^2\,\rd F_\omega^{\mu_0}(t)\,\rd\mu_0(\omega)
\right|
>\varepsilon
\Bigg)
\le
\mathsf{P}\big((\G_0+h)^2\notin\mathcal G_{B,w}\big)
\\
&+
\mathsf{P}\Bigg(
\sup_{g\in\mathcal G_{B,w}}
\left|
\int g(\omega,t)\,\rd F_{n,\omega}^{Q_n}(t)\,\rd P_n(\omega)
-
\int g(\omega,t)\,\rd F_\omega^{\mu_0}(t)\,\rd\mu_0(\omega)
\right|
>\varepsilon
\Bigg).
\end{align*}
The first term is bounded by $\eta$, and the second converges to $0$ by Lemma~\ref{ap:lem_DCT_P_hat_local}. Since $\eta>0$ is arbitrary, the quantity in \eqref{eq:local_cvm_aux_ineq_2} is $o_P(1)$. This concludes the proof.
\end{proof}



%----------------------------%
\subsubsection{Composite null hypothesis}\label{ap:sec_local_alternatives_composite}
%----------------------------%

This section contains the proof of Proposition \ref{prop:local_conv_roots_Qn}, which establishes the local convergence of $\hat{\btheta}_n$ and $\tilde{\btheta}_n$ to $\btheta_0$. Remark \ref{rem:local_conv_roots_Qn} shows a useful fact that will be applied in the proof.
%---------------------------%
\begin{remark}\label{rem:local_conv_roots_Qn}
%---------------------------%
By \ref{cp3}, as $\btheta\to\btheta_0$, one has
\begin{equation}\label{eq:U1_taylor_mu}
\mu_{\btheta_0}\psi_{\btheta} \;=\; \bV_{\btheta_0}(\btheta-\btheta_0)+o(\|\btheta-\btheta_0\|).
\end{equation}
Let $c\defin \frac{1}{2}\|\bV_{\btheta_0}^{-1}\|_2^{-1}$, with $\bV_{\btheta_0}$ as defined in \ref{cp3} and $\|\cdot\|_2$ the spectral norm. By \ref{cp3}, $\bV_{\btheta_0}$ is nonsingular, hence $c>0$. It holds that $\|\bV_{\btheta_0}\bv\|\ge 2c\|\bv\|$ for all $\bv\in\R^p$. 

In addition, by definition of $o(\|\btheta-\btheta_0\|)$, taking $U$ sufficiently small, it holds that
\begin{equation}\label{eq:U2_remainder_bound}
\big\|\mu_{\btheta_0}\psi_{\btheta} - \bV_{\btheta_0}(\btheta-\btheta_0)\big\|\le c\,\|\btheta-\btheta_0\|.
\end{equation}
\end{remark}
\begin{proof}[Proof of Proposition~\ref{prop:local_conv_roots_Qn}]
We divide the proof into its two parts, \ref{prop:local_conv_roots_Qn_i} and \ref{prop:local_conv_roots_Qn_ii}.

\emph{Proof of \ref{prop:local_conv_roots_Qn_i}.} By the reverse triangle inequality, one has
\begin{equation}\label{eq:mu_separation}
\|\mu_{\btheta_0}\psi_{\btheta}\|
\ge
\|\bV_{\btheta_0}(\btheta-\btheta_0)\|-\big\|\mu_{\btheta_0}\psi_{\btheta} - \bV_{\btheta_0}(\btheta-\btheta_0)\big\|
\ge
2c\|\btheta-\btheta_0\|-c\|\btheta-\btheta_0\|
=
c\|\btheta-\btheta_0\|
\end{equation}
in a neighborhood of $\btheta_0$.

By \eqref{eq:local_cp1_prop_local_conv}, for every $\btheta\in U$,
\begin{align*}
\|G\psi_{\btheta}\|
\le
\|G\psi_{\btheta_0}\| + \|G(\psi_{\btheta}-\psi_{\btheta_0})\| 
\le
\|G\psi_{\btheta_0}\| + G\|\psi_{\btheta}-\psi_{\btheta_0}\| 
\le
\|G\psi_{\btheta_0}\| + GL\|\btheta-\btheta_0\|.
\end{align*}
Since $GL<\infty$ by Cauchy--Schwarz and $U$ is bounded, it follows that
\begin{equation}\label{eq:MG_bound_local_conv}
M_G \defin \sup_{\btheta\in U}\|G\psi_{\btheta}\| < \infty.
\end{equation}

Since $Q_n\psi_{\tilde\btheta_n}=0$,
by definition of $Q_n$ one has
$$
\left(1-\frac{1}{\sqrt{n}}\right)\,\|\mu_{\btheta_0}\psi_{\tilde\btheta_n}\|
=
\frac{\|G\psi_{\tilde\btheta_n}\|}{\sqrt{n}}
\le
\frac{M_G}{\sqrt{n}}.
$$
Hence, for $n$ sufficiently large,
$$
\|\mu_{\btheta_0}\psi_{\tilde\btheta_n}\|\le \frac{M_G}{\sqrt{n}\left(1-\frac{1}{\sqrt{n}}\right)} \le \frac{2M_G}{\sqrt{n}},
$$
and combining with \eqref{eq:mu_separation} at $\btheta=\tilde\btheta_n$ gives $
\|\tilde\btheta_n-\btheta_0\| \to 0 $ as $n\to \infty$. 

\emph{Proof of \ref{prop:local_conv_roots_Qn_ii}.} For each $j=1,\ldots,p$, let $\cF_j\defin\{\psi_{\btheta,j} : \btheta \in U\}$. We show that $\cF_j$ satisfies the conditions of Corollary~\ref{cor:bernoulli}. Each class $\cF_j$ is $\mu_{\btheta_0}$- and $G$-Donsker by \eqref{eq:local_cp1_prop_local_conv} and \ref{cp2}, because $\mu_{\btheta_0}L^2<\infty$, $\mu_{\btheta_0}\|\psi_{\btheta_0}\|^2<\infty$, $GL^2<\infty$, and $G\|\psi_{\btheta_0}\|^2<\infty$ (see Example 19.7 in \cite{vandervaart1998}). Moreover, for every $\btheta\in U$,
\begin{align*}
\big|(G-\mu_{\btheta_0})\psi_{\btheta,j}\big|
\le
\big\|(G-\mu_{\btheta_0})\psi_{\btheta}\big\|
\le
\|G\psi_{\btheta}\| + \mu_{\btheta_0}\|\psi_{\btheta}-\psi_{\btheta_0}\|
\le
M_G + \mu_{\btheta_0}L\|\btheta-\btheta_0\|,
\end{align*}
so $\sup_{\btheta\in U}\big|(G-\mu_{\btheta_0})\psi_{\btheta,j}\big|<\infty$ because $U$ is bounded.



An application of Corollary~\ref{cor:bernoulli} to each $\cF_j$ and the continuous mapping theorem gives $\sup_{\btheta\in U}|(P_n-Q_n)\psi_{\btheta,j}|=O_{Q_n}(n^{-1/2})$. Since $p<\infty$,
\begin{equation}\label{eq:application_bernoulli}
\sup_{\btheta\in U}\|(P_n-Q_n)\psi_{\btheta}\|
\le
\sqrt{p}\,\max_{1\le j\le p}\sup_{\btheta\in U}|(P_n-Q_n)\psi_{\btheta,j}|
=O_{Q_n}(n^{-1/2}).
\end{equation}
Since $P_n\psi_{\hat\btheta}=0$, \eqref{eq:application_bernoulli} implies that, on the event $\{\hat\btheta\in U\}$,
\begin{equation}\label{eq:Qn_psi_hat_bound}
\|Q_n\psi_{\hat\btheta}\|
=
\|(P_n-Q_n)\psi_{\hat\btheta}\|
\le
\sup_{\btheta\in U}\|(P_n-Q_n)\psi_{\btheta}\|
=
O_{Q_n}(n^{-1/2}),
\end{equation}
and in particular $\|Q_n\psi_{\hat\btheta}\|=o_{Q_n}(1)$ on $\{\hat\btheta\in U\}$. 

By \eqref{eq:mu_separation},
\begin{equation}\label{eq:sep_ring_mu}
\inf_{\btheta \in U, \;\|\btheta-\btheta_0\|\ge\varepsilon}\|\mu_{\btheta_0}\psi_{\btheta}\|
\ge
c\varepsilon.
\end{equation}
Moreover,
$$
\|(Q_n-\mu_{\btheta_0})\psi_{\btheta}\|
=
\frac{1}{\sqrt{n}}\|(G-\mu_{\btheta_0})\psi_{\btheta}\|
\le
\frac{1}{\sqrt{n}}\big(\|G\psi_{\btheta}\|+\|\mu_{\btheta_0}\psi_{\btheta}\|\big)
\le
\frac{1}{\sqrt{n}}\Big(M_G+\sup_{\btheta\in U}\|\mu_{\btheta_0}\psi_{\btheta}\|\Big),
$$
so $\sup_{\btheta\in U}\|(Q_n-\mu_{\btheta_0})\psi_{\btheta}\|=O(n^{-1/2})=o(1)$. The fact that $\sup_{\btheta\in U}\|\mu_{\btheta_0}\psi_{\btheta}\|<\infty$ follows from \eqref{eq:local_cp1_prop_local_conv}.
Therefore, for $n$ sufficiently large,
\begin{equation}\label{eq:drift_small_ring}
\sup_{\btheta\in U}\|(Q_n-\mu_{\btheta_0})\psi_{\btheta}\| < \frac{c\varepsilon}{2}.
\end{equation}

Combining \eqref{eq:sep_ring_mu} and \eqref{eq:drift_small_ring}, for $n$ sufficiently large and all $\btheta\in U$ with $\|\btheta-\btheta_0\|\ge\varepsilon$,
\begin{equation*}
\|Q_n\psi_{\btheta}\|
\ge
\|\mu_{\btheta_0}\psi_{\btheta}\|-\|(Q_n-\mu_{\btheta_0})\psi_{\btheta}\|
>
c\varepsilon-\frac{c\varepsilon}{2}
=
\frac{c\varepsilon}{2}.
\end{equation*}
The preceding bound applied at $\btheta=\hat\btheta$ on the event $\{\hat\btheta\in U\}$ gives
$$
\mathsf P\big(\|\hat\btheta-\btheta_0\|\ge\varepsilon\big)
\le
\mathsf P(\hat\btheta\notin U)
+
\mathsf P\Big(\|Q_n\psi_{\hat\btheta}\|\ge \frac{c\varepsilon}{2},\ \hat\btheta\in U\Big)
\to0,
$$
because $\mathsf P(\hat\btheta\notin U)\to0$ and $\|Q_n\psi_{\hat\btheta}\|=O_{Q_n}(n^{-1/2})$ by \eqref{eq:Qn_psi_hat_bound} if $\hat\btheta\in U$. This implies $\mathsf P(\|\hat\btheta-\btheta_0\|\ge\varepsilon)\to0$ for every $\varepsilon>0$, i.e. $\hat\btheta \convp \btheta_0$.
\end{proof}

We introduce Lemmas \ref{lem:aux_P_n_Q_n_tilde} and \ref{lem:aux_P_n_Q_n_tilde_hat}, which are auxiliary results used in the proof of Theorem \ref{thm:local_alt_conv_process_correction}.
%----------------------%

%----------------------%
\begin{lemma}\label{lem:aux_P_n_Q_n_tilde_hat}
%----------------------%
Assume conditions \ref{cp1} and \ref{cp2}, $GL^2<\infty$, and $G\|\psi_{\btheta_0}\|^2<\infty$.
Let $\tilde\btheta_n\in\Theta$ be such that $Q_n\psi_{\tilde\btheta_n}=0$ and $\tilde\btheta_n\to\btheta_0$.  
Let also $\hat\btheta\in\Theta$ be such that $P_n\psi_{\hat\btheta}=o_P(n^{-1/2})$ and $\hat\btheta\convp\btheta_0$.
Under the local alternative defined in \eqref{eq:local_alt_Qn_def_composite}, it holds that 
$$
\left(P_n - Q_n\right)\left(\psi_{\hat{\btheta}} - \psi_{\tilde\btheta_n}\right) = o_P(n^{-1/2})
$$
\end{lemma}
\begin{proof}
Define, for a law $P$ and a measurable function $f$ taking values in $\R^p$, the pseudo-norms
$$
\rho_P(f) \defin \left(P\|f-Pf\|^2\right)^{\frac{1}{2}}, \quad    \|f\|_{P,2} \defin \left(P\|f\|^2\right)^{\frac{1}{2}}.
$$
Fix $\rho>0$ such that $B(\btheta_0,\rho)$ is contained in a neighborhood of $\btheta_0$ on which \ref{cp1} holds, and define
\[
A_n\defin \{\hat\btheta\in B(\btheta_0,\rho),\ \tilde\btheta_n\in B(\btheta_0,\rho)\}.
\]
Since $\hat\btheta\convp\btheta_0$ and $\tilde\btheta_n\to\btheta_0$, $\Prob{A_n}\to1$.
Thus, it holds that
\begin{align}\label{eq:relation_pseudo_norms}
   \rho_{\mu_{\btheta_0}}(f) \le \|f\|_{\mu_{\btheta_0},2} \le \Big(1-\frac{1}{\sqrt{n}}\Big)^{-1/2}\|f\|_{Q_n,2}.
\end{align}

On $A_n$, by \ref{cp1}, one has 
\begin{align*}
    \|\psi_{\hat{\btheta}}-\psi_{\tilde{\btheta}_n}\|_{Q_n,2} \le \|L\|_{Q_n,2}\|\hat{\btheta}-\tilde{\btheta}_n\| \convp 0,
\end{align*}
Since $\Probbig{A_n}\to1$, this bound holds with probability increasing to one as $n\to\infty$. Hence, for every $\delta>0$, one has
\begin{equation}\label{eq:bound_asympt_equicont}
\begin{aligned}
\|\G_{n,Q_n}^{Q_n}\big(\psi_{\hat{\btheta}}-\psi_{\tilde\btheta_n}\big)\|
\le &    
\sup_{\substack{\btheta_1,\btheta_2\in B(\btheta_0,\rho)\\ \|\psi_{\btheta_1}-\psi_{\btheta_2}\|_{Q_n,2}<\delta}}
\|\G_{n,Q_n}^{Q_n}\big(\psi_{\btheta_1}-\psi_{\btheta_2}\big)\| + o_P(1) \\
\le &
\sup_{\substack{\btheta_1,\btheta_2\in B(\btheta_0,\rho)\\ \|\psi_{\btheta_1}-\psi_{\btheta_2}\|_{Q_n,2}<\delta}}
\|\G_{n,\mu_{\btheta_0}}^{\mu_{\btheta_0}}\big(\psi_{\btheta_1}-\psi_{\btheta_2}\big)\|
\\
&+
2\sup_{\btheta\in B(\btheta_0,\rho)}
\left\|\G_{n,Q_n}^{Q_n}\psi_{\btheta} - \G_{n,\mu_{\btheta_0}}^{\mu_{\btheta_0}}\psi_{\btheta}\right\|
+o_P(1).
\end{aligned}
\end{equation}
%The o_P(1) appears because the probability that $\|\psi_{\hat{\btheta}}-\psi_{\tilde{\btheta}_n}\|_{Q_n,2}>\delta$ goes to zero for every $\delta>0$ by \eqref{eq:psi_hat_theta_tilde_theta_convp_0}.

To show $\G_{n,Q_n}^{Q_n}(\psi_{\hat{\btheta}}-\psi_{\tilde\btheta_n})\convp0$, we first show that
\begin{align}\label{eq:coupling_score_bound}
\sup_{\btheta\in B(\btheta_0,\rho)}\left\|\G_{n,Q_n}^{Q_n}(\psi_{\btheta}) - \G_{n,\mu_{\btheta_0}}^{\mu_{\btheta_0}}(\psi_{\btheta})\right\|
\le
\sqrt{p}\,\max_{1\le j\le p}\sup_{\btheta\in B(\btheta_0,\rho)}
\left|\G_{n,Q_n}^{Q_n}(\psi_{\btheta,j}) - \G_{n,\mu_{\btheta_0}}^{\mu_{\btheta_0}}(\psi_{\btheta,j})\right|
\convp 0.
\end{align} 
For each $j=1,\ldots,p$, let $\cF_{\rho,j}\defin\{\psi_{\btheta,j} : \btheta \in B(\btheta_0,\rho)\}$. The assumptions of Proposition~\ref{prop:bernoulli} hold for each $\cF_{\rho,j}$ because each $\cF_{\rho,j}$ is $\mu_{\btheta_0}$- and $G$-Donsker, and
\begin{align*}
\sup_{f\in\cF_{\rho,j}}|(G-\mu_{\btheta_0})f|
\le
\sup_{\btheta\in B(\btheta_0,\rho)}\|(G-\mu_{\btheta_0})\psi_{\btheta}\|
\le
2\Big(\|G\psi_{\btheta_0}\|+\rho\,(GL+\mu_{\btheta_0}L)\Big)<\infty.
\end{align*}
The fact that each $\cF_{\rho,j}$ is $\mu_{\btheta_0}$- and $G$-Donsker follows from \ref{cp1} and \ref{cp2}, $\mu_{\btheta_0}L^2<\infty$, $GL^2<\infty$, and $G\|\psi_{\btheta_0}\|^2<\infty$ \citep[see Example 19.7 in][]{vandervaart1998}.
By Proposition~\ref{prop:bernoulli} applied to the class $\bigcup_{j=1}^p \cF_{\rho,j}$, there exists a version of $\G_{n,\mu_{\btheta_0}}^{\mu_{\btheta_0}}$ such that the RHS of \eqref{eq:coupling_score_bound} converges to zero in probability. 

For the first term on the RHS of \eqref{eq:bound_asympt_equicont}, since each class $\cF_{\rho,j}$ is $\mu_{\btheta_0}$-Donsker, it is asymptotically equicontinuous under $\mu_{\btheta_0}$ \citep[see pp. 138--139 in][]{vandervaart1998}. Hence, for every $\varepsilon>0$ and every $j=1,\ldots,p$,
\begin{equation}\label{eq:asympt_equicont_mu_theta0_score}
\lim_{\delta \downarrow 0} \limsup_{n \to \infty} \mathsf{P}\Bigg(\sup_{\substack{\btheta_1,\btheta_2\in B(\btheta_0,\rho)\\ \rho_{\mu_{\btheta_0}}(\psi_{\btheta_1,j}-\psi_{\btheta_2,j})<\delta}}\left|\G_{n,\mu_{\btheta_0}}^{\mu_{\btheta_0}}\big(\psi_{\btheta_1,j}-\psi_{\btheta_2,j}\big)\right|>\varepsilon\Bigg) = 0.
\end{equation}
Moreover, applying \eqref{eq:relation_pseudo_norms}, one obtains
\begin{align*}
\sup_{\substack{\btheta_1,\btheta_2\in B(\btheta_0,\rho)\\ \|\psi_{\btheta_1}-\psi_{\btheta_2}\|_{Q_n,2}<\delta}}\left\|\G_{n,\mu_{\btheta_0}}^{\mu_{\btheta_0}}\big(\psi_{\btheta_1}-\psi_{\btheta_2}\big)\right\|
\le
\sqrt{p}\,\max_{1\le j\le p}\sup_{\substack{\btheta_1,\btheta_2\in B(\btheta_0,\rho)\\ \rho_{\mu_{\btheta_0}}(\psi_{\btheta_1,j}-\psi_{\btheta_2,j})<2\delta}}
\left|\G_{n,\mu_{\btheta_0}}^{\mu_{\btheta_0}}\big(\psi_{\btheta_1,j}-\psi_{\btheta_2,j}\big)\right|.
\end{align*}
Since \eqref{eq:bound_asympt_equicont} holds for every $\delta>0$, taking $\limsup_{n\to\infty}$ and then letting $\delta\downarrow0$ yields
$\G_{n,Q_n}^{Q_n}(\psi_{\hat{\btheta}}-\psi_{\tilde\btheta_n})\convp0$, and hence $\left(P_n - Q_n\right)\left(\psi_{\hat{\btheta}} - \psi_{\tilde\btheta_n}\right) = o_P(n^{-1/2})$.
\end{proof}
%----------------------%

%----------------------%
\begin{lemma}\label{lem:aux_P_n_Q_n_tilde}
%----------------------%
Assume conditions \ref{cp1} and \ref{cp2}, $GL^2<\infty$, and $G\|\psi_{\btheta_0}\|^2<\infty$.
Let $\tilde\btheta_n\in\Theta$ be such that $Q_n\psi_{\tilde\btheta_n}=0$ and $\tilde\btheta_n\to\btheta_0$.  
Under the local alternative defined in \eqref{eq:local_alt_Qn_def_composite}, it holds that 
\begin{align}\label{eq:p_n_q_n_O_n-12}
\left(P_n - Q_n\right)\psi_{\tilde{\btheta}_n} = O_P(n^{-1/2}).
\end{align}
\end{lemma}
\begin{proof}
Fix $\rho>0$ such that $B(\btheta_0,\rho)$ is contained in a neighborhood of $\btheta_0$ on which \ref{cp1} holds.
Since $\tilde{\btheta}_n\to\btheta_0$ by assumption, there exists $n(\rho)>0$ such that $\tilde{\btheta}_n\in B(\btheta_0,\rho)$ for all $n>n(\rho)$.
Observe that also $\|\sqrt{n}(P_n-Q_n)\psi_{\tilde{\btheta}_n}\| \le \sup_{\btheta\in B(\btheta_0,\rho)}\|\sqrt{n}(P_n-Q_n)\psi_{\btheta}\|$ for $n>n(\rho)$. 


For each $j=1,\ldots,p$, let $\cF_{\rho,j}\defin \{\psi_{\btheta,j} : \btheta \in B(\btheta_0,\rho)\}$.  The assumptions of Corollary~\ref{cor:bernoulli} hold by the same arguments as in Lemma \ref{lem:aux_P_n_Q_n_tilde_hat}. An application of Corollary~\ref{cor:bernoulli} to each $\cF_{\rho,j}$ gives $\sup_{\btheta\in B(\btheta_0,\rho)}|\sqrt{n}(P_n-Q_n)\psi_{\btheta,j}|=O_P(1)$. Hence, since $p<\infty$,
\begin{align*}
\sup_{\btheta\in B(\btheta_0,\rho)}\|\sqrt{n}(P_n-Q_n)\psi_{\btheta}\|
\le
\sqrt{p}\,\max_{1\le j\le p}\sup_{\btheta\in B(\btheta_0,\rho)}|\sqrt{n}(P_n-Q_n)\psi_{\btheta,j}|
=O_P(1),
\end{align*}
and hence $\left(P_n - Q_n\right)\psi_{\tilde{\btheta}_n} = O_P(n^{-1/2})$. This concludes the result.
\end{proof}
%----------------------%


Using Lemmas \ref{lem:aux_P_n_Q_n_tilde} and \ref{lem:aux_P_n_Q_n_tilde_hat}, we can now prove Theorem \ref{thm:local_alt_conv_process_correction}.



%-----------------------%
\begin{proof}[Proof of Theorem \ref{thm:local_alt_conv_process_correction}]
We divide the proof int its two parts, \ref{thm:local_alt_conv_process_correction_i} and \ref{thm:local_alt_conv_process_correction_ii}.

    \emph{Proof of \ref{thm:local_alt_conv_process_correction_i}.} Since $\btheta\mapsto \mu_{\btheta_0}\psi_{\btheta}$ is differentiable at $\btheta_0$ by \ref{cp3}, it holds that
\begin{align*}
    \mu_{\btheta_0}\psi_{\btheta} = \mu_{\btheta_0}\psi_{\btheta_0} + \bV_{\btheta_0}(\btheta-\btheta_0) + o(\|\btheta-\btheta_0\|).
\end{align*}
Since $\tilde\btheta_n\to\btheta_0$, evaluating this expansion at $\btheta=\tilde\btheta_n$ gives
$$
\mu_{\btheta_0}\psi_{\tilde{\btheta}_n} = \mu_{\btheta_0}\psi_{\btheta_0} + \bV_{\btheta_0}(\tilde{\btheta}_n-\btheta_0) + o(\|\tilde{\btheta}_n-\btheta_0\|).
$$
Together with $Q_n\psi_{\tilde{\btheta}_n}=0$, this gives
\begin{align*}
    0= \left(1-\frac{1}{\sqrt{n}}\right)\left(\bV_{\btheta_0}(\tilde{\btheta}_n-\btheta_0) + o(\|\tilde{\btheta}_n-\btheta_0\|)\right) + \frac{1}{\sqrt{n}} G \psi_{\tilde{\btheta}_n},
\end{align*}
and thus
\begin{align}\label{eq:tilde_theta_n_theta_0}
    \bV_{\btheta_0}(\tilde{\btheta}_n - \btheta_0) = - \frac{1}{\sqrt{n}\left(1-\frac{1}{\sqrt{n}}\right)} G \psi_{\tilde{\btheta}_n} + o(\|\tilde{\btheta}_n-\btheta_0\|).
\end{align}

Now, by \ref{cp1} and the assumption that $\tilde{\btheta}_n\to\btheta_0$, it holds that 
$$
\|G\psi_{\tilde{\btheta}_n}\|\le \|G\psi_{\btheta_0}\| + GL \|\tilde{\btheta}_n-\btheta_0\|
$$ 
for $n$ sufficiently large, where $GL<\infty$ because $GL^2<\infty$.
%the assumption of convergence of \tilde{\btheta}_n to \btheta_0 is used here to ensure that \tilde{\btheta}_n is in a neighborhood of \btheta_0 where \ref{cp1} holds.
Plugging into \eqref{eq:tilde_theta_n_theta_0} gives
\begin{align*}
    \left(1-\frac{1}{\sqrt{n}}\right)\|\bV_{\btheta_0}(\tilde{\btheta}_n-\btheta_0)\| \le \frac{1}{\sqrt{n}}\|G\psi_{\btheta_0}\| + \frac{GL}{\sqrt{n}}\|\tilde{\btheta}_n-\btheta_0\| + o(\|\tilde{\btheta}_n-\btheta_0\|).
\end{align*}
By Remark~\ref{rem:local_conv_roots_Qn}, $\|\bV_{\btheta_0}(\tilde{\btheta}_n-\btheta_0)\|\ge 2c\|\tilde{\btheta}_n-\btheta_0\|$. Therefore, for every $\varepsilon>0$, there exists $n(\varepsilon)>0$ such that, for every $n>n(\varepsilon)$, one has
\begin{align*}
    2c\left(1-\frac{1}{\sqrt{n}}\right)\|\tilde{\btheta}_n-\btheta_0\|  \le \frac{1}{\sqrt{n}}\|G\psi_{\btheta_0}\| + \frac{GL}{\sqrt{n}}\|\tilde{\btheta}_n-\btheta_0\| + \varepsilon\|\tilde{\btheta}_n-\btheta_0\|.
\end{align*}
Thus,
\begin{align*}
    \left(2c\Big(1-\frac{1}{\sqrt{n}}\Big) - \frac{GL}{\sqrt{n}} - \varepsilon \right)\|\tilde{\btheta}_n-\btheta_0\| \le \frac{1}{\sqrt{n}}\|G\psi_{\btheta_0}\|.
\end{align*} 
Take $\varepsilon=\frac{c}{2}$ to obtain
\begin{align}\label{eq:tilde_theta_theta_0_On-12}
    \|\tilde{\btheta}_n-\btheta_0\| \le \frac{2}{c\sqrt{n}}\|G\psi_{\btheta_0}\| = O\big(n^{-\frac{1}{2}}\big),
\end{align}
for $n$ sufficiently large. The bound $\|G\psi_{\btheta_0}\|<\infty$ follows from $G\|\psi_{\btheta_0}\|^2<\infty$.

Now, rewrite \eqref{eq:tilde_theta_n_theta_0} as
\begin{align}\label{eq:tilde_theta_n_theta_0_rewrite}
   \bV_{\btheta_0} (\tilde{\btheta}_n - \btheta_0) = - \frac{1}{\sqrt{n}\Big(1-\frac{1}{\sqrt{n}}\Big)} G \psi_{\btheta_0} - \frac{1}{\sqrt{n}\Big(1-\frac{1}{\sqrt{n}}\Big)} G \bigl(\psi_{\tilde{\btheta}_n} - \psi_{\btheta_0}\bigr) + o(\|\tilde{\btheta}_n-\btheta_0\|).
\end{align}
By condition \ref{cp1}, one has
\begin{align*}
   \frac{1}{\sqrt{n}\Big(1-\frac{1}{\sqrt{n}}\Big)} \|G(\psi_{\tilde{\btheta}_n} - \psi_{\btheta_0})\| \le \frac{1}{\sqrt{n}\Big(1-\frac{1}{\sqrt{n}}\Big)} GL \|\tilde{\btheta}_n-\btheta_0\| = O(n^{-1})
\end{align*}
by \eqref{eq:tilde_theta_theta_0_On-12}. Thus, \eqref{eq:tilde_theta_n_theta_0_rewrite} yields
\begin{align*}
    \tilde{\btheta}_n - \btheta_0 = - \frac{1}{\sqrt{n}\Big(1-\frac{1}{\sqrt{n}}\Big)} \bV_{\btheta_0}^{-1} G \psi_{\btheta_0} + o\big(n^{-\frac{1}{2}}\big),
\end{align*}
which gives \eqref{eq:thm_local_alt_conv_process_correction_i}, applying the assumption $\|G \psi_{\btheta_0}\|<\infty$.

\emph{Proof of \ref{thm:local_alt_conv_process_correction_ii}.} Now, we turn to \eqref{eq:thm_local_alt_conv_process_correction_ii}. Since $P_n\psi_{\hat\btheta}=o_P(n^{-1/2})$ and $Q_n\psi_{\tilde\btheta_n}=0$, one has
\begin{align}\label{eq:hat_theta_split}
    0 = \left(P_n - Q_n\right)\psi_{\hat{\btheta}} + Q_n\left(\psi_{\hat{\btheta}} - \psi_{\tilde\btheta_n}\right) + o_P(n^{-1/2}).
\end{align}
To control the first summand in \eqref{eq:hat_theta_split}, write 
\begin{align*}
    \left(P_n - Q_n\right)\psi_{\hat{\btheta}} = \left(P_n - Q_n\right)\psi_{\tilde{\btheta}_n} + \left(P_n - Q_n\right)\left(\psi_{\hat{\btheta}} - \psi_{\tilde{\btheta}_n}\right).
\end{align*}
By Lemma \ref{lem:aux_P_n_Q_n_tilde_hat}, it holds that $\left(P_n - Q_n\right)\left(\psi_{\hat{\btheta}} - \psi_{\tilde{\btheta}_n}\right) = o_P(n^{-1/2})$. 
Therefore, 
\begin{align}\label{eq:first_RHS_hat_theta_split}
    \left(P_n - Q_n\right)\psi_{\hat{\btheta}} = \left(P_n - Q_n\right)\psi_{\tilde{\btheta}_n} + o_P(n^{-1/2}).
\end{align}

Focus now on the second summand of \eqref{eq:hat_theta_split}, i.e, $Q_n\left(\psi_{\hat{\btheta}} - \psi_{\tilde\btheta_n}\right)$. By definition of $Q_n$, one has
$$
Q_n\left(\psi_{\hat{\btheta}} - \psi_{\tilde\btheta_n}\right) = \left(1-\frac{1}{\sqrt{n}}\right)\mu_{\btheta_0} \left(\psi_{\hat{\btheta}} - \psi_{\tilde\btheta_n}\right) + \frac{1}{\sqrt{n}} G \left(\psi_{\hat{\btheta}} - \psi_{\tilde\btheta_n}\right).
$$
By \eqref{eq:U1_taylor_mu} applied at $\btheta=\hat{\btheta}$ and $\btheta=\tilde{\btheta}_n$, \ref{cp1} and $GL<\infty$, one has
\begin{equation}\label{eq:second_RHS_hat_theta_split}
\begin{aligned}
   Q_n\left(\psi_{\hat{\btheta}} - \psi_{\tilde\btheta_n}\right) & = \left(1-\frac{1}{\sqrt{n}}\right)\left(\bV_{\btheta_0}(\hat{\btheta}-\tilde\btheta_n) + o_P(\|\hat{\btheta}-\btheta_0\|) + o(\|\tilde{\btheta}_n-\btheta_0\|)\right) + O_P\left(\frac{1}{\sqrt{n}}\|\hat{\btheta}-\tilde\btheta_n\|\right) \\
    &= \bV_{\btheta_0}(\hat{\btheta}-\tilde\btheta_n) +  O_P\left(\frac{1}{\sqrt{n}} \|\hat{\btheta}-\tilde\btheta_n\|\right) + o_P(\|\hat{\btheta}-\btheta_0\|) + o(\|\tilde{\btheta}_n-\btheta_0\|) \\
    &= \bV_{\btheta_0}(\hat{\btheta}-\tilde\btheta_n) + O_P\left(\frac{1}{\sqrt{n}} \|\hat{\btheta}-\tilde\btheta_n\|\right) + o_P(\|\hat{\btheta}-\btheta_0\|) + o_P(n^{-\frac{1}{2}}),
\end{aligned}
\end{equation}
where the last equality follows from \eqref{eq:tilde_theta_theta_0_On-12}.

Combining \eqref{eq:first_RHS_hat_theta_split} and \eqref{eq:second_RHS_hat_theta_split} into \eqref{eq:hat_theta_split} gives
\begin{align}\label{eq:pre_hat_tilde}
    0 = (P_n-Q_n) \psi_{\tilde{\btheta}_n} + \bV_{\btheta_0}(\hat{\btheta}-\tilde\btheta_n) +  O_P\left(\frac{1}{\sqrt{n}} \|\hat{\btheta}-\tilde\btheta_n\|\right) + o_P(\|\hat{\btheta}-\btheta_0\|) + o_P(n^{-\frac{1}{2}})
\end{align}
We show that $\|\hat{\btheta}-\tilde\btheta_n\| = O_P(n^{-1/2})$. Using Lemma \ref{lem:aux_P_n_Q_n_tilde} in \eqref{eq:pre_hat_tilde}, we obtain
\begin{align*} 
    \left\|\bV_{\btheta_0}(\hat{\btheta}-\tilde\btheta_n)\right\| &\le O_P(n^{-\frac{1}{2}}) + O_P\left(\frac{1}{\sqrt{n}} \|\hat{\btheta}-\tilde\btheta_n\|\right) + o_P(\|\hat{\btheta}-\btheta_0\|) + o_P(n^{-\frac{1}{2}}) \\
    &= O_P(n^{-1/2}) + o_P(\|\hat{\btheta}-\btheta_0\|). 
\end{align*} 
The last equality uses $\|\hat{\btheta}-\tilde\btheta_n\|=o_P(1)$. By Remark~\ref{rem:local_conv_roots_Qn}, $\|\bV_{\btheta_0}(\hat{\btheta}-\tilde{\btheta}_n)\|\ge 2c\|\hat{\btheta}-\tilde{\btheta}_n\|$. Then, applying the triangle inequality to $\|\hat{\btheta}-\btheta_0\|$ and using $\|\tilde{\btheta}_n-\btheta_0\|=O(n^{-1/2})$ from \eqref{eq:tilde_theta_theta_0_On-12} gives 
$$
2c\|\hat{\btheta}-\tilde{\btheta}_n\| \le O_P(n^{-\frac{1}{2}}) + o_P(\|\hat{\btheta}-\tilde{\btheta}_n\|).
$$
By Lemma \ref{lem:absorption_OP}, this implies that $\|\hat{\btheta}-\tilde\btheta_n\| = O_P(n^{-1/2})$. Observe also that since $\|\hat{\btheta}-\tilde\btheta_n\|=O_P(n^{-1/2})$ and $\|\tilde{\btheta}_n-\btheta_0\|=O(n^{-1/2})$, then the triangle inequality gives $\|\hat{\btheta}-\btheta_0\|=O_P(n^{-1/2})$. 
Plugging back into \eqref{eq:pre_hat_tilde} and using $Q_n \psi_{\tilde{\btheta}_n}=0$ gives 
\begin{align*} 
    \hat{\btheta} - \tilde\btheta_n = -\bV_{\btheta_0}^{-1} P_n \psi_{\tilde{\btheta}_n} + o_P(n^{-\frac{1}{2}}),
\end{align*} 
which combined with \eqref{eq:thm_local_alt_conv_process_correction_i} gives \eqref{eq:thm_local_alt_conv_process_correction_ii}.  
\end{proof}
%----------------------%


%----------------------%
\begin{proof}[Proof of Theorem \ref{thm:local_alt_Gn_composite}]
First, observe that under the local alternative \eqref{eq:local_alt_Qn_def_composite} we can write, for every
$(\omega,t)\in\Omega\times\R_{\ge0}$,
\begin{align}
F_{n,\omega}^{Q_n}(t)-F^{\mu_{\hat\btheta}}_\omega(t)
&=
\Big(F_{n,\omega}^{Q_n}(t)-F^{\mu_{\btheta_0}}_\omega(t)\Big)
-
\Big(F^{\mu_{\hat\btheta}}_\omega(t)-F^{\mu_{\btheta_0}}_\omega(t)\Big).
\label{eq:proof_E5_start}
\end{align}
The asymptotic behavior of the first term in the RHS of \eqref{eq:proof_E5_start} was given in \eqref{eq:local_alt_weak_limit_process} with $\mu_0$ replaced by $\mu_{\btheta_0}$. 

It remains to analyze $F^{\mu_{\hat\btheta}}_\omega(t)-F^{\mu_{\btheta_0}}_\omega(t)$. By Assumption \ref{d}, it holds that
\begin{equation}
\sqrt{n}\Big(F^{\mu_{\hat\btheta}}_\omega(t)-F^{\mu_{\btheta_0}}_\omega(t)\Big)
=
\sqrt{n} \dot F^{\btheta_0}_\omega(t)^\top(\hat\btheta-\btheta_0)
+
o_{P}\big(\sqrt{n}\|\hat\btheta-\btheta_0\|\big).
\label{eq:proof_E5_taylor_scaled}
\end{equation}
where the remainder is uniform in $(\omega,t)\in\Omega\times\R_{\ge0}$.

Recall that by Theorem \ref{thm:local_alt_conv_process_correction}, it holds that
\begin{equation}
\sqrt{n}(\hat\btheta-\btheta_0)
=
-\bV_{\btheta_0}^{-1}\sqrt{n}\,P_n\psi_{\tilde\btheta_n}
-\bV_{\btheta_0}^{-1}G\psi_{\btheta_0}
+o_{P}(1),
\label{eq:proof_E5_use_95}
\end{equation}
Since $Q_n\psi_{\tilde\btheta_n}=0$, we have $\sqrt{n}\,P_n\psi_{\tilde\btheta_n}=\sqrt{n}(P_n-Q_n)\psi_{\tilde\btheta_n}=\G_{n,Q_n}^{Q_n}(\psi_{\tilde\btheta_n}).$
Moreover,
\begin{align*}
\mathsf E\left\|\G_{n,Q_n}^{Q_n}(\psi_{\tilde\btheta_n}-\psi_{\btheta_0})\right\|^2
\le
Q_n\left\|\psi_{\tilde\btheta_n}-\psi_{\btheta_0}\right\|^2 
\le
Q_n L^2\,\|\tilde\btheta_n-\btheta_0\|^2
\to 0,
\end{align*}
because $Q_nL^2$ is bounded and $\tilde\btheta_n\to\btheta_0$. Hence, by Markov's inequality, $\G_{n,Q_n}^{Q_n}(\psi_{\tilde\btheta_n}-\psi_{\btheta_0})=o_P(1)$. Therefore,
\begin{equation}
\sqrt{n}P_n\psi_{\tilde\btheta_n}
=
\G_{n,Q_n}^{Q_n}(\psi_{\btheta_0})+o_{P}(1).
\label{eq:proof_E5_sum_to_theta0}
\end{equation}
Plugging \eqref{eq:proof_E5_sum_to_theta0} into \eqref{eq:proof_E5_use_95} gives
\begin{equation}
\sqrt{n}(\hat\btheta-\btheta_0)
=
-\bV_{\btheta_0}^{-1}\,\G_{n,Q_n}^{Q_n}(\psi_{\btheta_0})
-\bV_{\btheta_0}^{-1}G\psi_{\btheta_0}
+o_{P}(1).
\label{eq:proof_E5_theta_expansion_final}
\end{equation}
Since $Q_n\|\psi_{\btheta_0}\|^2$ is bounded, $\G_{n,Q_n}^{Q_n}(\psi_{\btheta_0})=O_P(1)$; hence \eqref{eq:proof_E5_theta_expansion_final} gives $\sqrt{n}\|\hat\btheta-\btheta_0\|=O_P(1)$, and the remainder in \eqref{eq:proof_E5_taylor_scaled} is $o_P(1)$ uniformly in $(\omega,t)\in\Omega\times\R_{\ge0}$.
Substituting \eqref{eq:proof_E5_theta_expansion_final} into \eqref{eq:proof_E5_taylor_scaled}, we obtain
\begin{align}
\sqrt{n}\Big(F^{\mu_{\hat\btheta}}_\omega(t)-F^{\mu_{\btheta_0}}_\omega(t)\Big)
&=
-\dot F^{\btheta_0}_\omega(t)^\top \bV_{\btheta_0}^{-1}\,\G_{n,Q_n}^{Q_n}(\psi_{\btheta_0})
-\dot F^{\btheta_0}_\omega(t)^\top \bV_{\btheta_0}^{-1}\,G\psi_{\btheta_0}
+o_{P}(1).
\label{eq:proof_E5_Fhat_expansion_like_3p2}
\end{align}

Combining \eqref{eq:proof_E5_start} with \eqref{eq:proof_E5_Fhat_expansion_like_3p2}, the rest of the proof follows analogously to that of Theorem~\ref{thm:convergence_G_n_composite}. Define
\begin{equation*}
g_{\omega,t}(x):=-\dot F^{\btheta_0}_\omega(t)^\top \bV_{\btheta_0}^{-1}\psi_{\btheta_0}(x),
\qquad
f_{\omega,t}(x):=y_{\omega,t}(x)-F^{\mu_{\btheta_0}}_\omega(t)-g_{\omega,t}(x),
\end{equation*}
so that $\mu_{\btheta_0}f_{\omega,t}=0$.

Combining \eqref{eq:proof_E5_start}, \eqref{eq:local_alt_Qn_def_composite}, and \eqref{eq:proof_E5_Fhat_expansion_like_3p2},
and using that $\G_{n,Q_n}^{Q_n}(g_{\omega,t})=
-\dot F^{\btheta_0}_\omega(t)^\top \bV_{\btheta_0}^{-1}\G_{n,Q_n}^{Q_n}(\psi_{\btheta_0})$, one obtains
\begin{align}
\sqrt{n}\Big(F_{n,\omega}^{Q_n}(t)-F^{\mu_{\hat\btheta}}_\omega(t)\Big)
&=
\G_{n,Q_n}^{Q_n}(y_{\omega,t})
-\G_{n,Q_n}^{Q_n}(g_{\omega,t})
+h(\omega,t)
+\dot F^{\btheta_0}_\omega(t)^\top \bV_{\btheta_0}^{-1}G\psi_{\btheta_0}
+o_{P}(1)
\nonumber\\
&=
\G_{n,Q_n}^{Q_n}(f_{\omega,t})
+h(\omega,t)
+\dot F^{\btheta_0}_\omega(t)^\top \bV_{\btheta_0}^{-1}G\psi_{\btheta_0}
+o_{P}(1),
\label{eq:proof_E5_representation}
\end{align}
where we used that $\G_{n,Q_n}^{Q_n}(F^{\mu_{\btheta_0}}_\omega(t))=0$.

It remains to prove the weak convergence of $\G_{n,Q_n}^{Q_n}$ to $\tilde{\G}_{0}$ in $\ell^\infty(\cF)$, 
where $\cF:=\{f_{\omega,t}:(\omega,t)\in\Omega\times\R_{\ge0}\}$.
As in the proof of Theorem~\ref{thm:convergence_G_n_composite}, write $\cF\subset \cY-\mathcal{G}$, where
$$
\cY:=\{y_{\omega,t}-F^{\mu_{\btheta_0}}_\omega(t):(\omega,t)\in\Omega\times\R_{\ge0}\},
\qquad
\mathcal{G}:=\{g_{\omega,t}:(\omega,t)\in\Omega\times\R_{\ge0}\}.
$$
The class $\cY$ is $\mu_{\btheta_0}$- and $G$-Donsker by Theorem~5.1 of \cite{DubeyChenMuller2024}. For the class $\mathcal{G}$, note that
$$
g_{\omega,t}(x)=\sum_{j=1}^p \alpha_j(\omega,t)\,\psi_{\btheta_0,j}(x),
\qquad
\boldsymbol{\alpha}(\omega,t)^\top:=-\dot F^{\btheta_0}_\omega(t)^\top \bV_{\btheta_0}^{-1},
$$
so $\mathcal{G}\subset\mathrm{span}\{\psi_{\btheta_0,1},\dots,\psi_{\btheta_0,p}\}$ is finite-dimensional. Moreover, by assumptions \ref{cp2}, \ref{cp3} and \ref{d}, and by $G\|\psi_{\btheta_0}\|^2<\infty$, the envelope of $\mathcal{G}$ is in $L^2(\mu_{\btheta_0})$ and $L^2(G)$.
Thus, $\mathcal{G}$ is $\mu_{\btheta_0}$-Donsker and $G$-Donsker. Moreover, $\sup_{h\in \cY\cup \mathcal{G}}\mu_{\btheta_0}h^2<\infty$ and $\sup_{h\in \cY\cup \mathcal{G}}Gh^2<\infty$, hence $\cY-\mathcal{G}$ is $\mu_{\btheta_0}$-Donsker and $G$-Donsker, and therefore so is $\cF$.
If $\sup_{f\in\cF}|(G-\mu_{\btheta_0})f|<\infty$ holds, then Corollary~\ref{cor:bernoulli} with $\mu_0$ replaced by $\mu_{\btheta_0}$ applies and yields
\begin{equation}
\G_{n,Q_n}^{Q_n} \convl \tilde{\G}_{0}
\qquad\text{in }\ell^\infty(\cF).
\label{eq:proof_E5_donsker_apply}
\end{equation}
Combining \eqref{eq:proof_E5_representation} and \eqref{eq:proof_E5_donsker_apply} and using Slutsky's theorem gives \eqref{eq:local_alt_composite_consistency}.


It is only left to verify the condition $\sup_{f\in\cF}|(G-\mu_{\btheta_0})f|<\infty$, in order to apply Corollary~\ref{cor:bernoulli}.
Since $0\le y_{\omega,t}\le 1$, one has $\sup_{\omega,t}|(G-\mu_{\btheta_0})y_{\omega,t}|\le 2$.
Moreover, $(G-\mu_{\btheta_0})F^{\mu_{\btheta_0}}_\omega(t)=0$ and $\mu_{\btheta_0}\psi_{\btheta_0}=0$, hence
$$
|(G-\mu_{\btheta_0})g_{\omega,t}|
=
\big|\dot F^{\btheta_0}_\omega(t)^\top \bV_{\btheta_0}^{-1}G\psi_{\btheta_0}\big|
\le
\sup_{\omega,t}\|\dot F^{\btheta_0}_\omega(t)\|\,\|\bV_{\btheta_0}^{-1}\|\,\|G\psi_{\btheta_0}\|.
$$
By \ref{d}, \ref{cp3}, and $G\|\psi_{\btheta_0}\|^2<\infty$, it holds that $\sup_{f\in\cF}|(G-\mu_{\btheta_0})f|<\infty$.
\end{proof}
%----------------------%

\begin{proof}[Proof of Theorem \ref{thm:consistency_local_composite}]
The convergence of the Kolmogorov--Smirnov statistic follows from Theorem \ref{thm:local_alt_Gn_composite} and the continuous mapping theorem.

For the Cramér--von Mises statistic, the result follows by the same argument as in Theorem~\ref{thm:local_alt_CvM_simple}, replacing $\mu_0$ by $\mu_{\btheta_0}$ and $\G_0+h$ by $\tilde{\G}_0(\omega,t)+h(\omega,t)+\dot F^{\btheta_0}_\omega(t)^\top \bV_{\btheta_0}^{-1}G\psi_{\btheta_0}$.
As in Theorem \ref{thm:local_alt_Gn_composite}, the only caveat is to verify boundedness and uniform continuity of the deterministic drift, 
which in this case is given by $h(\omega,t)+\dot F^{\btheta_0}_\omega(t)^\top \bV_{\btheta_0}^{-1}G\psi_{\btheta_0}.$

The uniform continuity and boundedness of $h$ was established in the proof of Theorem~\ref{thm:local_alt_CvM_simple}. Uniform continuity of $(\omega,t)\mapsto\dot F^{\btheta_0}_\omega(t)^\top \bV_{\btheta_0}^{-1}G\psi_{\btheta_0}$ with respect to $d_{\Omega\times\mathbb R_{\ge0}}$
follows from the assumed uniform continuity of $(\omega,t)\mapsto \dot F^{\btheta_0}_\omega(t)$. Boundedness follows from assumptions \ref{d}, the invertibility of $\bV_{\btheta_0}$ by \ref{cp3}, and $G\|\psi_{\btheta_0}\|^2<\infty$.
\end{proof}





%---------------------------%
\begin{lemma}\label{lem:absorption_OP}
%---------------------------%
For $n\ge1$, let $A_n$, $B_n$, and $C_n$ be nonnegative random variables such that $A_n \le B_n + C_n$ for all $n$. Assume that $C_n = o_P(A_n)$, i.e., for every $\epsilon>0$, $\mathsf{P}(C_n>\epsilon A_n)\to0$. Then, $A_n = O_P(B_n)$.
\end{lemma}

\begin{proof}
Since $C_n=o_P(A_n)$, one has $\mathsf{P}(C_n>\tfrac12 A_n)\to0$, that is, $\mathsf{P}\left(C_n \le \tfrac12 A_n\right)\to1$.
If $C_n \le \tfrac12 A_n$, then
$A_n \le B_n + C_n \le B_n + \tfrac12 A_n$, so
$A_n \le 2 B_n.$
Therefore, for every $n$,
$$
\mathsf{P}(A_n>2B_n)
\le
\mathsf{P}\left(A_n>2B_n,\ C_n \le \tfrac12 A_n\right)
+
\mathsf{P}\left(C_n > \tfrac12 A_n\right)
=
\mathsf{P}\left(C_n > \tfrac12 A_n\right),
$$
where the equality follows from $A_n>2B_n$. As $\mathsf{P}(C_n > \tfrac12 A_n)\to0$, we conclude that $\mathsf{P}(A_n>2B_n)\to0$, hence $A_n = O_P(B_n)$.
\end{proof}



%---------------------------%
\section{Proofs of Section~\ref{sec:bootstrap}}\label{ap:sec_proofs_multiplier_bootstrap}
%---------------------------%

\begin{proof}[Proof of Theorem~\ref{thm:bootstrap_processes}]
We divide the proof into its two parts, \ref{thm:bootstrap_processes_i} and \ref{thm:bootstrap_processes_ii}.

\emph{Proof of \ref{thm:bootstrap_processes_i}.} It was shown in Section~\ref{sec:bootstrap} that
$$
\G_n^{*}(\omega,t)
=
\frac{m}{\tau}\sqrt n(P_n^{*}-P_n)y_{\omega,t}.
$$
Since $\cY$ is $\mu_0$-Donsker, Theorem~2.6 in \cite{Kosorok2008} gives
$$
\G_n^{*}\stackrel[*]{P}{\to}\G_0
\quad\text{in }\ell^\infty(\cY).
$$
This concludes the result.

\emph{Proof of \ref{thm:bootstrap_processes_ii}.} First observe that \ref{cp2} and Theorem~2.6 in \cite{Kosorok2008}, applied to the finite-dimensional class generated by the components of $\bpsi_{\btheta_0}$, give
$$
\frac{m}{\tau}\frac{1}{\sqrt n}\sum_{i=1}^n
\left(\frac{\xi_i}{\bar\xi_n}-1\right)\bpsi_{\btheta_0}(X_i)
=O_{P_{*}}(1).
$$
Hence, by the bootstrap Bahadur representation \eqref{eq:weighted_bahadur_bootstrap}, it holds that $\sqrt n(\hat{\btheta}^{*}-\hat{\btheta})=O_{P_{*}}(1)$. Moreover, \eqref{eq:bahadur} gives $\sqrt n(\hat{\btheta}-\btheta_0)=O_P(1)$, and hence $\sqrt n(\hat{\btheta}^{*}-\btheta_0)=O_{P_{*}}(1)$. Therefore, by \ref{d},
\begin{align*}
\sqrt n\left(F_\omega^{\mu_{\hat{\btheta}}}(t)-F_\omega^{\mu_{\btheta_0}}(t)\right)
&=
\dot F_\omega^{\btheta_0}(t)^\top\sqrt n(\hat{\btheta}-\btheta_0)
+o_P(1),
\\
\sqrt n\left(F_\omega^{\mu_{\hat{\btheta}^{*}}}(t)-F_\omega^{\mu_{\btheta_0}}(t)\right)
&=
\dot F_\omega^{\btheta_0}(t)^\top\sqrt n(\hat{\btheta}^{*}-\btheta_0)
+o_{P_{*}}(1).
\end{align*}
The remainders $o_P(1)$ and $o_{P_{*}}(1)$ are both uniform in $(\omega,t)\in\Omega\times\R_{\ge0}$.
Subtracting these two expansions and using \eqref{eq:weighted_bahadur_bootstrap} in the definition of $\tilde{\G}_n^{*}$ gives
\begin{equation}\label{eq:bootstrap_comp_process_remainder}
\sup_{\omega\in\Omega, t\ge0}
\left|
\tilde{\G}_n^{*}(\omega,t)
-
\frac{m}{\tau}\sqrt n(P_n^{*}-P_n)f_{\omega,t}
\right|
=o_{P_{*}}(1).
\end{equation}

Since $\cF$ is $\mu_{\btheta_0}$-Donsker, %by Theorem~\ref{thm:convergence_G_n_composite},  
Theorem~2.6 in \cite{Kosorok2008} implies that
$$
\frac{m}{\tau}\sqrt n(P_n^{*}-P_n)
\stackrel[*]{P}{\to}\tilde{\G}_0
\quad\text{in }\ell^\infty(\cF).
$$
Let $\varphi\in \mathrm{BL}_1(\ell^\infty(\cF))$. Define
$$
D_n^{*}
\defin
\sup_{\omega\in\Omega, t\ge0}
\left|
\tilde{\G}_n^{*}(\omega,t)
-
\frac{m}{\tau}\sqrt n(P_n^{*}-P_n)f_{\omega,t}
\right|.
$$
By \eqref{eq:bootstrap_comp_process_remainder}, $D_n^{*}=o_{P_{*}}(1)$. Moreover,
for every $\varepsilon>0$,
\begin{align}\label{eq:bootstrap_comp_BL_comparison}
\left|
\mathsf{E}_{*}\varphi(\tilde{\G}_n^{*})
-
\mathsf{E}_{*}\varphi\left(\frac{m}{\tau}\sqrt n(P_n^{*}-P_n)\right)
\right|
&\le
\mathsf{E}_{*}\left|
\varphi(\tilde{\G}_n^{*})
-
\varphi\left(\frac{m}{\tau}\sqrt n(P_n^{*}-P_n)\right)
\right|
\\
&\le
\varepsilon\,\mathsf{P}_{*}(D_n^{*}\le\varepsilon)
+
2\,\mathsf{P}_{*}(D_n^{*}>\varepsilon)
\\
&\le
\varepsilon+
2\,\mathsf{P}_{*}(D_n^{*}>\varepsilon).
\end{align}
Here, we have used that, if $D_n^{*}\le\varepsilon$, the Lipschitz property of $\varphi$ gives
$$
\left|
\varphi(\tilde{\G}_n^{*})
-
\varphi\left(\frac{m}{\tau}\sqrt n(P_n^{*}-P_n)\right)
\right|
\le D_n^{*}\le\varepsilon.
$$
Also, if $D_n^{*}>\varepsilon$, then using $\|\varphi\|_\infty\le1$, this absolute difference can be bounded by $2$. Since $\mathsf{P}_{*}(D_n^{*}>\varepsilon)\convp0$ by \eqref{eq:bootstrap_comp_process_remainder}, and $\varepsilon$ is arbitrary, the LHS of \eqref{eq:bootstrap_comp_BL_comparison} converges to zero in probability.

It remains to verify the asymptotic measurability condition. Since
$$
\frac{m}{\tau}\sqrt n(P_n^{*}-P_n)
\stackrel[*]{P}{\to}\tilde{\G}_0
\quad\text{in }\ell^\infty(\cF),
$$
Theorem~2.6 in \cite{Kosorok2008} gives, for every $\varphi\in\mathrm{BL}_1(\ell^\infty(\cF))$,
$$
\mathsf E_{*}\overline{
\varphi\left(\frac{m}{\tau}\sqrt n(P_n^{*}-P_n)\right)}
-
\mathsf E_{*}\underline{
\varphi\left(\frac{m}{\tau}\sqrt n(P_n^{*}-P_n)\right)}
\convp0.
$$
Since $\varphi$ is bounded by one and Lipschitz with constant at most one,
$$
\left|
\varphi(\tilde{\G}_n^{*})
-
\varphi\left(\frac{m}{\tau}\sqrt n(P_n^{*}-P_n)\right)
\right|
\le
\min\{D_n^{*},2\}.
$$
Taking majorants and minorants gives
\begin{align*}
\mathsf E_{*}\overline{\varphi(\tilde{\G}_n^{*})}
-
\mathsf E_{*}\underline{\varphi(\tilde{\G}_n^{*})}
\le\;&
\mathsf E_{*}\overline{
\varphi\left(\frac{m}{\tau}\sqrt n(P_n^{*}-P_n)\right)}
-
\mathsf E_{*}\underline{
\varphi\left(\frac{m}{\tau}\sqrt n(P_n^{*}-P_n)\right)}
\\
&+
2\mathsf E_{*}\overline{\min\{D_n^{*},2\}}.
\end{align*}
Moreover, for every $\varepsilon>0$,
$$
\mathsf E_{*}\overline{\min\{D_n^{*},2\}}
\le
\varepsilon+2\mathsf P_{*}(D_n^{*}>\varepsilon)
\convp \varepsilon
$$
by \eqref{eq:bootstrap_comp_process_remainder}. Since $\varepsilon$ is arbitrary,
$$
\mathsf E_{*}\overline{\varphi(\tilde{\G}_n^{*})}
-
\mathsf E_{*}\underline{\varphi(\tilde{\G}_n^{*})}
\convp0.
$$
This concludes the proof.
\end{proof}

\begin{proof}[Proof of Theorem~\ref{thm:bootstrap_test_statistics}]
The proof is divided into \ref{thm:bootstrap_test_statistics_i} and \ref{thm:bootstrap_test_statistics_ii}.

\emph{Proof of \ref{thm:bootstrap_test_statistics_i}.} The asymptotic distribution of the Kolmogorov--Smirnov statistic follows from Theorem~\ref{thm:bootstrap_processes}\ref{thm:bootstrap_processes_i} by the same argument used in the proof of Theorem~\ref{thm:test_statistics_simp}.

For the Cramér--von Mises statistic under the simple null hypothesis, the difference with the proof of Theorem~\ref{thm:test_statistics_simp} is that the convergence $\G_{n,\mu}^{\mu_0}\convl\G_0$ is replaced by
$\G_n^{*}\stackrel[*]{P}{\to}\G_0$ in $\ell^\infty(\cY).$
The only point that is not verbatim is the application of Skorohod's construction. To justify it, consider the sequence $(\G_n^{*})_{n\ge1}$ as random elements in $\ell^\infty(\cY)$, and let $(n_j)_{j\ge1}$ be an arbitrary subsequence of positive integers. There exists a further subsequence $(n_{j_k})_{k\ge1}$ such that, for almost every realization of the data,
$$
\sup_{\varphi\in\mathrm{BL}_1(\ell^\infty(\cY))}
\left|
\mathsf E_{*}\varphi(\G_{n_{j_k}}^{*})
-
\mathsf E\varphi(\G_0)
\right|
\to0.
$$
Fix a realization of the data for which this convergence holds and for which Lemma~\ref{ap:lem_DCT_P_hat} applies along the same subsequence. Conditionally on this data sequence, the law of $\G_{n_{j_k}}^{*}$ converges weakly to the law of $\G_0$ in $\ell^\infty(\cY)$. Therefore, the same Skorohod construction used in the proof of Theorem~\ref{thm:test_statistics_simp} can be applied. From this point, the same argument as in the proof of Theorem~\ref{thm:test_statistics_simp} gives, conditionally on the fixed data realization, $T_{n_{j_k}}^{*\mathrm{CM}}\convl T_0$, where
$$
T_0
\defin
\int_{\Omega\times\R_{\ge0}}\G_0(\omega,t)^2\,\rd F_{\omega}^{\mu_0}(t)\,\rd\mu_0(\omega).
$$
Equivalently,
\begin{align}\label{eq:bootstrap_CvM_subseq_BL_convergence}
\sup_{\varphi\in\mathrm{BL}_1(\R)}
\left|
\mathsf E_{*}\varphi(T_{n_{j_k}}^{*\mathrm{CM}})
-
\mathsf E\varphi(T_0)
\right|
\to0.
\end{align}
This result holds almost surely with respect to the data. Since the original subsequence $(n_{j})_{j\ge1}$ was arbitrary, every subsequence of the bounded Lipschitz distance has a further subsequence converging to zero almost surely. Therefore, 
\begin{align*}
\sup_{\varphi\in\mathrm{BL}_1(\R)}
\left|
\mathsf E_{*}\varphi(T_{n}^{*\mathrm{CM}})
-
\mathsf E\varphi(T_0)
\right|
\convp0.
\end{align*}
Additionally, observe that conditionally on the data, $T_{n}^{*\mathrm{CM}}$ is a finite sum of measurable functions of the multipliers. Hence, for every $\varphi\in\mathrm{BL}_1(\R)$, the asymptotic measurability condition 
$$
\mathsf E_{*}\overline{\varphi(T_{n}^{*\mathrm{CM}})}
-
\mathsf E_{*}\underline{\varphi(T_{n}^{*\mathrm{CM}})}
=0
$$
is automatically satisfied. Hence,
$$
T_n^{*\mathrm{CM}}
\stackrel[*]{P}{\to}
\int_{\Omega\times\R_{\ge0}}\G_0(\omega,t)^2\,\rd F_{\omega}^{\mu_0}(t)\,\rd\mu_0(\omega).
$$

\emph{Proof of \ref{thm:bootstrap_test_statistics_ii}.} For the Kolmogorov--Smirnov statistic, the result follows from Theorem~\ref{thm:bootstrap_processes}\ref{thm:bootstrap_processes_ii} using the same argument as in the proof of Theorem~\ref{thm:test_statistics_comp}.


The asymptotic distribution of the Cramér--von Mises statistic follows by the same argument as in the proof of \ref{thm:bootstrap_test_statistics_i}, replacing Theorem~\ref{thm:test_statistics_simp} by Theorem~\ref{thm:test_statistics_comp}, and using $\tilde{\G}_n^{*}\stackrel[*]{P}{\to}\tilde{\G}_0$ in $\ell^\infty(\cF),$ which follows from Theorem~\ref{thm:bootstrap_processes}\ref{thm:bootstrap_processes_ii}. Conditionally on the data, $\hat{\btheta}^{*}$ is measurable as a function of the multipliers and, for every $(\omega,t)$, the map $\btheta\mapsto F_\omega^{\mu_{\btheta}}(t)$ is Borel measurable. Hence, $\tilde T_n^{*\mathrm{CM}}$ is a finite sum of measurable functions of the multipliers. Therefore, the asymptotic measurability condition
$$
\mathsf E_{*}\overline{\varphi(\tilde T_n^{*\mathrm{CM}})} - \mathsf E_{*}\underline{\varphi(\tilde T_n^{*\mathrm{CM}})} =0
$$
for every $\varphi\in\mathrm{BL}_1(\R)$ holds.
\end{proof}

%---------------------------%
\section{Proofs for specific parametric models}\label{ap:sec_specific_parametric_models}
%---------------------------%

%---------------------------%
\subsection{Proof of Proposition \ref{prop:bahadur_vmf}}
%---------------------------%
\label{ap:sec_bahadur_vmf}

\begin{proof}[Proof of Proposition \ref{prop:bahadur_vmf}]
We now show that the expansion \eqref{eq:bahadur_vMF} satisfies the requirements of Theorem 5.41 of \cite{vandervaart1998}. First, since $\bxi_0\ne0$, the map $\bxi \mapsto \psi_{\bxi}(\bx)$ is twice continuously differentiable in a neighborhood of $\bxi_0$, for every $\bx \in \mathbb{S}^q$.
Moreover, since $\mathsf{E}_{\bxi_0}(\bX)=A_q(\kappa_0)\bmu_0$, it follows that $\mathsf{E}_{\bxi_0}\left(\psi_{\bxi_0}(\bX)\right)=0$.

Second, for a given sample $\mathcal{L}_n = \{\bX_1, \ldots, \bX_n\}$, the MLE is a consistent estimator and satisfies
$$
\begin{aligned}%\label{eq:bahadur_vMF_cond_one}
\sum_{i=1}^n \psi_{\hat{\bxi}_n}(\bX_i) = 0.
\end{aligned}
$$

To see that $\mathsf{E}\left(\|\psi_{\bxi_0}(\bX)\|^2\right) < \infty$, observe that
\begin{align}\label{eq:vMF_psi_norm}
\mathsf{E}\left(\|\psi_{\bxi_0}(\bX)\|^2\right)
=
\mathsf{E}\left(\|\bX-A_q(\kappa_0)\bmu_0\|^2\right)
=
1-A_q(\kappa_0)^2\le 1<\infty.
\end{align}

Since $\psi_{\bxi}$ is the derivative of the log-likelihood function with respect to $\bxi$, it holds that $\mathsf{P}_{\bxi_0}(\dot{\psi}_{\bxi_0}) = -\mathcal{I}_{\bxi_0}$, where $\mathcal{I}_{\bxi_0}$ is the Fisher information matrix. For a $\mathrm{vMF}(\bmu, \kappa)$, this matrix is positive-definite (thus non-singular).

It is only left to check that there exists an integrable function $g(\bx)$ such that $\|\Ddot{\psi}(\bx)\| \le g(\bx)$ for every $\bxi$ in a neighborhood of $\bxi_0$. Observe that the second-order derivatives obtained by differentiating \eqref{eq:vMF_der_psi} with respect to $\bxi$ do not depend on $\bx$ and can therefore be bounded by a constant in a neighborhood of $\bxi_0$. To show that this constant is finite, we show the boundedness of the four quantities involved: $\|\bxi\|$, $A_q(\|\bxi\|)$, $A_q^\prime(\|\bxi\|)$, and $A_q^{\prime\prime}(\|\bxi\|)$. %where $A^{\prime\prime}$ is the second derivative of $A_q$.

Since $\|\bxi_0\|>0$, there exists a neighborhood $U$ of $\bxi_0$ and constants $0<c\le C<\infty$ such that $c\le \|\bxi\|\le C$ for every $\bxi\in U$. Since $A_q$, $A_q^\prime$, and $A_q^{\prime\prime}$ are continuous on $(0,\infty)$, they are bounded on $[c,C]$.
\end{proof}

%---------------------------%
\subsection{Proof of Proposition \ref{prop:verification_assumptions_specific_models}}
%---------------------------%
\label{ap:sec_verification_assumptions_specific_models}

\begin{proof}[Proof of Proposition \ref{prop:verification_assumptions_specific_models}]
In all the models considered, the map $\btheta\mapsto F_\omega^{\mu_{\btheta}}(t)$ is differentiable in a neighborhood of the true parameter, for every $(\omega,t)\in\Omega\times\mathbb R_{\ge0}$.
We apply the fundamental theorem of calculus. For $\btheta$ sufficiently close to $\btheta_0$, let $\btheta_s=\btheta_0+s(\btheta-\btheta_0)$, $s\in[0,1]$. Then
$$
F_{\omega}^{\mu_{\btheta}}(t)-F_{\omega}^{\mu_{\btheta_0}}(t)
-
\dot F_{\omega}^{\btheta_0}(t)^\top(\btheta-\btheta_0)
=
(\btheta-\btheta_0)^\top
\int_0^1
\left(\dot F_{\omega}^{\btheta_s}(t)-\dot F_{\omega}^{\btheta_0}(t)\right)
\rd s.
$$
Consequently, to verify Assumption~\ref{d}, it is enough to prove that, in a neighborhood of the true parameter, the map
${\btheta} \mapsto \dot F_{\omega}^{\btheta}(t)$ is locally Lipschitz in $\btheta$, uniformly over $(\omega,t)$. That is,
$$
\sup_{\omega,t}
\left\|
\dot F_{\omega}^{\btheta}(t)-\dot F_{\omega}^{\btheta_0}(t)
\right\|
\le L\|\btheta-\btheta_0\|,
$$
for some $L<\infty$.

% The normal distribution on $\mathbb{R}$.

Consider $X \sim \mathcal{N}(\mu,\sigma_0^2)$, where $\sigma_0>0$ is fixed and $\mu \in \R$ is unknown.
 For any $\omega \in \R$, we have
\begin{align}\label{eq:dot_F_normal_mu}
\dot{F}_{\omega}^\mu(t) \defin \frac{\partial}{\partial \mu} F_{\omega}^\mu(t) = - \frac{1}{\sigma_0}\phi\left(\frac{\omega - \mu + t}{\sigma_0}\right) + \frac{1}{\sigma_0}\phi\left(\frac{\omega - \mu - t}{\sigma_0}\right).
\end{align}

Let $\mu_s=\mu_0+s(\mu-\mu_0)$. By the expression for $\dot F_{\omega}^{\mu}(t)$ in \eqref{eq:dot_F_normal_mu} and the boundedness of $\phi^\prime$,
$$
\sup_{\omega\in\R, t\ge0}
\left|
\dot F_{\omega}^{\mu_s}(t)-\dot F_{\omega}^{\mu_0}(t)
\right|
\le L|\mu_s-\mu_0|
\le L|\mu-\mu_0|,
$$
where
$$
L\defin \frac{2}{\sigma_0^2}\sup_{u\in\R}|\phi'(u)|
=\frac{2\,\phi(1)}{\sigma_0^2}
=\frac{2}{\sigma_0^2\sqrt{2\pi e}}.
$$
Hence,
$$
\sup_{\omega\in\R, t\ge0}
\left|
F_{\omega}^{\mu}(t)-F_{\omega}^{\mu_0}(t)
-
\dot F_{\omega}^{\mu_0}(t)(\mu-\mu_0)
\right|
\le L|\mu-\mu_0|^2
=o(|\mu-\mu_0|),
$$
so Assumption~\ref{d} holds.

If $X \sim \mathcal{N}(\mu_0,\sigma^2)$, where $\mu_0 \in \mathbb{R}$ is fixed and $\sigma>0$ is unknown, let $\sigma_s=\sigma_0+s(\sigma-\sigma_0)$, and take a neighborhood of $\sigma_0$ with $\sigma\ge \underline\sigma>0$.
For any $\omega \in \R$, we have
\begin{align}\label{eq:dot_F_normal_sigma}
\dot{F}_{\omega}^\sigma(t) \defin \frac{\partial}{\partial \sigma} F_{\omega}^\sigma(t) = - \frac{\omega - \mu_0 + t}{\sigma^2}\phi\left(\frac{\omega - \mu_0 + t}{\sigma}\right) + \frac{\omega - \mu_0 - t}{\sigma^2} \phi\left(\frac{\omega - \mu_0 - t}{\sigma}\right).
\end{align}

Write $h(u)=u\phi(u)$, so that \eqref{eq:dot_F_normal_sigma} becomes
$$
\dot F_{\omega}^{\sigma}(t) = -\frac{1}{\sigma} h\left(\frac{\omega-\mu_0+t}{\sigma}\right) + \frac{1}{\sigma} h\left(\frac{\omega-\mu_0-t}{\sigma} \right).
$$
Moreover,
$$
\frac{\partial}{\partial\sigma}
\left\{\sigma^{-1}h\left(u\right)\right\}
=
-\sigma^{-2}
\left\{
h(u)+u h^\prime(u)
\right\},
\qquad
u=u(\sigma)\defin\frac{\omega-\mu_0\pm t}{\sigma}.
$$
Since $h$ and $u h^\prime(u)$ are bounded on $\R$, it follows that
$$
\sup_{\omega\in\R, t\ge0}
\left|
\dot F_{\omega}^{\sigma_s}(t)-\dot F_{\omega}^{\sigma_0}(t)
\right|
\le L|\sigma_s-\sigma_0|
\le L|\sigma-\sigma_0|,
$$
where
$$
L\defin \frac{2}{\underline\sigma^2}
\sup_{u\in\R}|h(u)+u h'(u)|
=\frac{2}{\underline\sigma^2}\sup_{u\in\R}|(2u-u^3)\phi(u)|<\infty.
$$
Therefore, Assumption~\ref{d} holds.

Consider $X \sim \mathcal{N}(\mu,\sigma^2)$, where both $\mu \in \mathbb{R}$ and $\sigma>0$ are unknown, and let $\btheta=(\mu,\sigma)^\top$.
Take a neighborhood $U$ of $\btheta_0=(\mu_0,\sigma_0)^\top$ on which $\sigma\ge \underline\sigma>0$.
We have already proved that the second-order partial derivatives $\partial_{\mu\mu}F_{\omega}^{\btheta}(t)$ and $\partial_{\sigma\sigma}F_{\omega}^{\btheta}(t)$ are uniformly bounded on $\R\times\R_{\ge0}\times U$.
It remains to control the mixed partial derivatives. Using \eqref{eq:dot_F_normal_mu} and differentiating with respect to $\sigma$, we obtain
$$
\partial_{\mu\sigma}F_{\omega}^{\btheta}(t)
=
\partial_{\sigma\mu}F_{\omega}^{\btheta}(t)
=
\frac{1}{\sigma^2}h'\left(\frac{\omega-\mu+t}{\sigma}\right)
-
 \frac{1}{\sigma^2}h'\left(\frac{\omega-\mu-t}{\sigma}\right).
$$
Hence,
$$
\sup_{\omega\in\R,\,t\ge0,\,\btheta\in U}
\left|\partial_{\mu\sigma}F_{\omega}^{\btheta}(t)\right|
\le
\frac{2}{\underline\sigma^2}
\sup_{u\in\R}\left|h'(u)\right|<\infty.
$$
Therefore, all entries of the Hessian matrix are uniformly bounded.
Consequently, using the same reasoning as in the previous two cases, there exists $L<\infty$ such that
$$
\sup_{\omega\in\R, t\ge0}
\left\|
\dot F_{\omega}^{\btheta_s}(t)-\dot F_{\omega}^{\btheta_0}(t)
\right\|
\le L\|\btheta-\btheta_0\|,
$$
and Assumption~\ref{d} follows.



% The von Mises--Fisher distribution on $\mathbb{S}^{q}$.

Consider the canonical parameter $\bxi=\kappa\bmu \in \R^{q+1}\setminus\{0\}$.
For any $(\bomega,t)$, compute $\dot{F}_{\bomega}^{\bxi}(t)$ as follows:
\begin{equation}\label{eq:F_dot_vmf}
\begin{aligned}
\dot{F}_{\bomega}^{\bxi}(t)
&=
\nabla_{\bxi}\E{y_{\bomega, t}(\bX)}
\\
&=
\int_{\mathbb{S}^q} y_{\bomega, t}(\bx)\nabla_{\bxi}f_{\bxi}(\bx) \, \rd \bx
\\
&=
\int_{\mathbb{S}^q} y_{\bomega, t}(\bx) f_{\bxi}(\bx) \nabla_{\bxi}\log\left(f_{\bxi}(\bx)\right) \, \rd \bx
\\
&=
\Ebig{y_{\bomega, t}(\bX) \psi_{\bxi}(\bX)}.
\end{aligned}
\end{equation}
The second equality in \eqref{eq:F_dot_vmf} follows by an application of the Dominated Convergence Theorem (DCT) \citep[see, e.g., Exercise 51 in][]{DeBarra2003}. To show that the DCT applies, it suffices to show that $\nabla_{\bxi} f_{\bxi}(\bx)$ is locally uniformly (in $\bxi$) dominated by an integrable function $g$ (say). Obviously, $|y_{\bomega, t}(\bx)| \le 1$ for all $\bx$. It holds that $\nabla_{\bxi}f_{\bxi}(\bx) = f_{\bxi}(\bx) \left(\bx - A_q(\kappa)\bmu\right)$, so $\|\nabla_{\bxi}f_{\bxi}(\bx)\| \le 2 f_{\bxi}(\bx)$. Take any $\bxi_0 \in \R^{q+1}\setminus\{0\}$. For $\bxi$ in a small neighborhood of $\bxi_0$, there exists $0<C<\infty$ such that $f_{\bxi}(\bx) \le C f_{\bxi_0}(\bx)$, so $g(\bx) = 2C f_{\bxi_0}(\bx)$, which is integrable. Thus, $\nabla_{\bxi}f_{\bxi}(\bx)$ is locally uniformly (in $\bxi$) bounded by $g$, and the DCT applies.

Let $\bxi_0$ be the true population parameter and let $K$ be a compact neighborhood of $\bxi_0$ contained in $\R^{q+1}\setminus\{0\}$. For $\bxi$ sufficiently close to $\bxi_0$, $\bxi_s=\bxi_0+s(\bxi-\bxi_0)$ belongs to $K$ for all $s\in[0,1]$. By the formula for $\dot F_{\bomega}^{\bxi}(t)$ in \eqref{eq:F_dot_vmf},
$$
\dot F_{\bomega}^{\bxi_s}(t)-\dot F_{\bomega}^{\bxi_0}(t)
=
\int_{\mathbb{S}^q}
y_{\bomega,t}(\bx)
\left(
\nabla_{\bxi}f_{\bxi_s}(\bx)-\nabla_{\bxi}f_{\bxi_0}(\bx)
\right)
\rd \bx.
$$
Since $K\times\mathbb{S}^q$ is compact and $f_{\bxi}(\bx)$ is twice continuously differentiable in $(\bxi, \bx)$ on $K \times \mathbb{S}^q$, one has
$$
\sup_{\bxi\in K, \bx\in\mathbb{S}^q}
\left\|
\nabla_{\bxi}^2 f_{\bxi}(\bx)
\right\|<\infty.
$$
Using also that $|y_{\bomega,t}(\bx)|\le1$, there exists $L<\infty$ such that
$$
\sup_{\bomega\in\mathbb{S}^q, t\ge0}
\left\|
\dot F_{\bomega}^{\bxi_s}(t)-\dot F_{\bomega}^{\bxi_0}(t)
\right\|
\le L\|\bxi_s-\bxi_0\|
\le L\|\bxi-\bxi_0\|.
$$
Therefore, Assumption~\ref{d} holds.
\end{proof}

%---------------------------%
\subsection{Proof of Proposition \ref{prop:covariance_specific_models}}\label{ap:sec_covariance_specific_models}
%---------------------------%

\begin{proof}[Proof of Proposition \ref{prop:covariance_specific_models}]
We divide the proof into its four parts, \ref{prop:covariance_specific_models_i}, \ref{prop:covariance_specific_models_ii}, \ref{prop:covariance_specific_models_iii}, and \ref{prop:covariance_specific_models_iv}.

% The normal distribution, mean and scale.

\emph{Proof of \ref{prop:covariance_specific_models_i}.} Consider $X \sim \mathcal{N}(\mu_0, \sigma^2)$ for a fixed $\sigma>0$. For any $\omega \in \R$, we have
\begin{align*}
\dot{F}_{\omega}^{\mu_0}(t) \defin \left.\frac{\partial}{\partial \mu} F_{\omega}^\mu(t)\right|_{\mu=\mu_0} = - \frac{1}{\sigma}\phi\left(\frac{\omega - \mu_0 + t}{\sigma}\right) + \frac{1}{\sigma}\phi\left(\frac{\omega - \mu_0 - t}{\sigma}\right),
\end{align*}
where $\phi$ is the PDF of a standard normal distribution.

The MLE for $X$ given a sample $\mathcal{L}_n$ is $\hat{\mu} = \bar{X}$, and $\psi_{\mu_0}(x) = \frac{x-\mu_0}{\sigma^2}$. Observe that $\mathsf{E}(\psi_{\mu_0}(X)) = 0$, $\mathsf{E}(\psi_{\mu_0}(X)^2) = \frac{1}{\sigma^2} = \mathcal{I}_{\mu_0}$.

Now, to compute the remaining addends in the covariance process \eqref{eq:cov_comp}, a simple calculation gives
\begin{align*}
\mathsf{E}\left(y_{\omega,t}(X)\psi_{\mu_0}(X)\right)
&=
\frac{1}{\sigma^2}
\int_{\frac{\omega-\mu_0-t}{\sigma}}^{\frac{\omega-\mu_0+t}{\sigma}} u \phi(u)\sigma \, \rd u
\\
&=
- \frac{1}{\sigma}\phi\left(\frac{\omega - \mu_0 + t}{\sigma}\right)
+ \frac{1}{\sigma}\phi\left(\frac{\omega - \mu_0 - t}{\sigma}\right)
\\
&=
\dot{F}_{\omega}^{\mu_0}(t).
\end{align*}
Therefore,
$$
\begin{aligned}
\Cov{\G(f_{\omega_1,t_1})}{\G(f_{\omega_2,t_2})} = \mathcal{C}_{(\omega_1,t_1), (\omega_2, t_2)} -
\sigma^2 \dot{F}_{\omega_1}^{\mu_0}(t_1) \dot{F}_{\omega_2}^{\mu_0}(t_2).
\end{aligned}
$$


\emph{Proof of \ref{prop:covariance_specific_models_ii}.} Consider $X \sim \mathcal{N}(\mu, \sigma_0^2)$ for a fixed $\mu\in\R$. For any $\omega \in \R$, we have
\begin{align*}
\dot{F}_{\omega}^{\sigma_0}(t) \defin \left.\frac{\partial}{\partial \sigma} F_{\omega}^\sigma(t)\right|_{\sigma=\sigma_0} = - \frac{\omega - \mu + t}{\sigma_0^2}\phi\left(\frac{\omega - \mu + t}{\sigma_0}\right) + \frac{\omega - \mu - t}{\sigma_0^2} \phi\left(\frac{\omega - \mu - t}{\sigma_0}\right),
\end{align*}

The MLE for $\sigma$ given a sample $\mathcal{L}_n$ is $\hat{\sigma} = \sqrt{\frac{1}{n} \sum_{i=1}^n (X_i-\mu)^2}$, and $\psi_{\sigma_0}(x) = \frac{(x-\mu)^2 - \sigma_0^2}{\sigma_0^3}$. Observe that $\mathsf{E}(\psi_{\sigma_0}(X)) = 0$, $\mathsf{E}(\psi_{\sigma_0}(X)^2) = \frac{2}{\sigma_0^2}$.

Doing some algebra, one obtains
$$
\begin{aligned}
\mathsf{E}\left(y_{\omega,t}(X)\psi_{\sigma_0}(X)\right) &= \frac{1}{\sigma_0} \int_{\frac{\omega-\mu-t}{\sigma_0}}^{\frac{\omega-\mu+t}{\sigma_0}} (u^2-1) \phi(u) \, \rd u
\\
&=
- \frac{\omega - \mu + t}{\sigma_0^2}\phi\left(\frac{\omega - \mu + t}{\sigma_0}\right) + \frac{\omega - \mu - t}{\sigma_0^2} \phi\left(\frac{\omega - \mu - t}{\sigma_0}\right) = \dot{F}_{\omega}^{\sigma_0}(t).
\end{aligned}
$$
Therefore,
$$
\begin{aligned}
\Cov{\G_0(f_{\omega_1,t_1})}{\G_0(f_{\omega_2,t_2})} = \mathcal{C}_{(\omega_1,t_1), (\omega_2, t_2)}
-\frac{\sigma_0^2}{2} \dot{F}_{\omega_1}^{\sigma_0}(t_1) \dot{F}_{\omega_2}^{\sigma_0}(t_2).
\end{aligned}
$$

\emph{Proof of \ref{prop:covariance_specific_models_iii}.} Consider $X \sim \mathcal{N}(\mu_0, \sigma_0^2)$, where both $\mu_0$ and $\sigma_0$ are unknown. The Fisher information matrix is given by $\mathcal{I}_{\mu_0, \sigma_0} = \mathrm{diag}\left(\frac{1}{\sigma_0^2}, \frac{2}{\sigma_0^2}\right)$, a diagonal matrix with $\frac{1}{\sigma_0^2}$ and $\frac{2}{\sigma_0^2}$ in the diagonal. Writing
$$
\dot F_\omega^{\mu_0,\sigma_0}(t)
\defin
\begin{pmatrix}
\dot F_\omega^{\mu_0}(t)\\
\dot F_\omega^{\sigma_0}(t)
\end{pmatrix},
$$
and $\psi_{\mu_0,\sigma_0}\defin(\psi_{\mu_0},\psi_{\sigma_0})^\top$, one has $\mathsf{E}(y_{\omega,t}(X)\psi_{\mu_0,\sigma_0}(X))=\dot F_\omega^{\mu_0,\sigma_0}(t)$ and $\bV_{\mu_0,\sigma_0}=-\mathcal I_{\mu_0,\sigma_0}$. The covariance of the limiting Gaussian process is given by
$$
\begin{aligned}
\Cov{\G_0(f_{\omega_1,t_1})}{\G_0(f_{\omega_2,t_2})}
=
\mathcal{C}_{(\omega_1,t_1), (\omega_2, t_2)}
-
\big(\dot{F}_{\omega_1}^{\mu_0, \sigma_0}(t_1)\big)^\top
\mathcal{I}_{\mu_0, \sigma_0}^{-1}
\dot{F}_{\omega_2}^{\mu_0, \sigma_0}(t_2).
\end{aligned}
$$

% The vMF distribution on $\mathbb{S}^q$, canonical parameter $\bxi = \kappa \bmu$.

\emph{Proof of \ref{prop:covariance_specific_models_iv}.} We provide the explicit form of the covariance process \eqref{eq:cov_comp} for $\bX \sim \mathrm{vMF}(\bmu_0, \kappa_0)$ distribution on $\mathbb{S}^q$. Consider the estimation of the canonical parameter $\bxi_0 = \kappa_0 \bmu_0$ through maximum likelihood. The log-likelihood (for one observation) is given by
$$
\ell(\bxi, \bx) \defin \log\left(c_q^{\mathrm{vMF}}(\|\bxi\|)\right) + \bxi^\top \bx, \ \text{ where } \ c_q^{\mathrm{vMF}}(\|\bxi\|) = \frac{\|\bxi\|^{\frac{q-1}{2}}}{(2\pi)^{\frac{q+1}{2}}\mathcal{I}_{\frac{q-1}{2}}(\|\bxi\|)},
$$

The score function is
\begin{align*}
  \psi_{\bxi}(\bx) \defin \frac{\partial}{\partial \bxi} \ell(\bxi, \bx) = \nabla_{\bxi} \log\left(c_q^{\mathrm{vMF}}(\|\bxi\|)\right)+\bx = \bx - A_q(\|\bxi\|) \frac{\bxi}{\|\bxi\|} = \bx - \mathsf{E}_{\bxi}(\bX),
\end{align*}
where $A_q(\|\bxi\|) = \frac{\mathcal{I}_{\frac{q+1}{2}}(\|\bxi\|)}{\mathcal{I}_{\frac{q-1}{2}}(\|\bxi\|)}$. Because the $\mathrm{vMF}$ is an exponential family with canonical parameter $\bxi$, it holds that $\mathsf{E}_{\bxi}(\psi_{\bxi}(\bX)) = 0$. The likelihood equation for the MLE is
$$
A_q(\|\hat{\bxi}\|)\frac{\hat{\bxi}}{\|\hat{\bxi}\|}=\bar{\bX};
$$
equivalently, when $\bar{\bX}\ne0$,
$$
\hat\bmu=\frac{\bar{\bX}}{\|\bar{\bX}\|},
\qquad
\hat\kappa=A_q^{-1}(\|\bar{\bX}\|),
\qquad
\hat{\bxi}=A_q^{-1}(\|\bar{\bX}\|)\frac{\bar{\bX}}{\|\bar{\bX}\|}.
$$

The first term in \eqref{eq:cov_comp}, that is, $\mathcal{C}_{(\bomega_1, t_1), (\bomega_2, t_2)}$, can be calculated directly from the expression of the distance profile for the $\mathrm{vMF}$ distribution under the chordal and geodesic distance (see Section \ref{subsec:dp_vmf}).

The second quantity of interest is $\dot{F}_{\bomega_2}^{\bxi_0}(t_2)^\top \mathsf{E}\left(y_{\bomega_1, t_1}(\bX) \psi_{\bxi_0}(\bX)\right)$. Observe that
\begin{equation}\label{eq:expectation_y_score_vmf}
\begin{aligned}
\mathsf{E}\left(y_{\bomega, t}(\bX) \psi_{\bxi_0}(\bX)\right)
&=
\mathsf{E}(y_{\bomega, t}(\bX) \bX) - A_q(\kappa_0)\bmu_0 F_{\bomega}^{\bxi_0}(t)
\\
&=
- A_q(\kappa_0)\bmu_0 F_{\bomega}^{\bxi_0}(t)
+ F_{\bomega}^{\bxi_0}(t) \mathsf{E}\left(\bX \mid d(\bomega,\bX)\le t\right)
\\
&=
F_{\bomega}^{\bxi_0}(t)
\left(
\mathsf{E}\left(\bX \mid d(\bomega,\bX)\le t\right) - A_q(\kappa_0)\bmu_0
\right),
\end{aligned}
\end{equation}
with %$f_{\bxi}(\bx)$ the density of $\bX \sim \mathrm{vMF}(\bmu, \kappa)$ as in \eqref{eq:f_vmf}, $\mathcal{S} = \{\bx \in \mathbb{S}^q: \; d(\bomega_1, \bx) \le t_1\}$ and
$A_q(\kappa_0) = \frac{\mathcal{I}_{\frac{q+1}{2}}(\kappa_0)}{\mathcal{I}_{\frac{q-1}{2}}(\kappa_0)}$.
By \eqref{eq:F_dot_vmf} and \eqref{eq:expectation_y_score_vmf},
\begin{equation*}
\dot{F}_{\bomega}^{\bxi_0}(t)
=
F_{\bomega}^{\bxi_0}(t)
\left(
\mathsf{E}\left(\bX \mid d(\bomega,\bX)\le t\right)
-
A_q(\kappa_0)\bmu_0
\right).
\end{equation*}


It is only left to calculate the Fisher information matrix $\mathcal{I}_{\bxi_0}$:
\begin{equation}\label{eq:fisher_info_vmf}
\begin{aligned}
\mathcal{I}_{\bxi_0}
&=
\E{\psi_{\bxi_0}(\bX) \psi_{\bxi_0}(\bX)^\top}
=
\E{(\bX-\E{\bX}) (\bX-\E{\bX})^\top}
=
\mathsf{Var}(\bX)
\\
&=
A_q^\prime(\kappa_0) \bmu_0 \bmu_0^\top
+ \frac{A_q(\kappa_0)}{\kappa_0}\left(\bI_{q+1} - \bmu_0 \bmu_0^\top\right)
\\
&=
\frac{A_q(\kappa_0)}{\kappa_0} \mathbf{I}_{q+1}
+ \left( 1 - A_q(\kappa_0)^2 - \frac{q+1}{\kappa_0} A_q(\kappa_0) \right) \boldsymbol{\mu}_0 \boldsymbol{\mu}_0^\top.
\end{aligned}
\end{equation}
Here, $A_q^\prime(\kappa_0)$ is the derivative of $A_q(\kappa)$ with respect to $\kappa$ evaluated at $\kappa_0$, which satisfies $A_q^\prime(\kappa_0) = 1 - A_q(\kappa_0)^2 - \frac{q}{\kappa_0} A_q(\kappa_0)$. See Section \ref{ap:sec_variance_vmf} for a brief note on the variance of a vMF distribution.

Finally, combine \eqref{eq:cov_simp},  \eqref{eq:expectation_y_score_vmf}, \eqref{eq:F_dot_vmf}, and  \eqref{eq:fisher_info_vmf} with $\bV_{\bxi_0}=\mathsf{E}(\dot\psi_{\bxi_0}(\bX))=-\mathcal I_{\bxi_0}=-\mathsf{Var}(\bX)$ to produce
\begin{align*}
\Cov{\G(f_{\bomega_1, t_1})}{\G(f_{\bomega_2, t_2})}
&=
\mathsf{P}\left(d(\bomega_1, \bX) \le t_1, d(\bomega_2, \bX) \le t_2\right)
\\
&\quad -
F_{\bomega_1}^{\bxi_0}(t_1) F_{\bomega_2}^{\bxi_0}(t_2)
\left(1 + \bm_1^\top \mathsf{Var}(\bX)^{-1} \bm_2 \right),
\end{align*}
where $\bm_i = \mathsf{E}(\bX \mid d(\bX, \bomega_i) \le t_i) - A_q(\kappa_0)\bmu_0$, $i=1,2$, and $\mathsf{Var}(\bX) = \mathcal{I}_{\bxi_0}$ is given in \eqref{eq:fisher_info_vmf}.
\end{proof}




%---------------------------%
\section{The covariance process under the simple null}\label{ap:sec_exact_covariance_simple}
%---------------------------%
\begingroup


\noindent\textcolor{red}{\textbf{DRAFT}}

This section complements Proposition~\ref{prop:covariance_specific_models} by recording exact formulas for the covariance process \eqref{eq:cov_simp} of the limiting Gaussian process $\G$ under the simple null in some tractable cases. For the normal example, we restrict attention to the one-dimensional setting. For the vMF and HvMF examples, we work on $\mathbb{S}^2$ and $\mathbb{H}^2$, respectively; in these cases the conditioning argument reduces the joint probability in \eqref{eq:cov_simp} to a one-dimensional integral involving an angular probability on the circle.

%---------------------------%
\begin{proposition}\label{prop:exact_covariance_simple_examples}
%---------------------------%
Let $\mathcal{C}$ denote the covariance process in \eqref{eq:cov_simp}. The following formulas hold:
\begin{enumerate}[label=(\textit{\roman*}), ref=(\textit{\roman*})]
\item\label{prop:exact_covariance_simple_examples_i} If $X \sim \mathcal{N}(\mu,\sigma^2)$ on $\R$, then
\begin{equation}\label{eq:exact_covariance_simple_normal}
\mathcal{C}_{(\omega_1,t_1),(\omega_2,t_2)}
=
1_{\{L \le U\}}
\left[
\Phi\left(\frac{U-\mu}{\sigma}\right)
-
\Phi\left(\frac{L-\mu}{\sigma}\right)
\right]
-
F_{\omega_1}(t_1)F_{\omega_2}(t_2),
\end{equation}
where
$$
L \defin \max(\omega_1-t_1,\omega_2-t_2), \qquad U \defin \min(\omega_1+t_1,\omega_2+t_2),
$$
and
$$
F_{\omega}(t)
=
\Phi\left(\frac{\omega+t-\mu}{\sigma}\right)
-
\Phi\left(\frac{\omega-t-\mu}{\sigma}\right).
$$

\item\label{prop:exact_covariance_simple_examples_ii} Let $\bX \sim \mathrm{vMF}(\bmu,\kappa)$ on $\mathbb{S}^2$. Consider either the chordal distance or the geodesic distance, and let $0\le t_1,t_2\le 2$ in the chordal case and $0\le t_1,t_2\le \pi$ in the geodesic case. Define
$$
a_i \defin
\begin{cases}
1-\frac{t_i^2}{2}, & \text{for the chordal distance},\\
\cos(t_i), & \text{for the geodesic distance},
\end{cases}
\qquad
U \defin \bomega_1^\top \bX,
$$
where $f_U$ is the density of $U$ given by \eqref{eq:density_t_vmf}, and
\begin{equation}\label{eq:exact_covariance_simple_vmf}
\mathcal{C}_{(\bomega_1,t_1),(\bomega_2,t_2)}
=
\int_{a_1}^{1} \Probbig{\bomega_2^\top \bX \ge a_2 \mid U=u} f_U(u) \, \rd u
-
F_{\bomega_1}(t_1)F_{\bomega_2}(t_2),
\end{equation}

\item\label{prop:exact_covariance_simple_examples_iii} Let $\bX \sim \mathrm{HvMF}(\bmu,\kappa)$ on $\mathbb{H}^2$, endowed with the geodesic distance, and assume that $\rho \defin d_{\mathbb{H}^2}(\bomega_1,\bomega_2)>0$. Let
$$
R \defin d_{\mathbb{H}^2}(\bomega_1,\bX),
$$
where $f_R$ is the density of $R$ given by \eqref{eq:density_zr_hvmf}. Then
\begin{equation}\label{eq:exact_covariance_simple_hvmf}
\mathcal{C}_{(\bomega_1,t_1),(\bomega_2,t_2)}
=
\int_{0}^{t_1} \Prob{d_{\mathbb{H}^2}(\bomega_2,\bX)\le t_2 \mid R=r} f_R(r) \, \rd r
-
F_{\bomega_1}(t_1)F_{\bomega_2}(t_2),
\end{equation}
\end{enumerate}
\end{proposition}

\begin{proof}
We divide the proof into three parts.

\emph{Proof of \ref{prop:exact_covariance_simple_examples_i}.}
The event $\{|X-\omega_i|\le t_i\}$ is the interval $[\omega_i-t_i,\omega_i+t_i]$, $i=1,2$. Therefore,
$$
\{|X-\omega_1|\le t_1,\ |X-\omega_2|\le t_2\}
=
\{L\le X\le U\}
$$
when $L\le U$, and it is empty otherwise. This gives
$$
\Prob{|X-\omega_1|\le t_1,\ |X-\omega_2|\le t_2}
=
1_{\{L \le U\}}
\left[
\Phi\left(\frac{U-\mu}{\sigma}\right)
-
\Phi\left(\frac{L-\mu}{\sigma}\right)
\right].
$$
The expression for $F_\omega(t)$ is immediate. Combining both identities with \eqref{eq:cov_simp} yields \eqref{eq:exact_covariance_simple_normal}.

\emph{Proof of \ref{prop:exact_covariance_simple_examples_ii}.}
Set
$$
\rho \defin \bomega_1^\top \bomega_2, \qquad m_1 \defin \bmu^\top \bomega_1, \qquad \lambda(u) \defin \kappa\sqrt{(1-u^2)(1-m_1^2)}.
$$
For both the chordal and the geodesic distance,
$$
\{d(\bomega_i,\bX)\le t_i\} = \{\bomega_i^\top \bX \ge a_i\}, \qquad i=1,2.
$$
Hence,
$$
\Prob{d(\bomega_1,\bX)\le t_1,\ d(\bomega_2,\bX)\le t_2}
=
\Prob{\bomega_1^\top \bX \ge a_1,\ \bomega_2^\top \bX \ge a_2}.
$$
If $\rho=1$, then $\bomega_2=\bomega_1$, so
$$
\Prob{\bomega_1^\top \bX \ge a_1,\ \bomega_2^\top \bX \ge a_2}
=
\Prob{U \ge \max(a_1,a_2)}.
$$
If $\rho=-1$, then $\bomega_2=-\bomega_1$, so
$$
\Prob{\bomega_1^\top \bX \ge a_1,\ \bomega_2^\top \bX \ge a_2}
=
\Prob{a_1 \le U \le -a_2}.
$$
In both cases, \eqref{eq:exact_covariance_simple_vmf} holds.

Assume now that $\rho\in(-1,1)$, and define
$$
\bbe \defin \frac{\bomega_2-\rho \bomega_1}{\sqrt{1-\rho^2}}.
$$
Then $\bomega_2=\rho \bomega_1+\sqrt{1-\rho^2}\,\bbe$. Write
$$
\bX = U\bomega_1 + \sqrt{1-U^2}\,\bY,
$$
where $\bY\in\mathbb{S}^1\cap \bomega_1^\perp$. By \eqref{eq:density_t_vmf} with $q=2$ and $\bomega=\bomega_1$, the variable $U$ has density $f_U$.

If $|m_1|<1$, let
$$
\widetilde{\bmu}_1 \defin \frac{\bmu-m_1\bomega_1}{\sqrt{1-m_1^2}}.
$$
The proof of Proposition~\ref{prop:dp_examples} shows that, conditionally on $U=u$, the density of $\bY$ on $\mathbb{S}^1\cap \bomega_1^\perp$ is proportional to
$$
\exp\left\{\lambda(u)\widetilde{\bmu}_1^\top \bY\right\}.
$$
Choose $\bbe_\perp\in\mathbb{S}^1\cap \bomega_1^\perp$ orthogonal to $\bbe$ so that
$$
\widetilde{\bmu}_1 = \cos(\alpha)\,\bbe + \sin(\alpha)\,\bbe_\perp.
$$
Writing $\bY=\cos(\varphi)\,\bbe+\sin(\varphi)\,\bbe_\perp$, one gets that $\varphi\mid U=u$ follows a von Mises distribution with density
$$
\frac{1}{2\pi\mathcal{I}_0(\lambda(u))}e^{\lambda(u)\cos(\varphi-\alpha)}, \qquad \varphi\in[-\pi,\pi).
$$
Moreover,
$$
\bomega_2^\top \bX
=
\rho U + \sqrt{1-\rho^2}\sqrt{1-U^2}\,\bbe^\top \bY
=
\rho U + \sqrt{1-\rho^2}\sqrt{1-U^2}\cos(\varphi).
$$
Define
$$
b(u) \defin \frac{a_2-\rho u}{\sqrt{1-\rho^2}\sqrt{1-u^2}}.
$$
Therefore, conditionally on $U=u$,
$$
\{\bomega_2^\top \bX \ge a_2\}
=
\{\cos(\varphi)\ge b(u)\}.
$$
Hence, if $|m_1|<1$ and $-1<b(u)<1$,
\begin{equation}\label{eq:conditional_prob_vmf_simple}
\Probbig{\bomega_2^\top \bX \ge a_2 \mid U=u}
=
\frac{1}{2\pi\mathcal{I}_0(\lambda(u))}
\int_{-\arccos(b(u))}^{\arccos(b(u))}
e^{\lambda(u)\cos(\varphi-\alpha)} \, \rd \varphi.
\end{equation}
If $|m_1|<1$ and $b(u)\le -1$, then $\Prob{\bomega_2^\top \bX \ge a_2 \mid U=u}=1$, whereas if $|m_1|<1$ and $b(u)\ge 1$, then $\Prob{\bomega_2^\top \bX \ge a_2 \mid U=u}=0$.

If $|m_1|=1$, then $\lambda(u)=0$, so $\varphi\mid U=u$ is uniform on the circle. Therefore, if $-1<b(u)<1$,
\begin{equation}\label{eq:conditional_prob_vmf_simple_uniform}
\Probbig{\bomega_2^\top \bX \ge a_2 \mid U=u}
=
\frac{\arccos(b(u))}{\pi},
\end{equation}
whereas the conditional probability is again equal to $1$ when $b(u)\le -1$ and equal to $0$ when $b(u)\ge 1$. Combining \eqref{eq:conditional_prob_vmf_simple} and \eqref{eq:conditional_prob_vmf_simple_uniform} with the density $f_U$ given in \eqref{eq:density_t_vmf} and with the definition of $\mathcal{C}$ in \eqref{eq:cov_simp} yields \eqref{eq:exact_covariance_simple_vmf}.

\emph{Proof of \ref{prop:exact_covariance_simple_examples_iii}.}
Set
$$
\rho \defin d_{\mathbb{H}^2}(\bomega_1,\bomega_2), \qquad \rho_\mu \defin d_{\mathbb{H}^2}(\bomega_1,\bmu), \qquad \lambda_H(r) \defin \kappa\sinh(r)\sinh(\rho_\mu).
$$
For completeness, if $\rho=0$, then $\bomega_1=\bomega_2$, and the joint event is $\{R\le \min(t_1,t_2)\}$. This boundary case is omitted from the statement.

Assume from now on that $\rho>0$. Choose a hyperbolic transformation $\bA$ such that
$$
\bA\bomega_1=(1,0,0)^\top, \qquad \bA\bomega_2=(\cosh(\rho),\sinh(\rho),0)^\top.
$$
Write
$$
\bA\bX = \big(\cosh(R),\sinh(R)\cos(\varphi),\sinh(R)\sin(\varphi)\big)^\top.
$$
If $\rho_\mu>0$, write also
$$
\bA\bmu
=
\left(
\cosh(\rho_\mu),
\sinh(\rho_\mu)\cos(\alpha_H),
\sinh(\rho_\mu)\sin(\alpha_H)
\right)^\top.
$$
Then, by the Minkowski inner product,
$$
(\bA\bmu,\bA\bX)
=
-\cosh(\rho_\mu)\cosh(R)
+ \sinh(\rho_\mu)\sinh(R)\cos(\varphi-\alpha_H).
$$
Using the polar-coordinate representation from the proof of Proposition~\ref{prop:dp_examples}, it follows that the joint density of $(R,\varphi)$ is proportional to
$$
e^{\kappa(\bmu,\bomega_1)\cosh(R)}
e^{\lambda_H(R)\cos(\varphi-\alpha_H)}
\sinh(R),
$$
and therefore the marginal density of $R$ is $f_R$, while $\varphi\mid R=r$ follows a von Mises distribution with density
$$
\frac{1}{2\pi\mathcal{I}_0(\lambda_H(r))}e^{\lambda_H(r)\cos(\varphi-\alpha_H)},
\qquad \varphi\in[-\pi,\pi).
$$
If $\rho_\mu=0$, then $\lambda_H(r)=0$ and $\varphi\mid R=r$ is uniform on the circle.

Next, by the distance formula on $\mathbb{H}^2$,
$$
\cosh\big(d_{\mathbb{H}^2}(\bomega_2,\bX)\big)
=
-(\bA\bomega_2,\bA\bX)
=
\cosh(\rho)\cosh(R)-\sinh(\rho)\sinh(R)\cos(\varphi).
$$
Define
$$
b_H(r) \defin \frac{\cosh(\rho)\cosh(r)-\cosh(t_2)}{\sinh(\rho)\sinh(r)}.
$$
Hence, conditionally on $R=r$,
$$
\{d_{\mathbb{H}^2}(\bomega_2,\bX)\le t_2\}
=
\{\cos(\varphi)\ge b_H(r)\}.
$$
Hence, if $\rho_\mu>0$ and $-1<b_H(r)<1$,
\begin{equation}\label{eq:conditional_prob_hvmf_simple}
\Prob{d_{\mathbb{H}^2}(\bomega_2,\bX)\le t_2 \mid R=r}
=
\frac{1}{2\pi\mathcal{I}_0(\lambda_H(r))}
\int_{-\arccos(b_H(r))}^{\arccos(b_H(r))}
e^{\lambda_H(r)\cos(\varphi-\alpha_H)} \, \rd \varphi.
\end{equation}
If $\rho_\mu>0$ and $b_H(r)\le -1$, then $\Prob{d_{\mathbb{H}^2}(\bomega_2,\bX)\le t_2 \mid R=r}=1$, whereas if $\rho_\mu>0$ and $b_H(r)\ge 1$, then $\Prob{d_{\mathbb{H}^2}(\bomega_2,\bX)\le t_2 \mid R=r}=0$.

If $\rho_\mu=0$, then $\lambda_H(r)=0$, so $\varphi\mid R=r$ is uniform on the circle. Therefore, if $-1<b_H(r)<1$,
\begin{equation}\label{eq:conditional_prob_hvmf_simple_uniform}
\Prob{d_{\mathbb{H}^2}(\bomega_2,\bX)\le t_2 \mid R=r}
=
\frac{\arccos(b_H(r))}{\pi},
\end{equation}
whereas the conditional probability is again equal to $1$ when $b_H(r)\le -1$ and equal to $0$ when $b_H(r)\ge 1$. Combining \eqref{eq:conditional_prob_hvmf_simple} and \eqref{eq:conditional_prob_hvmf_simple_uniform} with the density $f_R$ in \eqref{eq:density_zr_hvmf} and with the definition of $\mathcal{C}$ in \eqref{eq:cov_simp} yields \eqref{eq:exact_covariance_simple_hvmf}.
\end{proof}

\endgroup

%---------------------------%
\section{The variance of the vMF distribution}\label{ap:sec_variance_vmf}
%---------------------------%
\cite{Mardia1999} (Equation 10.3.16) show that the variance of $\bX \sim \mathrm{vMF}(\bmu, \kappa)$ can be written as
$$
\mathsf{Var}(\mathbf{X}) = \mathcal{I}_{(\kappa, \bmu)} =
\begin{pmatrix}
A_p^{\prime}(\kappa) & 0 \\
0 & \dfrac{A_p(\kappa)}{\kappa}\left(\mathbf{I}_{q+1} - \boldsymbol{\mu} \boldsymbol{\mu}^T\right)
\end{pmatrix}
$$
Observe that, since $\bmu \in \mathbb{S}^q \subset\R^{q+1}$, this matrix has dimensionality $(q+2) \times (q+2)$. This is because the Fisher information matrix is expressed with respect to $(\kappa, \bmu)$.

The expression for the variance shown in \eqref{eq:fisher_info_vmf} can be derived from the properties of exponential families. For a $\mathrm{vMF}(\bmu,\kappa)$ distribution, this gives
$$
\mathcal{I}_{\bxi} = \frac{\partial}{\partial \bxi^\top} A_q(\kappa)\bmu =  \frac{\partial}{\partial \bxi^\top} A_q(\|\bxi\|)\frac{\bxi}{\|\bxi\|} = A_q^\prime(\kappa) \bmu \bmu^\top + \frac{A_q(\kappa)}{\kappa}\left(\bI - \bmu \bmu^\top\right).
$$
%Using 

\textcolor{blue}{
Diego: Since I am not an expert in vMF distributions, I think it is a good practice to compare this result with \cite{Hillen2017} as a sanity check. \cite[][Lemma 3.1]{Hillen2017}  show that
\begin{equation}\label{eq:sanity_check_variance}
\mathsf{Var}(\bX) = \frac{1}{2}\left(1-\frac{\mathcal{I}_2(k)}{\mathcal{I}_0(k)}\right) \bI_2+\left(\frac{\mathcal{I}_2(k)}{\mathcal{I}_0(k)}-\left(\frac{\mathcal{I}_1(k)}{\mathcal{I}_0(k)}\right)^2\right) \bmu\bmu^T.
\end{equation}
It is easy to see that \eqref{eq:sanity_check_variance} equals \eqref{eq:fisher_info_vmf} when $q=1$, by observing that $\frac{\mathcal{I}_2(\kappa)}{\mathcal{I}_0(\kappa)} = 1 - \frac{2\mathcal{I}_2(\kappa)}{\kappa\mathcal{I}_0(\kappa)} = 1 - \frac{2}{\kappa}A_1(\kappa)$. This follows from the recurrence relation for modified Bessel functions of the first kind $\mathcal{I}_{\mu-1}(\kappa) - \mathcal{I}_{\mu+1}(\kappa) = \frac{2\mu}{\kappa} \mathcal{I}_{\mu}(\kappa)$ for $\mu \in \mathbb{Z}_{\ge 0}$.
}

\end{document}
